% =====================================================================
%  GENERATED; DO NOT EDIT BY HAND.  Built by ebook-method/moderne-filosoffer/
%  rebuild.sh (draft.py --emit, collate/check.py --apply, emarks.py,
%  draft.py --post) from template.tex.  Edit the template, patches.py,
%  decisions.tsv, pagepass.tsv, italics*.tsv, spaced.tsv or last.tsv and
%  rebuild.
%
%  Harald Høffding, MODERNE FILOSOFER.  København: Gyldendalske
%  Boghandel, Nordisk Forlag, 1904.  X + 157 pp.  (Verso: "Trykt i 2200
%  Ekspl."; Aarhus Stiftsbogtrykkeri.)  TRANSCRIPTION (Danish) of the 1904
%  print, complete: title page, Indhold, Indledning, the three groups,
%  Noter.
%
%  (1) SOURCE.  `bibliotek/Høffding, Harald/1904-moderne-filosofer.pdf`
%  (SCANS.tsv, sha256 5a44cfd6...c474): Google's PDF of the University of
%  Minnesota copy (call no. 184 H67mo; HathiTrust umn.319510019950770),
%  as uploaded to the Internet Archive (modern-filosofer), 180 pp., 600 ppi
%  bitonal, with Google's OCR as text layer.  A second physical copy
%  exists (HathiTrust wu.89084975002, Wisconsin; not consulted).
%  PAGE MAP: PDF = printed + 16 (p. 1 = PDF 17, p. 157 = PDF 173); title
%  PDF 9, verso PDF 10, Indhold pp. V-VI = PDF 11-12, Indledning
%  pp. VII-X = PDF 13-16, Noter pp. 148-157 = PDF 164-173.  Every page turn
%  is marked "% --- p. N ---" + \opage{N} at the first word of the page; a
%  word divided over a page turn ends in %.
%
%  (2) CONVENTIONS.  Each text page has a RUNNING HEAD (folio and a short
%  summary of the page, e.g. "Wundts Overgang til Filosofien"): not set,
%  but kept in a comment "% running head: ..." after the page marker, as
%  the layer read it (not proof-read).  Sheet signatures ("H. Høffding:
%  Moderne Filosofer. N") and folios are not transcribed.  Headings:
%  \grouphead (FØRSTE/ANDEN/TREDIE GRUPPE + name), \chapterhead /
%  \chapterheadb (the chapter numeral on the line or above it), \sectionhead
%  (the lettered sections of the Guyau and Nietzsche chapters; the
%  Nietzsche titles are italic).  The notes are ENDNOTES (Noter, small
%  type), marked in the text by a raised figure and a parenthesis, set
%  \textsuperscript{N}); the labels in the Noter ("1) S. 8.") are raised
%  figures in the print, set plain here.  Italic -> \textit; names in SMALL
%  CAPITALS -> \textsc (normal case inside); letterspacing (only in the
%  Noter: the names of cited authors) -> \emph; Greek as Greek (the
%  subsection labels α, β, γ of the Wundt chapter are italic, as is the
%  sign α in Avenarius's formulas, pp. 103-106); » « as printed; the long
%  dash ---.  Raised edition figures (System.\textsuperscript{1},
%  Psychol.\textsuperscript{4}) as printed.
%  THE Ø RULE (Hans, 2026-10-08): where the scan shows a plain o and the
%  o-form is no Danish word, the text has ø (religiøs, Nødvendighed,
%  Følger, religiøse p. IX); where both forms are words (fores/føres) the
%  image decided.  "Iagttagelse": both OCRs read the capital I as l; set I
%  throughout (read at the image, p. 87 and 5 more).
%  PRINTED AS IS (misprints kept): p. VIII "hegge" (begge); p. IX
%  "Nietzche", "kømme"; p. 4 "1880-—83" (hyphen and dash); p. 12 "knnnet";
%  p. 15 "Iagttagelscn"; p. 18 "holde" (holdt); p. 34 "piû" (più); p. 35
%  "hved" (hvad); pp. 57 and 70 "ban" (han); p. 73 "derved." (a point
%  before lower-case "at"); p. 135 "Kærligked"; p. 151 "hau"; p. 152
%  "Absolutc"; p. 155 "mouvemeut", "peuseur"; Indhold "I26" (126) and
%  "Indledning ... I" (the Indledning begins on p. VII).  Also read at the
%  image and kept (the translators' odd spots, 2026-10-14): p. 20
%  "mendens", "alene. ---"; p. 21 "Forestillens og Villens"; p. 22
%  "übereinsstimmt"; p. 25 "Aderitisme"; p. 27 "Ødelanden"; p. 34 "Ide.
%  der"; p. 36 "tutti"; p. 39 "bl. A."; p. 41 "egef"; p. 48 "mærkes",
%  "Ledelsen"; p. 60 "d'étre"; p. 61 "hvilke"; p. 63 "utrumqve"; p. 65
%  "Villen"; p. 67 "systêmes"; p. 72 "hekæmper"; p. 73 "Forestillen"; p. 74
%  "Primitivt,"; p. 77 "abstrakts."; p. 79 "Bontroux", "egentligt"; p. 86
%  "erklærede. at"; p. 90 "Forbindelse. nogen"; p. 91 "Zitren"; p. 99
%  "Begriben"; p. 106 "Avenarins"; p. 110 "cicade"; p. 121 "Villle"; p. 123
%  "Dag. med"; p. 128 "jevne"; p. 129 "Nietsche"; p. 130 "Førster"; p. 131
%  "Delfi's"; p. 137 "Fix Ide"; p. 145 "Barbarina"; p. 148 "dan"; p. 149
%  "Smlgn" (no point); p. 150 "Villen"; p. 151 "an 19e", "Kamp", "perdi",
%  "pua" (twice), "peusiero"; p. 153 "Mensches"; p. 155 "Questione",
%  "Carstanjev"; p. 156 "Arklepios", "Fröliche", "Unzeitgemeine"; p. 157
%  "ndtrykke".
%
%  (3) METHOD AND HOW ACCURACY WAS ESTABLISHED (2026-10-14; HOEFFDING-QUEUE
%  #8; TRANSCRIPTION-PLAYBOOK §00 and §6).  No one typed the text.
%  (a) DRAFT by script from Google's layer (draft.py: lines rebuilt from the
%  word boxes, running heads cut, hyphens joined, paragraphs from indents,
%  endnote marks placed where the Noter assign them).  Twelve places where
%  the layer moved a dash and the words after it onto the line above
%  (pp. 2, 46, 57, 100, 104, 126, 137, 139, 141, 142 twice, 149) were
%  rebuilt at the image (patches.py, pagepass.tsv).
%  (b) COLLATION word by word with punctuation against an INDEPENDENT
%  tesseract reading of the images (collate/check.py, ~1070 disagreements
%  over 51,500 words).  The SAGA / Lindhardt og Ringhof e-book (2020,
%  ISBN 9788726363005; edition-checked: the 1904 text, itself an OCR text)
%  voted as a THIRD witness (rules.py: the reading it agrees with, where it
%  agrees with exactly one; the dashes the layer drops; spacing and marks);
%  never as the base, never for page markers.  The 338 disagreements it
%  did not settle were read BLIND at the image (crop sheets b1, b2: the
%  reader sees neither OCR reading; 60 planted controls, 60 right).  An
%  AUDIT of 78 random e-book votes for the layer's reading (where the two
%  OCRs that share the layer's images could share a fault) found 1 wrong
%  ("Sandbed", p. 133): about 1 in 80, i.e. perhaps 5 such faults among
%  the ~400 unaudited votes.
%  (c) E-BOOK DIFFERENCES: the collated text was compared word by word
%  with the e-book (ebookdiff.py: 217 differences); 139 were read blind at
%  the image (sheets b2, XE items), the rest in the main session; the
%  remaining differences are the e-book's own faults (read at the image),
%  the small capitals it sets as full capitals, and its quotation marks.
%  The e-book's paragraphs were compared too (parcmp.py): 22 disputed
%  paragraph starts read at the image.
%  (d) ITALICS from the e-book's <i> (emarks.py), checked against the
%  stroke slant of every word (check.py --italics): 64 disagreements read
%  blind (sheets i1, 12/12 controls), 15 italic spans added from the image.
%  The e-book's "underline" spans proved NOT to be letterspacing (measured
%  letter gaps, spacecmp.py; and at the image): discarded.  LETTERSPACING
%  occurs only in the Noter (the names of cited authors): every Noter line
%  with a word whose letter gaps measure as spaced was read blind (sheets
%  l1, 45 lines, 10/10 controls; 43 spaced runs).
%  (e) The FRONT MATTER and the INDLEDNING (pp. VII-X, which check.py does
%  not collate) were read whole, blind, at the image and compared word by
%  word (12 dashes the layer dropped, 3 words mended).
%  (f) SAMPLED PAGE CHECK (pagecheck.py): pp. 12, 35, 52, 73, 97, 99, 128,
%  147, 153 (9 pages, one or two per section, one of the Noter; 3053
%  words) transcribed blind at the image and compared word by word: 2 word
%  errors ("Ardigos" for Ardigòs, "Sé ailles"), 1 punctuation error, 3
%  italic omissions, and the Noter's letterspacing (then still to be set):
%  all mended.  About 1 error per 1000 words before the check (95 % upper
%  bound ~3 per 1000).
%  (g) TRANSLATION as a second reading: the six translators listed 117 odd
%  spots; the plausible ones (about 70) were read at the image: 10
%  transcription faults mended (p. 10 the label γ, p. 133 Sandhed, p. 141
%  "Zahl- und", p. 150 hundertjährigen and a false white line, p. 153 the
%  lost "1900.", p. 156 notes 45 and 46 run together, p. 111 a dash set in
%  a heading, p. 123 four lost points and a false paragraph); the rest are
%  the print's own misprints (listed under (2)).
%  The scripts, the decisions and every reading are in
%  ebook-method/moderne-filosoffer/ (README entry in ebook-method/README.md).
% =====================================================================
\documentclass[12pt,a4paper,openany]{book}
\usepackage[utf8]{inputenc}
\usepackage[T1]{fontenc}
\usepackage{libertinus}
\usepackage{libertinust1math}
\usepackage{textalpha}
\usepackage{amsmath}
\usepackage[danish]{babel}
\usepackage{setspace}
\usepackage{fancyhdr}
\setlength{\headheight}{28pt}
\addtolength{\topmargin}{-13.5pt}
\pagestyle{fancy}
\fancyhf{}
\fancyhead[LE,RO]{\thepage}
\fancyhead[RE]{\textit{Harald Høffding: Moderne Filosofer}}
\fancyhead[LO]{\textit{\leftmark}}
\setstretch{1.3}
\newcommand{\whiteline}{\par\vspace{\baselineskip}\par}
\newcommand{\secrule}{\par\begin{center}\rule{2cm}{0.4pt}\end{center}\par}
% the headings (moderne-filosoffer): a group (FØRSTE GRUPPE. + its name), a chapter on one line (I. WILHELM WUNDT.)
% or with its numeral above (II. / ROBERTO ARDIGÒ.), and the lettered sections of the Guyau and Nietzsche chapters
\newcommand{\grouphead}[2]{\clearpage\begin{center}{\Large #1}\\[8pt]{\large #2}\end{center}\addcontentsline{toc}{part}{#1 #2}\bigskip}
\newcommand{\chapterhead}[2]{\bigskip\begin{center}{\large #1 #2}\end{center}\addcontentsline{toc}{chapter}{#1 #2}\markboth{#2}{#2}\medskip}
\newcommand{\chapterheadb}[2]{\bigskip\begin{center}{\large #1}\\[4pt]{\large #2}\end{center}\addcontentsline{toc}{chapter}{#1 #2}\markboth{#2}{#2}\medskip}
\newcommand{\sectionhead}[1]{\par\medskip\begin{center}#1\end{center}\addcontentsline{toc}{section}{#1}\smallskip}
\newcommand{\supplied}[1]{\textlangle#1\textrangle}
\usepackage{graphicx}
\DeclareUnicodeCharacter{0254}{\reflectbox{c}}   % ɔ: (id est), p. 103
\usepackage{hyperref}
\hypersetup{
  pdftitle={Moderne Filosofer},
  pdfauthor={Harald Høffding},
  hidelinks
}
\usepackage{marginnote}
\newcommand{\opage}[1]{\marginnote{\footnotesize\textit{[#1]}}}
\newcommand{\apage}[1]{\marginnote{\footnotesize\textit{[#1]}}}

\begin{document}

\frontmatter
% --- front matter (read blind at the image, reader K, 2026-10-14) ---
\begin{titlepage}
\centering
\vspace*{2cm}
{\Huge MODERNE FILOSOFER}\par
\vspace{1.2cm}
{\footnotesize AF}\par
\vspace{0.8cm}
{\large HARALD HØFFDING}\par
\vfill
{[\textit{publisher's device}]}\par
\vfill
{\large KØBENHAVN}\par
\smallskip
{\small\emph{GYLDENDALSKE BOGHANDEL}}\par% letterspaced in the print (\emph, the house convention)
{\footnotesize \textsc{Nordisk Forlag}}\par
\smallskip
{\footnotesize 1904}\par
\end{titlepage}
\thispagestyle{empty}
\vspace*{\fill}
\begin{center}
\textit{\textbf{Trykt i 2200 Ekspl.}}\par
\vspace{2cm}
{\footnotesize \textsc{Aarhus Stiftsbogtrykkeri.}}
\end{center}
\clearpage
% --- the Indhold (pp. V-VI, PDF 11-12), set as printed (entries, numerals, italics; the dot leaders as \dotfill) ---
\chapter*{INDHOLD.}
\markboth{Indhold}{Indhold}
\opage{V}%
\begin{small}
\newcommand{\ix}[3]{\par\noindent\hspace*{#1}#2\dotfill #3}
\newcommand{\iz}[2]{\par\noindent\hspace*{#1}#2}
\hfill{\footnotesize Side}
\ix{0em}{\textit{Indledning}}{I}
\par\medskip\centerline{FØRSTE GRUPPE.}
\ix{1em}{Objektiv-systematisk Retning}{1}
\iz{1em}{I. \textit{Wilhelm Wundt.}}
\ix{2em}{1. Udviklingsgang}{1}
\ix{2em}{2. Det psykologiske Problem}{6}
\iz{3em}{a. Fysiologi og Psykologi.}
\iz{3em}{b. Sjælelivets Ejendommelighed.}
\iz{3em}{c. Sjælelivets Elementer.}
\ix{2em}{3. Erkendelsesproblemet}{14}
\iz{3em}{a. Naiv og kritisk Realisme.}
\iz{3em}{b. Iagttagelse, Forstand. Fornuft.}
\ix{2em}{4. Tilværelsesproblemet}{19}
\iz{3em}{a. Metafysik og Empire.}
\iz{3em}{b. Ideerne.}
\iz{3em}{c. Kritiske Bemærkninger.}
\ix{2em}{5. Det etiske Problem}{24}
\iz{3em}{a. Historie og Etik.}
\iz{3em}{b. Samfund og Individ.}
\iz{3em}{c. Etik og Metafysik.}
\iz{1em}{II. \textit{Roberto Ardigò.}}
\ix{2em}{1. Italiens Filosofi efter Renaissancen}{29}
\ix{2em}{2. Ardigòs Udviklingsgang}{31}
\ix{2em}{3. Udviklingslære og Erkendelseslære}{35}
\ix{2em}{4. Psykologi}{37}
\ix{2em}{5. Etik}{40}
\iz{1em}{III. \textit{Francis Bradley.}}
\ix{2em}{1. Idealismen i den nyeste engelske Filosofi}{43}
\ix{2em}{2. Bradley og Filosofien}{44}
% --- p. VI ---
\opage{VI}%
\ix{2em}{3. Fænomen og Realitet}{48}
\ix{2em}{4. Karakteristik og Kritik}{51}
\iz{1em}{IV. \textit{Alfred Fouillée og den franske Filosofi i vore Dage.}}
\ix{2em}{1. Indledning (Taine. Renan)}{55}
\ix{2em}{2. Fouillées Udviklingsfilosofi}{64}
\ix{2em}{3. Diskontinuitetsfilosofi}{68}
\ix{2em}{4. Charles Renouvier}{69}
\ix{2em}{5. Émile Boutroux}{77}
\par\medskip\centerline{ANDEN GRUPPE.}
\ix{1em}{Erkendelsesteoretisk-biologisk Retning}{81}
\ix{1em}{Indledning}{81}
\iz{1em}{I. \textit{Filosoferende Naturforskere.}}
\ix{2em}{1. James Clerk Maxwell}{82}
\ix{2em}{2. Ernst Mach}{87}
\ix{2em}{3. Heinrich Hertz}{92}
\ix{2em}{4. Wilhelm Ostwald}{95}
\iz{1em}{II. \textit{Problemernes Naturhistorie.}}
\ix{2em}{Richard Avenarius}{98}
\par\medskip\centerline{TREDIE GRUPPE.}
\ix{1em}{Vurderingsfilosofi}{107}
\ix{1em}{Indledning}{107}
\iz{1em}{I. \textit{Jean Marie Guyau.}}
\ix{2em}{1. Kritik af den engelske Etik}{110}
\ix{2em}{2. Etik}{111}
\ix{2em}{3. Æstetik}{114}
\ix{2em}{4. Religionsfilosofi}{115}
\iz{1em}{II. \textit{Friedrich Nietzsche.}}
\ix{2em}{1. Karakteristik og Biografi}{118}
\ix{2em}{2. Nietzsches Skrifter}{I26}% PRINTED AS IS: "I26" (126)
\ix{2em}{3. Kulturhistorisk Udgangspunkt}{131}
\ix{2em}{4. Historiens Maal og den sociale Dualisme}{133}
\ix{2em}{5. Slaveopstanden i Moralen}{138}
\ix{2em}{6. Filosofisk Begrundelse af den sociale Dualisme}{141}
\ix{2em}{7. Det sidste Ja --- og Zarathustras Død}{143}
\ix{0em}{\textit{Noter}}{148}
\par\bigskip\centerline{\rule{2cm}{0.4pt}}
\end{small}
\clearpage

% --- p. VII ---
\chapter*{INDLEDNING.}
\addcontentsline{toc}{chapter}{Indledning}
\markboth{Indledning}{Indledning}
\opage{VII}I min Fremstilling af den nyere Filosofis Historie stansede jeg
ved Aaret 1880. Forskellige Grunde gjorde, at jeg satte
dette Aar som Grænse. For det Første var der ved dette Tidspunkt
naat en vis foreløbig Afslutning i Debatten og Vexelvirkningen
mellem de to store Tankeretninger i det 19. Aarhundrede,
Romantiken og Positivismen. Begge havde udfoldet deres Konsekvenser
og tildels fundet deres Korrektiver. I Lotzes og Spencers
karakteristiske Standpunkter og i den paany fremtrængende
kritiske Filosofi var der vundet foreløbige Hvilepunkter. Ved
Siden deraf var der vel fremtraadt nye Bestræbelser, men disse
havde endnu ikke formet sig til klare og karakteristiske Retninger.
For det Andet er Fremstilling og Vurdering af umiddelbart
samtidige Tankeretninger underkastet andre Vilkaar end Behandlingen
af filosofiske Fænomener, der foreligge historisk afsluttede.
En psykologisk og biografisk Belysning bliver ikke saa
let mulig; og dog er en saadan Belysning nu engang nødvendig
paa det filosofiske Omraade paa Grund af det inderlige Vexelforhold
mellem Tænkning og Personlighed. Behandlingsmaaden
maa derfor blive en anden end overfor Retninger, der mere afgjort
høre Fortiden til. Og for det Tredie begynde mine egne
mere selvstændige filosofiske Arbejder omtrent ved 1880; ogsaa
dette gør, at det er mindre let for mig at staa objektivt overfor
fremmede Arbejder fra dette Tidspunkt af.

Naar jeg nu forsøger en Karakteristik af de filosofiske Bestræbelser
i det sidste Fjerdedelsaarhundrede, som jeg anser for
% --- p. VIII ---
% running head: Indledning.
\opage{VIII}de betydningsfuldeste, er jeg klar over, at den personlige Faktor
vil gøre sig mere gældende her end i det tidligere Værk.
Baade i Valget af Repræsentanter, i Fremstilling og i Vurdering
vil denne Faktor virke. Det er Grunden til, at jeg lader
denne Bog fremtræde som et selvstændigt Værk, og ikke uden
videre som tredie Del af den nyere Filosofis Historie. ---

Skal jeg strax fra Begyndelsen af fremhæve en Ejendommelighed
ved den nyeste Tids filosofiske Tankeforsøg, maa det være
den, at de ere endnu vanskeligere at klassificere end tidligere
Tiders Filosofi. Jo grundigere man studerer Filosofi, des mere
viser det sig, hvorlidt de sædvanlige Rubriker, alle de mange
-ismer, slaa til. Baade Tanken og Tilværelsen have for stor Mangesidighed
i Væsen og Virken, til at en saadan udvortes Ordning
skulde være mulig. Særligt viser det sig atter og atter, hvorledes
en filosofisk Bevægelse baade er et Tankearbejde og et Tidens
Tegn. Filosofien er baade Problembehandling og Symptom,
og den fremtræder i vore Dage i udpræget Form paa hegge
disse Maader. Dels viser Tilværelsen sig os stadigt fra flere Sider,
i større Dybde og Mangfoldighed, dels bliver man mere og mere
opmærksom paa subjektive Faktorers Medvirken og Medklingen
ved al Tænken. I begge Henseender --- man kunde kalde dem
den objektive og den subjektive Side ved Filosofien --- fremtræder
vor Tids Filosoferen med stærkere Accent end tidligere
Tiders. Man betoner stærkere end før Nødvendigheden af Iagttagelse,
Analyse, Kritik, objektiv Sammenhæng, og man betoner
stærkere end før det subjektive Valg af Synspunkter og Udgangspunkter
og af afsluttende Tanker. Og dog (dette er maaske
det mest Karakteristiske) mener man egentligt ikke, at der her
er en uforenelig Modsætning til Stede; snarere er der en stille
Overbevisning om, at den objektive og den subjektive Tankebevægelse
tilsidst ville falde sammen.

De filosoferende Personligheder ledes snart af mere objektive,
snart af mere subjektive Grunde i deres Forsken. Wilhelm
Wundt og Ernst Mach ere fra Naturvidenskaben gaaede over til
Filosofien, og selv Naturforskere som Clerk Maxwell og Heinrich
Hertz, der ville forblive indenfor deres Fags Grænser, føle
dog Trang til at gøre sig Rede for de første Forudsætninger,
% --- p. IX ---
% running head: Indledning.
\opage{IX}de bygge paa, og kømme derved i Berøring med Filosofien.
Roberto Ardigò føres fra katolsk Teologi til kritisk og positivistisk
Filosoferen, uden skarpt Brud, men saaledes at et langt
Tankearbejde umærkeligt fører ham over til et ejendommeligt
positivistisk Standpunkt. Friedrich Nietzche begynder som Filolog
og Kulturhistoriker, men ved Kulturproblemet, saaledes som
det stillede sig allerede indenfor en snevrere historisk Ramme,
vækkes hos ham et lidenskabeligt Tankearbejde, der fører ham
gennem voldsom Kritik og Polemik til en ideal Konstruktion,
hvor Fantasi og Følelse virke sammen med Tanken. Hos Mænd
som Francis Bradley, Richard Avenarius og Jean Marie Guyau
synes den filosofiske Drift at være vaagnet mere direkte, skønt
den ogsaa hos dem fremtræder under individuelt forskellige Former. ---
Vi filosofere egentligt Alle, skønt det --- paa Grund af
de forskellige indre og ydre Forhold --- kun er Faa, hos hvem
et sammenhængende Tankearbejde kommer til Udvikling. »Blandt
hundre, som gruble, til Tænker blir én.«

En Inddeling af de vigtigste filosofiske Fænomener fra det
sidste Fjerdedelsaarhundrede vil man kunne faa ved at skelne mellem
tre Grupper eller Retninger.

\textit{Den systematiske Retning} repræsenteres af en Gruppe af
Tænkere, der især gaa ud paa at give en Belysning af Tilværelsesproblemet,
altsaa arbejde paa Udviklingen af en sammenhængende
Verdensanskuelse. Saadanne Mænd ere Wundt, Ardigò,
Bradley og Fouillée. --- Hos en anden Retning, den \textit{erkendelsesteoretisk-biologiske,}
staar Erkendelsesproblemet i første Række,
og Bestræbelsen gaar ud paa at finde de simpleste Veje, ad
hvilke Tankelivets Krav kunne fyldestgøres, idet selve Tankelivet
betragtes som en speciel Art af Liv, følgende Livets almindelige
Love. Herhen høre Mænd som Naturforskerne Maxwell,
Hertz og Ernst Mach; sin mest karakteristiske Udførelse faar
denne Tendens i Richard Avenarius' Forsøg paa at give en Problemernes
Naturhistorie. --- En tredie Retning holder sig væsentligt
til Vurderingsproblemet. Denne \textit{Vurderingsfilosofi} tager
de etiske og religiøse Grundproblemer op og søger at belyse eller
løse dem fra nye Synspunkter. Hos Forfatterne af denne Retning,
Guyau og Nietzsche, er den subjektive Faktor fremherskende. ---

% --- p. X ---
% running head: Indledning.
\opage{X}Indenfor alle tre Grupper spille psykologisk Iagttagelse og
Analyse en væsentlig Rolle, saa at den Firhed af Problemer,
som jeg i »Den nyere Filosofis Historie« har fundet ad historisk
Vej, og som jeg i Afhandlingen »Filosofiske Problemer«
har søgt nærmere at bestemme og begrunde, bekræfter sig ogsaa
her. ---

Gennem Karakteristiken af de nævnte tre Grupper af moderne
Tænkere haaber jeg at kunne orientere i vor Tids Tankeverden.

\mainmatter

% --- p. 1 ---
\grouphead{FØRSTE GRUPPE.}{OBJEKTIV-SYSTEMATISK RETNING.}

\opage{1}\chapterhead{I.}{WILHELM WUNDT.}

1. Wundt staar som et stort Exempel paa Tyskernes kolossale
Arbejdsevne. Der er ingen nulevende Filosof, der overgaar
ham i Kundskabsfylde, i Evnen til at sammenfatte sin mangfoldige
Viden under almindelige Synspunkter og i utrættelig
Bestræbelse for Nøjagtighed og Klarhed i det Enkelte saavel
som i det Store. Han har ikke Lotzes Finhed og Inderlighed;
han staar mere bredt og solidt paa Virkelighedens Grund og
vover sig ikke saa langt ud i Metafysiken som Lotze. I Sammenligning
med Spencer er han mere Idealist. Fra Erfaringens
brede Grundvold søger han op imod den Tanketinde, som Romantikens
Filosofi mente at have besteget, men som efter Wundt,
der her er besjælet af Kants kritiske Aandsretning, kun er
et højeste Punkt i Synskredsen, som vi benytte til Orientering.

Wundt er født i Nærheden af Mannheim den 16. August
1832. Han studerede Medicin, og det var gennem sine fysiologiske
Studier, at han kom til Filosofien. Hans \textit{Beiträge zur
Theorie der Sinneswahrnehmung} (1859---62) undersøge især
Øjenbevægelsernes Indflydelse paa Rumsopfattelsen. Paa Naturforskermødet
i Speier 1861 holde han Foredrag om \textit{den fysiologiske
Tid,} d. v. s. om den Tid, der medgaar til Opfattelse og
Gengivelse af et Sanseindtryk. I \textit{Vorlesungen über die Menschen-
und Thierseele} (1863, 3. Udg. 1897) anvendes den etno%
% --- p. 2 ---
% running head: Wundts Overgang til Filosofien
\opage{2}logiske Metode ved Siden af den experimentalfysiologiske. Efter
i en Række af Aar at have virket som extraordinær Professor
i Fysiologi i Heidelberg blev han 1874 Albert Langes Efterfølger
som Professor i induktiv Filosofi i Zürich, men overtog
allerede det følgende Aar et filosofisk Professorat i Leipzig,
hvor han endnu virker. Wundts Liv er hans Studiers Gang.
Hans Personlighed er ikke meget fremtrædende i hans Værker;
overfor den, der træder i Berøring med ham i Livet, virker
den med mild Varme og aaben Hjertelighed og med stille intellektuel
Begejstring.

Kort førend Wundt drog til Leipzig, udgav han sin \textit{Physiologische
Psychologie} (1874), et Hovedværk i den experimentale
Psykologi, næst efter Fechners »Elemente der Psychophysik«.
Den fremstiller især Metoderne og det fysiologiske Grundlag
for Psykologien, og kun de Afsnit af Psykologien, der særligt
kunne belyses ad denne Vej, behandles udførligt. Beskrivelse
og Analyse træde i Skyggen for Experimentet. --- En mere
ligeligt udførlig Fremstilling af Psykologien har Wundt senere
givet i sin \textit{Grundriss der Psychologie} (1896). --- I Leipzig oprettede
han et experimentalpsykologisk Laboratorium, den første
Universitetsanstalt af den Art, og begyndte Udgivelsen af Tidsskriftet
\textit{Philosophische Studien} (20 Bind, fra 1883---1903).

Om den Betydning, experimentalpsykologiske Studier havde
for Wundt ved Overgangen fra Fysiologien til Filosofien, har
han udtalt sig paa en interessant Maade i Afhandlingen \textit{Ueber
psychische Causalität und das Princip des psychophysischen
Parallelismus} (Philos. Studien 1894. p. 121---124). De overbeviste
ham om den Psykologiens Selvstændighed i Forhold til
Fysiologi og Metafysik, der allerede var hævdet af den engelske
Skole og af Kant. »Naar man vilde spørge mig«, siger han, »hvori
for mig den experimentale Iagttagelses psykologiske Værdi har
bestaaet og endnu bestaar, saa vilde jeg svare, at den hos mig
har fremkaldt og stedse mere befæstet en helt ny Anskuelse
om de sjælelige Processers Natur og Sammenhæng. Da jeg for
første Gang traadte overfor psykologiske Problemer, delte jeg
den almindelige Fordom, der ligger saa nær for Fysiologen, at
Sanseiagttagelsernes Tilblivelse blot er et Værk af Sanseorganer%
% --- p. 3 ---
% running head: ad to Veje .
\opage{3}nes fysiologiske Egenskaber. Jeg lærte først ved Synssansens
Virken [især hvad Rumsopfattelsen angaar] at forstaa den Akt
af skabende Syntese, der efterhaanden blev mig en Fører til
ogsaa at afvinde de højere Forstands- og Fantasifunktioners Udvikling
en psykologisk Forstaaelse, som den gamle Psykologi
ikke havde kunnet hjælpe mig til. Da jeg derpaa gik over til
Forestillingernes Forløb i Tiden, aabnede der sig for mig et
nyt Indblik i Villiesfunktionernes Udvikling [nemlig ved den
Betydning, Forberedelse og Anspændelse have for den fysiologiske
Tids Forkortelse], de ydre Villiesfunktioners Udvikling
af de indre, de sammensattes af de simple, --- og tillige et Indblik
i den nøje Sammenhæng mellem alle de sjælelige Funktioner,
der ellers skilles ad ved kunstige Abstraktioner og Navne
(Forestilling, Følelse og Villie), med andre Ord i det sjælelige
Livs Udelelighed og indre Ensartethed paa alle Trin«.

Wundt var ad denne Vej kommen til det Syntesebegreb,
som Kant i sin Tid havde vundet ved Analyse af Forstandsfunktionerne og af Rums-
og Tidsopfattelsen. Ved »skabende
Syntese« forstaar han en Sammenfatten og Sammensætten, hvis
Produkt faar Egenskaber, som intet af Elementerne hvert for
sig havde. Ikke sjeldent overser Wundt, at dette Begreb ikke
giver nogen Forklaring, saa godt det end egner sig til Beskrivelse
af det for sjælelige Processer Ejendommelige. Han betoner
særligt Aktivitetsmomentet i Syntesen, og ledes derved
efterhaanden til (i de senere Udgaver af sin Psykologi) at betegne
Villien som det Centrale i Sjælelivet.

Psykologiske Studier vare dog ikke den eneste Vej, ad
hvilken Wundt førtes over fra Naturvidenskaben til Filosofien.
Som Naturforsker blev han opmærksom paa, at al Videnskab
om Naturen hviler paa visse Forudsætninger. Spørgsmaalet opstaar
da uvilkaarligt, hvorledes disse Forudsætninger, der begrunde
Alt for os, selv kunne begrundes. Ad denne Vej opstaar
Erkendelsesteorien. I sit Skrift \textit{Die physikalischen Axiome
und ihre Beziehung zum Causalprincip.} Ein Capitel aus einer
Philosophie der Naturwissenschaften (Erlangen 1866) drøftede
Wundt dette Spørgsmaal. Han søgte her at udlede Fysikens
Grundsætninger af den Sætning, at Alt har sin Grund, i For%
% --- p. 4 ---
% running head: Wundts Logik og Etik.
\opage{4}bindelse med den anden Sætning, at alle Naturprocesser ere
Bevægelser. Denne anden Sætning begrunder han igen ved,
at Bevægelsen er den simpleste Forandring, idet en Ting, der
bevæger sig, kun forandrer sin Stilling i Forhold til andre Ting,
uden at behøve at miste nogen af sine Egenskaber. --- Wundt
er her endnu ikke klar over de erkendelsesteoretiske Vanskeligheder
ved denne Opfattelse, der i Sammenligning med hans
senere Erkendelsesteori har en noget dogmatisk Karakter.

Ad to Veje førtes Wundt saaledes ind paa Filosofien, og i
sine senere Skrifter har han videre uddybet og udvidet sit Begreb
om Filosofi.

Drøftelsen af Erkendelsesproblemet førte ham til at udarbejde
sit store Værk over Logik \textit{(Logik,} eine Untersuchung
der Principien der Erkenntniss und der Methoden wissenschaftlicher
Forschung. 1880-{}---83. 2. Udgave i 3 Bind 1893---95).
Hvad der giver dette Værk dets store Værdi, er især det rige
Stof fra de specielle Videnskaber, som er indarbejdet deri.
Wundt betragter snarere Videnskaben som en historisk given
Kendsgerning i det menneskelige Aandsliv end som en Bestræbelse,
hvis Mulighed og Berettigelse skal undersøges. Dette er
et for Wundts hele Filosofi karakteristisk Træk. Hans logiske
Værk faar derved snarere Karakteren af en Encyklopædi end
af en Drøftelse af Principspørgsmaal. I formel og principiel
Henseende ere \textsc{Sigwarts}, \textsc{Schuppe's} og \textsc{Benno Erdmanns}
logiske Værker maaske at foretrække; men intet af dem har
den reale Fylde som Wundts.

Wundts Fremstilling af \textit{Etiken} (1886. 3. Udg. 1902) har
en Karakter, der er analog med Karakteren af hans Logik. Han
betragter det Etiske, ligesom Videnskaben, som en stor Kendsgerning,
de Enkelte skulle anerkende og fortsætte, og han gaar
ikke ind paa de Vanskeligheder, Fordringen om denne Anerkendelse
indeholder, naar der kræves en bestemt Begrundelse af
den. Etiken bygger hos Wundt paa Folkepsykologien: den viser
hvorledes de etiske Grundsætninger udspringe af den historiske
og sociale Sammenhæng, indenfor hvilken ethvert Menneskeliv
føres. --- Om Etikens Plads i Wundts hele Filosofi kommer
jeg senere til at tale. ---

% --- p. 5 ---
% running head: Filosofiens Udspring.
\opage{5}Wundts filosofiske Hovedværk er \textit{System der Philosophie}
(1889, 2. Udgave 1897). Det er tillige det Værk, hvor hans
Fremstillingskunst er størst. Han giver her sin afsluttende Opfattelse
af Filosofien som Helhed og af dens Problemer.
--- Til dette Værk slutter sig \textit{Einleitung in die Philosophie} (1901),
som dog bringer forskellige Nuancer i Opfattelsen i Forhold til
de tidligere Skrifter. Den største Plads indtager i dette Skrift
Karakteristiken af de forskellige Retninger i Filosofiens Historie.
---

Hvad der fører til Filosofi, er ifølge Wundt en intellektuel
Trang, som mere og mere fremtræder i sin Forskellighed fra
den religiøse Trang, uden at den dog behøver at ophæve denne.
Filosofi vil sige: en Verdens- og Livsanskuelse, som tilfredsstiller
baade Forstandens og Følelsens Krav. I nyere Tid er
Filosofien ikke længere, som hos Grækerne, ene om at imødekomme
hin intellektuelle Trang. De enkelte Videnskaber have
udviklet sig og have delt Erfaringen mellem sig. De kunne
deles i tre Grupper. Matematiken er Læren om Tal, Rum og
Bevægelse. Naturvidenskaben er Læren om Naturprocesserne
og om Naturgenstandene. Aandsvidenskaben deles i Psykologi,
Filologi og Historie. Forskellen mellem disse tre Grupper bestaar
dog ikke blot i deres forskellige Genstande, men især i
de forskellige Synspunkter, de anlægge. Udviklingen af alle
disse Videnskaber har mægtigt bidraget til den intellektuelle
Interesses Væxt. Filosofien kan ikke, som den romantiske Spekulation
mente, skyde dem til Side; men dens Opgave bliver
at forbinde de af dem vundne Erkendelser paa modsigelsesfri
Maade, saaledes at intet af de specielle Synspunkter ensidigt
kommer til at raade, og saaledes at der tilvejebringes Harmoni
mellem de intellektuelle Interesser og de etiske Fornødenheder.
Historisk kan der, ifølge Wundt, paavises en Udviklingsgang af
vore Begreber fra antik Filosofi gennem Specialvidenskaberne
tilbage til Filosofien.

Filosofien inddeler Wundt i Erkendelseslære og Principlære.
Under \textit{Erkendelseslæren} hører Læren om Tænkningen
(den formale Logik), og Erkendelsens Historie og Teori (real
Erkendelseslære). \textit{Principlæren} har en almindelig Del, som
% --- p. 6 ---
% running head: Etikens og Psykologiens Plads.
\opage{6}Wundt kalder Metafysiken, og om hvilken han bemærker, at
den er Filosofiens Slutning, ikke dens Begyndelse, --- og specielle
Dele, nemlig Naturfilosofi og Aandsfilosofi.

Herefter skulde der kun existere to filosofiske Problemer,
Erkendelsesproblemet og Tilværelsesproblemet. --- Etiken faar
sin Plads i Principlæren som en af Aandsfilosofiens specielle
Grene (sideordnet med Æstetik og Religionsfilosofi), men repræsenterer
ikke en selvstændig Problembehandling. Dog er der
her en Vaklen hos Wundt. I sin Fremstilling af de filosofiske
Retninger (i »Einleitung in die Philosophie«) skelner han udtrykkeligt
mellem tre Problemer, idet han stiller det etiske
Problem ved Siden af det erkendelsesteoretiske og det metafysiske;
det erkendelsesteoretiske Problem har Mennesket som
erkendende Væsen, det etiske Problem som villende Væsen
til Genstand, og det metafysiske Problem angaar begge Sider
i hans Væsen. Og dette stemmer ogsaa med hans Opfattelse af
Filosofien som den, der skal bringe Overensstemmelse mellem
de intellektuelle Interesser og de etiske Fornødenheder; ti denne
Opgave forudsætter det etiske Problems Selvstændighed.
--- Psykologien har sin Plads mellem de specielle Videnskaber, ved
Siden af Filologien og Historien, og regnes ikke til Filosofien.
Men indenfor Wundts egen Psykologi fremtræder dog tydeligt
en Forskel mellem de specielle, experimentale Undersøgelser og
de Synspunkter for Sjælelivet, der blive af Betydning for »Aandsfilosofien«,
særligt for dennes almindelige Del, der skal give
»en sammenhængende Grundopfattelse af aandelig Væren og
Vorden«. Han indrømmer fremdeles, at Psykologien faar en
Undtagelsesstilling mellem Enkeltvidenskaberne derved, at den
træder i mere direkte Forhold til den filosofiske Erkendelseslære
end nogen anden Videnskab, idet enhver Erkendelsesakt
er en empirisk given aandelig Proces, der som saadan hører
under Psykologiens Forum. Heraf følger, at det vil være berettiget
at tale om et særligt psykologisk Problem, --- og vi
kunne saaledes anvende de fire tidligere omtalte filosofiske Problemer
til Inddeling af vor Fremstilling af Wundts Filosofi.

2. Det maa overlades Experimentalpsykologiens Historie
at fremstille Wundts Betydning for denne Disciplins Metodelære
% --- p. 7 ---
% running head: Psykologi og Fysiologi .
\opage{7}og Organisation. Wundt staar her som den betydeligste Forsker
næst efter selve Fagets Grundlægger, Fechner. Men medens
Fechner (saaledes som Wundt har fremhævet i sin Mindetale
over ham) egentligt ikke havde en udpræget psykologisk Interesse,
idet han kun brød sig om, hvad der angik Bevidsthedens
Tærskel og derved formentligt kunde kaste Lys over Forholdet
mellem det Aandelige og det Materielle, er Wundts psykologiske
Interesse mere omfattende, og det er desuden paa en
hel Række forskellige Punkter, at Psykologien for ham bliver
af Betydning for Filosofien i det Hele. Jeg skal her dvæle
ved tre saadanne Punkter, nemlig Forholdet mellem Psykologi
og Fysiologi, Sjælelivets Ejendommelighed og Sjælelivets Elementer.


a. Forholdet mellem Psykologi og Fysiologi faar nødvendigvis
filosofisk Betydning, da det bliver afgørende for Bestemmelsen
af Forholdet mellem det Aandelige og det Materielle.
Foruden i sin Psykologi og i sit System har Wundt ogsaa behandlet
dette Spørgsmaal i sin allerede ovenfor omtalte Afhandling
»Ueber psychische Causalität und das Princip des psychophysischen
Parallelismus«.

I Modsætning til Fechner anser Wundt ikke den psykofysiske
Parallelisme for en fuldstændig Løsning af Problemet.
Forskellen mellem Videnskaberne beror jo mere paa Forskelle
i Synspunkter end paa Forskelle i Genstande, og en og samme
Genstand kan derfor forekomme i meget forskellige Videnskaber.
Fysiologi og Psykologi betragte samme Genstand fra forskellige
Synspunkter. Men det, at forskellige Synspunkter ere nødvendige,
beviser ikke, at der er en Forskel mellem selvstændige
Genstande til Stede. Og naar vi Skridt for Skridt følge det
Aandeliges Udvikling fra lavere til højere Trin, nødes vi til
den Antagelse, at denne Trinrække er forberedt i den ubevidste
Natur, saa at Naturen staar som Aandens Selvudviklingsproces.
Wundt hylder altsaa tilsidst en afgjort idealistisk Tydning.
»Parallelismen« er for ham kun en Hjælpehypotese, ikke et
definitivt Standpunkt. Men som foreløbig Hypotese er den nødvendig,
fordi vi ikke kunne antage Andet, end at der maa være
Ensartethed mellem Aarsag og Virkning, saa at en Indvirken
% --- p. 8 ---
% running head: Wundt og Parallelismen.
\opage{8}af det Aandelige paa det Materielle, eller omvendt, vilde staa
som et Under. Han lægger i denne Sammenhæng ogsaa Vægt
paa Loven om den fysiske Energis Bestaaen, der kræver fysiske
Ækvivalenter for enhver Forøgelse eller Forminskelse af fysisk
Energi, hvorved den sædvanlige Vexelvirkningslære udelukkes.

Gentagne Gange har Wundt maattet forsikre, at han stadigt
fastholder denne Synsmaade, selv om han tilsyneladende udtrykker
sig i Modstrid med den. Han erklærer, at han, hvor
det ikke drejer sig om principielle Spørgsmaal, benytter den
populære Udtryksmaade med samme Ret som den, hvormed en
Kopernikaner lader Solen staa op og gaa ned. Jeg tror nu
dog, at den Vanskelighed, man her finder hos Wundt, ikke blot
stammer fra Udtryksmaaden. Der er nemlig to Punkter, hvor
han udtrykkeligt statuerer Undtagelser fra Parallelismen. For
det Første skal Forbindelsen mellem Bevidsthedselementerne
ikke kunne belyses fra det fysiologiske Synspunkt; kun de
enkelte Elementer, ikke deres Forbindelse, have fysiologiske
Korrelater. Og for det Andet skal heller ikke den Værdi, der
knytter sig til de sjælelige Fænomener have et saadant Korrelat.
--- Den første Begrænsning af Parallelismen beror paa en
uberettiget Anvendelse af Distinktionen mellem Form og Stof,
Forbindelse og Element. Begrebet Element betegner stedse kun
en Tilnærmelse til et blot Givet, og det er derfor kun i meget
grove Træk og rent relativt muligt at skelne mellem Forbindelsen
og Elementerne. Desuden er Forbindelsen ligesaa vel
et faktisk Givet, som Elementerne er det. Den anden Begrænsning
beror paa en uberettiget Sondring mellem Værdien og det,
den er knyttet til. At Noget føles som værdifuldt, er en psykisk
Proces, som enhver anden. Eller, om man vil, Værdi er ligesaa
vel et psykisk Element eller en psykisk Kvalitet, som Farve
eller Tone er det. Det synes saaledes at være inkonsekvent
at gøre disse Undtagelser. At Wundt er kommen til at gøre
dem gældende, skyldes sikkert hans Iver for at hævde Psykologiens
Selvstændighed. I tidligere Tid var --- mener han ---
denne Selvstændighed især truet ved Overgreb fra spiritualistisk
Side; men i vore Dage er det materialistiske Overgreb, som især
ere at befrygte\textsuperscript{1}).

% --- p. 9 ---
% running head: Det aktuelle Sjælebegreb .
\opage{9}En særlig Vanskelighed for ret Forstaaelse af Wundts Teori
har det voldet, at han vel antager, at den fysiske Energi i Verden
er konstant, men taler om en Væxt af Energien paa det aandelige
Omraade. Dog ses det af Hovedstederne for denne Del af
hans Lære, at han ved Væxt af aandelig Energi dels mener de
nye psykiske Kvaliteters Opstaaen, dels Tilblivelsen af nye
Værdier\textsuperscript{2}). Det forekommer mig at være rigtigst her at tale
om Koncentration og Organisation af Energien i Stedet for om
dens Væxt. Psykisk Værdi forudsætter en saadan Koncentration,
hvorimod den ikke forudsætter en Tiltagen af Energien i
det Hele. Og paa den anden Side er netop Principet om den
psykiske Energis Bestaaen nødvendigt til Forstaaelse af saadanne
sygelige Tilstande, der skyldes Ensidighed og Delthed i
Bevidstheden.

b. Gennem alle Sjælelivets Trin og Ytringsformer fremtræder
der ifølge Wundt Egenskaber, der stille det i Modsætning
til det materielle Naturliv, skønt dette er dets ydre, sanseligt
anskuelige Side. I sine forskellige Fremstillinger opregner
Wundt disse Ejendommeligheder paa noget forskellig Maade.
Men de vigtigste Træk ere dog stedse de tre, der i det Følgende
skulle omtales.

\textit{α.} Ethvert sjæleligt Indhold er en Proces, en Virken, i
stadig Vexelvirken med andre Processer og bestemt ved forudgaaende
Processer. Sjælen er ingen »Ting«, ingen »Substans«,
men Virksomhed, Aktivitet. Dette Sjælebegreb kalder Wundt
det aktuelle, i Modsætning til det substantielle Begreb om Sjæl,
som baade Spiritualismen og Materialismen, hver paa sin Maade,
hylde\textsuperscript{3}). Sjælens Væsen lægger sig for Dagen i selve den kontinuerlige
Sammenhæng af al sjælelig Skeen. Det er denne,
der gør det umuligt at opfatte Sjælen i Analogi med et materielt
Atom, som Spiritualismen stedse har Tilbøjelighed til.
Paa alle Punkter, hvor vi kunne iagttage Sjælelivet, finde vi
stedse Kontinuiteten, jo nøjagtigere vore Undersøgelser anstilles.
Det vil da være berettiget at antage, at den gælder ogsaa ud
over det Omraade, hvor psykologisk Iagttagelse er mulig, idet
det Ubevidste bruges som Hjælpebegreb.

\textit{β.} En anden Ejendommelighed ved Sjælelivet er dets Evne
% --- p. 10 ---
% running head: Skabende Syntese og forholdssættende Analyse.
\opage{10}til ved Sammenfatten af givne Elementer at frembringe et kvalitativt
nyt Indhold. Det er denne Ejendommelighed, der kan betegnes
ved Begrebet skabende Syntese. Den ytrer sig i enhver
Sanseiagttagelse, tydeligst ved Rumsopfattelsen, idet denne opstaar
ved Sammenarbejden af Syns-, Berørings- og Bevægelsesfornemmelser,
og ved Tonernes Klangfarve, der opstaar ved
Sammenarbejden af Over- og Undertoner. Først er den bleven
opdaget ved de højere sjælelige Processer, ved Fantasibilleders,
Begrebers og Ideers Tilblivelse. Men det er en Ejendommelighed,
der fremtræder i de mest elementære psykiske Processer
saavel som i de højeste. Ogsaa fra denne Side viser Kontinuiteten
i Sjælelivet sig.

I Udtrykket skabende Syntese lægger Wundt ganske særligt
Vægt paa Ordet skabende. Han fremhæver (som allerede
berørt) ikke det Problem og den Grænse for Forstaaelse, som
heri ligger. Ti det kvalitative Nye, der kommer til --- f. Ex.
naar Rumsformen, ifølge den genetiske Teori, opstaar ved Syntese
af forskelligartede Fornemmelser, eller naar en genial Ide
springer frem paa Grundlag af Motiver, den bringer i helt ny
Sammenhæng --- finder dog ikke sin Forklaring ved at kaldes
Værk af en skabende Syntese. Dette Udtryk kan med Rette
bruges til at beskrive, og det egner sig godt til at betone det
Nye i Produktet i Sammenligning med Faktorerne; men det
lader et stort Problem staa tilbage og maa derfor ikke bruges
som Afslutning en Gang for alle.

\textit{γ.} Til den nys nævnte Ejendommelighed slutter sig meget
nøje et tredie Særkende for Sjælelivet. Intetsteds ytrer den
sjælelige Aktivitet sig ifølge Wundt tydeligere end i den Maade,
paa hvilken det, der staar som Genstand for Bevidsthed, leddeles,
sondres i enkelte Bestanddele ved en Proces, der kan
betegnes som forholdsbestemmende Analyse (beziehende Analyse).
Denne Leddeling sker gennemgaaende ikke saaledes, at
Delene derefter optræde som for sig bestaaende Enheder, men
saaledes, at de vedblive at staa i Forhold til Helheden og have
deres Betydning ved deres bestemte Plads indenfor denne. Som
oplysende Exempel bruger Wundt den Maade, hvorpaa vi indenfor
hele Synsfeltet (Blickfeld) formaa at lade et enkelt Punkt
% --- p. 11 ---
% running head: Begrebet Apperception.
\opage{11}(Blickpunkt) træde særligt frem, ved at lade det Indtryk, der
svarer til det, falde paa Stedet for den tydeligste Seen paa Nethinden.
Den sondrende Opmærksomhed, der ytrer sig i den
forholdssættende Analyse, kalder Wundt Apperception. Ved dens
Hjælp træde Delene af et Helhedsbillede successivt frem for
os, uden at miste Sammenhængen med Helheden. Saaledes staar
for den Talende Tanken som et Hele, før de enkelte Dele artikuleres;
og for Kunstneren staar Værket som et Hele, før de
enkelte Dele skelnes og udføres. Det er en fremskridende Klarhed
og Tydelighed, som derved naas.

c. Wundt blev først opmærksom paa den aktive Side i Bevidsthedslivet
ved Undersøgelsen af Sansning og Forestilling.
Denne Undersøgelse førte ham til Begrebet Apperception, som
jeg allerede foreløbigt har nævnt --- et af de vanskeligste Begreber
i Wundts Filosofi.

Apperceptionen er ikke blot en Aktivitet, der ytrer sig paa
Forestillingernes Omraade; den ytrer sig ogsaa i al Følelse, idet
Lyst og Ulyst bestemmes ved Forestillingsindholdets Forhold til
den psykiske Virken, --- og den ytrer sig i, hvad man i snevrere
Forstand kalder Villen. Naar Wundt som oplysende Exempel
benytter Forskellen mellem Synskreds (Blickfeld) og Fixeringspunkt
(Blickpunkt), saa maa det dog mærkes, at Apperceptionen
efter hans Teori ikke blot er skelnende, som naar et enkelt
Punkt fixeres indenfor en hel Sammenhæng, men ogsaa forbindende,
saaledes ved al Begrebs- og Domsdannelse. Den staar
i Modsætning til de overvejende passive og uvilkaarlige Processer
--- saaledes bl. A. til den blotte Forestillingsassociation.
Wundt udtaler sig imod den psykologiske Teori, der vil føre
Alt i Bevidstheden tilbage til udvortes Association mellem selvstændige
Forestillinger, den saakaldte association-psychology.
Han hævder, at der i al Association virker en »Apperception«,
saa at Sondringen mellem Association og Apperception egentligt
beror paa »en Abstraktion, til hvilken Virkeligheden mere eller
mindre kan nærme sig«. Men det er ikke let at bringe alle
Wundts Udtalelser i Overensstemmelse hermed. Ti Wundt omtaler
simple Associationsprocesser, ved hvilke Villien aldeles
ingen Indflydelse skal have; de karakteriseres »durch die völ%
% --- p. 12 ---
% running head: Kritik af Apperceptionsbegrebet.
\opage{12}lige Einflusslosigkeit des Willens auf die Art ihres Eintritts«.
Og naar han antager, at der i Hjernen findes et eget »Apperceptionscentrum«,
saa synes ogsaa deri at ligge, at Apperceptionen
maa være en selvstændig Proces i Forhold til Associationen
og andre mere passive Processer\textsuperscript{4}). Men Forholdet bliver
endnu mere indviklet ved, at der igen indenfor Apperceptionen
skelnes mellem passiv og aktiv Apperception. Aktiv Apperception
er karakteriseret ved, at den er bestemt ved en forudgaaende
Forestilling og kender mere end én Mulighed, medens
den passive Apperception bestemmes ved et Sanseindtryk og
gaar i én bestemt Retning.

Jeg har ikke knnnet overbevise mig om, at der var nogen
Nødvendighed for at indføre dette Begreb i Psykologien. Begrebet
Opmærksomhed indeholder i Virkeligheden Alt, hvad Erfaringen
frembyder. Som uvilkaarlig Opmærksomhed (svarende
nogenlunde til Wundts passive Apperception) medvirker den
ved enhver Sansning (f. Ex. naar et enkelt Punkt i Synskredsen
fixeres) og ved enhver Forestillingsassociation (idet det stedse
for en væsentlig Del beror paa den herskende Følelse og Interesse,
hvilken Forestilling der fremkaldes). Som vilkaarlig Opmærksomhed
(svarende nogenlunde til Wundts aktive Apperception)
fremtræder den, naar der forud for Opmærksomhedsakten
er en Forventning om eller en Stræben efter at opfatte eller
forestille noget vist Bestemt. Det er da Psykologiens Opgave
at efterspore alle specielle Grader og Former for Opmærksomheden,
fra de mest elementære til de allerhøjeste.

Det Væsentlige i Wundts Apperceptionsteori bestaar i, at
den fremhæver Betydningen af de indre, centrale Betingelser
for psykisk Virksomhed i Modsætning til de ydre og periferiske.
Særligt betoner han, at hele Bevidsthedslivets forudgaaende Udviklingshistorie
bestemmer den Maade, hvorpaa vi selvvirksomt
bearbejde og opfatte Sanseindtryk og Forestillinger. Som jeg
helst vilde udtrykke det: Opmærksomhedens Grad og Retning
er bestemt ved Livets Trang og Stræben paa det givne Udviklingstrin,
og denne Trang og Stræben er igen bestemt ved Livets
Forhistorie.

Vanskelighederne ved Apperceptionsteorien hos Wundt hæn%
% --- p. 13 ---
% running head: Wundt og Voluntarismen.
\opage{13}ger sandsynligvis sammen med den Omstændighed, at han er
bleven foranlediget til at danne Begrebet Apperception dels ad
experimental Vej, navnlig ved Forsøg, som angik forberedt og
bevidst Opmærksomhed, dels efter historiske Forbilleder (især
Leibniz og Kant), hos hvem ogsaa udpræget bevidste Funktioner
laa til Grund for Apperceptionslæren. Derved er Forholdet
mellem Sjælelivets klart bevidste og halvt bevidste, vilkaarlige
og uvilkaarlige Funktioner kommet til at staa hen i en vis
Uklarhed, ogsaa efterat Wundt senere hen havde udvidet Apperceptionsbegrebet
til at gælde for hele Sjælelivet. Begrebet kom
ikke til at undergaa de Ændringer, det større Omfang nødvendigvis
fordrede. Det er Wundts Stræben, som tydeligt fremtræder
gennem Værkernes og Udgavernes Række, at give Begrebet
Aktivitet, som Udtryk for en indefra bestemt Virken og
Stræben, en fremragende Plads i Psykologien. Villiesbegrebet
kommer mere og mere til at staa for ham som det typiske,
centrale Begreb, i Analogi med hvilket alt Andet i Sjælelivet
skal tænkes. I anden Udgave af sin Logik betegner han sin
psykologiske Opfattelse som Voluntarisme, et Udtryk, han efter
sin egen Udtalelse har laant fra \textsc{Friedrich Paulsen}, hos
hvem det stilles i Modsætning til Intellektualismen i den ældre
Psykologi\textsuperscript{5}). Dog er Wundts Psykologi end ikke i sine seneste
Former (i \textit{Grundriss,} i tredie Udgave af \textit{Menschen- und Thierseele}
og i femte Udgave af \textit{Physiologische Psychologie)} fremstillet
som en Villiespsykologi, der følger Villien i dens Udvikling
fra Trin til Trin og ser de andre Bevidsthedselementer
i deres Forhold til den. Saadanne Forsøg, Forsøg af ikke ringe
Interesse, ere gjorte af \textsc{Friedrich Jodl}, af \textsc{Alfred Fouillée}
og af G. F. \textsc{Stout}. Ja, Wundt henregner end ikke Villien til
Bevidsthedslivets Elementer. Han behandler Villiesfænomenerne
som de mest sammensatte og specielle Former for Bevidsthedsliv,
og regner kun Fornemmelser og Følelser for psykiske »Elementer«.
Dette forstaas kun ved, at Wundt har begyndt sine
psykologiske Studier med Sansningen og de til den sig knyttende
intellektuelle Funktioner og arbejdet sig frem til at se
den psykiske Aktivitets principielle Betydning, uden at det er
lykkedes ham at gøre denne til det Alt betingende Hovedelement.
% --- p. 14 ---
% running head: Erkendelsesproblemet.
\opage{14}Forholdet mellem de elementære og de sammensatte psykiske
Fænomener kommer hos ham stedse til at staa uklart hen\textsuperscript{6}).

De Vanskeligheder, der saaledes frembyde sig ved Studiet
af Wundts psykologiske Værker, vidne efter denne Forklaring
netop om hans uafbrudte og utrættelige Arbejde gennem 40
Aar, ud fra det Punkt, paa hvilket han fra Naturforskningen
var kommen over paa Filosofiens Omraade. Paa denne Vej har
han ved sin Alsidighed og Kundskabsfylde kastet Lys over en
Mængde Fænomener, som en mere retliniet og skarpt begrænset
Fremstilling maaske ikke havde formaaet.

3. Hvad Wundts Behandling af Erkendelsesproblemet angaar,
holder jeg mig i denne Fremstilling hovedsageligt til hans
\textit{System der Philosophie,} i hvilket det forekommer mig, at de
afgørende Synspunkter fremtræde med størst Klarhed og Nøjagtighed.


\textit{a.} Allerede i den sædvanlige Erkendelse, der endnu ikke
er videnskabeligt udformet, fremtræde de vigtigste Tankeprocesser
tydeligt. Erkendelsen er dels beskrivende, dels fortællende:
den gaar dels ud paa at angive Egenskaber, dels paa at fastslaa
Begivenheder. Til Beskrivelsen svarer i den strenge Tænken
Paavisningen af Identitet, til Fortællingen Godtgørelsen af
et logisk Afhængighedsforhold. Hist fremtræder det Konstante
eller det Statiske, her det Skiftende eller det Dynamiske. Afhængighedsforholdet
fremtræder med fuld Klarhed i Slutningen,
i den Maade, paa hvilken en Sætning udledes af andre Sætninger.
Derfor søger Tanken at føre alle Forhold i Tid og
Rum tilbage til logiske Afhængighedsforhold, saa at de blive af
samme Art, som det Forhold, der i en Slutning er mellem Præmisser
og Konklusion. Man bliver da ikke staaende ved den
blotte Konstateren af Egenskaber, men varierer dem for at finde
deres indbyrdes Afhængighedsforhold. Denne Bestræbelse er
naturlig, da den logiske Afhængighed er den eneste Art Afhængighed,
som umiddelbart er Et med Tankens egen Lovmæssighed;
Afhængigheden bliver her ikke en ydre; Tankens Frihed
og Nødvendigheden ere her umiddelbart Et. Alle reale Afhængighedsforhold
søges derfor reducerede til logiske Afhængighedsforhold.
Denne Bestræbelse finder sit Udtryk i det Prin%
% --- p. 15 ---
% running head: Naiv og kritisk Realisme.
\opage{15}cip, at Alt har sin Grund (Sætningen om Grunden). Dette Princip
forbinder vore forskellige Tankeakter indbyrdes paa den inderligste
Maade.

Erkendelsen begynder stedse med Forvisningen om Forestillingernes
reale Gyldighed --- som \textit{naiv Realisme.} Der skelnes
endnu ikke mellem selve Erkendelsen og dens Genstand; en
saadan Distinktion skyldes en senere Reflexion. Filosoferne
overse ofte, at Reflexionen stedse betegner et senere Stadium,
ikke det oprindelige Stadium. Der maa tilstrækkelige Grunde
til at opgive det umiddelbart Givnes Realitet. Saadanne Grunde
udspringe især af Modsætningen mellem det Skiftende og det
Konstante. Hvad der har Gyldighed, kan kun være det, der
bestaar trods al Skiften, --- og det findes ved Tankearbejde,
ikke ved umiddelbar Opfattelse. Naar Tankearbejdet danner
Objektsbegreber, der kunne afløse den umiddelbare Opfattelses
skiftende Indhold, gaar den naive Realisme successivt over til
\textit{kritisk Realisme.} Denne er kun Fortsættelse af den Proces,
der allerede begynder indenfor den naive Realisme, naar denne
skelner mellem Iagttagelse, Erindring og Fantasi. --- Om Forholdet
mellem naiv og kritisk Realisme har Wundt skrevet en
fortræffelig Afhandling, der findes i 12. og 13. Bind af »Philosophische
Studien«.

Der er paa Vejen fra naiv til kritisk Realisme tre Stadier
at skelne imellem: Iagttagelsens, Forstandserkendelsens og
Fornufterkendelsens. Disse tre Stadier svare henholdsvis til det
praktiske Livs Standpunkt, til Specialvidenskabernes og til Filosofiens
Standpunkt.

\textit{b.} Indenfor \textit{Iagttagelscn} træde Tid og Rum frem som to
Former, i hvilke Ordningen af det Givne sker, --- som det Konstante
i Modsætning til de skiftende Kvaliteter,
--- og som Udtryk for et stadigt Tankearbejde. Rummets og Tidens Love
kunne fixeres begrebsmæssigt, og Tanken har derfor en naturlig
Tendens til at holde sig til dem og betragte Kvaliteterne
som noget Subjektivt i Sammenligning med dem.

\textit{Forstandserkendelsen} begynder, naar den strenge Fordring
om konstant Bestaaen fører bort fra den umiddelbare Anskuen
og erstatter denne med Ord og Symboler. Kun det, der kan
% --- p. 16 ---
% running head: Materiebegrebet.
\opage{16}bestemmes begrebsmæssigt, kaldes her Objekt. Vi komme her
nødvendigvis ind paa Hypotesernes Omraade, men det er logiske
Grunde, der nøde dertil, idet Tanken stræber efter at finde
en modsigelsesfri Sammenhæng mellem alle Dele af Erfaringsindholdet.
Til Tankelovene gaar Alt her tilbage. Men ren
Tænken have vi paa intet Punkt saalidt som ren Erfaring. Derved
afvises Apriorismen, saavel som Empirismen. --- Baade i
Matematik, Naturvidenskab og Psykologi er det Forstandserkendelsen,
der virker. --- Et godt Exempel er det naturvidenskabelige
Begreb om Materien eller om den materielle Substans.
Dette Begreb dannes ved, at alle materielle Egenskaber føres
tilbage til Forhold i Rummet, --- til Bevægelser og Stillinger.
Og denne Reduktion motiveres netop ved, at det kun derved
bliver muligt paa en simpel Maade og i Kraft af Tankens egne
Love at udlede den ene materielle Forandring af den anden.
De empiriske Kvaliteter benyttes i Naturvidenskaben kun som
Hjælpemidler for at slutte tilbage til Objekternes Forhold i Rum
og Tid. Det er ikke Opgaven for en Teori om Materien at udvikle
Forestillinger, der ligne de empiriske Legemers Fænomener,
ja, ikke engang overhovedet at forme anskuelige Forestillinger,
men at fastslaa Begreber, ved hvis Hjælp empiriske
Fænomener kunne udledes af tidligere givne Fænomener. Hvad
der gælder for Materiebegrebet, gælder ogsaa for Begreber som
Inerti og Energi. Det er stadigt Fordringen om modsigelsesfri
Forbindelse i Overensstemmelse med Sætningen om Grunden,
der ligger til Grund ved disse Begrebsdannelser. --- Paa det
psykologiske Omraade er der ifølge Wundt ingen Grund til at
foretage en Begrebskonstruktion i Analogi med dem, der fører
til Materiebegrebet i Naturvidenskaben. Ti her godtgør Iagttagelsen
klart, at det mest Konstante, der gives i os, er vor
Villen, og at den er en aldrig hvilende Virksomhed, en stadig
Skeen; der bliver da hverken Mulighed for eller Grund til at
danne et Begreb, der svarer til, hvad Stillingen i Rummet er
ved de materielle Fænomeners Virkemaade. (Smlgn. ovenfor om
Wundts »aktuelle« Sjælebegreb). De Konstruktioner og Hypoteser,
som vi paa det psykologiske Omraade have Brug for,
angaa dels Forholdet mellem de psykiske Elementer, der sam%
% --- p. 17 ---
% running head: Forstand og Fornuft.
\opage{17}virke i de indre Processer, dels Forholdet mellem det Psykiske
og det Materielle. --- Skønt Naturvidenskaben fører alle Forandringer
tilbage til Forandringer i Stilling og Bevægelse, negter
den ikke, at Tingene i Verden have indre Egenskaber, som
ikke udtrykkes ved deres ydre Relationer; det er kun ikke dens
Sag at drøfte dem. Men uden saadanne indre Egenskaber vilde
det være umuligt at forstaa, hvorledes Livet, især Sjælelivet har
kunnet opstaa i Verden. Og overalt, hvor en Række af Begivenheder
fremtræder som en Helhed, indenfor hvilken hvert Led
staar paa sit bestemte Sted som medvirkende til et afgørende
Resultat, er det berettiget at anvende Formaalsbegrebet, selv
om man ikke tænker sig, at Formaalet existerer i en Forestilling
forud for Resultatet. Det er den Maade, hvorpaa det enkelte
Led viser sig at være bestemt ved sit Forhold til og Betydning
for Helheden, der gør den teleologiske Opfattelse berettiget, uden
at den paa noget Punkt udelukker den rent mekaniske. Kun ved
en saadan Opfattelse bliver det ifølge Wundt berettiget at tillægge
Begrebet Udvikling Betydning ud over det organiske
Omraade.

\textit{Fornufterkendelsen} fører ud over Erfaringen, idet den søger
en ubegrænset Sammenhæng, medens i Erfaringen al Sammenhæng
er begrænset. Ordet Fornuft tager Wundt i snevrere Betydning,
om den Tendens til Enhed, der fører til at danne
Totaliteter paa Grundlag af det brudstykkeagtigt Givne. Fornuften
bliver til, naar Bevidstheden vaagner om, at Tankevirksomheden
ifølge sin egen Lov stedse gaar videre, idet enhver Begrænsning
i Begrundelsen maa kunne overskrides eller dog tænkes
overskreden. Der opstaar da Ideen om en Totalitet af alle Afhængighedsforhold.
Fornuftens Ideer opstaa ved en Fortsættelse
af den Proces, der førte Forstanden til at danne sine Begreber.
Medens Forstanden blot søgte at forbinde det Givne paa modsigelsesfri
Maade, søger Fornuften --- i Kraft af de samme
Tankelove, der ledede Forstanden --- at bringe et sammenhængende
System tilveje.

Denne Opgave stiller sig altsaa for Wundt med en Tankenødvendigheds
Magt. Forgæves ville Empirisme og Skepticisme
negte Tilstedeværelsen af de Problemer, her frembyde sig. Wundt
% --- p. 18 ---
% running head: Real og imaginær Transscendens.
\opage{18}indrømmer, at man ad logisk Bevisførelses Vej ikke kan gendrive
de empiriske og skeptiske Standpunkter; men han mener,
at de gendrives ved selve Tænkningens Existens: ti det er selvmodsigende
at anvende Tanken til at forbinde Enkeltheder indbyrdes
og dog at afvise den Opgave at bringe de saaledes
vundne Forbindelser i indbyrdes Sammenhæng. Sætningen om
Grunden fører stedse videre. Empirismens og Skepticismens
videnskabelige Betydning beror paa, at de holde Spekulationen
indenfor visse Skranker og sørge for, at den ikke volder Forvirring
ved at sammenblandes med Forstandserkendelsen.

Fornufterkendelsen kan gaa ud over Erfaringen eller blive
»transscendent« paa to Maader. Den kan fortsætte de Rækkedannelser,
der ere begyndte indenfor Erfaringen, i samme Retning.
Saaledes i Talrækken, i Rummets Dimensioner, i Tidsrækken,
i Spinozas to Attributer (Aand og Materie). Her tænke
vi det som stadigt fortsat Kontinuitet, hvad der i Erfaringen
stedse kun er givet som Brudstykke. Denne Overskriden
af Erfaringen kalder Wundt \textit{den reale Transscendens.} Men Fornufterkendelsen
kan ogsaa supplere Erfaringen ved at antage
andre Rækker end dem, der fremtræde i Erfaringen. Medens
der ved den reale Transscendens kun er Tale om kvantitativ
Uendelighed, indføres der her en kvalitativ Uendelighed, idet
man antager flere og andre Sider eller Egenskaber ved Tilværelsen
end de givne. Et Exempel fra Filosofiens Historie
frembyder Spinozas Antagelse af, at Tilværelsen har uendeligt
mange »Attributer«, ikke blot de to, Erfaringen viser os (Aand
og Materie). Denne Art Overskriden af Erfaringen kalder Wundt
\textit{den imaginære Transscendens,} fordi den antager nye Kvalitetsrækker
i Tilværelsen, ligesom de imaginære Tal kunne tydes
som Udtryk for nye Retninger i Forhold til den Retning, der
betegnes ved de positive og negative reelle Tals Række,
--- »sidelængs Enheder«, som en berømt Matematiker har kaldt
dem. Ogsaa en Hypotesedannelse af denne Art er berettiget,
naar den tilfredsstiller Fornuftens Enhedstrang, og naar den
ikke krænker Erfaringsvidenskaben.

Kun ad disse to Veje kan der efter Wundt naas en almindelig
Verdensanskuelse. Al Metafysik, der skal have Betyd%
% --- p. 19 ---
% running head: Trangen til Metafysik.
\opage{19}ning, opstaar enten ved Fortsættelse eller ved Supplering af Erfaringserkendelsen.
Ad disse to Veje føres vi tillige over fra
Erkendelsesproblemet til Tilværelsesproblemet.

4. \textit{a.} Metafysiken begynder egentligt allerede indenfor de
specielle Videnskaber, nemlig forsaavidt disse opstille definitive
Principer eller Hypoteser, hvis Indhold ikke kan paavises empirisk,
skønt de ere uundværlige for Forbindelsen af Erfaringsindholdet.
Maaske kunne de ikke engang konstrueres fuldstændigt
i Analogi med Genstande, der ere givne i Erfaringen. De
gøre Erfaringen begribelig, men ligge selv hinsides Erfaringen.
Dette gælder f. Ex. om de naturvidenskabelige Hypoteser om
Materiens Beskaffenhed og Grundegenskaber. Wundt lægger
Vægt paa at beholde Udtrykket Metafysik om saadanne Antagelser,
forat de ikke skulle forvexles med Kendsgerninger.
Enhver definitiv Hypotese er metafysisk, og enhver Metafysik
er hypotetisk.

Motivet til Metafysik er Trangen til Sammenhæng og dermed til Forstaaelse af Erfaringsindholdet,
--- tilsidst Trangen
 til
en sammenhængende Verdensanskuelse. Hvert enkelt Erfaringsomraade,
og som Følge deraf hver enkelt Videnskab har en
vis Tendens til at betragte sig selv som absolut, som en Nøgle
til hele Tilværelsen. Derfor behøves der her en stadig Kritik\textsuperscript{7}).

Wundts Forhold til Positivisme, Kriticisme og Romantik
fremtræder her paa karakteristisk Maade. Han nærmer sig
Positivismen ved den Vægt, han lægger paa Erfaringen som
Maalestok. Med stor Forsigtighed gaar han Skridt for Skridt
henimod Erkendelsens Grænsespørgsmaal, og den Maade, paa hvilken
han baade anerkender de Problemer, der her fremstille sig,
og Vanskeligheden ved deres Løsning, viser ham som Tilhænger
eller Fortsætter af den kritiske Filosofi. Men Forsigtigheden
raader kun hos ham, medens han endnu er ved at nærme
sig til Grænsen. Naar han staar, eller mener at staa ved selve
Grænsen er han ikke saa varsom, og han gennemtænker ikke
altid de sidste Konsekvenser af sine afsluttende Tanker. Han
bliver her let tung og dogmatisk. Enhedsprincipet betager ham,
og den kritiske Reflexion træder tilbage. Man mærker ikke
noget til den Brænding, der naturligt opstaar paa Tankens Kyster.
% --- p. 20 ---
% running head: Kosmologiske Ideer.
\opage{20}Paa dette Punkt er der Mere at lære af den engelske Filosof
Francis Bradley, som vi senere ville beskæftige os med. Han
supplerer forsaavidt Wundt paa dette Punkt, mendens han mangler
Wundts brede empiriske og experimentale Basis. --- Hidtil
have vi dog væsentligt lært den empiriske og den kritiske Side
hos Wundt at kende; ved Fremstillingen af hans metafysiske
Standpunkt vil det Romantiske i hans Synsmaade komme for
Dagen. ---

b. Der gives, som den metafysiske Tænknings Historie
viser, tre Grupper af Ideer, d. v. s. af Begreber, der ved real
eller imaginær Transscendens bringe en Afslutning paa Verdensanskuelsen.
Wundt kalder dem de kosmologiske, de psykologiske
og de ontologiske Ideer.

\textit{α.} Ideen om den fysiske Verden som Helhed opstaar saa
meget des naturligere indenfor Naturvidenskaben, som det egentligt
kun er indenfor afsluttede Helheder, at man kan paavise
bestemte Love. Enhver Undersøgelse forudsætter en vis Isolation
af det, der skal undersøges; man maa saavidt muligt tænke
sig det uden Paavirkning udefra. Skridt for Skridt kan man
saa maaske udvide Betragtningen. Det gælder ikke blot, naar
Talen er om den fysiske Verden som Helhed, men ogsaa, naar
det drejer sig om den fysiske Verdens mindste Dele. Tid og
Rum kunne vi ved blot Fortsættelse tænke udvidede ud over
Erfaringen; her er der kun real Transscendens til Stede. Men
enhver Hypotese om den fysiske Verden som Helhed eller om
dens mindste Dele bliver imaginært transscendent, da der ved
enhver Grænse, der er naat, er Mulighed for, at noget kvalitativt
Nyt kunde gælde paa den anden Side. Ogsaa saadanne
Hypoteser have dog deres Berettigelse, naar de ikke komme i
Strid med Erfaringserkendelsen\textsuperscript{8}).

Hvis man bygger sin Metafysik paa Grundlag af de kosmologiske
Ideer alene. --- altsaa alene paa Begreber som Rum,
Tid, Materie, Mekanik, --- faar denne Metafysik Materialismens
Præg.

\textit{β.} Wundt bemærker, at næsten al Strid paa det metafysiske
Omraade drejer sig om de psykologiske Ideer. Vi kende
allerede Wundts Stilling til det psykologiske Problem, forsaavidt
% --- p. 21 ---
% running head: Psykologiske og ontologiske Ideer.
\opage{21}dette stilles paa Erfaringens Grund. Aktivitet og Enhedsstræben
ere for ham Sjælelivets Karaktermærker, og Begrebet Sjælesubstans
indeholder for ham, selv hvor det opstilles med spiritualistisk
Motivering, en skjult Materialisme. Det er ham tillige
--- kunne vi her føje til --- for individualistisk, da det
isolerer de enkelte sjælelige Væsener fra hverandre. Forskellen
mellem Sjæl og Legeme existerer kun i vor Opfattelse. Hvad
Erfaringen viser os, er en aandelig Organisation, der i sig selv
er Et med en legemlig Organisation. Men Erfaringen fører ikke
til noget afsluttende Begreb her, og derfor bliver ethvert Sjælebegreb,
det aktuelle, saavel som det substantielle, tilsidst imaginært
transscendent. Men naar man vil afslutte sine Ideer om
dette Spørgsmaal, maa man ifølge Wundt ikke glemme, at det
enkelte Sjæleliv i vor Erfaring stedse fremtræder som Led af
et Samfund, af en aandelig Helhed, indenfor hvilken dets Motiver,
dets Tendenser og dets Indhold ene blive forstaaelige.
Wundt lægger derfor stor Vægt paa Folkepsykologien. Som den
højeste psykologiske Ide fremtræder da tilsidst Tanken om en
omfattende Enhedsgrund (ein Uebergeistiges), der ligger til Grund
for alle sjælelige Væsener og for deres indbyrdes Sammenhæng.
Imaginært transscendent bliver denne Ide, da for os
alle psykologiske Bestemmelser kun kunne anvendes paa individuelle
Væsener, hvorfor vi i vor Erfaring mangle Udtryk
for det, der skal udtrykke al Forestillens og Villens Totalitet.

Dersom man bygger sin Metafysik paa de psykologiske
Ideer alene, faar den Idealismens Præg.

\textit{γ.} Ville vi ikke følge Materialismens og Idealismens mere
eller mindre ensidige Veje, men stræbe efter et mere omfattende
Begreb om Tilværelsen, saa maa vi ifølge Wundt forbinde
de kosmologiske og de psykologiske Ideer. Derved komme
vi over til, hvad han kalder de ontologiske Ideer.

Denne Forbindelse kan ske paa den Maade, at den kosmologiske
Analyse suppleres ved den psykologiske. Den kosmologiske
Undersøgelse viser os Tilværelsen som en Sammenhæng
af Elementer, om hvis sidste Natur den intet Bestemt kan sige;
den psykologiske Undersøgelse har idetmindste vist os Villien
som vort egentlige Væsen. Kombinere vi nu begge, faa vi
% --- p. 22 ---
% running head: Den afsluttende Ide.
\opage{22}Ideen om Tilværelsen som en Helhed af stræbende og villende
Væsener.

Man kan, mener Wundt, ikke som Spinoza blive staaende
ved Ideen om et Væsen, hvis Beskaffenhed ikke nærmere bestemmes.
Der maa stedse spørges, om Enhedsprincipet i Tilværelsen
ikke stemmer med noget af vore givne Begreber,
--- om det ligger det Materielle eller det Aandelige nærmest. Enten
maa man --- hævder han --- tænke Verden som en materiel
Enhed eller som en aandelig Enhed: et Tredie gives ikke! ---
Og her er Valget ikke tvivlsomt for ham. Den eneste Virksomhed,
som er os umiddelbart given, er og bliver vor Villen.
En uendelig Totalvillie (Gesammtwille) bliver da den sidste, afsluttende
Ide. Denne Ide, Gudsideen, er imaginært transscendent,
og dens Indhold er ubestemmeligt. Ved dens Hjælp bliver
det dog muligt at opfatte den kosmiske Mekanisme som det
ydre Hylster for aandelig Virken og Stræben, og vor egen
psykofysiske Existens staar for os som en Verden i det Smaa,
en Mikrokosme. Herved er paa én Gang Fornuftens Enhedsfordring
tilfredsstillet, og vi kunne se vore menneskelige Idealer
som Følger, der udspringe af selve Verdensgrunden. Den tomme
og trøstesløse Verdensanskuelse, der, støttet til den blotte Forstandserkendelse,
ser Tingenes Væsen udtømt i deres ydre
Forhold og Ordninger, er forladt.

\textit{c.} Den definitive Stilling, Wundt indtager til Grænsespørgsmaalene,
foranlediger naturligt nogle kritiske Bemærkninger, dels
med Hensyn til Metoden, dels med Hensyn til Motiverne for
det Standpunkt, han ender med at indtage.

Hvad Metoden angaar, mener han selv udelukkende at foretage
en saadan Fortsættelse og Supplering af Erfaringserkendelsen,
som Tankens Enhedstrang kræver. Han afviser derimod
udtrykkeligt den Mulighed, at man her skulde kunne støtte sig
til Analogi: ingen Analogi vilde slaa til! --- Men naar han
skal begrunde den idealistiske Vending, han tilsidst giver sin
Filosofi, støtter han sig jo dog afgjort til Analogien mellem
Mikrokosme og Makrokosme. Han spørger, hvilket af vore givne
Begreber, Enhedsprincipet i Tilværelsen nærmest \textit{stemmer med}
(übereinsstimmt). Der kan nu her, mener han fremdeles, kun
% --- p. 23 ---
% running head: Kritik af Wundts afsluttende Ide .
\opage{23}være to Muligheder: »Enten maa Verden tænkes af os som en
materiel Enhed eller som en aandelig Enhed, naar den overhovedet
er at tænke som en Enhed; --- et Tredie gives ikke«.
(System.\textsuperscript{1} p. 411). Det er rigtigt, at vi erfaringsmæssigt kun
kende aandelige og materielle Fænomener; men derfor have vi
(som Spinoza allerede har set) ingen Ret til at antage, at Tilværelsens
Væsen er udtømt ved dem, og at et Tredie ikke gives.
Wundt synes paa dette Punkt at glemme sin egen Teori om
imaginær Transscendens, som viser hen til, at Tilværelsen kunde
være mere mangfoldig i sit Væsen, end Metafysikerne mene.
I hvert Tilfælde er det afgjort, at han støtter sig til en Analogi,
og at han endog vælger mellem de to ifølge hans Mening eneste
mulige Analogier. Hvad der bestemmer hans Valg, er Hensynet
til, hvilken af de to Arter af Fænomener, der er os mest umiddelbart
kendt. (Smlgn. System.\textsuperscript{1} p. 434: --- dem gleichend, das
wir in uns selber erleben). Hans Tankegang minder her afgjort
om Leibniz' og Lotzes, kun at disse med fuld Bevidsthed støtte
sig til en Analogi. Der savnes da en nøjere Undersøgelse af Analogiens
Berettigelse i Tænkningen, og af de forskellige Anvendelser
af Analogi, som Videnskaben kender\textsuperscript{9}).

Hvad Motiverne angaar, da er det klart, at Wundts afsluttende
Tanke ikke blot skyldes teoretisk Trang. Han erklærer,
at den Verdensanskuelse, til hvilken vi vilde være henviste,
naar vi ikke vilde gaa den Vej, han gaar, er »tom og trøstesløs«.
Og rent erfaringsmæssigt set staa vi, siger han fremdeles,
med vore etiske Idealer, saa værdifulde og uundværlige de end
maatte være, ved en Afgrund, over hvilken der ikke fører
nogen Bro.

Der savnes her en Række Undersøgelser af psykologisk og
etisk Natur til nærmere Belysning af den Trang, ved hvilken
den ontologiske Enhedsidé saaledes skulde blive praktisk nødvendig
og berettiget. I første Udgave af Wundts »System«
forekom den Sætning, at Filosofien vel ikke kan gøre Tro til
Viden, men at den kan bevise Troens Nødvendighed. Denne
Sætning er falden bort i anden Udgave, maaske fordi Wundt
har set, at Troens psykologiske Nødvendighed ikke kan bevises.
--- Der savnes desuden en Drøftelse af de Vanskeligheder, en%
% --- p. 24 ---
% running head: Etik og Folkepsykologi .
\opage{24}hver Monisme, og især en paa etiske Motiver grundet Monisme
har at kæmpe med overfor de Disharmonier, Erfaringen frembyder\textsuperscript{10}). ---


Ved Slutningen af Wundts Behandling af det etiske Problem
ville vi vende tilbage til nogle af de her berørte Punkter.

5. Om \textit{det etiske Problems} noget uklare Stilling i Wundts
Filosofi har jeg allerede talt ovenfor. Den finder for en Del
sin Forklaring ved Karakteren af Wundts Etik, idet denne dels
tenderer til at blive Et med Folkepsykologi, dels munder ud i
en metafysisk eller religiøs Betragtning.

\textit{a.} Folkepsykologien er ifølge Wundt Etikens Forhal. Han
har i sin Alderdom paabegyndt et stort og bredt anlagt Værk,
der skal behandle Folkepsykologiens vigtigste Emner: Sprog,
Myte og Sæd (Sitte); kun den første Del foreligger endnu, og
en sammenhængende Drøftelse af Wundts Folkepsykologi kan
derfor endnu ikke anstilles. Men hans Hovedtanke er, at den
individuelle Bevidsthed gennem Sprog, Religion og fælles Livsvaner
og Handlevaner hænger sammen med Folkets, ja med
hele Menneskehedens Liv. Den individuelle Villie finder sig
selv som Element i en Totalvillie (Gesammtwille), af hvilken
den bestemmes baade i de Motiver, af hvilke den ledes, og i
de Formaal, henimod hvilke den stræber. Kultur og Historie
udgøre et sandt Fællesliv og ere ikke blotte Resultater af talløse
Enkeltbestræbelsers Sammentræf. Med stor Uret betragter
den Individualisme, under hvilken hele den nyere Tid lider,
Enkeltvillien som den eneste Realitet. Der existerer intet oprindeligt
isoleret individuelt Menneske. Individualiseringen foregaar
successivt, ud fra en Tilstand af social Enhed og uden
nogen fuldstændig Løsrivelse fra Totalvillien.

Samfundets Existens er den vældigste af alle historiske
Kendsgerninger. Af det bæres den Enkelte, selv hvor han
tror at røre sig mest selvstændigt; ved det bestemmes Sympatien
og Pieteten, den sociale Følelses Grundlag. Og Delagtigheden
i Totalaanden (Gesammtgeist) er størst hos Menneskehedens
store Aander, de førende Aander; de formaa at drage
saadan Næring af Totalaanden, at de igen kunne anvise den
nye Opgaver, nye Retninger.

% --- p. 25 ---
% running head: Formaalenes Udvidelse og Forskydning.
\opage{25}Det vilde være Aderitisme, om Værdien af Menneskehedens
Historie skulde bero paa, hvorvidt den Enkeltes eller
den enkelte Menneskegruppes Vel fremmes. Individer og Folkeslag
ere forgængelige, og de ere underlagte Lidenskaber, Fordomme
og Svagheder. Men Historiens Aand er uforgængelig
og har altid Ret. Den historiske Udvikling følger Love, som
det enkelte Individ og det enkelte Folk, der staa midt i Udviklingen,
ikke kunne overse. Af største Vigtighed er her den
Kendsgerning, at menneskelige Handlingers Virkninger stedse
mere eller mindre række ud over Individernes bevidste Motiver
og Hensigter. Heri ytrer sig en Formaalenes Metamorfose
(die Heterogonie der Zwecke kalder Wundt det), som gør nye
subjektive Motiver mulige, da de uforudsete Virkninger kunne
fremkalde nye Følelser og Drifter. Nydannelsen af Motiver ud
af givne Virkninger er den vigtigste Udviklingslov, som gælder
for den etiske Bevidsthed\textsuperscript{11}). Den gør, at vi ikke kunne blive
os de højeste Formaal for Udviklingen bevidste; den Retning,
i hvilken de ligge, kunne vi vel skønne, men Overbevisningen
om, at alle Stadier tjene en stadigt fremskridende Udvikling,
kan tilsidst kun være en Tro, ingen Viden. ---

Den Vanskelighed, som heri ligger for en videnskabelig
Etik, fremhæver Wundt ikke tilstrækkeligt. Det synes dog
klart, at Muligheden af Nydannelser maa gøre al Etik mere
empirisk, end Wundt med sin Appel til Historiens Aand vil
indrømme. Vi staa her atter ved et Punkt, hvor han minder
om Romantikens Filosofi. Og han vedkender sig det aabent,
at hans Etik i visse Grundtanker staar den spekulative Idealisme
nær. Allerede Hegel havde anerkendt en real sædelig
Magt i Totalvillien. Hegels Fejl var ifølge Wundt, at han kun
betragtede Totalvillien som objektiv sædelig Magt, medens den
individuelle Villie blot skulde være den ubevidste Bærer og
Fuldbringer. I sin Lære om Forholdet mellem Individ og Samfund
søger Wundt at raade Bod paa denne Ensidighed.

\textit{b.} Den Enkelte bæres vel af Samfundet, men virker ved
sin egen Tanke- og Villiesretning tilbage paa det. Enkeltbevidstheden
er skabende, Samfundsbevidstheden bevarende. Det Nye
kommer fra de Enkelte, men Samfundet gør det frugtbringende
% --- p. 26 ---
% running head: Totalvillien og Enkeltvillierne.
\opage{26}for en
senere Udvikling og virker derved for det aandelige
Livs Kontinuitet. Ikke Alle ere her produktive i samme Grad.
Det er kun de førende Aander, der afgørende bestemme Totalvilliens
Retning, og allermest gælder dette om en saadan etisk
Genius, som Historiens Aand maaske kun frembringer én Gang
i Aarhundreder eller i Aartusinder, og som fører det etiske
Liv videre ved at vække hidtil slumrende Drifter til Live.

Naar Wundt saaledes, i Modsætning til Hegel, lægger Produktiviteten
over i de enkelte Individer, kommer han i Modsigelse
med sin egen stærke Fremhæven af »Totalvillien«. Konsekvent
maatte han tillægge Individualismen en langt større Betydning,
end han gør, ti Udviklingen af kraftige Individualiteter
maa være et etisk Hovedformaal. I denne Henseende
kunde han have lært Meget af den engelske Skole, som han ser
ned paa formedelst dens Empirisme. Og ikke blot som Udgangspunkter,
ogsaa som Endepunkter, som Formaal maatte han
konsekvent stille de enkelte Individer højere, end han gør.
»Den Enkelte«, siger han, »være nok saa rigt udstyret af Lykken
og nok saa fuldkommen, --- han er dog en Draabe i Livets
Hav. Hvad kunne hans Lykke og hans Smerte betyde for
Verden?« --- Jeg kan ikke besvare dette Spørgsmaal. Men jeg
formaar ikke at forbinde nogen Mening med Ord som Formaal
eller Værdi, naar de slet ikke belyses ved et Forhold til de
Livsvilkaar, som Væsener med Evne til at føle Lyst og Smerte
ere underlagte. Wundts Etik ender i en mystisk Dualisme,
idet Formaalene ligge i »Totalvillien«, Midlerne i »Enkeltvillierne«.
Og for de tragiske Konflikter, der kunne opstaa ved Sammenstød
mellem Enkeltvillien og en historisk udformet Totalvillie
har han ingen Plads i sin Etik\textsuperscript{12}).

Karakteristisk for denne Side ved Wundts Etik er ogsaa
den Sætning, at Retfærdighed ikke er en individuel, men en
offentlig Dyd, fordi den forudsætter Magt til at stifte Ret og
paalægge Pligt. Han overser (maaske under Indflydelse af tysk
Bureaukratisme), at ethvert Individ gennem sin Stilling i Familie,
Samfund og Stat har en vis Magt, og at det i sin Overbevisning,
i sine Domme om Andre og i sin Færd overfor dem
har Adgang til at vise Retfærdighed. Selv Trællen har det i
% --- p. 27 ---
% running head: Etik og Metafysik.
\opage{27}Forhold til Herren. Magttanken har gjort Filosofen blind for
Betydningen af de mange selvstændige Udgangspunkter for
Dømmen og Handlen.

Og dog vise flere Ytringer, at Wundt har Øje for den
nuværende sociale Ordnings Misligheder. Han anker over, at
man bedømmer Tyven strengere end Ødelanden, Aagerkarlen
og Hazardspilleren. Han finder, at det nu bestaaende Forhold
mellem Besiddelse og Arbejde frembringer to modsatte Anledninger
til Immoralitet: Besiddelse uden Kald avler Nydelsessyge,
Besiddelsesløshed uden Kald avler Misundelse. Disse Onder
kunne kun afhjælpes ved en ny Retsorden. Vor nuværende
Samfundstilstand lider under Modsætningen mellem en forældet
Retsopfattelse og nye Kulturelementer, der ikke kunne bringes
ind under de gamle Begreber.

\textit{c.} Wundt hævder Etikens Selvstændighed overfor Spekulation
og Metafysik, netop af den Grund, at den afgiver et af de
vigtigste Bidrag til Fundamentet for en almindelig Verdensanskuelse.
Og som vi have set, fortoner Etiken sig ogsaa tilsidst
over i Spekulation eller Tro. Ja, Etiken behøver ifølge
Wundt mere end andre Omraader en metafysisk Supplering.
I Kraft af denne Opfattelse (der staar i Modsigelse med Antagelsen
af Etikens Selvstændighed overfor Metafysiken) gaar
Etiken for Wundt over til Religionsfilosofi. Naar Idealerne
gaa ud over det, der kan naas ved menneskelig Stræben, faa
de en religiøs Karakter. Filosofien formaar her kun at give
ubestemte Antydninger; de positive Religioner give derimod
konkrete Symboler. Den religiøse Hovedtanke er »den Fordring,
at alle aandelige Værker (geistige Schöpfungen) besidde
en absolut, eller uforgængelig Værdi«. Jo mere Religionen
hæver sig over Naturreligionens Standpunkt, des mere stemmer
den med Videnskaben; en Gud, der gør Mirakler er en
Naturgud, ikke den etiske Religions Gud. Kristendommens
Udvikling undgik ikke Tilbagefald til Naturreligionen. Men Kristendommens
sluttelige Opgave maa --- i Henhold til dens Stifters
egne Udsagn --- være Overvindelse af alle saadanne Elementer
i den religiøse Tro, der hæmme de religiøse Ideers
etiske Indhold. Kristus vil, netop naar man opfatter ham ikke
% --- p. 28 ---
% running head: Wundts Objektivisme.
\opage{28}som Gud, men som Menneske, vedblive at staa som etisk Forbillede,
og tillige som Vidne om Verdens uendelige, uudgrundelige,
men med det etiske Ideal identiske Grund og Formaal. ---

Wundts Filosoferen, saaledes som jeg nu har søgt at karakterisere
den, kan siges at være typisk for vor Tids Tænkning.
Hvor den trænger til Korrektiver, antyder den selv de Synspunkter
og Metoder, der maa anvendes for at føre den videre.
Navnligt er dette Tilfældet, hvor han vil »supplere« Erfaringsvidenskaben
ved metafysiske og religiøse Elementer. Saa energisk
han søger at fastholde og gennemføre et rent objektivt
Standpunkt, gør der sig dog subjektive Faktorer gældende, der
kræve en bestemtere Undersøgelse. Wundts Objektivisme, der
paa visse Punkter bliver til Mystik, er baade hans Styrke og
hans Begrænsning.

% --- p. 29 ---
\chapterheadb{II.}{ROBERTO ARDIGÒ.}

\opage{29}1. I min Fremstilling af den nyere Filosofis Historie har
jeg kun for Renaissancetidens Vedkommende havt Lejlighed til
at omtale italiensk Filosofi. Det var kun i Renaissancen, at
den italienske Tænkning fremtraadte med en Originalitet og
med en Dristighed, der kunde give den Betydning for Filosofiens
almindelige Udvikling. Den Aand, der i Oldtiden havde
besjælet Pythagoræer og Eleater, og som havde inspireret Lucretius
til hans vældige Læredigt, levede op igen i Pomponazzi
og Telesio, i Bruno og Galilei. Men den italienske Renaissancebevægelse
blev slaaet til Jorden efter Brunos Baal og
Galileis aftvungne Mened. I de følgende Aarhundreder er
der kun enkelte Navne af filosofisk Interesse at nævne, som
\textsc{Giambattista Vico} († 1744), Forløberen for den moderne
Sociologi. I Slutningen af det 18. Aarhundrede fik den franske
Filosofi stor Indflydelse, og om Psykologen og Sociologen
\textsc{Romagnosi} († 1835) samlede sig en Gruppe begejstrede Disciple.
Men ved Midten af det 19. Aarhundrede fik den filosofiske
Bevægelse en anden Karakter, idet religiøs, national og
filosofisk Begejstring indgik en nøje Forbindelse. Friheds- og
Enhedsbestræbelserne i Italien rørte sig mægtigt hos en Skare
af unge Gejstlige, der drømte om en stor Harmoni mellem Religion
og Tænkning og mellem Kirke og Stat og som troede
paa Muligheden af, at Kirken vilde stille sig i Spidsen for den
% --- p. 30 ---
% running head: Blik paa italiensk Filosofi i det 19. Aarhundrede.
\opage{30}nationale Reisning. Ligesom paa Renaissancens Tid Begejstringen
for Italiens Frihed laa bagved Machiavellis politiske Filosofi,
saaledes besjælede den nu \textsc{Rosmini's} († 1855) og \textsc{Gioberti's}
(† 1852) Spekulationer. Deres Filosofi var en Art
Platonisme, baaren af Troen paa en evig, over al Erfaring ophøjet
Sandhed. Filosofien var for dem dels Indledning til Religionen,
en Art Logoslære, dels Midlet for Patriotismen. Forskellen
mellem de to moderne Platonikere beror især paa, at
Gioberti antager en umiddelbar Anskuen af den ideale Sandhed,
medens Rosmini hævder (paa lignende Maade som Schelling
i sin senere Lære), at Tænkningen kun fører til et System
af Muligheder, og forsaavidt staar den kritiske Filosofi nærmere.
\textsc{Mamiani} († 1885) udviklede den italienske Idealisme i en
Retning, der gav Erfaringen større Indflydelse, og i denne mæglende
Form raadede den en Tidlang paa de italienske Universiteter.


Den Rytme, der saa ofte viser sig i de filosofiske Ideers
Udvikling, ytrede sig paa karakteristisk Maade i den Rolle, en
positivistisk Retning kom til at spille i Italien i den sidste
Menneskealder. Forskellige Aarsager samvirkede ved denne
Ændring i de filosofiske Interesser. Italien havde opnaaet sin
Enhed og Frihed, saa at der var Plads for specielle videnskabelige
Interesser; den ideale Begejstring kunde afløses af
det reale Arbejde. Samtidigt kom der en stærk Paavirkning af
den nyere franske og engelske Filosofi gennem Comtes og Mill's
Skrifter; \textsc{Villari} anvendte Positivismens Grundtanker paa Historieopfattelsen,
\textsc{Angiulli} især paa Psykologiens og Pædagogikens
Omraade. Dertil kom den moderne Naturvidenskabs
Indflydelse. Og paa den anden Side stillede den katolske Kirke
sig nu mere og mere skarpt overfor enhver Filosofi, der ikke
strengt gik i den middelalderlige Tænknings Spor. Rosminis og
Giobertis Disciple indenfor Gejstligheden bleve forfulgte. Man
udtalte endog, at det nu var Opgaven at føre Civilisationen tilbage
til den katolske Opfattelse, fra hvilken den i de sidste tre
Aarhundreder var afvegen! 1864 udsendte Pave Pius IX en
»Syllabus«, en Fortegnelse over Tidens Vildfarelser, hvortil regnes
Samvittighedsfrihed, Naturalisme og Rationalisme, og ganske
% --- p. 31 ---
% running head: Ardigòs Overgang fra Katolicisme til Positivisme.
\opage{31}særligt den Antagelse, at den skolastiske Teologis Metode og
Principer ikke skulde stemme overens med vor Tids Krav og
Videnskabens Resultater. I en Encyclica af 1879 erklærede
senere Pave Leo XIII, at Thomas Aquinas' Filosofi var et guddommeligt
Hjælpemiddel, som af alle Lærere bør lægges til
Grund og benyttes til Gendrivelse af de moderne Vildfarelser.
Senere har han i en Rundskrivelse til de franske Biskopper
(1899) særligt fordømt den kritiske Filosofi\textsuperscript{13}).

Saaledes vare ikke blot Motiverne til den intellektuelle Harmoni,
de italienske Filosofer ved Aarhundredets Midte havde
haabet og tilstræbt, faldne bort, men Modsætningen mellem Kirken
og Videnskaben var skærpet, dels ved Tænkningens Udvikling,
dels ved Kirkens Forsøg paa at skrue Tiden tilbage.

Under disse Forhold udviklede \textsc{Roberto Ardigò} sig i
Klosterets Ensomhed fra barnligt troende og patriotisk begejstret
Katolik til energisk Positivist. Hans Udvikling og hans Standpunkt
har ikke blot Interesse som Tidssymptom, men er af
Betydning for Problembehandlingen ved hans Tankekraft og
hans psykologiske Talent, saavel som ved den ejendommelige
Ændring, Positivismen undergaar i hans Hænder.

2. Ardigò er født i Nærheden af Cremona i Aaret 1828.
Faderen, der var Landmand, flyttede til Mantova, for at Sønnen
kunde studere. Moderen var en from Katolik, hvem Ardigò
oftere i sine Skrifter omtaler med inderlig Taknemlighed.
Det skyldtes hendes Indflydelse, at han blev Præst. Efter
Forældrenes Død tog Biskop Martini sig af ham og gjorde ham
til Kannik ved Domkirken i Mantova. Han studerede i sin Ensomhed
med Iver skolastisk Filosofi og moderne Videnskab, inderligt
overbevist om, at de »moderne Vildfarelser« maatte kunne
gendrives. Men langsomt rejste der sig --- tydeligere for hans
Omgivelser end for ham selv --- en helt ny Tankebygning i
hans Indre, og det gik da tilsidst i et bestemt Øjeblik op for
ham, at han ikke længere hyldede Kirkens dualistiske Lære, men
at han troede paa en stor Kontinuitet i Alt. I et af sine Skrifter
\textit{(La morale dei positivisti} II, 3, 2, i et Kapitel, der handler
om Muligheden af Moral uden Religion) udtaler han sig
om dette Gennembrud paa en interessant Maade.

% --- p. 32 ---
% running head: Ardigòs egen Skildring.
\opage{32}»Ved min Moders, en enfoldig og fattig Bondekones, Exempel
og Lærdom, opstod og næredes i min Sjæl religiøs Tro og
religiøst Liv. Endnu den Dag idag kan jeg ikke tænke paa
den ophøjede Enfold i min Moders religiøse Følelse, uden at
derved den største Begejstring og den ømmeste Stemning vækkes
hos mig. Det tro Billede af denne Følelse er endnu til i min
Søster, som hverken skylder Skole eller dannede Kredse sin
Udvikling, og jeg føler overfor den Maade, hvorpaa det fremtræder
hos hende, den største Ærefrygt, som for en Helligdom.
Denne min barnlige Religiøsitet styrkedes senere i høj Grad
ved Samlivet med Mgr. Martini, der strakte sig over mere end
tyve Aar. Han havde taget mig til sig efter mine Forældres
Død; han gav mig det Brød, jeg manglede, og gjorde mig den
videnskabelige Løbebane mulig ... Han, der ansaa Videnskab
og oprigtigt Sindelag for Religionens Grundvold og Sjæl,
elskede mig, fordi han saa mig begejstret for Studiet og fandt
en Karakter uden Svig hos mig. Han haabede, at jeg engang
skulde blive til Nytte for Kirken imod den uvidende,
overtroiske, bigotte og hykleriske Art af Religion, som han afskyede
.... Jeg hengav mig (foruden til Studier af Naturvidenskaberne
og af Filosofien, som jeg aldrig har opgivet) med
hele min Sjæl til Teologien, især den dogmatiske og apologetiske.
Jeg samlede mig et Bibliotek af Kirkefædre og Teologer, idet
jeg paa Studiet af dem, især af St. Thomas, anvendte mine
bedste Ungdomsaar. Tilsidst forfattede og offentliggjorde jeg
en Bog om Skriftemaalet imod Protestanterne. Men Udgangen
af mit Studium blev ganske modsat den, jeg arbejdede
hen imod og ventede. Lidt efter lidt kom Tvivlen, der allerede
havde rejst sig paa alle Sider i mine tidligere Aar, men som
jeg med uafbrudt Eftertanke og Studium havde bekæmpet og i
lang Tid havde trot overvundet, tilsidst til at staa uden Modstander
og fremtraadte en skøn Dag til min største Forundring
som afgjort Overbevisning og ubestridelig Vished. Forunderligt!
Ligetil den Dag havde jeg hengivet mig til Bestræbelsen
for at staa fast ved min gamle religiøse Tro, og saa havde
dog i mit Indre, uden mit Vidende, det positivistiske System
udfoldet sig fuldstændigt under det System af religiøse
% --- p. 33 ---
% running head: Ideel Kontinuitet trods Bruddet.
\opage{33}Ideer, der var Frugter af saa meget Arbejde og af saa lang
Tid. Det nye System fandt jeg til min største Forbavselse allerede
fuldt færdigt og urokkeligt fæstnet i min Aand i selvsamme
Øjeblik, i hvilket jeg --- siddende paa en Sten under
en Busk i den lille Have, jeg havde indrettet ved min Kannikebolig
--- mærkede, hvorledes den sidste Reflexion sønderbrød
den sidste Traad, der endnu holdt mig bunden til Troen.
Det forekom mig da pludseligt, at jeg aldrig nogensinde i hele
mit Liv havde trot, og at jeg aldrig havde gjort Andet end
beflitte mig paa at udvikle den rent videnskabelige Tendens
hos mig. Og dette udsprang, tror jeg, af den Iver, hvormed
jeg fra først til sidst havde søgt saavidt muligt at kende alle
de Grunde, der stred mod Religionen, for at være i Stand til
at tro med god Samvittighed og til at værne Troen mod alle
Angreb.« --- Ardigò tilføjer, at det Skridt, han nu gjorde, var
forbundet med Smerte og Vemod paa Grund af Erindringen
om hans Moder og paa Grund af Forholdet til hans Velgører.
Men Villien blev styrket i denne Kamp, efter hvilken Religionen
kom til at staa som en poetisk Erindring. Og han erfarede,
at den moralske Idealitet ikke svækkedes; tvertimod stod det
fra nu af endnu mere afgjort for ham, at de sande Goder ere
de, der erhverves ved at følge Samvittighedens Lov i nyttig
Virksomhed og ved at hengive sig til Tankens Arbejde.

Der kom nu nogle trange Aar for Ardigò. Han maatte
ernære sig ved Skoleundervisning, og først efter adskillige Aars
Forløb udnævnte en liberal Undervisningsminister ham til filosofisk
Professor i Padova (1881). Hans Retning stred ikke blot
mod Teologien, men ogsaa mod den Retning, som under Mamianis
Indflydelse herskede paa de italienske Universiteter.
Hans Undervisning har øvet stor Indflydelse og vakt stor Begejstring,
hvad der fik offentligt Udtryk i et Festskrift i Anledning
af hans 70. Fødselsdag (1898).

Ardigòs Hovedproblem hænger nøje sammen med hans
personlige Udvikling. Allerede i en Tale, han holdt 1869 om
Pietro Pomponazzi, i hvem han ser en Forgænger, skildrer han
Tanken som en Kraft, der bliver til, uden at det mærkes, og
som, naar den er modnet, bryder frem paa uimodstaaelig Maade.
% --- p. 34 ---
% running head: Ardigò og Positivismen.
\opage{34}Det var det, han selv havde erfaret, og det var ham et Exempel
paa, hvorledes al Udvikling gaar for sig. I sin Tiltrædelsestale
i Padova (1881) hentyder han til sin egen Udviklingsgang
som den, der havde stillet ham et Problem, der kommer igen
for større Forholds Vedkommende, ligesom en Molekules Bevægelse
kan minde om en Verdensklodes Rotation. Al Udvikling
bestaar i en Overgang fra det Ubestemte til det Bestemte, dog
saaledes at Kontinuiteten bevares: den Helhed, der paa mere
ubestemt Maade fremtræder paa det første Stadium, ligger stadigt
til Grund som bærende Grundlag for de Forskelligheder,
der successivt fremtræde. Den første videnskabelige Opgave,
han stillede sig, var den at forfølge den menneskelige Tankes
naturlige Udvikling (la formazione naturale del pensiero). Denne
Plan blev ikke udført, men det Materiale, han havde samlet,
findes anvendt rundt om i hans forskellige Skrifter. Men Udgangspunktet
fik afgørende Betydning for hans Standpunkt.

Han hævder, at han har naat sit Resultat uafhængigt af
den franske og engelske Positivisme, som han først senere har
lært at kende. Naar hans Opfattelse af Udviklingsbegrebet minder
om Spencers, bliver dog den Forskel tilbage, at medens
Spencer væsentligt bygger paa Analogien med biologisk Udvikling,
bygger Ardigò væsentligt paa Analogien med Tankens Udvikling,
denne vidunderligste af alle naturlige Formationer (la
piû miragliosa delle formazione naturale). Ardigò kalder sig
Positivist. Men det Væsentlige i Positivismen er for ham Udgangspunktet
i Erfaringen, ikke Afslutningen. Positivisten har
ingen Hast med at slutte; han vil ikke danne sig en Ide. der
kan bruges som Partifane, men gaar frem Skridt for Skridt,
efterhaanden som Sandheden aabner sig for ham. Det gælder
om at holde Horizonten aaben. Denne Tankegang kommer særligt
frem hos Ardigò overfor Problemet om Tilværelsens Enhed,
i Modsætning dels til Materialismen, dels til Tilbøjeligheden
til at blive staaende ved en Mængde spredte Fakta. I sit seneste
Skrift \textit{(L'unità della coscienza} 1898) udtaler han sig særligt i
denne Retning, hvad det psykologiske Problem angaar. En af
hans Disciple, \textsc{Marchesini}, har i et særligt Skrift, som han
har kaldet »Positivismens Krisis« \textit{(La crisi del positivismo.} To%
% --- p. 35 ---
% running head: Det Usondrede og Sondringerne.
\opage{35}rino 1898), fra dette Synspunkt drøftet Forholdet mellem Ardigòs
Filosofi og Positivismen. En vis Lighed mellem Ardigòs
Standpunkt og Wundts kommer frem paa dette Punkt. Allerede
Villari og Angiulli stode iøvrigt paa lignende Maade kritisk
overfor Positivismens ældre Former.

3. De Grundtræk, Ardigò havde fundet i Tankelivets Udvikling,
genfinder han i al Naturudvikling. Tankens Udvikling
er for ham Exempel paa en Verdenslov. I sit Skrift \textit{La formazione
naturale nel fatto del sistema solare} (1877) har han
nærmere udviklet denne Lov gennem Analyse af Kants og
Laplaces bekendte Hypotese, hvilken han betragter som Exempel
paa en videnskabelig Forklaring. Egentligt skulde det nævnte
Skrift kun have været et enkelt Kapitel af det paatænkte Værk
om de menneskelige Forestillingers Udvikling.

Ifølge den nævnte Hypotese skyldes Solsystemets nuværende
Tilstand en Sondringsproces (\textit{distinzione}), idet der har
dannet sig mindre Dele eller Helheder indenfor den store, usondrede
Masse. Men Helheden ophæves derved ikke; Helheden
--- det Usondrede (l'Indistinto) --- bestaar stadigt, og kun derved
bliver det forstaaeligt, hvorledes der kan være Vexelvirkning
mellem de udsondrede Dele (Kloderne): disse høre sammen
nu, ligesaa vel som før de udsondredes. Det Usondredes
Bestaaen er Grunden til Solidariteten. --- Spørger man,
hvorledes Forskellighederne kunde opstaa af det Usondrede,
svarer Ardigò, at de allerede vare til Stede i Begyndelsesstilstanden
som Muligheder eller i latent Tilstand (forza latente, o
virtuale); den specielle Form naas gennem successiv Udfoldelse.
Han indrømmer, at denne Forklaring blot er en Tilbageslutning
fra Erfaringen. Vi kunne ikke paa Forhaand vide, hvilke specielle
Former og Dele der ville udvikle sig af den usondrede
Begyndelsestilstand. Kun Iagttagelse og Forsøg kan paavise
dem. Hverken Matematiken eller Metafysiken kan konstruere Naturen.
Og Grunden hertil er, at den usondrede Helhed ikke
specialiseres »af sig selv«, men kræver ydre, for os uforudseelige
og \textit{forsaavidt} »tilfældige« Betingelser for at faa sine specielle
Former og Dele. Men hvad der saaledes successivt udfoldes,
er dog kun Ækvivalentet for, hved der allerede bestod som
% --- p. 36 ---
% running head: Udviklingslære og Erkendelseslære .
\opage{36}koexisterende i den usondrede Tilstand. I Verdenssubstansen
opmagasineres det af de forudgaaende Aarsager udførte Arbejde,
og derved betinges Muligheden af de fremtidige Virkninger\textsuperscript{14}).
Den Kontinuitet, der saaledes fremtræder mellem l'Indistinto og
de udformede Dele (distinti), er det, vi udtrykke ved Ordet
\textit{Natur.} (La natura è la continuità di una cosa con tutti le altre.
\textit{Form. nat.} p. 205).

Ardigò ser nu dog klart, at der maa spørges, hvori hine
Muligheder eller »latente Kræfter« egentligt bestaa. Mulighed
er for ham kun en Virkelighed af anden Beskaffenhed end den
givne, og den bestaar i en Aktivitet: l'essere è attività \textit{(L'unità
della coscienza} p. 479). Kontinuiteten er en stadig Energi. Derfor
viser enhver usondret Helhed, ethvert Indistinto tilbage til
et endnu mere omfattende Hele, hvoraf den har udsondret sig.
Forskellen mellem indistinto og distinto, saavel som den mellem
Mulighed og Virkelighed, er derfor stedse kun relativ. En
Afslutning er her ikke mulig. Vi komme ind i en uendelig
Række. Men Videnskabens sidste Ord angaar Grundforholdet
mellem distinto og indistinto og udsiger, at alle Forskelligheder,
hvorsomhelst de frembyde sig, udspringe af et Hele og omfattes
af et Hele. ---

Ardigòs Udviklingslære og hans Erkendelseslære hænge
nøje sammen. Enhver Forklaring er en Sondren, en Analyse.
Hvad der vedbliver at staa usondret, staar som uforklaret. Tanken
drives stedse videre frem, skønt den har en Tilbøjelighed
til at stanse ved et distinto finito; --- og det er selve Naturens
Uendelighed, der saaledes driver Tanken frem. Ja, det er
den uendelige Natur, der er selve Energien i den logiske Tænkens
egen Lov\textsuperscript{15}).

At Tankens eller Forklaringens Lov saaledes er den samme
som Naturens Lov, kan ikke forundre, da selve Tanken er Natur
eller en naturlig Formation, ligesom alle andre Ting. Men
naar Tanken er en Erfaringskendsgerning som enhver anden,
og selv kun er et Exempel paa naturlig Formation, saa er det
umuligt at forklare hele Naturen ved at udlede den af Tanken,
saaledes som Metafysiken og Teologien ville. Selve Tanken forstaas
jo kun ved den almindelige Lov om Udvikling fra in%
% --- p. 37 ---
% running head: Psykologiens Selvstændighed.
\opage{37}distinto til distinti. Vi ere henviste til at forklare den ene
Kendsgerning ved den anden, uden at kunne faa nogen absolut
Afslutning. Naturen er en umaadelig Cirkel, hvis Midtpunkt
er overalt, og hvis Omkreds er intetsteds. ---

Ardigò har i for høj Grad ladet Erkendelseslæren blive
absorberet af Udviklingslæren. Psykologisk eller biologisk betragtet
er Tankens Udvikling ganske vist kun et enkelt Exempel
paa Udviklingens almindelige Love. Men Erkendelsesproblemet
melder sig, saasnart der spørges om, hvorledes vi begrunde
Gyldigheden af de almindelige Love, vi mene at finde
for Tanken saavel som for andre Fænomener. Dette Problem
falder ikke bort ved blot Henvisning til Udviklingens Lov som
fælles for alle Fænomener. Der er jo det Særegne ved Tanken,
at det er ved og gennem den, vi opfatte og udtrykke al
Udvikling og al Lov, baade dens egen og andre Tings. Og
Spørgsmaalet er da, hvilken Gyldighed denne Opfattelse og Udtryksform
har\textsuperscript{16}).

4. Ardigò har med Iver virket for \textit{Psykologiens} Selvstændighed
som Erfaringsvidenskab, og herved kom baade hans
Talent for psykologisk Beskrivelse og Analyse og hans Kendskab
til Naturvidenskaben ham til Gode. Hans psykologiske Interesse
og Begavelse gaar i to Retninger. Dels søger han at
finde større Sammenhæng i Sjælelivet, end der for foreløbig
Iagttagelse synes at være; han stræber at finde Kontinuitet
mellem de bevidste Tilstande indbyrdes og mellem de bevidste
og de ubevidste Tilstande. Dels søger han at paavise finere
Nuancer og Forskelligheder, end den sædvanlige Iagttagelse
kender; han paastaar saaledes, at tilsyneladende ensformige sjælelige
Tilstande i Virkeligheden bestaa af rytmiske Bevægelser.
Nogle af hans mest fremragende Disciple, som Fysiologen Giulio
Fano og Kriminalisten Enrico Ferri have i Festskriftet til den
gamle Filosof især dvælet ved disse Ejendommeligheder ved
hans Undervisning. Det er jo de to Retninger, i hvilke al
Forsken gaar og maa gaa; men Ardigò har --- som ogsaa hans
Lære om Indistinto og Distinti og om Forholdet mellem dem
viser --- havt særligt Øje for dem og søgt at forene dem.

Hans psykologiske Hovedskrifter ere \textit{La psicologia come}
% --- p. 38 ---
% running head: >>Den psykofysiske Realitet<< .
\opage{38}\textit{scienza positiva} (1870) og det Skrift, han selv har betegnet
som sit filosofiske Testamente: \textit{L'unità della coscienza} (1898).

Psykologiens Opgave er for Ardigò Studiet af vore indre
Tilstande, hvis fælles Ejendommelighed udtrykkes ved Begrebet
Sjæl. Det er et subjektivt Faktum, Psykologien gaar ud fra;
men det maa belyses ad objektiv, fysiologisk Vej. Denne dobbelte
Belysning viser os Sjælefænomenernes Forskellighed fra
alle andre Fænomener. Saa stor Betydning fysiologiske Undersøgelser
end have for Psykologien, saa kan Fysiologien dog
aldrig træde i Stedet for Psykologien. Forholdet er dette, at
de psykiske og fysiologiske Fænomener ere Udtryk for en og
samme »Substans«, den psykofysiske Substans, eller som Ardigò
hellere kalder den: den psykofysiske Realitet (realtà psicofisica).
I hans senere Skrifter (hvor hans almindelige Udviklingslære
fremstilles) er den psykofysiske Realitet Et med l'indistinto\textsuperscript{17}).
Det er en naturlig Totalitet, som gaar forud for og ligger til
Grund for de Forskelligheder, vi betegne som Sjæl og Materie.
Begrebet Indistinto udtrykker her som overalt Enheden og Solidariteten.
Baade Materialismen og Spiritualismen grunde sig
paa blotte Abstraktioner. Fysiologisk set hænger den aandelige
Virken sammen med de fysisk-kemiske Processer. Psykologisk
set er det Materielle ligesaa vel som det Sjælelige kun givet
for os som Fornemmelse og Forestilling; baade Bevægelse og
Tanke kende vi kun som psykiske Akter. I Stedet for at spørge
om, hvorledes Materie kan gaa over til Sjæl, var det rigtigere
at spørge om, hvorledes vore oprindeligt mere ubestemte Forestillinger
efterhaanden kunne differentieres, saa nogle af dem
komme til at staa som Udtryk for et Jeg, andre som Udtryk
for et Ikke-Jeg. Den populære, materialistiske Opfattelse skriver
sig ifølge Ardigò fra, at man ikke indser, at Egenskaber
som Udstrækning og Bevægelse ligesaa vel skyldes psykiske
Funktioner som Farve og Lugt og de andre saakaldte sekundære
Egenskaber.

Ideen om den psykofysiske Realitet giver dog --- det indskærper
Ardigò bestemt --- ingen Forklaring. Dens Betydning
er kun at hævde en Sammenhæng, som vi i vore Abstraktioner
ere tilbøjelige til at opløse. Han sammenligner den med
% --- p. 39 ---
% running head: Den psykologiske Grundtanke.
\opage{39}Ideen om Gravitationen, der jo heller ikke giver nogen Forklaring.
Det maa overlades Fremtiden, om den kan finde en
virkelig Forklaring. Positivisten har ingen Hast med at slutte.
Videre end til et foreløbigt Begreb kunne vi ikke komme. Forholdets
indre Natur kunne vi ikke bestemme. Den Gaade, der
bliver tilbage, staar dog for Ardigò ikke alene. Hvorledes en
Tanke kan være Ækvivalent for en Bevægelse, er for ham ikke
nogen større Gaade end, hvorledes en Billardkugle kan sætte
en anden i Bevægelse ved et Stød. Overalt i Naturen kende
vi kun Samtidigheds- og Successionsforholdene ved Erfaring,
ikke gennem Indsigt i Leddenes indre Væsen. --- Ved denne
sidste Bemærkning overser Ardigò dog det Faktum, at Videnskaben
overalt, hvor det er muligt, stræber at omsætte den
ydre Sammenhæng mellem de forskellige Fænomener til Forbindelse
af ensartede Led, og at derved naas en større Grad af
Kontinuitet. Jo større Modstand denne Bestræbelse støder paa,
des mere gaadefuldt maa det Fænomen siges at være, som vi
stanse ved. ---

I sit filosofiske Testamente, Skriftet om Bevidsthedens Enhed,
udtaler Ardigò sin Beundring for Kants geniale Anelse i
Læren om Bevidsthedsenheden. Vel har han fejlet ved at antage
for stor Modsætning mellem Erkendelsens Stof og dens
Form. Men han har grebet Psykologiens videnskabelige Grundtanke.
Der foregaar i Bevidsthedslivet en uafbrudt Sammensætnings-
og Forbindelsesproces, hvorved alle nedarvede eller
erhvervede Dispositioner og alle nye Elementer virke sammen
i bestemte Retninger. Der er fra først til sidst en Solidaritet
mellem alle psykofysiske Funktioner, en stadig Tendens til at
forenes i én Strøm. Denne Tendens til »aandelig Samstrømmen«
(confluenza mentale) ligger bl. A. til Grund for Forestillingernes
Association. Ofte opdages den kun, naar man tager
Hensyn til ubevidste Begyndelsespunkter eller Mellemled, og
det er særligt dette, der gør, at Fysiologiens Hjælp er saa uundværlig
for Psykologien. Den Enhed, der er Bevidsthedslivets
Særkende, kan ikke forklares som blot Produkt af mangfoldige
Elementers Samvirken, ti disse Elementer opdage vi først ved
en Distinktion, der stedse forudsætter en forud given Helhed.
% --- p. 40 ---
% running head: Den sociale Idealitet.
\opage{40}Og det er fra Iagttagelsen af Bevidsthedslivet, vi overhovedet
hente vore Begreber om Enhed og Flerhed. Overalt, hvor vi
tale om Enhed og Mangfoldighed i Verden, have vi abstraheret
disse Begreber og Ideen om deres Forbindelse fra Bevidsthedslivet.
Mikrokosme og Makrokosme belyse gensidigt hinanden.

5. Ardigòs \textit{Etik} er baaren af de samme Grundtanker,
som karakteriserer de andre Dele af hans Filosofi. Individet
udvikler sig indenfor Samfundet, og staar i samme Forhold til
dette, som det sondrede Element (distinto) stedse staar til Helheden
(indistinto). Og Samfundet udvikler sig --- ligesom Solsystemet
og ligesom Tanken --- ved en naturlig Proces (formazione
naturale), som det er Sociologiens Opgave at beskrive,
og som Etiken forudsætter. Det er Etikens Opgave (som nomologia,
i Modsætning til Sociologien som nomografia og nomogonia)
at skelne mellem de Elementer i Samfundslivet, der have
mistet deres Betydning, og dem, der kunne gøres frugtbringende
paa samme eller ny Maade. Derved baner den Vejen for Fremtidens
praktiske Moral.

Gennem den oprindelige og stadige Sammenhæng med Samfundslivet
vil der hos Individet udvikle sig en antiegoistisk
Følelse. (Ardigò foretrækker dette Ord for Ordet Altruisme).
De menneskelige Forestillinger faa deres første Indhold fra de
sociale Omgivelser. Disses Domme og Interesser blive foreløbigt
gældende for Individet. Og alle Forestillinger have fra
først af en impulsiv Karakter, en Tendens til umiddelbart at
gaa over til Handling; det skyldes en senere Differentiation
(distinzione), naar Tanke og Handling skilles ad. Begge Dele,
det uvilkaarligt optagne sociale Forestillingsindhold og Forestillingernes
oprindelige impulsive Karakter, sammenfatter Ardigò
i Begrebet social Idealitet (idealità sociale). Ved Idealitet forstaar
han Evne til at bestemmes ved Tanker, der pege ud over
Øjeblikket. Den sociale Idealitet udvikles allerede i Familien
som varigt Samfund, der plejer og beskytter Spirerne til Fremtidens
Slægt. Familiefølelsen bliver en Arneild, fra hvilken
Varmen kan straale ud i større Kredse. Den sociale Idealitet
har mange Grader og Former. Den indøves ved Efterligning
og Gentagelse. --- I hele denne Lære udvikler Ardigò en Tanke%
% --- p. 41 ---
% running head: Heroisk Handlen.
\opage{41}gang, der i Oldtiden især fandtes hos Stoikerne, i nyere Tid
hos Adam Smith, og i Nutiden hos Tarde, Leslie Stephen og
Baldwin. Denne sidstes Begreb social Arvelighed vilde passe
godt ind i Ardigòs Filosofi.

Den sociale Idealitet er ikke i lige Grad udviklet hos Alle,
skønt Evnen til uegoistisk Handlen maa siges at ligge i den
menneskelige Natur. Familiekærligheden vidner herom, ligesom
den uvilkaarlige Medfølelse, --- Ærestrangen, der egentligt er
meningsløs hos en Egoist, --- og Harmen over Overgreb og
Retskrænkelser. Paa sit Højdepunkt bliver det etiske Sindelag
en Art helligt Raseri (furore santo), som fører til Selvopofrelse
uden Tanke om egen Fordel, i Troen paa, at der af det Menneskeliges
tragiske Undergang vil hæve sig et Evigt og Guddommeligt
(eterno divino che sorge delle ruine tragiche dell' umano,
et Udtryk, Ardigò laaner fra en aandsbeslægtet italiensk Forfatter).
Det Sande i den teologiske Forestilling om Naaden er,
at der gives en uvilkaarlig Trang, en lidenskabelig Higen, som
fører Mennesket til uinteresserede Handlinger.

En saadan heroisk Handlen er mulig uden Religion, ligesom
Etiken idethele er uafhængig af Religionen, naar man ved
Religion ikke forstaar et Forhold til det Uendelige (efter Max
Müllers Definition), men et Forhold til det Overnaturlige. Det
videnskabelige Begreb om det Uendelige staar i skarp Modsætning
til det Overnaturlige: det betegner den kontinuerlige Funktion,
hvis begrænsede og enkelte Form det Endelige er. Det
Uendelige er Loven, ved hvilken Individets egef Væsen betinges;
derved forstaas den dybe Tilfredsstillelse, Individet føler,
naar dets Higen mod det Uendelige tilfredsstilles. Positivisten
finder saaledes det Uendelige i sig selv, naar han bliver sig
Loven for sit Væsen bevidst. Atter her gentager Forholdet
mellem Element (distinto) og Helhed (indistinto) sig. Forestillingen
om det Overnaturlige bliver derimod til, naar Tingenes
Lov gøres til en fra selve Tingenes Virken forskellig Ting, der
saa kun paa udvortes Maade kan sættes i Forhold til det, der
foregaar i Naturen.

Dog betegner Begrebet det Overnaturlige kun den teoretiske
Side ved Religionen, udgør ikke hele dens Væsen. Det
% --- p. 42 ---
% running head: Ardigò om Religion.
\opage{42}Væsentlige er det Afhængighedsforhold, i hvilket Mennesket med
Hensyn til sit hele Ve og Vel føler sig at staa til det Overnaturlige.
Det Hellige var fra først af det Frygtelige; det er
Frygten, der skaber Guder. I hvert Tilfælde er Følelsen ved
det Hellige ikke, som man har paastaaet, en aldeles enkelt og
uopløselig Følelse\textsuperscript{18}). Gradevis udskilles Frygtelementet ved Overgangen
fra Religion til den videnskabelige Bevidsthed, som overalt
sætter Loven i Stedet for Forestillingen om et hemmelighedsfuldt
Væsen som det, der bestemmer Menneskets Skæbne.
Religionen forefinder stedse en vis social Idealitet og optager
den i sig; deri ligger, at Moralen er uafhængig af Religionen.
Mange Overgangsformer ere her mulige; men tilsidst gælder
det et Enten-Eller. Kunstig Opretholdelse af Religionen fører
kun til, at Folket staar uforberedt, naar dens Tid for Alvor er
forbi. ---

Det er en vel elementær Religionspsykologi, Ardigò bygger
paa, og en omfattende Drøftelse af det religiøse Problem har
han ikke anstillet. Hans Standpunkt er væsentligt bestemt ved
hans intellektuelle Interesse, og vel ogsaa ved en Kontrastvirkning,
som den skarpe Modsætning mellem hans tidligere og
hans senere Livsanskuelse naturligt maatte medføre. Forundring
maa det dog her vække, hvorledes Religionen efter hans Omslag
kunde staa for ham som en »poetisk Erindring«, naar han
anser Frygten, som Kærnen i al Religion --- fra de laveste til
de højeste Former. Den energiske Tænker har sikkert her ikke
benyttet sine egne virkelige Erfaringer.

% --- p. 43 ---
\chapterheadb{III.}{FRANCIS BRADLEY.}

\opage{43}1. Agter man paa Aand og Retning i den seneste Tids
engelske Filosoferen, kan man ikke undgaa at komme til det
Resultat, at den klassiske engelske Skole --- den, der begynder
med Locke, og hvis seneste Repræsentant Spencer er, --- har
ophørt at existere. Den kan siges at have fuldført sin Mission,
som bestod i at hævde Erfaringens Ret, at fordre Problemerne
stillede paa Virkelighedens Grund, og at forberede praktisk,
reformatorisk Stræben. Den begyndte hos Locke med at lægge
Individets Erfaring til Grund, og den ender hos Spencer med
at lægge Slægtens Erfaring til Grund. Hvad der her er indskærpet,
er blevet overleveret til andre filosofiske Retninger og
har virket til at berigtige og omforme disse. Det er ikke længere
den enkelte Retnings Særkende\textsuperscript{19}). Men den engelske
Skole havde paa sin Side visse Mangler, der traadte grelt frem,
jo mere den kom i Vexelvirkning med andre Retninger. Herhen
hører den mekaniske atomistiske Opfattelse, som den fra
Naturvidenskaben overførte til Aandsvidenskaben, og som førte
den til at betragte Sjælelivet som Produkt af selvstændige psykiske
Elementer og Samfundet som ydre Forbund mellem selvstændige
Individer. I Modsætning hertil er i den seneste Tid,
dels som Reaktion, dels ved Indflydelse fra tysk Tænkning,
Helhedsproblemet kommet stærkere frem. Der spørges om den
sidste Sammenhæng mellem Elementerne, der maa forudsættes,
% --- p. 44 ---
% running head: T. H. Green. F. H. Bradley.
\opage{44}forat de kunne komme til at danne en Helhed. Der var vel
ogsaa tidligere i det 19. Aarhundrede en Bevægelse, der gik
imod den »insulære« Filosofi; Navne som Coleridge, Carlyle og
Hamilton vidne derom. Men et metodisk Studium af den tyske
Tænkning bliver først i den sidste Fjerdedel af det 19. Aarhundrede
for Alvor foretaget. Det er her Oxforderuniversitetet,
som gaar i Spidsen. Her havde der stedse tidligere været Modstand
mod den nationale engelske Filosofi (mod Hobbes, mod
Locke, mod Hume); dog havde det mere været af teologiske
end af filosofiske Grunde; det er først nu, at en virkeligt filosofisk
Modsætning gør sig gældende.

Som den, der først vakte en større aandelig Bevægelse af
denne Art maa nævnes \textsc{Thomas Hill Green} (f. 1836, fellow
of Balliol College 1860, Professor i Moralfilosofi 1878, død
1882). Stærkt paavirket af Wordsworth og Carlyle udviklede
han en ejendommelig religiøs Idealisme, hvis filosofiske Grundlag
skyldtes Studiet af Kant og Hegel. Han havde en overordentligt
stor Indflydelse paa den studerende Ungdom i Oxford,
især ved sin Personligheds Idealitet og Begejstring. Han
virkede ikke blot i videnskabelig, men ogsaa i socialreformatorisk
og i fri religiøs Retning. Næsten hele den yngre Generation
af sociale Tænkere i Oxford skylder ham sin Vækkelse.
Hans Hovedværker ere \textit{Introduction to Hume} (1874). en skarpsindig
Kritik af den ældre engelske Skoles Principer, og \textit{Prolegomena to Ethics} (1883), en ejendommelig Kombination af Erkendelsesteori
og Etik, hvorved der tilstræbes en ikke-empirisk
Begrundelse af det Etiske.

2. Englands betydeligste Tænker i den nyeste Tid er utvivlsomt
\textsc{Francis Herbert Bradley}, der er født 1846 og lever
som fellow i Merton College i Oxford. Under sin Studiegang
paavirkedes han af Green og af Hegels og Lotzes Skrifter. Hans
Hang til Ensomhed og hans Tilbageholdenhed forstærkedes ved
Svagelighed. Han er en spinozistisk Natur; kun mangler han
Spinozas realistiske Blik paa psykologiske og sociale Fænomener.
Med stor Tankeenergi fordyber han sig i en eneste Ide, der
atter og atter sætter hans Reflexion i Bevægelse og snart mere
nærmer ham til Skepticismen, snart til Mystiken. Selv skarpe
% --- p. 45 ---
% running head: Kritik af den atomistiske Psykologi .
\opage{45}Kritikere anerkende den Vækkelse og den Tanketugt, som især
Studiet af hans Hovedværk bevirker. Dog bliver hans Skarpsindighed
i det Enkelte ofte Spidsfindighed, og ved hans specielle
Analyser føles ofte Manglen af empirisk Materiale. I
energisk og utrættelig Reflexion staar han maaske øverst blandt
vor Tids Tænkere, i hvert Tilfælde, hvad de yderste Grænseproblemer
angaar. Han danner et Modstykke til Wundt. Medens
denne forsigtigt nærmer sig Tankegrænserne, men slutter noget
hastigt og dogmatisk af paa Grænsen, iler Bradley for hastigt
til den yderste Grænse, men bevæger sig saa der med aarvaagen
Kritik og drøfter Problemerne fra alle Sider.

Det første betydelige Skrift af Bradley er hans \textit{Ethical
Studies} (1876). Han optræder her mod den engelske Psykologis
Atomisme. Bevidstheden, hævder han, kan ikke beskrives som
en blot Samling (collection) af Elementer; ti det vilde være
umuligt at forstaa, hvorledes en saadan Samling kunde blive
sig selv bevidst (aware of itself). Dette Problem havde (hvad
Bradley ikke har lagt Mærke til) Stuart Mill selv opstillet i en
af de senere Udgaver af »Examination of Sir William Hamiltons
philosophy«. Det førte Bradley til Studiet af den tyske
Filosofi, som jo særligt har betonet Bevidsthedens Enhed og
indre Sammenhæng. Han dømmer strengt om den engelske
Tankegang, idet han dadler den for Ensidighed og Dogmatisme
og for Mangel paa Dristighed, især hvad det religiøse Problem
angaar. »Vi bo paa en Ø, og vor nationale Tænkemaade vil,
hvis vi ikke udvide den, faa en øagtig (insular) Karakter«.
--- Denne Dom er dog selv ensidig; Bradley glemmer, at den kritiske
Filosofi egentligt allerede er grundlagt af Locke, og at
den engelske Erfaringsfilosofi med sin Fordring om Paavisning
af Forestillingernes Udspring har været et mægtigt Vaaben mod
Dogmatismen.

Menneskets etiske Stræben gaar ifølge Bradley ud paa at
virkeliggøre dets Selv, udvikle sig til en Helhed, et harmonisk
og afsluttet Hele. Ti det er Selvets Væsen at være Totalitet,
--- paa én Gang afsluttet og rigt. Paa lignende Maade gaar
den teoretiske Stræben ud paa at tænke Tilværelsen som et
sammenhængende og afsluttet Hele (a consistent whole). Naar
% --- p. 46 ---
% running head: Bradleys Grundtanke.
\opage{46}vi ikke kunne blive et Hele selv, maa vi gøre os til Del af et
mere omfattende Hele, ligesom vi, hvor en mindre Helhed frembyder
Modsigelse, maa forme Tanken om en større Helhed. Der
er altsaa Overensstemmelse mellem vor praktiske og vor teoretiske
Natur. Og vi have i vor egen Natur en Maalestok for,
hvad der skal kaldes højere og lavere: det beror paa Graden
af Selvvirkeliggørelse, altsaa paa den Harmoni og den Selvstændighed,
der er i vor Tanke og i vort Liv. Smerter ved
Modsigelse kunde Mennesket ikke føle, hvis det ikke selv var
et Hele og anede, at det var det. Modsigelse fremkaldes dels
ved Mangel paa Harmoni i det Indre, dels ved Mangel paa
Overensstemmelse med de ydre Forhold.

Den Maalestok for teoretisk og praktisk Fuldkommenhed,
som Bradley her havde opstillet, giver Grundtanken i hele hans
Filosofi. Den indeholder Spiren til en hel Livs- og Verdensanskuelse,
saaledes som han senere (1893) udviklede den i sit
Hovedværk \textit{Appearance and Reality,} efterat han i sine \textit{Principles
of Logic} (1883) havde givet en Fremstilling af Erkendelsens
Principer\textsuperscript{20}).

Skriftet om »Fænomen og Realitet« betegner Bradley som
en metafysisk Afhandling, og han forklarer dette Udtryk ved,
at der tilsigtes en Undersøgelse af de Principer, der ligge til
Grund ved Anerkendelse af Realitet i Modsætning til Fænomen
eller Tilsyneladelse (appearance), eller den Maalestok, ved hvis
Hjælp vi kunne skelne mellem højere og lavere Grader af Realitet.
Et System vil han ikke give.

Bradley tænker sig forskellige Indvendinger mod det Forsøg,
han her vil gøre. --- Man anser maaske Forsøget for haabløst.
Men paa Forhaand kan dette ikke vides --- uden netop
i Kraft af en metafysisk Indsigt, d. v. s. en Indsigt i, hvilken
Maalestok der maa anvendes for at kunne skelne mellem Realitet
og Tilsyneladelse. --- Eller man mener, at Resultatet vil
være værdiløst. Men selv om Resultatet er ufuldkomment, vil
det have sin Værdi, naar det tjener til Belysning af, hvad Realitet
er. Ja, selv om vi endte i fuldstændig Skepticisme, saa vilde
derved tilvejebringes en gavnlig Modvægt mod de dogmatiske
Retninger --- den teologiske Ortodoxi paa den ene Side, den
% --- p. 47 ---
% running head: Det filosofiske Motiv .
\opage{47}vulgære Materialisme paa den anden Side, der ellers dele Sindene
mellem sig. --- Efterat have fremført disse Bemærkninger
føjer Bradley til: »Der er endnu en Grund, som for mit
eget Vedkommende maaske har den største Vægt. Jeg tror,
at vi alle mere eller mindre føle os dragne ud over de sædvanlige
Faktas Omraade. Hver paa sin Maade mene vi at staa
i Berøring og i Samfund med Noget, der ligger ud over den
synlige Verden. Paa forskellig Vis finde vi noget Højere, der
baade bærer os og ydmyger os, baade revser og begejstrer os.
Og for nogle Naturer er den intellektuelle Stræben efter Forstaaelse
af Tilværelsen den vigtigste Vej, ad hvilken de erfare
Guddommen. Ingen, der ikke har følt dette, hvor forskelligt
han end udtrykker det, har bekymret sig meget om Metafysik.
Men hvor det føles stærkt, er det sin egen Retfærdiggørelse.«
Til nærmere Forstaaelse af denne Ytring, som aabner et Indblik
i hans indre Liv, og for at undgaa Misforstaaelse, tilføjer
Bradley: »Jeg har været nødt til at tale om Filosofien som en
Tilfredsstillelse af, hvad man kunde kalde den mystiske Side
af vor Natur, en Tilfredsstillelse, der af nogle Individualiteter
ikke kan naas paa anden Maade. Det kunde maaske se ud,
som om jeg mener, at Metafysikerne indvies i Noget, der er
langt højere, end hvad den store Hob kan besidde. En saadan
Lære vilde dog hvile paa en beklagelig Vildfarelse, nemlig paa
den Overtro, at den blotte Forstand er den højeste Side af vor
Natur, og paa den falske Forestilling, at det intellektuelle Arbejde,
der anvendes paa højere Genstande, derfor er højere Arbejde
... Intetsomhelst Kald, ingensomhelst Bestræbelse er en
privat Vej til Guddommen, og den Vej, der fører gennem Spekulation
over de sidste Sandheder, staar ikke over andre Veje. Der
er ingen Synd, som Filosofien saa lidet kan retfærdiggøre, som
aandeligt Hovmod, hvor tilbøjelig Filosofen end maatte være
til den.«

Ved disse Udtalelser har Bradley ikke blot givet Bidrag
til sin egen Karakteristik, men ogsaa til Forstaaelse af den inderste
Kærne i al filosofisk Stræben, naar denne er i den
store
Stil. Den inderlige Trang til Forstaaelse og den stadige Erfaring
om Tankens Grænse, --- den lidenskabelige Anspændelse
% --- p. 48 ---
% running head: Naturvidenskab og Metafysik.
\opage{48}af Tankeenergi og den stadige Følelse af, at det ikke er os, der
tænker, men at »det tænker i os«, --- den intensive Hengivelse
i Tankelivet, som om det var det eneste værdifulde Liv, og den
stadige Erfaring om, at der ogsaa er anden Livstrang i os, der
kræver Tilfredsstillelse, --- alt dette kan læses ud af den engelske
Tænkers Ord. Disse Ord ere værd at mærkes overfor den
haandværksmæssige Karakter, det videnskabelige Arbejde paa
mange Omraader synes i Færd med at antage.

Efter at have betragtet Bradleys almindelige Forhold til
Filosofien, gaar jeg nu over til at fremdrage Grundtankerne i
hans Hovedværk.

3. Det gælder en kritisk Undersøgelse af de Forestillinger,
ved hvis Hjælp man har søgt at forstaa Tilværelsen.

Saadanne Begreber som Materie, Rum, Tid, Energi, med
hvilke Naturvidenskaben arbejder, egne sig fortræffeligt til at
bestemme begrænsede Fænomeners indbyrdes Forhold, men føre
til Modsigelser, naar de skulle udtrykke Tilværelsens sande
Væsen. De ere Forholdsbegreber, som karakterisere Tingene i
Sammenligning og Sammenhæng med hverandre, men ikke sige
Noget om selve de enkelte Ting, der staa i disse Forhold. De
føre ind i uendelige Rækker, idet der stedse kan spørges om
Forholdet mellem selve Leddene og de Forhold, hvori de staa,
og idet Undersøgelsen af selve Ledelsen fører til at konstatere
Forhold indenfor ethvert af dem. Saadanne Begreber ere intellektuelle
Konstruktioner, som kunne være hensigtsmæssige
og nødvendige indenfor den enkelte Videnskab, men ikke oplyse
Tilværelsens inderste Væsen. De ere Arbejdstanker (working
ideas), hvis Betydning er teknisk, ikke rent teoretisk. Derfor
kan Naturvidenskab ikke være Metafysik; men derfor kan
der heller ikke være nogen Strid mellem Naturvidenskab og
Metafysik. Grunder man en Metafysik udelukkende paa de
naturvidenskabelige Grundbegreber, faar den Materialismens
Præg.

Hvis man nu tager sin Tilflugt til Aandsvidenskabens Grundbegreb,
Begrebet Sjæl, har man den Fordel, at vi paa dette Omraade
finde en inderligere Sammenhæng af Enhed og Mangfoldighed
end paa det Omraade, hvor de naturvidenskabelige
% --- p. 49 ---
% running head: Begrebet Erfaring.
\opage{49}Grundbegreber finde Anvendelse. Enhed og Mangfoldighed
staa i Sjælelivet ikke i det udvortes Forhold til hinanden som
i den fysiske Natur. Derfor er psykologisk Erfaring den højeste
Erfaring, vi have. Men heller ikke den egner sig til at udtrykke
absolut Realitet. Indenfor Selvet viser Erfaring og Analyse
os Modsætninger, Forskelle og Relationer, og Selvet som
Helhed undergaar Forandringer. Det er derfor umuligt at sammenfatte
dets Væsen i en eneste Anskuelse, at danne et afsluttet
Begreb om det. Derfor kunne vi ligesaa lidt bygge en
Metafysik paa Psykologien alene, som paa Naturvidenskaben
alene. Idealismen kan ligesaa lidt som Materialismen udtrykke
den fulde Sandhed. Psykologien er en speciel Videnskab, og
enhver speciel Videnskab finder kun Halvsandheder og opererer
med hensigtsmæssige Fiktioner. Begrebet Sjæl er en Abstraktion,
ligesaa vel som Begrebet Legeme. Realiteten kan
hverken være »Sjæl« eller »Legeme«. Hvad vi have givet, er
stedse kun Begivenheder, der for Betragtningen frembyder to
Sider. ---

Men Forudsætningen for hele denne Tankegang, der fører
til at afvise baade Idealismen og Materialismen, er dog, at vi
have en Maalestok (standard) for, hvad Realitet og Sandhed er.
En saadan Maalestok ligger i \textit{Begrebet Erfaring.} I dette Begreb
ligger to Ting i nøje indbyrdes Forbindelse: et Omfang
og en Harmoni. Erfaring forudsætter en given Mangfoldighed
og en indre Overensstemmelse, et konsekvent og harmonisk
Forhold mellem de mangfoldige Elementer. Den fuldkomne Erfaring
vilde bestaa i et altomfattende Indhold forbundet med
fuld Konsekvens og Harmoni til en Helhed. Enhver Erfaring,
vi virkeligt have, indeholder en Tilnærmelse i Retning af dette
Ideal.

Maalestokken for Realitet og for Værdi er den samme.
I vort praktiske Ideal maa enhver Side i vor Natur være tilfredsstillet
i Harmoni med de andre Sider. Enhver utilfredsstillet
Drift er en ufuldbyrdet Tanke; i enhver Smerte ytrer
der sig en Disharmoni og ligger der en Spore til at arbejde sig
ud af denne Disharmoni.

Denne Maalestok viser os, hvorvidt vi staa i Forbindelse
% --- p. 50 ---
% running head: Modsætning mellem Omfang og Harmoni.
\opage{50}med Realiteten. Vi kunne ikke danne os nogen Forestilling om
det, der fuldstændigt tilfredsstiller Maalestokken. »We cannot
construe the one absorbing experience to ourselves«. Det Højeste
maa være altomfattende og absolut samstemmende. Men for
os er der en stadig Modsætning mellem Omfang (extension) og
Samstemmen (harmony) eller Selvoverensstemmelse (self-consistency).
Og Grunden er, at vort Omfang er for ringe og vor
Harmoni for ufuldstændig. Indre Disharmonier skyldes Begrænsning
og kunne kun ophæves gennem et mere omfattende
Indhold, der ophæver Afhængigheden af ydre Forhold og derved
gør det muligt at opnaa indre Sammenhæng. Det er Forstyrrelserne
udefra, der volde den indre Uorden. Derfor kan
kun et uendeligt Væsen være fuldstændigt harmonisk. De to
Bestanddele i vor Maalestok for Realitet (og for Værdi) hænge
saaledes nøje sammen. Paa det teoretiske Omraade viser det
sig deri, at vi kun kunne ophæve de Modsigelser, Erfaringen
frembyder, ved at udvide vor Erfaring. Tidsforholdet staar især
som en Hindring for at opnaa harmonisk Totalitet. Det staar
altsaa i Strid med Kriteriet for Realitet, og derfor kan det kun
have fænomenal Betydning. Om Fremskridt og Tilbagegang kan
der kun være Tale ved begrænsede og ufuldkomne Realiteter;
det Absolute kan ingen Historie have, skønt det indeholder historiske
Processer uden Tal; det har ingen Aarstider. Intet
Fuldkomment, intet sandt Realt forandres.

Vor Tanke stræber derfor stedse efter Noget, som er Mere
end Tanke, vor Personlighed efter Noget, som er Mere end
Personlighed, vor Moralitet efter Noget, som er højere end al
Moral. Der er ingen Modsigelse i, at et Væsen attraar en Fuldendelse,
i hvilken det selv taber sig. Floden strømmer til Havet,
og Selvet taber sig selv i Kærligheden. Det Højere maa stedse
være mere omfattende end det Lavere, maa indbefatte eller
være dette og Mere til. ---

Filosofien fører efter Bradley til en sund Skepticisme, for
hvilken Videnskaben er en ringe Ting i Forhold til Tilværelsens
Rigdom. Selv den højeste Sandhed, vi kunne finde, betinges
ved Noget, vi ikke kende. Vi vide ikke, hvilke andre Arter af
Erfaring der kan existere foruden vor egen. Hvad vi kunne
% --- p. 51 ---
% running head: Det Endeliges Lod.
\opage{51}naa, er kun en Begrundelse af vor Ret til at kalde Noget højere
eller lavere end noget Andet.

Det gaar overfor de sidste Gaader Religionen, som det gaar
Filosofien. Ogsaa den er nødt til at udtrykke det Højeste ved
Forestillinger, der ere hentede fra Erfaringen. Og den anstiller
end ikke, saaledes som Filosofien gør det, en nøjagtig Undersøgelse
af saadanne Forestillingers Væsen og Gyldighed. Fra
denne Side set, som Erkendelse, staar Filosofien over Religionen.
Men i en anden Henseende --- som Stræben efter at lade Anerkendelsen
af det Godes Realitet gennemtrænge enhver Side
af vort Væsen (the attempt to express the complete reality of
goodness through every aspect of our being) --- staar Religionen
højest. ---

Er dette Resultat nu ikke utilfredsstillende? Paa dette
Spørgsmaal svarer Bradley med det Modspørgsmaal: Hvem siger,
at vi skulle have fuld Tilfredsstillelse af al vor Trang? Ufuldstændighed,
Uro og utilfredsstillet Idealitet er nu engang det
Endeliges Lod! Og han føjer til, at ethvert Forsøg paa at forestille
sig en fuld Tilfredsstillelse stedse foretager et Udvalg af
vore Fornødenheder, et Udvalg, hvis Berettigelse ikke kan godtgøres.


4. Bradley's Filosofi er en Gennemtænken af det Problem,
Wundt og Ardigò nærmede sig til, men ikke tog op til
udtrykkelig Behandling, --- det, som Kant havde stillet ved sin
Lære om »Ideerne« som Totalitetsbegreber, og som har sit Udspring
i selve Tankens Væsen som sammenfattende Virken (Syntese).
Selv mener Bradley at staa i stor Gæld til Hegel, og han
er, som flere engelske Filosofer i den nyeste Tid, imponeret af
Hegels dristige Dialektik. Men i Virkeligheden vilde han, hvis
han endeligt skulde klassificeres, langt snarere være at kalde
Kantianer end Hegelianer. Paa enkelte Punkter minder han
om William Hamilton. Med fast Haand drager han paa selve
Tankens Grænse Konsekvensen af Tankens egne Love. Det er
især to Love, han støtter sig til: Forholdsloven (the relational
way of thought) og Virkelighedskriteriet (complete experience
som standard of reality). De føre ham begge til samme Resultat,
--- til Umuligheden af at stanse paa noget Punkt, skønt vi
% --- p. 52 ---
% running head: Bradleys Skepticisme.
\opage{52}Skridt for Skridt kunne naa frem til nærmere Bestemmelse af
Realiteten.

En vis Skepsis karakteriserer Bradley, idet han lægger
mere Vægt paa Umuligheden af en Afslutning end paa de nærmere
Bestemmelser af Realiteten, som ere mulige. Derved kommer
han til at tillægge den specielle og empiriske Viden for
ringe positiv Betydning for Verdensanskuelsen. De Synspunkter,
med hvilke der arbejdes paa de specielle Erfaringsomraader,
affærdiger han som »unyttige Fiktioner«, »blotte praktiske Kompromiser«,
uden at tage Hensyn til, at de dog ingen »Nytte«
vilde være til, hvis de ikke i nogen Forstand førte os nærmere
til Realiteten. Han glemmer Goethes Ord, at Naturen hverken
har Kærne eller Skal.

Han er ogsaa Skeptiker i den Forstand, at der er flere
enkelte Problemer, han erklærer for uløselige. Vi kunne ikke
udlede Tilværelsens Mangfoldighed, de mange endelige Centra,
udfra hvilke der gøres Erfaringer, alt det Fragmentariske, ved
hvilket vi maa stanse, af ét Princip, selv om Mangfoldighedens
Faktum ikke strider mod, at der gives en omfattende Enhed.
Hvorfor der skulde være »Tilsyneladelser« (appearances) og
ikke blot »Realitet«, kunne vi ikke sige, men det hindrer os ikke
i at fastholde Realitetsbegrebet. Særligt er det Uorganiske en
Grænse for vor Erkendelse. Det Organiske tilfredsstiller med
stor Tilnærmelse vort Begreb om sand Realitet som Harmoni
og Sluttethed, skønt der ogsaa paa det organiske Omraade findes
Delthed og Disharmoni; men det Uorganiske staar dette Begreb
om Realitet langt fjernere. Bradley er tilbøjelig til at
mene, at Antagelsen af et absolut Uorganisk i Naturen kun
skyldes vor Uvidenhed. Han indrømmer, at den specielle Videnskab
af praktiske Hensyn maa skelne mellem Organisk og
Uorganisk, men tør ikke tillægge denne Distinktion absolut Betydning.
Her fremtræder tydeligt hans Afstand fra Hegel; ti
den romantiske Naturfilosofi arbejdede dristigt med Tanken om
hele Naturen som en stor Organisme. --- Ogsaa Problemet om
Forholdet mellem Sjæl og Legeme anser han for uløseligt. Det
er umuligt at danne et Begreb om, hvorledes disse to Existensformer
forholde sig til hinanden. Men heri finder han en Be%
% --- p. 53 ---
% running head: Bradleys Mysticisme.
\opage{53}kræftelse paa sin almindelige Teori om Realiteten; ti Vanskeligheden
stammer fra, at vi ved at stille Sjæl og Legeme i Modsætning
til hinanden, gøre to Abstraktioner til virkelige Væsener
og altsaa behandle »appearances«, som om de vare Realiteter.


Dog vilde Skepticisme ikke være det rette Ord om Bradleys
Standpunkt. Han stanser ikke ved et Svælg mellem Arbejdet
og Maalet, Fænomenet og Realiteten. Paa ethvert Trin
er det Højeste til Stede, og ethvert Trin har sin Sandhed; der
gives mange Grader og Stadier, men de ere alle uundværlige.
Vi kunne ikke finde nogen Provins i Verden saa ringe, at det
Absolute ikke skulde ibo det. --- Snarest bliver han at kalde
Mystiker, og han er det afgjort, naar hans Tanke hviler, og
naar han polemiserer mod Tidsbegrebet, og mod Aktivitetens
Betydning. Han gaar her over til en ubevægelig Skuen, et
hvilende Blik, en Betragten sub specie æterni. Det er gaaet
ham her, som det gik Spinoza. Ti Spinozas »Substans« er jo
egentligt den som et fuldendt Væsen anskuede Maalestok for
Realitet, Virkelighedskriteriet som existerende Ideal.

Jeg kan ikke se, at vi her kunne komme ud over den
stadige Mulighed af nye Processer, ny Aktivitet. Og hvorfor
skulde det ikke være det Højeste at udvikle sig igennem Tiden?
Tilmed, da det stedse vil vise sig umuligt at eliminere Tidsbegrebet
af en Virkelighedserkendelse. Af den stadige Inkongruens
mellem Tanken og Realiteten og af den stadige Nødvendighed
af nyt Tankearbejde --- for den Enkelte og for Slægten
--- kunde man maaske være berettiget til at drage den
»metafysiske« Konsekvens, at Tilværelsen selv ikke er afsluttet,
ikke er færdig. Og den \textit{kan} ikke være færdig, naar Tanken,
der jo dog er en Del af Tilværelsen, ikke er færdig. Denne
Opfattelse fører til et inderligere Forhold mellem Tanken og Realiteten
end det, som Bradley havner i. Ti naar Tid og Aktivitet
ikke ere Former, der blot høre Fænomenerne, »Tilsyneladelsen «
til, saa vil Tanken netop som stræbende og arbejdende kunne
føle sig Et med det Inderste i Tilværelsen\textsuperscript{21}). --- Men jeg vil
ikke her gaa nærmere ind paa, hvad jeg har udviklet i det
% --- p. 54 ---
% running head: Bradleys engelske Kritikere.
\opage{54}erkendelsesteoretiske Parti af min Religionsfilosofi og i mine
»Filosofiske Problemer«.

I England selv er Bradley dels bleven kritiseret fra et
empirisk og kritisk Synspunkt (saaledes af \textsc{James Ward} i »Mind«
1894), dels fra et Standpunkt, der betoner vore Grundbegrebers
økonomiske og tekniske Karakter og grunder en idealistisk
Verdensanskuelse paa personlig Trang (saaledes i det af otte
Oxforder Akademikere udgivne Skrift \textit{Personal Idealism.} London
1902). Denne sidste Retning mener at være en Fortsættelse
af den Tankeretning, der i de sidste tredive Aar har
raadet i Oxford. Den har mest Ret i sin Kritik, hvor den bestrider
Bradleys negative Holdning overfor Sandheder, der gælde
indenfor bestemte Betingelser, og hævder, at vi kunne have
gyldig Erkendelse, selv om vi ikke kunne naa en udtømmende
Forstaaelse af Tilværelsen. Derimod forekommer det mig ikke,
at den engelske Kritik tilstrækkeligt har anerkendt den virkelige
Dybsindighed og Tankeenergi i Bradley's Værk.

% --- p. 55 ---
\chapterheadb{IV.}{ALFRED FOUILLÉE OG DEN FRANSKE\\ FILOSOFI I VORE DAGE.}

\opage{55}1. I den franske Filosofis Historie staar Auguste Comte
(† 1857) som den største Skikkelse i det 19. Aarhundrede. Den
Indflydelse, han har havt videnom, fik Udtryk ved Afsløringen
af hans Statue i Maj 1902. Positivismen er den mest ejendommelige
og betydningsfulde Tankeretning, som Frankrig i
det sidste Hundrede Aar har frembragt. Den havde i den
største Del af Aarhundredet til Modstander en populær Spiritualisme,
der repræsenteredes af Victor Cousin og hans Disciple
og begunstigedes af Regeringerne. Dog øvede omtrent fra Midten
af Aarhundredet Taine og Renan en stor Indflydelse paa
den yngre Slægt, ikke blot i filosofisk, men ogsaa, og nok saa
meget i literær Henseende. Ved Siden af dem, men som en
mere ensom Skikkelse staar Charles Renouvier, Repræsentant
for en streng kritisk Filosofi. --- Der er paa dette Punkt et
Hul i min Fremstilling af den nyere Filosofis Historie, idet jeg
for Frankrigs Vedkommende er stanset ved Comte. Dette Hul
ser jeg mig heller ikke nu i Stand til at udfylde. Jeg kan kun
give en kort Karakteristik af Taine og Renan som Indledning
til den franske Form for Udviklingsfilosofi, der repræsenteres
af Fouillée, og senere en lignende Karakteristik af Renouvier i
Sammenhæng med den Rolle, Diskontinuitetsprincipet kommer
% --- p. 56 ---
% running head: Hippolyte Taine .
\opage{56}til at spille i den nyeste franske Tænkning. En Monografi
(eller flere) om disse tre Mænd som Filosofer vilde være af
ikke ringe Interesse til Belysning af den aandelige Udviklings
Gang fra det 19. Aarhundredes Midte af. Det danske Videnskabernes
Selskab har gjort Forsøg paa at fremkalde en Monografi
over Renouvier, hvis omfattende Forfattervirksomhed frembyder
særegne Vanskeligheder for Overblik og Karakteristik;
men der kom ingen Besvarelser af den stillede Opgave.

\textsc{Hippolyte Taine} (1828---1893) er mest bekendt som
Literatur- og Kunsthistoriker og som Kritiker i »realistisk« Aand.
Vor Literatur besidder en værdifuld Karakteristik af ham fra
dette Synspunkt i \textsc{Georg Brandes'} Afhandling \textit{Den franske
Æstetik i vore Dage} (1870). Taine stræbte som Kritiker først
og fremmest at forstaa Kunstværkerne og deres Frembringere
ved at efterforske de ydre Forhold (le milieu), under hvilke de
udviklede sig, den Situation (le moment), i hvilken de arbejdede,
den Race, af hvilken de vare fremgaaede, den herskende Evne,
der ytrede sig i deres Værk. Dette sidste Element var væsentligt
betinget ved de første tre Elementer.

Fra sin første Ungdom arbejdede Taine under Tryk og
Modstand, men med usvækket Energi og Begejstring. Nyligt
udgivne Breve fra hans unge Dage give et smukt Billede af
ham fra denne Tid. Senere fik han som Professor i Kunsthistorie
ved École des beaux-arts Lejlighed til at virke paa
større Kredse og kom til at arbejde under gunstigere Forhold,
en Lejlighed, han benyttede med den Jernflid, som var hans
Særkende hele hans Liv igennem. Karakteristisk for Taine som
Taler, Forfatter og Forsker var især hans Evne til at frembringe
et Helhedsbillede gennem Ophoben af enkelte Træk.
Paa hans Forelæsninger kunde man --- efter min personlige
Erfaring --- til en Begyndelse føle sig frastødt ved den tørre,
objektive Beskrivelse, der gik fra Del til Del, fra Egenskab til
Egenskab. Men der var en saadan Rigdom i det Stof, der stod
til hans Raadighed, og en saadan Energi i den Maade, hvorpaa
det blev ført frem, at Totalbilledet inden Timens Slutning stod
klart og levende for Fantasien. Jeg erindrer saaledes en Forelæsning
over græsk Skulptur. Den var ikke ledsaget af Bille%
% --- p. 57 ---
% running head: Taines Psykologi.
\opage{57}der, men Taine sad under den lange Beskrivelse, han gav, og
tegnede i Luften, og tilsidst var det, som om der stod en
Billedstøtte foran ham paa Bordet. Hans Fremstillingskunst gik
i modsat Retning af Julius Langes, hvis Styrke bestod i at
analysere et Helhedsbillede, som fremkaldtes ved Afbildning
eller ved strax at vække Tilhørernes Fantasi. Taine gik fra Del
til Helhed, Lange fra Helhed til Del.

Et rent filosofisk Værk ydede Taine i \textit{De l'intelligence}
(1870), en Psykologi i den engelske Skoles Aand. Den røber
stærk Indflydelse af St. Mill, Bain og Spencer. Af særlig Interesse
i denne glimrende skrevne Bog er for det Første Forklaringen
af Erkendelsens Udvikling gennem en Kamp for Tilværelsen
mellem de enkelte psykiske Elementer. Der opstaar
stadigt Elementer med Virkelighedspræg, en Slags normale
Hallucinationer; de brydes med hverandre, og det sejrende er
en Fornemmelse, en Sansning, som derfor kan defineres som
une hallucination vraie. Under denne Brydning medvirker dernæst
Erkendelseselementernes oprindelige Bevægelsestendenser.
Enhver Fornemmelse eller Forestilling er forbunden fra først
af med en Tilskyndelse til Bevægelse. Efterhaanden afslibes
denne motoriske Tendens, og først da bliver rent teoretisk Sansning
og Forestilling til. Endeligt er det ogsaa karakteristisk
for Taines Psykologi, at han i saa stor Udstrækning gør Brug
af sygelige Tilstande til at oplyse Bevidsthedslivets Natur.
--- Dette sidste Træk er siden vedblevet at være karakteristisk for
den franske Psykologi, maaske fordi den franske Nervepatologi,
især hos \textsc{Charcot} og hans Skole, ydede saa stort Materiale. \textsc{Ribot},
hvis første Skrift, \textit{La psychologie anglaise contemporaine,} udkom
1870, er i sine senere Skrifter især gaaet denne Vej; ligeledes
\textsc{Alfred Binet} og \textsc{Pierre Janet}. Den moderne franske
Psykologi har herved faaet sit Særpræg, ligesom den engelske
Psykologi (hos Ward, James og Stout) har sit i Analysen og
den tyske (hos Fechners og Wundts Disciple) sit i Experimentet.


Tidligere (især i Skriftet \textit{Les philosophes français du 19e
siècle)} havde Taine skarpt kritiseret den herskende Spiritualisme
(Cousins Skole) og vist hen til Comte. Positivist var ban
% --- p. 58 ---
% running head: Taines Estetik og Historieforskning.
\opage{58}dog kun i dette Ords videre Betydning. Positivismen var for
ham, ligesom for Ardigò, Udgangspunkt, men ikke Slutning.
Dels fordrede han Analyserne af Begreber og af Tilstande førte
videre, end Positivisterne af den strenge Skole fandt nødvendigt,
dels fastholdt han Problemet om en afsluttende Verdensanskuelse
og troede forsaavidt paa en Metafysik. Han ender
med Henblik paa dette Problem sit Skrift De l'intelligence saaledes:
»Je vois les bornes de mon esprit; je ne vois pas les
bornes de l'esprit humain«. I Taines berømte literatur- og
kunsthistoriske Arbejder spiller Psykologien en stor Rolle, ofte
for stor en Rolle, idet han er tilbøjelig til at gaa deduktivt
frem og udlede et kunstnerisk Fænomen udelukkende af et
enkelt Element (en herskende Evne eller et milieu). Det hindrer
ham undertiden, f. Ex. i hans Behandling af Shakespeare,
i at faa de finere Nuancer og de inderligere Sjæleytringer med,
som kun aabenbare sig ret for længere fortsat Iagttagelse og
taalmodigere Analyse. Disse Arbejder ligge alle forud for 1870.
Det havde nu været hans Hensigt at fortsætte sit rent psykologiske
Arbejde og supplere Bogen om Erkendelsen med en
Bog om Villien. Men hans Fædrelands forfærdelige Ulykker
førte ham ind paa historiske Studier. Han vilde forstaa sit
Lands Skæbne af dets Fortid og gennem Forstaaelsen vinde
nyt Haab. Deraf udsprang det store Værk \textit{Les origines de la
France contemporaine} (1876---1894). Af filosofisk Interesse ere
især det første Bind (L'ancien régime), hvor han dog paa væsentlige
Punkter følger Tocqueville, og de to sidste Bind (Le régime
nouveau), som Ingen, der vil forstaa det 19. Aarhundrede, maa
lade ulæste. Det er vigtige Bidrag til Sociologien. Om hans
Behandling af den franske Revolution (i de mellemliggende
Bind) har man med Rette sagt, at den har grundlagt den videnskabelige
Betragtning af denne Begivenhed, i Modsætning til
den deklamatoriske og agitatoriske. Der er derimod indvendt
imod den, at han for meget holder sig til de almindelige Aarsager
og undervurderer Indflydelsen af den specielle og øjeblikkelige
politiske Situation. Hans Had til Frasen har tillige
ført ham saa vidt, at han miskender den Begejstring, de store
% --- p. 59 ---
% running head: Ernest Renan.
\opage{59}Forhaabninger vakte. Man forstaar efter hans Fremstilling ikke,
at det var den Tid, da Marseillaisen blev til.

Taines Livsanskuelse var bestemt ved hans Temperament.
Han var efter sit eget Udsagn un homme naturellement triste.
Statskupet 1851 og Nederlaget 1870 maatte forøge hans Mismod
overfor de ydre Forhold. Sit bedste Liv levede han i
Arbejdet og i Tankens Verden. Han følte Trang til at fordybe
sig i Tilværelsens store, nødvendige Sammenhæng. Marcus
Aurelius og Spinoza vare hans Yndlingsforfattere. I religiøs
Henseende sympatiserede han med Protestantismen, og efter
sit Besøg i England haabede han paa en fri religiøs Bevægelse.
Men han forblev Stoiker til det Sidste og læste endnu i sine
sidste Dage Marcus Aurelius »som en Slags Liturgi«\textsuperscript{22}).

\textsc{Ernest Renan} (1823---1892) frembyder i sin Udviklingsgang
en usædvanlig lang Række af Stadier, og den biografiske
Kunst, der skal kunne finde den indre Sammenhæng
her, er endnu ikke bleven anvendt. Sin Barndoms- og sin Ungdomstid
har han selv skildret i den af hans Bøger, som staar
højest ved Stemningens Inderlighed og Dybde \textit{(Souvenirs d'enfance
et de jeunesse).} Det var efter hans egen Paastand historisk
Kritik, der førte ham til at forlade Præsteseminariet og
den katolske Tro. Imidlertid udviklede der sig snart hos ham
en begejstret Tro paa Videnskaben og paa den Betydning, dens
Resultater vilde faa for den store Mængde. Denne filosofiske
og demokratiske Optimisme udtalte han i et Skrift om »Videnskabens
Fremtid« \textit{(L'avenir de la science),} som han forfattede
i Aarene 1848 og 1849, men som han først udgav langt senere,
efterat han havde mistet det Haab, der udtaltes deri. Det var
Statskupet og Reaktionen, der indgød ham en Pessimisme, som
han aldrig forvandt. Senere kom Nederlaget og Kommuneopstanden
til: Som Videnskabsmand blev han især berømt ved
sine sprog- og religionshistoriske Værker, som det ikke her er
Stedet at omtale nærmere. Under den tredie Republik fik han
en høj videnskabelig Stilling, som Direktør for Collège de France,
og alt imedens han lagde den sidste Haand paa sine Arbejder,
syslede hans Reflexion paa en let og spøgende Maade med de
store Problemer, der havde optaget ham saa stærkt i hans unge
% --- p. 60 ---
% running head: Den gamle Forskers Tankeleg .
\opage{60}Dage. Man lyttede i store Kredse til den skeptiske Causeur,
om hvem man --- idetmindste paa tredie og fjerde Haand ---
vidste, at han var en fremragende Mand. Han var i sine Udtalelser
og Skrifter fra denne Tid snart stemningsfuld, snart
kaad, --- snart sublim, snart blaseret; og for Mange var denne
stadige Svingen, --- en Svingen, hvis Midtpunkt ikke var let at
opdage, om der var noget, --- et Udtryk for den sande Aandsfrihed
og Aandsrigdom. Tankearbejdet var vel nu for ham selv
snarest en Leg, hvorved han adspredte sig fra sine strengere
Studier. Og han var bleven betænkelig ved enhver bestemt
Stilling overfor Grænsespørgsmaalene, fordi han selv havde indtaget
saa mange forskellige Stillinger. Derved bedrog han sig
selv og Andre for et værdifuldt Resultat af sin lange Stræben
og Forsken. Man har med Rette kunnet sige om ham: Il fut
dupe de la peur d'étre dupe.

Renans Betydning som Tænker i de senere Aar af hans
Liv ligger mere i, at han staar som Symptom paa en vis Retning
i Tiden, end i, at han skulde have øvet et virkeligt Arbejde
i Problemernes Tjeneste. Hemmeligheden var maaske
den, at hans Hjerte hang ved hans Barndoms Religion, den,
han havde levet sig ind i, medens han blev »opdragen af Kvinder
og Præster« (en Opdragelse, i hvilken han selv mente, at
Forklaringen af hans Fortrin og hans Fejl maatte søges). Hvad
der senere hen besjælede ham, var en Efterklang deraf --- at
ligne med Klangen af den i Havet nedsænkede Klokke, der
endnu i stille Vejr skal kunne høres ved Bretagnes Kyst. En
almindelig idealistisk Tro er vel mulig; men den forslaar kun,
saalænge vi endnu leve under de gamle Vaners Herredømme.
Dette vil ikke vare ved. Vi leve --- siger Renan flere Steder
--- af Skyggen af en Skygge; hvad vil man leve af efter os?
Til en fast og mandig Stilling overfor dette Spørgsmaal kom
Renan ikke: det var kun i visse Stemninger, han følte dets
Brod.

Det er to Arbejder af Renan, som især have Interesse for
hans filosofiske Karakteristik i hans senere Aar.

\textit{Dialogues et fragments philosophiques} forfattedes i Foraaret
1871, da Renan under Kommuneopstanden opholdt sig i
% --- p. 61 ---
% running head: Demokrati og Videnskab.
\opage{61}Versailles, men blev først udgiven 1876. Det er Samtaler, der
føres mellem »tre Filosofer af den Skole, hvis Grundprinciper
ere Kultus af Idealet, Negtelse af det Overnaturlige og Udforsken
af Virkeligheden«. De drøfte det Spørgsmaal, hvad det
ideale Maal for Verdensudviklingen er. Ti et Maal maa der
være; Universet er ikke en tom Bølgegang, der ender i Intet.
Maalet for Udviklingen kan ikke søges i den store Masse af
Væsener, som har et Minimum af Nydelse og af Dannelse, men
et Maximum af Arbejde. Maalet maa være Fornuftens Herredømme,
og kun en lille Kreds kan blive delagtig i den. Den
store Masses Betydning kan kun være den, at den er den Jordbund,
af hvilke Genierne voxe frem. Genierne --- og de Intelligente,
som forstaa at vurdere dem --- ere Historiens Endemaal.
Det er beklageligt, at den store Masses Uvidenhed er
Betingelse for, at dette Maal naas; men Naturen bekymrer sig
ikke om Sligt. Det vilde føre til Degeneration, til Sænkning af
Niveauet, om man vilde gaa Demokratiets Veje; ad disse Veje
frembringes ingen Guder. Demokrati og Videnskab ere Antipoder.
Udviklingen betegnes ved, at Instinktet afløses af Reflexionen,
Religionen og Kunsten af Videnskaben.

Renan ser heri Noget, som allerede Middelalderen paa sin
Maade havde Øje for. Han fandt en Sandhed i den middelalderlige
Ordning, at der var Mennesker, der bad for dem, der
ikke havde Tid til at bede. Videnskabsmændene skulle være
den nye Tids Præster. Og det er aabenbart en Efterklang af
Præsten i Renan, der her kommer frem og --- i hvert Tilfælde
for et Øjeblik --- lader ham havne i denne Løsning af Kulturproblemet,
der iøvrigt --- som vi senere skulle se --- omtrent
samtidigt udtaltes af en tysk Forfatter med mere Fynd og
Klem, skønt med større Inkonsekvens. For Renan stod det dog
saaledes, som om de Vidende, der ere Udviklingens Maal, meget
vel kunde elske Folket, saa at deres Herredømme kunde komme
dette til Gode. Nietzsche forsøger paa at negte dette.

Nogle Aar før sin Død gav Renan et Slags filosofisk Testamente
i sin \textit{Examen de conscience philosophique} (1888). (Trykt
i »Feuilles détachées«). En for Renan karakteristisk Udtalelse
forekommer strax i Begyndelsen af denne Afhandling. »Det
% --- p. 62 ---
% running head: Renans filosofiske Testamente.
\opage{62}ærlige Menneskes første Pligt«, hedder det, »er ikke at øve
Indflydelse paa sine egne Meninger, men lade Virkeligheden
genspejle sig i ham som i Fotografens Mørkekammer og staa
som Tilskuer ved de indre Kampe, Tankerne føre i Dybet af
Bevidstheden ... Vi bør være passive overfor de indre Forandringer
paa vor intellektuelle Nethinde.« --- Det er naturligvis
sandt, at naar vi skulle gøre vort intellektuelle Regnskab op,
maa vi ikke gribe ind, men se, hvad der virkeligt er tænkt i
os. Men ved Tilblivelsen af de enkelte Poster paa Regnskabet
medvirker en Stræben efter bestemte Maal, en uvilkaarlig
Higen, ligesom vort Øje i hvert Øjeblik rettes saaledes, at det,
der tiltrækker sig Opmærksomheden, kan opfattes af den centrale
Del af Nethinden. Det vil derfor være en Illusion, naar
man mener at være rent passiv Tilskuer af, hvad der foregaar
paa den aandelige Nethinde. I hver videnskabelig Hypotese
saavel som i hver Livsanskuelse er der Noget af et Vovestykke.
Grundtankernes og Grundværdiernes Kamp følge vi ingenlunde
som uinteresserede Tilskuere. Og prøve vi det, bliver det ingenlunde
til Fordel for Tænkningen.

To Tanker ligge til Grund i Renans filosofiske Testamente.
Den ene gaar ud paa, at vi baade, naar vi gaa i det Store, og
naar vi gaa i det Smaa, støde paa det Uendelige. Saa bestemt
vi end indenfor vor Verden (eller vor Del af Verden) maa
holde os til Erfaringen, saa liden Ret have vi dog til at mene,
at dens Resultater have absolut Gyldighed. Fra Uendelighedens
Synspunkt er Intet umuligt, og den uoverskuelige Fremtid faar
mange Vanskeligheder til at forsvinde. Lad os da Intet benegte,
Intet bekræfte, men lad os haabe! --- Den anden Tanke supplerer
den første. Midt i alt det Gaadefulde og Usikre, vi omgives
af, fremtræde de fire ophøjede Autoriteter, Kærligheden, Religionen,
Poesien og Dyden, som Egoisten negter, og som dog
lede Verden fremad. Gennem dem høre vi Universets eller,
om man vil, Guds Stemme. Det er de himmelske Sfærers Harmoni,
det Uendeliges Toner. I dem ytrer sig den dybe Stræben
(le nisus profond), som arbejder sig frem gennem Verdensudviklingen
under Modstand af det stridige Stof (særligt stridigt
maaske i vor Del af Verden). Den vil vinde frem i Klarhed
% --- p. 63 ---
% running head: Svingen mellem Muligheder som det Højeste !
\opage{63}og i Magt i en Grad, vi ikke kunne gøre os nogen Forestilling
om. ---

Det filosofiske Testamente er mindre legende og skeptisk
end de tidligere Udtalelser. Teorien om de store Mennesker
som Verdensmaalet forekommer ikke i det. Men om han har
taget den tilbage, kan man ikke se. Heller ikke kan man se,
om han endnu fastholder, hvad han faa Aar forud havde udtalt
i Slutningen af sin Afhandling om Amiel (der ogsaa er
trykt i Feuilles détachées). Amiel havde bebrejdet ham den
legende og ironiske Maade, paa hvilken han talte om de alvorligste
Ting. Renan svarer, at Ironien maa være Filosofiens
sidste Ord. Hvad der er Maalet for Verdensudviklingen, vide
vi ikke, end ikke, om der er et Maal. Maaske er det Hele une
mauvaise farce. Og lad os dog saa for enhver Pris ikke blive
narrede! Ved at svinge mellem de to Muligheder --- at Livet
er Alvor, eller at det er en daarlig Spøg --- undgaar man dog
at blive \textit{helt} bedragen, hvorimod man ved definitivt at vælge
den ene eller den anden, udsætter sig for at lide en fuldstændig
Skuffelse. Det er især for Andres Skyld, man maa være
forsigtig. For sin egen Skyld kunde man indlade sig paa store
Vovestykker; men man har ikke Lov at forlede Andre til at
udsætte sig for fuldstændigt Skibbrud. Derfor gælder det at
være ad utrumqve paratus, og det opnaar man ved at svinge
--- svinge mellem Tro og Tvivl, mellem Optimisme og Ironi.

\textsc{Charles Renouvier} har hertil med Rette sagt, at ved
saadanne Svingninger mister man let sin Adkomst i begge Retninger,
og han har henvist til den Mulighed, at \textit{selve Valget}
kunde være det Væsentlige\textsuperscript{23}). Det kunde jo være, at den personlige
Energi, der vælger og fastholder sit Stade og tager
Konsekvenserne deraf, stod den egentlige Realitet i Verden nærmere
end de to Muligheder, Renan skildrer.

En psykologisk Forklaring af Renans Standpunkt i hans
senere Dage faar man maaske ved hans Forhistorie og hans
Race. Kontrasten mellem Barndomstroens Fylde og Kritikens
tilsyneladende Tomhed havde grebet dybt ind i hans Liv, og
det stod for ham, som om det eneste Positive i ham egentligt
kun var en Efterklang, hvis Toner han ikke kunde hindre i at
% --- p. 64 ---
% running head: Alfred Fouillée .
\opage{64}dø hen. Ogsaa for Ardigò stod i hans senere Tid Religionen
som en poetisk Drøm; men han byggede sig med al Tankens
og Villiens Energi et nyt, fast Stade. Renans Katolicisme har
vel ogsaa været af en anden Art end Ardigòs. Som ung Teolog
var han bleven vant til at faa sine Livserfaringer paa
passiv Maade, udefra, ikke gennem energisk indre Arbejde, ikke
gennem aktiv Deltagelse i Livet. Derfor faldt det ham ogsaa
senere lettere at staa som blot Tilskuer ved Tankernes Kamp,
saavel som ved Livets Kamp. Han lod Stemningerne følge efter
hverandre paa kaleidoskopisk Vis; der foregik ingen indre Krystallisationsproces
i ham. Og hans galliske Temperament havde
ikke Evne til at lade de modsatte Stemninger forene sig til en
Helhed med en bestemt Klangfarve.

2. Skulde man nævne en Forsker, som paa det filosofiske
Omraade kan siges at fortsætte Taines Arbejde, maatte først
og fremmest \textsc{Alfred Fouillée} (f. 1838) nævnes. Taines stærke
Fremhæven af den lovbestemte Sammenhæng genfindes hos
Fouillée, der tillige træder i bestemt Modsætning til Renans betydningsfulde
filosofiske Dilettantisme i hans senere Aar. Men desuden
ere Fouillées filosofiske Skrifter Vidnesbyrd om den idealistiske
Retning, den filosofiske Tænken gennemgaaende slog ind
paa henimod Aarhundredets Slutning.

Fouillée har virket som Professor i Filosofi i Bordeaux og
i Paris, men senere af Helbredshensyn trukket sig tilbage til
Mentona.

Hans første Skrifter angaa den græske Filosofis Historie.
Gennem Studiet af Platon har han arbejdet sig frem til sin
filosofiske Opfattelse. \textit{(La philosophie de Platon.} 1869). For
Platon stod Ideverdenen og Erfaringsverdenen i en skarp Modsætning
til hinanden, saaledes at kun hin var den sande Realitet.
Den moderne Naturalisme hævder derimod Erfaringsverdenen,
det vil sige: den i Erfaringen givne Natursammenhæng,
som den eneste Realitet. Da Fouillée fra Studiet af Platon
vendte tilbage til Nutidens Problemer, maatte det blive
hans Opgave, »at føre Platons Ideer tilbage fra Himlen til Jorden
og derved forsone Idealismen og Naturalismen«. Saaledes
betegner han selv sin Bestræbelse. \textit{(Le mouvement idéaliste et}
% --- p. 65 ---
% running head: Begrebet idée-force.
\opage{65}\textit{la reaction contre la science positive.} 1896. p. XXI). Denne
Forsonings Mulighed bygger han paa det Faktum, at Tanke
(idée) kan føre til Handling. Særligt betoner han, at Tanken
om Friheden kan vække Stræben og udløse Kraft, hvad enten
vi ere »frie« i indeterministisk Betydning eller ikke. \textit{(La liberté
et le determinisme.} 1872). Her staa vi allerede ved Grundtanken
i Fouillée's Psykologi saavel som i hele hans Verdensanskuelse.
Denne Grundtanke har han selv udtrykt ved Ordet
Tankekraft, \textit{idée-force.} Dette Begreb gør det efter hans Opfattelse
muligt at tilvejebringe en Tilnærmelse mellem de forskellige
filosofiske Retninger. Han har første Gang brugt Udtrykket
i Revue philosophique 1879. Med Tankerne i Bevidstheden
ere Bevægelsestendenser i Hjernen umiddelbart forenede. Det
er paa én Gang et fysiologisk og et psykologisk Faktum, vi her
have for os, saa at der her aabenbarer sig et Ideal, der ytrer
sig i selve Naturens Determinisme. Allerede forud for den bevidste
Tanke virker det i den organiske Udviklings- og Udvidelsestendens
saavel som i den instinktmæssige Sympati. Der
rører sig en Villie i hele Naturen, baade i ydre Bevægelse og
i indre Fornemmelse. Det Aandelige og det Materielle, Bevidstheden
og Livet, det Individuelle og det Sociale, Friheden og
Solidariteten virke her som Et. \textit{(La science sociale contemporaine.}
1880).

I det betydeligste af sine Arbejder, \textit{La psychologie des
idées-forces} (1893), har Fouillée udviklet sin Grundtanke fra
den rent psykologiske Side. Det er den bedste Fremstilling af
Voluntarismens Psykologi. Fouillée definerer ligefrem Psykologien
som Studiet af Villien. Han ser klart, at vor Villen ikke er Genstand
for umiddelbar Iagttagelse; men i Stedet for deraf at
drage den Slutning, som Nogle drage, at saa er Villien Intet,
drager han --- efter min Opfattelse\textsuperscript{24}) med fuld Ret --- den
modsatte Slutning, at vor Villen er Et med os selv. Psykologien
har ifølge Fouillée hidtil lidt af Intellektualisme. Man har
ikke set, at de sjælelige Fænomener stedse ere Ytringer af
Trang eller Drift (appétition), der ledsages af Lyst eller Smerte,
alt eftersom de begunstiges eller hæmmes. Det er her ligegyldigt,
om vi betragte Fænomenerne fra den fysiologiske eller
% --- p. 66 ---
% running head: Fouillées Etik.
\opage{66}fra den psykologiske Side. Naar vi stille disse to Sider i Modsætning
til hinanden, maa vi erindre, at vi i dem egentligt kun
have to Abstraktioner (deux extraits d'une réalité unique et
totale). Enhver Skelnen (discernement), selv den mest elementære,
forudsætter allerede et Valg, en Foretrækken (préférence,
choix pratique rudimentaire). Skelnen og Foretrækken gaa i
simpleste Tilfælde ud i Et, f. Ex. ved Dyrets Skelnen mellem
Lystfølelsen ved at spise og Ulystfølelsen ved Sult. Kun det
sanses, der er af Betydning i Kampen for Tilværelsen, og det
er Villien (dette Ord taget i videste Betydning), der driver
Sansningen fremad og betinger dens mere differentierede Former.
Hvad der gælder for Sansning, gælder ogsaa for Genkendelse
og Erindring. Hvad der genkendes og erindres, er først
og fremmest, hvad der har praktisk Interesse. Skridt for Skridt
voxer Interessen --- i Vexelvirkning med Erfaringen --- ud
over det Øjeblikkelige og Nærliggende. Selv de abstrakte logiske
Principer have formet sig ad denne Vei. Enhver Tanke
eller Ide betegner en mere eller mindre bevidst Retning af vort
stræbende og sansende Liv. Der ytrer sig her en stor Kontinuitet
mellem alle sjælelige Processer. Enhver Bevidsthedstilstand
kan kaldes Tanke (idée), forsaavidt den bestaar i en Skelnen,
--- Kraft (force), forsaavidt den bestaar i en Foretrækken. ---

Fouillées Etik hænger nøje sammen med hans Psykologi.
Han betoner særligt, at jeg ikke kan blive mig bevidst uden
at have Bevidsthed om andre Væsener: jeg opfatter dem i
Analogi med mig selv og opfatter mig i Forhold til dem. Med
denne Relativitet i vor Opfattelse er allerede en vis Solidaritet
og en altruistisk Tendens given. Det er mig allerede rent intellektuelt
umuligt at gøre mig selv til det eneste Væsen. Begrænsningen
af det intellektuelle Hovmod medfører en Begrænsning
af den praktiske Egoisme. Her have vi atter en Ide, der
tillige er Kraft, idet den er Et med en uvilkaarlig Stræben,
der --- naar den bliver bevidst --- optræder som et Imperativ.
Idealet ligger i Forlængelsen af den Retning, vor Stræben uvilkaarligt
bevæger sig i, og former sig (formedelst Sammenhængen
mellem Relativitet og Solidaritet) som Ideen om et Frihedens,
Lighedens og Retfærdighedens Rige. --- Denne Opfattelse
% --- p. 67 ---
% running head: Muligheden af en Metafysik.
\opage{67}af Etiken har Fouillée ikke hidtil udført i en sammenhængende
Fremstilling, men fremsat forskellige Steder i sine Skrifter. Den
ligger til Grund i et glimrende kritisk Arbejde, af stor Interesse
for den moderne Etiks Historie. \textit{(Critique des systêmes
de morale contemporains.} 1883). ---

I sin almindelige Verdensanskuelse lægger Fouillée især
Vægt paa, at der ingen Grund er til at tyde Tilværelsen blot
efter de mest elementære Fænomener eller de mest abstrakte
Synspunkter. Materialismen gør sig skyldig heri, naar den gør
Bevægelsen til Alt; Bevægelse er kun det Synspunkt for Virkeligheden,
som støtter sig til Syns- og Berøringsfornemmelserne.
Idealismen gør sig skyldig deri, naar den gør Tanken til Alt;
Tanken er ligesaa vel en Abstraktion, som Bevægelsen er det.
De enkelte Videnskaber have overhovedet en Tendens til hver
for sig at gøre deres Synspunkter til det eneste Udtryk for
Virkeligheden. Men ethvert specielt Synspunkt er ensidigt. En
Verdensanskuelse maa opfatte Virkeligheden som en Totalitet,
indenfor hvilken de specielle Synspunkter hvert for sig have
deres Berettigelse og Betydning. Og da den psykiske Existens
er den eneste umiddelbart givne Virkelighed, vi kende, er det
berettiget at tyde Tilværelsen som Helhed i Analogi med den.
Siden Kants Tid har man altfor ofte opfattet Tanken som det,
der skiller os fra Virkeligheden, i Stedet for at opfatte den som
det, der forener os med Virkeligheden og selv er en Art af
Virkelighed. Og Metafysiken er henvist til at bygge paa den
eneste umiddelbare Erfaring, vi have. Den Erfaring, at Smerte
skys, er mere umiddelbar end den, at Stød bevæger. --- Sine
metafysiske Ideer har Fouillée især udviklet i \textit{L'avenir de la
metaphysique} (1889).

Al Metafysik er hypotetisk og bygger paa Analogi. Den
er en højeste Expansion, udsprungen af det personlige Livs
Trang til Enhed med det universelle Liv. Dette universelle
Liv opfattes som et stort Hele, et Samfund af stræbende Kræfter.
Det er altsaa ikke blot psykologiske, men ogsaa sociologiske
Analogier, der ligge til Grund ved de metafysiske Ideer. Ordet
Gud --- der
altid mere eller mindre udtrykker en fra menneskelige
Forhold hentet Forestilling --- betegner den inderste Grund
% --- p. 68 ---
% running head: Systemernes Kamp for Tilværelsen.
\opage{68}og Higen i Tilværelsens universelle Samfund. Men udover et
hypotetisk Skema kunne vi her ikke naa. Den sidste Syntese,
vor Tanke forsøger, kan ikke forestilles af os paa samme positive
Maade som en speciel Sammenhæng mellem Fænomener,
om hvilke vi have en bestemt og afsluttet Erkendelse. Vi behøve
ikke med Kantianerne at indføre nye Dogmer i Moralens
Navn; ligesaa lidt ere de afsluttende Ideer Genstand for Ironi
eller Spot, som i Renans Dilettantisme. I selve Trangen til at
handle efter vore Idealer trods Usikkerheden af deres Sejr
ligger der en idealistisk Forudsætning skjult: c'est une spéculative
en pensée et en acte sur le sens du monde et de la vie!
Formuleringen af denne Forudsætning i bestemte Sætninger er
stedse hypotetisk. Der staar en stadig Kamp for Tilværelsen
mellem de metafysiske Systemer. Det gælder, hvilket af dem
der bedst trives i den raadende videnskabelige Atmosfære. Mulighederne
for metafysisk Opfattelse ville begrænses mere og
mere, alt eftersom Erkendelsen skrider frem. Den Opfattelse,
der vil sejre, vil være den, der baade lader Analysen og Syntesen
komme til deres Ret paa fuldkomneste Maade. Men Isis'
Slør vil aldrig blive hævet, og der vil stedse staa et Spillerum
aabent for Fantasien og for den religiøse Symbolisme. --- Denne
sidste Tankegang er ikke blot rettet mod Renan, men ogsaa
mod Albert Lange og mod en Tænker, der iøvrigt staar Fouillée
meget nær, og som vi senere komme til at beskæftige os med,
nemlig Guyau.

3. Taine og Fouillée, saavel som Wundt og Ardigò ere afgjorte
Kontinuitetsfilosofer. Deres Stræben gaar ud paa at
finde saa stor Sammenhæng som muligt indenfor Erfaringen,
og at supplere Erfaringen, saaledes at Kontinuitetsprincipet saavidt
muligt fastholdes. I Modsætning til denne Art Forsøg
træder en Række Bestræbelser, der dels af Erfaringsgrunde,
dels af etiske Grunde hævde Diskontinuitetens Betydning. Saadanne
Bestræbelser ere karakteristiske for den idealistiske Retning
ved Overgangen til det 20. Aarhundrede, medens den
Idealisme, der raadede ved Begyndelsen af der 19. Aarhundrede,
væsentligt var en Kontinuitetsfilosofi. I den franske
filosofiske Literatur fremtræder Diskontinuitetsfilosofien paa en
% --- p. 69 ---
% running head: Dilkontinuitetsfilosofien i Frankrig.
\opage{69}særligt interessant og energisk Maade. Der kan skelnes mellem
tre forskellige Motiver, som ledende for Diskontinuitetsfilosofi.
Erfaringen frembyder Kvalitetsforskelligheder, som det hverken
var lykkedes Spekulationen eller Udviklingslæren at reducere.
Comte's Positivisme havde udtrykkeligt anerkendt den Kløft,
der skilte de forskellige Naturomraader fra hverandre. Enhver
ny Videnskab betød for Comte en særegen, irreduktibel Gruppe
af Fænomener. --- Dertil kommer, at selv indenfor hver enkelt
Gruppe af Fænomener kan Aarsagssætningen kun ufuldkomment
finde sin Bekræftelse. Man drager derfor Hume frem
igen og stiller hans Empirisme mod Kants og Evolutionismens
Forsøg paa at overvinde den. --- Endeligt henviser man til Bevidstheden
om Initiativ, om Evnen til i sin Tænken og Handlen
at sætte noget Nyt ind i Verden, og man betoner stærkt
denne Evnes moralske Betydning. --- Hos de to Tænkere, vi
ville fremføre som Repræsentanter for Diskontinuitetsfilosofien
virke disse tre Motiver sammen paa noget forskellig Vis.

4. \textsc{Charles Renouvier} er Filosofiens Nestor i vor Tid.
Han er født 1818, og endnu 1903, altsaa 85 Aar gammel, har
han udgivet et vel gennemtænkt Skrift. Han fik sin videnskabelige
Uddannelse som Elev i den polytekniske Skole og
var i sin Ungdom ivrig Saint-Simonist og Republikaner. Hele
sit Liv igennem har han levet som Privatmand, men har øvet
en ikke ringe Indflydelse gennem sine Bøger og gennem et
Tidsskrift, han en Tidlang udgav. En energisk og kundskabsrig
Tænker og en alvorlig Karakter fremtræder for os gennem
hans Skrifter; dog er han, som saa mange franske Filosofer,
stærkest som Kritiker. Han optraadte baade i Logikens og i
Moralens Navn mod al Art Kontinuitetslære, hvad enten den
støttede sig til teologisk Mystik, til metafysisk Spekulation eller
til naturvidenskabelig Hypotesedannelse. Hans Virksomhed som
Kritiker og Tænker begynder længe før det Tidsrum, som nærværende
Skrift behandler. Der er her, som allerede bemærket,
et Hul i min Fremstilling af den nyere Filosofis Historie for
det moderne Frankrigs Vedkommende, et Hul, jeg heller ikke
nu ser mig i Stand til at udfylde. Jeg forbigaar her hans tidligere
Værker (af hvilke det vigtigste er \textit{Essais de Critique}
% --- p. 70 ---
% running head: Charles Renouvier.
\opage{70}\textit{générale.} 1854---1869) og holder mig til hans Skrifter fra det
sidste Fjerdedelsaarhundrede. Til Orientering angaaende hans
Standpunkt vil det være af Vigtighed at studere \textit{Classification
systématique des Doctrines philosophiques} (1885), \textit{Les dilemmes
de la métaphysique} (1901) og \textit{Le personalisme} (1903),
foruden forskellige Afhandlinger i det af Renouviers Ven og
Medarbejder Pillon udgivne Aarsskrift \textit{L'année philosophique.}

Renouvier har selv givet en Fremstilling af sin Udviklingsgang,
idet ban lægger Vægt paa, at ethvert System --- som
Værk af et Menneske, der lever og tænker under bestemte
indre og ydre Forhold --- er betinget ved personlige Faktorer,
der faa Indflydelse paa det Synspunkt, der vælges. (Se Afhandlingen
\textit{Comment je suis arrivé à cette conclusion} i anden Del
af det store Værk om de filosofiske Teoriers Klassifikation).

I sin Ungdom havde han, som allerede omtalt, en Periode,
i hvilken han var ivrig Socialist af Saint Simons Skole, og samtidigt
førte hans matematiske Studier ham ind paa Syslen med
Uendelighedsbegrebet. Ad begge Veje var han kommen ind
paa en Kontinuitetsfilosofi: som Socialist betonede han Solidariteten
paa Individualitetens Bekostning, og som Tænker lod
han alle endelige Forskelle fortone sig gennem uendelige Overgange
og stadigt fortsatte Processer. Men i begge Henseender
vaktes Tvivl, som næredes ved ivrigt Studium af Descartes og
Kant. Han saa nu, at ethvert Forsøg paa at konstruere et System,
der skal gælde for Tilværelsen som Helhed, ender i Modsigelser.
Kants Lære om Antinomierne virkede især stærkt
paa ham, skønt han ikke kunde give Kant Ret i, at hverken
Antiteser eller Teser kunde gendrives; efter hans Overbevisning
havde Teserne, der hævde Verdens Endelighed, Ret overfor
Antiteserne, der hævde dens Uendelighed. De Modsigelser,
Spekulationen tilsidst hilder sig i, hjemsøge ikke blot de rent
filosofiske Forsøg, men ogsaa ethvert teologisk Forsøg paa at
forene Guds Uforanderlighed med hans Skabervirksomhed og
hans Uendelighed og Forudviden med menneskelig Frihed. Renouvier's
Tænken er ikke blot baaren af rent teoretisk Interesse,
men ogsaa af Bestræbelsen for at bane Vej for en rationel
Reform af de religiøse Forestillinger. --- Efterhaanden førtes
% --- p. 71 ---
% running head: Néokriticismen.
\opage{71}Renouvier til en ejendommelig Form af kritisk Filosofi, der i
Frankrig kaldes le néocriticisme, og som jeg nu skal karakterisere
i dens Hovedtræk.

De første Principer for vor Verdensanskuelse bestemmes
ifølge Renouvier ved et Valg mellem de Modsætninger, Tænkningen
stiller os overfor. --- Der maa saaledes vælges mellem
Uendelighed og Endelighed. Begrebet Uendelighed indeholder
en Modsigelse, dersom man derved mener Andet og Mere end
en Mulighed for Tanken af stedse at gaa videre (f. Ex. af at
føje stedse nye Dele til mit Fantasibillede af Rummet). Som
virkeligt given maa Verden være endelig, skønt vor Erfaring
ikke kan omfatte den. Kants Teser have Ret imod hans Antiteser.
--- Ved dette Valg er Kontinuiteten allerede brudt; ti
med Uendelighedsbegrebet falder Muligheden af uendeligt mange
Overgange mellem Fænomenerne. Deraf følger igen, at Aarsagssætningen
kun kan have begrænset Gyldighed. Renouvier
slutter sig her til Hume og hævder med ham Aarsagssætningens
Ubevislighed. Der kan indtræde nye Begyndelser, Aarsager,
som ikke ere Virkninger. Det bliver af Betydning for
Moralen, idet Villiens Frihed derved bliver mulig. Filosofien
kan derfor ikke blive Konstruktion; den kan kun blive en Lære
om, hvad der sker forud og efter de rene Begyndelser, og
om hvorledes Friheden bruges paa rette Maade. Friheden selv
er en første Sandhed i Erkendelsens Verden; den er en Forudsætning,
som egentligt selve Deterministen benytter sig af,
naar han vælger at gøre Nødvendigheden til første Princip, ---
kun at han derved hilder sig i en Modsigelse, som Frihedsfilosofen
undgaar\textsuperscript{25}). --- Det er saaledes en udpræget voluntaristisk
(eller »arbitrær«) Erkendelsesteori, Renouvier hylder, ---
ikke blot fordi Principerne vælges frit, men ogsaa fordi det,
der vælges, er Frihedsprincipet, Diskontinuitetens Princip. De
Modsigelser og Vanskeligheder, Renouvier støder paa eller forefinder
i Tænkningens Historie, føre ham ikke til at høre op at
tænke, ej heller til at lege med Tankerne, men til at fastholde
den Antagelse, der baade logisk og etisk staar for ham som
den rette, med en Villiesakt. Han er Renans afgjorte Antipode.
--- Sin Kamp mod Kontinuiteten fører Renouvier ikke blot,
% --- p. 72 ---
% running head: Ronouvier mod historisk Fatalisme.
\opage{72}hvad Individets, men ogsaa hvad Slægtens Udvikling angaar.
Der gives, paastaar han, ikke nogen altomfattende historisk Lov,
men mange Love, hver gældende for sin Fase i Historien, idet
hver enkelt Fase indledes ved en ny Begyndelse, en Frihedsakt,
der havde kunnet falde anderledes ud, end den faktisk
gjorde. I de store Menneskers Optræden ser han saadanne nye
Begyndelser, og han bestrider den Opfattelse, der betragter
saadanne Personligheders Opstaaen som Produkt af lange Tiders
langsomme Udvikling. Det Skabende hos dem forstaas
derved ikke. Han ser en stor Fare i den historiske Betragtningsmaade,
der betragter Alt, hvad der faktisk sker, som det
eneste Mulige; derved kommer man til at retfærdiggøre Alt.
Han hekæmper Romantiken, som især har indført denne historiske
Opfattelse, der rehabiliterer al Fortid uden Hensyn til de
Muligheder for en helt anden Udviklingsretning, som paa ethvert
Punkt ere til Stede. I sit Skrift \textit{Uchronie} (L'Utopie dans
l'histoire) (1901) har han forsøgt at vise, at den antike Kultur
ingenlunde behøvede at gaa til Grunde for at afløses af Overtro
og Barbari, og han har skildret en Reformation af Romerriget
og af den antike Religion, som laa indenfor Mulighedernes
Grænser og vilde have ført os udenom Middelalderen.
Ved denne »europæiske Civilisationshistorie, saaledes som den
ikke har været, men som den har kunnet være«, vil Renouvier
styrke Frihedsbevidstheden og Ansvaret og modarbejde de reaktionære
Tendenser, der støtte sig til en Kultus af Historien.

Gennem sin Kritik af Uendelighedsbegrebet og af Kontinuitetsprincipet
kom Renouvier til at indse Forholdslovens eller
Relativitetsprincipets Betydning. Begrebet Forhold (relation) blev
nu for ham Grundlaget i vor Erkendelse, Grundbegrebet, der
lægger sig for Dagen i alle de Kategorier, hvormed Erkendelsen
arbejder. Han slutter sig her til William Hamilton, der
først havde paavist dette. Enhver Erkendelse naas gennem at
udfinde de Forhold, i hvilke Tingene staa, og enhver Genstand,
vi kende, staar derfor ogsaa for os som et System af Relationer.
Relativitetsprincipet er en Metode (la méthode des relations),
som udspringer af selve vor Bevidstheds Væsen. Bevidstheden
selv er et Forhold, men forskellig fra alle andre
% --- p. 73 ---
% running head: Forholdsloven og dens Konsekvenser.
\opage{73}Forhold derved. at den er et Forhold mellem mit Jeg som Subjekt
og mit Jeg som Objekt (relation de soi à soi), altsaa et Forhold
mellem identiske Led. Derved er den Grundlaget for alle andre
Forhold, der kun blive tydelige ved at henføres til dette Grundforhold.
--- Af Forholdsloven følger, at al Erkendelse kun angaar
Fænomener, idet Alt, hvad vi skulle erkende, maa fremtræde
for Bevidstheden og bestemmes ved dens Love. Vi kende
kun Tingene som Genstande for Forestillen. Vil man ud over
Fænomenerne, kan det kun ske ved religiøse Postulater. Heller
ikke Postulaterne kunne dog frigøres fra Forholdsloven; Gud
kan ligesaa lidt være uendelig, som Verden kan det. Men ad
Postulatets Vej komme vi dog ud over Erfaringens Grænser.
--- I sin tidligere Tid var Renouvier tilbøjelig til at overføre
sin Republikanisme paa Teologien og mente, at Pluralisme (i
religiøs Form Polyteisme) var en Konsekvens af Relativismen.
Senere fandt han dog saa megen Sammenhæng i Verden, at
Monoteismen stod for ham som nødvendig, skønt han ikke kunde
antage Gud for et uendeligt Væsen. Det fulgte --- mente han
nu --- netop af Forholdsloven, at man maa ende ved en første
Aarsag, altsaa ved en Skabelse, som et Punkt, hvor Tanken
stanser (le point d'arrêt de la pensée). Ellers komme vi ind i
en uendelig Proces, der ingen Løsning af Gaaderne giver os.
Renouvier mener, at mange Tænkere forvexle den stadige Forskydning
af Problemerne (le reculement indéfini des problèmes)
med deres virkelige Løsning. Men naar man saaledes ender i
Forestillingen om en Verdensaarsag, der selv er endelig, saa er
--- da nu det definitive Slutningspunkt for Tanken skal være
naat --- enhver Grublen over denne Verdensaarsags Oprindelse
resultatløs. Her begynder den absolute Uvidenhed.

Det sidste Valg, som bestemmer vor Erkendelses Resultater,
staar mellem Begreberne Ting og Person, mellem en Verdensanskuelse,
der underordner det Personlige og Subjektive
under det Upersonlige og Objektive, og en Verdensanskuelse,
der opfatter Alt i Analogi med det Personlige. Det har kostet
et langt og stort Arbejde for Mennesketanken at naa til en
objektiv Betragtningsmaade fra Mytologiernes kaotiske Personifikationer.
Det gjaldt om at gennemføre Tanken om den lov%
% --- p. 74 ---
% running head: Monadelære.
\opage{74}mæssige, objektivt givne Sammenhæng, som danner Forudsætning
og Genstand for de personlige Væseners Arbejde. Men
Begrebet om Objekt (Tingen, Naturen, Substansen, »Ideen«) faar
da let en saa overvældende Magt over Tanken, at denne glemmer,
at hvad der er realt, kun er givet som Genstand \textit{for den} som Fænomen,
og ikke indser, at en forstaaelig Verdensopfattelse kun faas
ved at lægge vore umiddelbare Erfaringer til Grund ved Tydningen
af den objektive Sammenhæng. Ad disse to Veje, den
erkendelsesteoretiske Paavisning af Subjektets logiske Prioritet
og den paa Analogi grundede metafysiske Tydning af Naturen,
kommer Renouvier til en Monadelære, han har fremstillet i
sine to seneste Værker, \textit{La nouvelle Monadologie} (1901) og
\textit{Le personalisme} (1903). Det er hans Overbevisning, at efterat
den videnskabelige Forstaaelse af Natursammenhængen er sikret
gennem streng og metodisk Videnskab, er Tiden kommen for
en ny Anvendelse af Personlighedsbegrebet i Verdensanskuelsen,
en Anvendelse, der nu kan ske med større Kritik og med
klarere Bevidsthed end i de mytologiske Fantasier og i de teologiske
Spekulationer. Det er her særligt Villiesbegrebet, han
støtter sig til. Medens han er enig med Hume i, at Aarsagsforholdets
objektive Gyldighed ikke kan godtgøres, saa at vi
maa lade os nøje med paa saa mange Punkter i Naturen som
muligt at paavise uforanderlig Rækkefølge, skiller han sig fra
Hume ved at hævde Aarsagsbegrebets uopløselige Sammenhæng
med vor Bevidstheds Natur som en af de Kategorier, hvormed
den ifølge sit Væsen maa arbejde. Og han skiller sig endnu
mere fra Hume ved at lære, at Aarsagsbegrebet tilsidst har sit
Grundlag i vor Bevidsthed om vor Villen. Intet Forhold mellem
Aarsag og Virkning er os forstaaeligt (pénétrable), naar det
ikke tydes som noget med vor Villiesvirksomhed Analogt. Begrebet
Villie gør os Begrebet Kraft forstaaeligt, men kan selv
ikke defineres ved at føres tilbage til noget mere Primitivt,
Ved dette Begreb maa Tanken stanse --- i Kraft af Forholdsloven.


Det bliver under disse Forudsætninger let for Renouvier
at tænke sig Sjælelivets Opstaaen og Udvikling i Verden. Spirerne
til Sjælelivet maa --- i Kraft af Monadologien --- have
% --- p. 75 ---
% running head: Renouviers Mytologi.
\opage{75}været til Stede fra først af, saa at der blot behøves gunstige
Betingelser for deres Udfoldelse. For at forklare sig det Ondes
Opstaaen antager han en meget tidligt opstaaet Disharmoni, en
Art Syndefald i Monadernes Verden, som Guddommen, der selv
er begrænset og har begrænset sig, ikke har kunnet hindre. I
sine teologiske Forestillinger lægger han Vægt paa Nødvendigheden
af at tænke Guddommen i Analogi med Mennesket og
derfor som endelig og begrænset. Kun da kan en etisk Verdensanskuelse
fastholdes, medens Uendelighedslæren fortoner
alle Modsætninger, ogsaa de etiske.

Medens flere Tænkere (Bruno, Böhme, Spinoza) overfor den
Udvidelse af Verdensanskuelsen, som den kopernikanske Lære
gav Stødet til, fandt Tilhold i Ideen om Uendeligheden og omfattede
den med Begejstring, indskærper Renouvier, at man,
hvis man vil fastholde de etiske Ideers ubetingede Betydning,
maa gøre Mennesket til Verdens Centrum, ligesaa vel som
Aristoteles og Bibelen gjorde det: l'anthropocentrisme est le
point de vue moral de l'univers. Den ydre Verden er i Størrelse
og fysisk Kraft os langt overlegen. Men vor Jord kunde
dog maaske være Sædet for noget i aandelig Henseende langt
Betydeligere end den øvrige Verden, og naar det gælder aandelig
Magt, vilde det da snarere være os, der med vor Overlegenhed
(supériorité d'agents intelligents) kunde knuse Universet,
end omvendt. Fuldkommenhed bestaar stedse i Begrænsning, og
det er derfor et Fantasteri, som fører bort fra den sande Storhed,
naar man fortaber sig i den materielle Verdens Ubegrænsethed.
--- Denne Hentydning til en bekendt Udtalelse af Pascal
er karakteristisk for Renouviers Standpunkt, og den fremkommer
unegteligt hos ham med større Konsekvens end den
tilsvarende Sætning hos Pascal. Overhovedet er det karakteristisk
for Renouvier, at selv hvor han mere optræder som Mytolog
eller Teolog end som Filosof, former han sine Anskuelser
med skarp Logik og gaar ikke paa Akkord. Han er tillige klar
over, at han bevæger sig paa de dristige Gisningers og Hypotesers
Omraade. »Jeg dogmatiserer ikke«, siger han, »men
søger at begribe. Hvad jeg i Virkeligheden har til Øjemed i
mine dristige Gætninger er at belyse de ophøjede Muligheder
% --- p. 76 ---
% running head: Kritik af Renouvier.
\opage{76}. . . . Hvor det gælder at efterspore Sandheder, som ligge ud
over Erfaringen, ja, som først skulle forklare selve Erfaringen,
er Hypotesen, uden anden Kontrol end Logiken og Moralen,
Filosofens eneste Udvej. Naar han forbliver i Overensstemmelse
med den rene Fornufts Love, har han Ret til at hævde
Sandsynligheden af de Sætninger om vor Verdens Oprindelse
og Natur, der stemme med hans Forstands Lov og med hans
Følelse ved Livet«. ---

Kun en enkelt, men ganske vist meget indgribende kritisk
Bemærkning har jeg at gøre, inden jeg forlader vor Tids filosofiske
Patriark. Han støtter Antagelsen af en første Aarsag,
og dermed hele sin ejendommelige Religionsfilosofi, til Forholdsloven
(loi de relativité). Denne Lov kræver, mener han, at vi
sætte en Grænse for den stadige Tilbageskriden fra Virkning
til Aarsag. Men det Led, vi stanse ved, kommer saa til at betegne
en Grænse for Forholdsloven. Det skal vel staa i Forhold
til det, der udledes af det, men ikke selv paa den anden
Side være udledet af noget Andet, saaledes som det dog gælder
om alle de tidligere Led i Rækken. Forholdsloven ophæves
altsaa, naar vi ere komne tilbage til det første Led. Dette er
en afgjort Selvmodsigelse, naar man som Renouvier opfatter
Forholdsloven som en for vor Tanke nødvendig Lov, dybt grundet
i selve Bevidsthedens Natur. Forholdsloven kræver, at vi
paa ethvert Punkt knytte Led til Led, og hvor vi høre op med
det, hører ogsaa vor Forstaaelse op. Forholdsloven fører forsaavidt
snarest til en Uendelighedslære, idet det ikke paa noget
Punkt kan tillades os at mene, at vi have naat et absolut Udgangspunkt.
Renouvier har Ret i at sige, at Henvisningen til
en uendelig Proces ikke giver nogen Forklaring. Men Sagen
er, at vi, hvor vi staa overfor en Række, der i Kraft af sin
Lov stadigt kan fortsættes, ingen anden Forstaaelse kunne naa
end den, der ligger i selve Loven\textsuperscript{26}). Der er her en mytologisk
Tendens hos Renouvier, som faar Overhaand over hans
Kriticisme. I Stedet for at fordybe sig i den indre Sammenhæng
mellem Fænomenerne vil han havne i en første Begyndelse
og giver os saa kun Valget mellem at antage en »Ting«
eller en »Person« som denne første Begyndelse. I Kraft af
% --- p. 77 ---
% running head: Emile Boutroux.
\opage{77}selve Forholdsloven maa den kritiske Filosofi afvise dette Dilemma.
Kontinuitetsfilosofien vil, netop naar Trangen til Forstaaelse
betones, stedse beholde et Fortrin for Diskontinuitetsfilosofien.


5. Medens Renouvier arbejder med skarpt tilskaarne Postulater
og med Energi forkynder de store Brud paa Kontinuiteten,
søger \textsc{Emile Boutroux} (f. 1845, Professor ved Sorbonne)
mere en rent teoretisk Begrundelse af Diskontinuitetsfilosofien
og vender især sin Opmærksomhed mod de smaa
Differencer. I Boutroux's Filosoferen spores Fortsættelsen af
en Retning, der i en tidligere Periode af fransk Filosofi repræsenteredes
af Maine de Biran (som er omtalt i andet Bind af
den nyere Filosofis Historie).

I to betydningsfulde Afhandlinger \textit{(De la contingence des
lois de la nature} 1875, og \textit{De l'idée de la loi naturelle dans
la science et dans la philosophie contemporaines} 1895) drøfter
Boutroux Aarsagssætningen i de forskellige Former, i hvilke
de forskellige Videnskaber formaa at anvende den. Hans Hovedformaal
er at godtgøre, hvor Meget der mangler i dens Gennemførelse,
naar den opfattes i den strenge Form, i hvilken den
mekaniske Naturopfattelse anvender den. Især ved Overgangen
fra et Erfaringsomraade til et andet staa vi overfor Diskontinuitetsforhold
--- saaledes ved Overgangen fra Logik og Matematik
til Mekanik, fra Mekanik til Fysik, og saaledes videre gennem
Kemi, Fysiologi, Psykologi, Sociologi. De mere konkrete
Videnskabers Grundlove kunne ikke reduceres til de mere abstraktes;
hvert nyt Erfaringsomraade kræver nye, specielle Principer,
der ikke ere indeholdte i det foregaaendes. Boutroux er
her i mærkelig Overensstemmelse med Comte, der allerede
stærkt havde betonet Principernes Diskontinuitet ved Overgangen
fra en Videnskab til en anden og bestemt havde advaret
mod at forsøge at udlede en mere konkret Videnskabs Principer
af en mere abstrakts. Boutroux har selv været opmærksom
paa, hvorledes han her i en vis Forstand mødes med Positivismens
Grundlægger. I en aandfuld Karakteristik af Comte,
der tjente som Indledning til en Diskussion i det franske filosofiske
Selskab, har han desuden særligt fremdraget de Ele%
% --- p. 78 ---
% running head: Variations contingentes
\opage{78}menter i Comtes Filosofi, der pege ud over det strengt empiriske
Standpunkt, det var Comtes Bestræbelse at indtage.
\textit{(Bulletin de la société française de philosophie.} Janvier 1903).

Ejendommeligt for de konkrete Videnskaber er det ifølge
Boutroux, at de angaa Helheder, der tænkes at bestaa af Elementer,
uden at der dog kan paavises et Ækvivalensforhold
mellem en saadan Helhed og de Elementer, af hvilke den antages
at bestaa. Der er f. Ex. intet Ækvivalensforhold mellem
et Menneske som Helhed og de Elementer, der kunne paavises
i dets Natur. Derfor ere vi her stedse henviste til Erfaring;
kun den kan belære os om, hvilke Helheder der dannes ved
Elementernes Vexelvirkning. Her fremtræder det da klart, at
Aarsagssætningen er en abstrakt Lov, der ikke udtømmer hele
den konkrete Realitet. Under den faktiske Udvikling optræder
der uforudseelige, og forsaavidt tilfældige Forandringer (variations
contingentes). Tingenes Historie indeholder for os Nøglen
til deres Væsen, men kan ikke udledes af deres Natur. Det
dynamiske Synspunkt faar, jo mere vi gaa fra de abstrakte til
de konkrete Videnskaber, Overvægten over det statiske, og Kvaliteterne
raade i Stedet for Kvantiteterne. Naturlovene angive
kun de konstante Forhold, i hvilke Forandringerne finde Sted,
saa at sige de Vaner, Tingene lægge for Dagen. Forandringen
er Principet, Permanensen er kun et Resultat. Det Konstante
i Naturen, de saakaldte Naturlove, er Flodlejet, i hvilket Begivenhedernes
Strøm flyder, --- og som Strømmen selv har uddybet,
skønt den nu følger det. Dersom der i selve Tingene
indtraadte dybere gaaende Variationer, vilde dette Flodleje kunne
ændre sin Form og sin Retning; selve de objektive, logiske Forhold
vilde da undergaa Forandring.

Naar Ændringerne, gennem hvilke Udviklingen fra de elementære
til de højere Tilværelsesformer foregaar, ikke kunne
bringes i Ækvivalensforhold til de forudgaaende Tilstande, raader
forsaavidt Tilfældigheden. Men denne kan tydes som det
ydre, empiriske Tegn paa Frihed. Friheden selv kan ikke umiddelbart
erfares, men den ytrer sig i den menneskelige Handlen,
hver Gang en Vane overvindes. Paa det psykologiske Omraade
angaa Lovene kun det, Sjælen overlader Vanen at udrette. Va%
% --- p. 79 ---
% running head: og deres Tydning.
\opage{79}nen har her, som overalt i Naturen, den Betydning, at der ved
Mekaniseringen af de mest nødvendige Funktioner kan frigøres
Kraft til nyt Arbejde.

Det er altsaa for Bontroux den frie Udvikling gennem
Nydannelser, der er det egentligt Arbejdende og Virkende i
Tilværelsen, den egentlige Existens. De konstante Former og
Love ere kun saakaldte Resultater. Hvis han, da han skrev
sine aandfulde Afhandlinger, havde kendt Hugo de Vries'
»Mutationsteori«, vilde han have kunnet tage den til Indtægt
for sin Anskuelse, idet den nævnte Teori antager springvis
Varieren og Opdukken af nye Typer. Lamarcks Teori sympatiserer
han øjensynligt med, idet den jo antog en indre Trang
og Stræben hos de levende Væsener, der føre dem til at tillempe
sig efter Forholdene og saaledes udvikle faste Former.

Men i sit Skrift om Begrebet Naturlov i den moderne
Videnskab lægger Boutroux dog ogsaa Vægt paa et andet Synspunkt
for det Konstante i Naturen. Erkendelsesteoretisk set
er det, vi kalde Naturens Love, Indbegrebet af de Metoder, vi
have fundet til at gøre Tingene tilgængelige for Forstanden
(assimiler les choses à notre intelligence) og benytte dem til
at opnaa vore Formaal. Og hvad der saaledes gælder Begrebet
Naturlov, gælder alle de Begreber, vi danne forat forstaa
Tingene, særligt de Distinktioner, vi foretage forat kunne se
klarere. Naar vi saaledes skelne saa skarpt mellem Tanke og
Bevægelse, og maaske antage disse to Omraader adskilte ved
en Afgrund, saa have vi ingen Ret til at tillægge denne Distinktion
absolut Betydning. Vore Begreber og Begrebsforskelle
udtrykke snarere vor Maade at betragte Tingene paa end Tingenes
egen Maade at være paa. --- Ved disse sidste Bemærkninger
nærmer Boutroux sig til den økonomiske Erkendelsesteori,
vi ville støde paa indenfor det følgende Afsnit. Havde
han gennemført den erkendelsesteoretiske Betragtningsmaade,
han her antyder, vilde han maaske være bleven overbevist om,
at det »Tilfælde«, som han tyder som Udtryk for »Frihed«,
egentligt kun betegner en negativ Grænse for Forskningen,
der maaske ikke er definitiv, men stiller nye Opgaver. Diskontinuitetsfilosofien
bliver dogmatisk, naar den antager de
% --- p. 80 ---
% running head: Variation et Problem .
\opage{80}Forskelle eller Spring, ved hvilke vi foreløbigt stanse, for absolute.
Det maa bestandigt staa som et Problem, hvorledes de
nye Variationer ere blevne til. At de for vor Iagttagelse indtræde
pludseligt, have de fælles med mange Begivenheder,
--- maaske endog, naar man ser nøje til, med alle Begivenheder
i Verden.

% --- p. 81 ---
\opage{81}\grouphead{ANDEN GRUPPE.}{ERKENDELSESTEORETISK-BIOLOGISK\\ RETNING.}

Erkendelsesteoretiske Overvejelser vare ingenlunde fremmede
for den Gruppe af Tænkere, vi i det Foregaaende beskæftigede
os med. Det drejede sig dog for dem tilsidst om det Spørgsmaal,
hvorvidt noget af de specielle Synspunkter, de enkelte
Videnskaber anvende hver paa sit Omraade, egne sig til at
være Princip for en sammenhængende og foreløbigt afsluttende
Verdensanskuelse. Men de Vanskeligheder dette Spørgsmaals
Besvarelse stødte paa, maatte naturligt vække Trangen til en
fornyet Undersøgelse af de videnskabelige Synspunkter og
Principer, deres Udspring og Betingelserne for deres Gyldighed.
Karakteren af det Bidrag, Naturvidenskaben kan give til
en almindelig Verdensanskuelse, blev --- ikke mindst indenfor
selve Naturforskernes Kreds --- tagen op til fornyet Drøftelse.
Den saakaldte mekaniske Naturopfattelse, der er grundlagt af
Galilei og Descartes, og som førte Alt tilbage til Tryk, Stød og
Tiltrækning, havde man været tilbøjelig til at betragte som en
højeste Sandhed, som Mere end en Hypotese. Der kom nu nye
Fænomener frem (navnligt ved den elektromagnetiske Lysteori),
som det var vanskeligt at bringe ind under den mekaniske Naturopfattelse,
og desuden blev det ogsaa indskærpet, at en Teori
ikke tør anses for absolut rigtig, fordi den bekræftes ved Er%
% --- p. 82 ---
% running head: Filosoferende Naturforskere .
\opage{82}faring, saalænge Muligheden af en anden Teori, hvis Konsekvenser
ogsaa kunde bekræftes ved Erfaring, ikke er udelukket.

I samme Retning virker den Vægt, der lægges paa Udviklingen,
paa den stadige Udvidelse af Erfaringens Grænser og
en mulig Ændring i vor Aands indre Struktur. Det Dynamiske
faar derved Overhaand over det Statiske, og vor Erkendelseslære
vil kun kunne have Gyldighed for et vist Udviklingstrin.

Naar endeligt Villieselementet saa stærkt betones, særligt
ved selve Erkendelsen af Principerne, saa ligger det nær at
drøfte det Spørgsmaal, hvilket Maal det egentligt er, vi stille os
i vor Erkendelse, og hvilke Midler vi raade over til at nærme
os det. ---

Ad disse forskellige Veje føres vi over til den Gruppe af
Tænkere, vi have betegnet som den erkendelsesteoretisk-biologiske.
Til denne Gruppe regner jeg først nogle berømte Naturforskere,
der dels ved deres specielle Forskninger, dels ved Studiet af
deres Videnskabs Historie ere førte til Eftertanke over Erkendelsens
Forudsætninger, dels en Filosof, som har foretaget en
mere gennemført Undersøgelse af Problemernes Opstaaen og
Forsvinden under Livets og Udviklingens Gang.

\chapterhead{I.}{FILOSOFERENDE NATURFORSKERE.}

1. Den udmærkede Fysiker \textsc{James Clerk Maxwell} (f.
i Edinburgh 1831, død som Professor i Fysik i Cambridge
1879), især berømt som den, der i matematisk Form har udviklet
den elektromagnetiske Lysteori, faar Betydning for Filosofiens
Historie ved den Interesse. han har næret for at gøre
sig klart Rede for sin Videnskabs første Principer. Det var
William Hamilton, hvis Forelæsninger han hørte som Student i
Edinburgh, der grundlagde hans filosofiske Interesse; han fik
her (efter hvad hans Biografer sige) en Impuls, der aldrig tabte
sin Virkning. Hamilton fra sin Side følte ogsaa stor Interesse
for den unge skarpsindige og videbegærlige Student. Maxwells
Opfattelse af Tænkningens Væsen og Betydning minder endnu
senere hen, da han kunde bygge den paa sin egen dybtgaaende
% --- p. 83 ---
% running head: James Clerk Maxwell.
\opage{83}Forsken, paa afgørende Punkter om Hamiltons Filosofi. De Arbejder
af Maxwell, som især ere af Betydning, findes i andet
Bind af hans \textit{Scientific Papers,} som samledes og udgaves efter
hans Død.

Man har paastaaet, siger Maxwell, at metafysisk Spekulation
kun skulde høre Fortiden til, og at Fysiken skulde have
udryddet den. Der er dog ingen Fare for, at Drøftelsen af
Tilværelsens Grundbegreber skal afsluttes vor Tid. Øvelse
af den spekulative Evne vil vedblive at være ligesaa tiltrækkende
for enhver frisk Aand, som den var det i Thales' Tid!
--- Men det gælder om at udvikle Tanker, med hvilke vi kunne
arbejde. De exakte Videnskabers Fremskridt beror paa Opdagelse
og Udvikling af hensigtsmæssige (appropriate) og nøjagtige
Forestillinger, ved hvis Hjælp vi kunne give en aandelig
Genfremstilling (representation) af Kendsgerningerne, --- en
Genfremstilling, som paa den ene Side er tilstrækkeligt almindelig
til at kunne repræsentere (stand for) ethvert specielt Tilfælde,
og paa den anden Side tilstrækkeligt nøjagtig til at danne
Grundlag for matematiske Slutninger. Denne Vej er det, man
er gaaet, lige fra Euklid til Faraday.

Naar der spørges, hvorledes vi fra først af falde paa at
opstille og anvende et Princip paa et bestemt Omraade, henviser
Maxwell til Analogiens Betydning for Videnskaben. Ved
\textit{fysisk Analogi} forstaas den partielle Lighed mellem to Erfaringsomraaders
Love, der gør det muligt, at det ene Omraade
kan belyse det andet. Al Anvendelse af Matematik i Videnskaben
er saaledes grundet paa en Analogi mellem de fysiske
Størrelsers Love og Tallenes Love, og de exakte Videnskabers
Bestræbelse gaar ud paa at føre Naturens Problemer tilbage
til Bestemmelse af Størrelser formedelst Taloperationer\textsuperscript{27}). Og
af lignende Beskaffenhed er egentligt hele den saakaldte mekaniske
Naturopfattelse. Den hviler paa en Analogi mellem
Lovene for de kvalitative Forandringer i Naturen og Lovene
for Bevægelse. Derved bliver det muligt at beskrive Legemernes
Forandringer, som om de blot bestod i Smaadeles Bevægelser.
En saadan Beskrivelse af et Fænomen, der gør det muligt
at se dette som Exempel paa et Princip, der ogsaa kan anven%
% --- p. 84 ---
% running head: Hvilen forstaas ved Bevægelsen.
% CHECK p. 84: indent after a broken word?
\opage{84}des paa andre Fænomener, er en Forklaring, og Erfaring viser,
at der med enhver saadan Forklaring følger en ejendommelig
aandelig Tilfredsstillelse. Naar blot Analogien gælder, behøver
man ikke at gøre sig bestemte Forestillinger om hine Smaadeles
Natur, end ikke tillægge dem Udstrækning, Form, Haardhed
eller Blødhed. Hvad man maa forudsætte om Materiens sidste
Dele, er kun Konstans og Lovmæssighed af deres Bevægelser.
Det er Bevægelsens Konstans og Kontinuitet, der ligger til
Grund for Materiens Begreb.

Paa dette sidste Punkt slutter Maxwell sig til Tanker, der tidligere
vare gjorte gældende af Faraday og W. Thomson (Lord
Kelvin). \textit{Kun ved Bevægelse,} hævder han, \textit{forstaar man Hvilen
og Ligevægten;} derfor maa Kinematiken (Læren om Bevægelsen,
afset fra Beskaffenheden af det, der bevæger sig) gaa
forud for Statiken (Læren om Ligevægtsforhold). Kinematiken
adskiller sig blot fra Geometrien ved, at Tiden udtrykkeligt
indføres som maalelig Kvantitet. Dog er i Grunden selve Geometrien
afledet i Forhold til Bevægelseslæren, eller er en Del
af denne; ti Geometrien handler egentligt om den Proces, hvorved
Figurer i Rummet frembringes. En Linie er ikke fra først
af en Streg paa Tavlen, der ligesaa godt kan kaldes A B som
B A, men er Banen for en Bevægelse fra A til B. Forestillingen
om Bevægelse ligger til Grund for Forestillingen om
Form. At dette er saa, viser ogsaa den psykologiske Lære om
Opmærksomheden: denne er til enhver Tid begrænset til en
enkelt udelt Forestilling, men bevæger sig frem gennem en
kontinuerlig Række af Forestillinger. Man holder sig i Reglen
for meget til de færdige Tegn og Figurer og overser Sansningens
og Tænkningens genetiske Proces.

Erkendelsens Principer udspringe altsaa af Analogi og anvendes
paa Forandring før end paa Hvile. Et tredie Punkt er
allerede antydet: Principernes Sandhed bestaar i deres Gyldighed,
og deres Gyldighed i deres \textit{Anvendelighed,} deres Evne til
at sætte Forskningen i Bevægelse. Et Princips Betydning bestaar
i, at det fører til at stille aldeles bestemte Spørgsmaal.
Saaledes bestaar Energisætningens Betydning i, at alle fysiske
Fænomener kunne betragtes som Exempler paa Energiomsæt%
% --- p. 85 ---
% running head: Gyldighed = Anvendelighed.
\opage{85}ning, saa at man ved etvert nyt Fænomen kan spørge: hvilken
Energiomsætning foregaar her? hvorfra og hvortil? hvilke ere
Omsætningens Betingelser? Saaledes har Antagelsen af Fjernvirkning,
selv om den --- hvad Maxwell selv tror --- skulde
være urigtig, havt den store videnskabelige Betydning, at man,
forat besvare det Spørgsmaal, hvorledes Legemer kunne virke
paa hverandre i Afstand, med Iver har undersøgt det mellemliggende
Mediums Egenskaber. --- Man mindes her William
Hamiltons Sætning, at en levende Usandhed er bedre end en
død Sandhed. --- Maxwell anvender endogsaa denne Synsmaade
paa selve Atomlæren, den mekaniske Naturopfattelses Forudsætning.
Efterat have vist, at ingen af de Aarsager, vi kalde
naturlige, kunne forklare Atomernes Existens og navnligt deres
Ensartethed, udtaler han: »Er det da saa, at vore videnskabelige
Spekulationer allerede nu ere naaede ned under Tingenes
synlige Fremtræden, hvor Opstaaen og Forgaaen raade, og staa
ved Indgangen til den Verden af Orden og Fuldkommenhed,
der den Dag idag er, som den skabtes, fuldkommen i Tal,
Maal og Vægt? Vi tage maaske fejl. Ingen har endnu set
eller befølt et Atom, og vor Atomhypotese vil maaske blive
fortrængt af en ny Teori om Materiens Konstitution. Dog er
Forestillingen om utalte enkelte Ting, der ere ens og uforanderlige,
ikke opstaaet i den menneskelige Bevidsthed uden at bære
Frugt.« Altsaa: selv om Atomlæren ikke er »rigtig«, kan den
meget vel for et langt Tidsrum have været nyttig. Paa analog
Maade udtaler Maxwell sig om den Slutning, at da der stadigt
foregaar en Spredning af fysisk Energi, maa der have været
en Begyndelsestilstand, som ikke kan være bleven til paa naturlig
Maade. Det er ikke sagt, bemærker han, at denne Hypotese
er Videnskabens sidste Ord, --- og desuden betyder spredt
Energi egentligt kun Energi, vi ikke kunne tænke os nogen
Brug for. Det gaar her, tilføjer han, som naar vi tale om Forstyrrelser
(disturbations), der hindre en Naturlov i at træde
klart frem: det er i Virkeligheden Omstændigheder, vi ikke
kende, eller have overset, eller maa gemme til senere Undersøgelse;
»Forstyrrelse« er en Indbildning, ikke et Naturfaktum;
% --- p. 86 ---
% running head: Om Grænser for Videnskaben.
\opage{86}i Naturens Virken er der ingen Forstyrrelse. --- Man tror her
at høre Spinoza tale.

Dog er der andre Steder, hvor Maxwell ikke er saa metodisk
og konsekvent. Som hos ikke faa Naturforskere ligge undertiden
hos ham de teologiske Forestillinger saaledes paa Springet,
at de træde op, naar han mener at staa ved en absolut
Grænse, eller maaske rettere, at de forlede ham til at mene at
staa ved en saadan Grænse. Ligesom Descartes udledte de
højeste Naturlove af Guds Villie, --- ligesom Newton erklærede.
at Solsystemets Sammenhæng ikke kunde have nogen naturlig
Oprindelse, --- og ligesom Linné og Cuvier erklærede det Samme
om de organiske Typer, og Lord Kelvin i den nyeste Tid om
selve Livets Opstaaen\textsuperscript{28}), saaledes lærer Maxwell nogle Steder
(som ikke kunne forenes med det ovenfor anførte Sted), at de
ensartede og uforanderlige Atomer, der ganske have Præget af
Fabrikata, maa være skabte; vi ere da her ved Grænsen for al
Videnskab, tilføjer han, ti Videnskaben kan ikke diskutere Skabelsen
af Intet. Tilblivelsen af et Atom er en Begivenhed, der
ikke kan have fundet Sted indenfor den Naturorden, vi nu leve
under. Men en saadan Begivenhed maa have fundet Sted, da
det vilde være absurd at tænke sig absolut ensartede Elementer existerende fra Evighed af. --- Hvori dette Absurde egentligt
skulde ligge, siger Maxwell ikke, og hans Atomlære har
her aabenbart en mere dogmatisk Karakter end i andre Udtalelser,
der bedre stemmer med denne fine Forskers intellektuelle
Energi.

Ogsaa paa et andet Punkt fandt Maxwell en Grænse for
Videnskaben.

Skønt man, siger han, længe har kaldt Sjælen immateriel,
laa dog lige til den nyeste Tid den Forudsætning til Grund, at
man vel ikke ad anatomisk eller kemisk Vej (ved Dissektion
eller ved Vejning) vilde kunne støde paa Sjælen ved en Undersøgelse
af Hjernen, men dog ved at gaa tilstrækkeligt langt
tilbage vilde støde paa en materiel Bevægelse, der ikke havde
nogen materiel Aarsag, men direkte var bevirket af en Sjæl.
Denne Antagelse strider ifølge Maxwell mod Loven om Energiens
Bestaaen, og han paastaar, at man ikke vilde stole paa
% --- p. 87 ---
% running head: Ernst Mach.
\opage{87}et Experiment, der viste en betydelig Forskel mellem det af
et levende Væsen udførte Arbejde og den modtagne og forbrugte
Energi. Men hvorledes kan Bevidstheden da være bleven
til? Det nytter ikke at tillægge Atomerne psykiske Egenskaber,
ti Bevidstheden kan ikke forklares som et Produkt af
psykiske Atomer. Naar vi komme til Personlighedens Dybder,
ere vi ude over Videnskabens Grænse. De objektive Fakta høre
op. Den allerførste Sætning, som gælder om Subjektet: »Jeg
er !« kan ikke tages af to Personer i ganske samme Betydning
og kan derfor slet ikke blive Genstand for Videnskab.

Til denne sidste Udtalelse maa bemærkes, at naar man
kan opdage Forskellen mellem to Personers Brug af samme
Tanke, maa man kunne iagttage og sammenligne deres Tanker,
og derved bliver da allerede en empirisk Psykologi indledet.
Ogsaa her stanser Maxwell for tidligt, hvad der er saa meget
des mærkeligere, som han selv i en tidligere anført Udtalelse
beraaber sig paa psykologisk Iagttagelse.

2. Hos \textsc{Ernst Mach} (f. 1838) virkede en meget tidligt
vakt filosofisk Interesse sammen med naturvidenskabelig Forsken
og især med Studiet af Naturvidenskabens Historie til at stille
Erkendelsesproblemet paa en ny og ejendommelig Maade. Denne
Udviklingsgang lægger sig i det Ydre for Dagen i hans Lærervirksomhed,
idet han først en Aarrække var Professor i Fysik
(i Prag) og senere blev Professor i Filosofi (i Wien). Hans
Værker angaa især Naturvidenskabens Historie \textit{(Die Geschichte
und die Wurzel des Satzes von der Erhaltung der Arbeit.} 1872.
--- \textit{Die Mechanik in ihrer Entwickelung.} 2. Aufl. 1889. ---
\textit{Die Principien der Wärmelehre.} Historisch-kritisch erläutert.
1896). I \textit{Die Analyse der Empfindungen} (1886, 4. Aufl. 1903)
har han givet et rent filosofisk Værk.

Mach var ikke mere end 15 Aar, da han i sin Faders
Bibliotek fik fat i Kants' »Prolegomena«. Dette Skrift gjorde,
efter hvad han senere har udtalt, et vældigt Indtryk paa ham,
og han betragtede det som en Lykke, at han saa tidligt var
bleven bekendt med dets Tankegang. Fuldstændig Kantianer
vedblev han ikke at være; saaledes fandt han »Ding an sich«
overflødig. Desuden virkede senere Studiet af Spinoza, Herbart
% --- p. 88 ---
% running head: 888 Videnskab som økonomisk Beskrivelse.
\opage{88}og Fechner paa ham. Det gjaldt for ham, som han har udtrykt
det, at indtage \textit{et Standpunkt, han ikke behøvede at opgive, naar
han gik over fra Fysiken til Psykologien.} Alle Videnskaber
maa jo dog tilsidst udgøre et Hele. Og Mach saa klart, at
Atomlæren, saaledes som den i Almindelighed opfattes, ikke
muliggør Bevarelsen af samme Standpunkt ved Overgangen fra
Fysik til Psykologi. Heldigere stillet er man, naar man fra
første Færd opfatter alle Legemer som Tankesymboler for Fornemmelseskomplexer.
Verden bestaar da ikke af gaadefulde
Væsener i Vexelvirkning med et gaadefuldt Jeg; Farver, Toner,
Rums- og Tidsforhold o. s. v. ere kun de sidste Elementer,
hvis Sammenhæng det gælder at udforske, og som vi sammenfatte
og afgrænse paa den i praktisk og teoretisk Henseende
mest hensigtsmæssige Maade. Der gives kun en stor sammenhængende
Masse af Fornemmelser, stærkest forbunden i det,
vi kalde vort Jeg, som derved staa i Modsætning til det, vi kalde
Verden. Men Jeg, Verden, Aand, Materie og lignende Begreber
ere kun Tankesymboler, som betegne Grænser og Forskelle, det
har vist sig mere eller mindre hensigtsmæssigt at fastholde.

Denne Opfattelse har Mach allerede udtalt i sine Forelæsninger
over Psykofysik i 1863. Sansefysiologiske Studier fik
overhovedet en stor Indflydelse paa hans Tankegang. De overbeviste
ham saaledes om, at Rumsanskuelsen er bunden til
Sanserne og derfor kun gælder for, hvad der kan sanses, saa
at man ikke har Ret til at tillægge Atomerne, der ikke kunne
sanses, rumlige Egenskaber.

Videnskabens Opgave er at give en \textit{økonomisk} Fremstilling
af Kendsgerningerne, d. v. s. en Fremstilling, i hvilken
man kun benytter de strengt nødvendige og de mest hensigtsmæssige
Forestillinger. I denne Opfattelse af Videnskabens Opgave
erkender Mach, at han har en Forgænger i Maxwell, hvorimod
han mener at have Prioritet for Kirchhoff, der hyldede
en beslægtet Opfattelse\textsuperscript{29}).

Denne Betragtningsmaade af Videnskabernes Natur stemmer
godt med Udviklingshypotesen. Mach gør opmærksom paa,
at Spencer allerede før Darwin betragtede de sjælelige Processer
som \textit{Tillempelser} efter det Givne. Selv har han, støttet til
% --- p. 89 ---
% running head: Det Usædvanlige forstaas ved det Sædvanlige.
\opage{89}Videnskabernes Historie, udviklet dette biologiske Synspunkt
for Erkendelsen i sin Tale \textit{Ueber Umbildung und Anpassung
im naturwissenschaftlichen Denken} (1884). Tillempelsen bestaar
dels i Sammenfatten, dels i Reduktion: dels i en Forbindelse
af det Mangfoldige under ét Synspunkt, dels i en Udskillelse
af det, der ikke er nødvendigt til Forstaaelse. Hvilke
Kendsgerninger der lægges til Grund, beror paa Bekvemmelighed,
Overlevering eller Vane. Men selve de Kendsgerninger,
der lægges til Grund for Orienteringen, kunne ikke »forstaas«.
Naar f. Ex. mekaniske Forhold synes os mere indlysende end
andre, kommer det blot af, at vi ere mere vante til dem. Vor
Forstaaelse bestaar stedse i, at usædvanlige Uforstaaeligheder
føres tilbage til sædvanlige Uforstaaeligheder. Naturligvis skal
i Videnskaben enhver tilvant Dom, vi tillægge Gyldighed paa
et andet Omraade end det, hvor den oprindeligt dannedes, prøves
med Hensyn til, om den passer paa det nye Tilfælde; ellers
bliver den en Fordom. Vor Erkendelse arbejder sig stadigt
frem gennem Kamp mellem Domme og Fordomme. Ligesom
de levende Væsener ikke strax danne et nyt Organ til en ny
Funktion, der er bleven nødvendig (som naar et Hvirveldyr
skal flyve eller svømme), men forsøge at benytte de Organer,
de allerede have, --- saaledes benytter Videnskaben en Omdannelse
af de ældre Ideer til Fremstilling af nye Erfaringer.
I Newtons »Tiltrækning« stikker der endnu Noget af den gamle
Forestilling om at søge det naturlige Sted. Hvad det gælder
om, er at tænke det Nye paa simpleste Maade og efter Sparsommelighedens
Princip.

Aritmetiken sparer den direkte Tællen, idet selve Tallet
udtrykker, at to Ordningsmaader kunne være identiske, selv
om det, der ordnes, er forskelligt. Allerede Maxwell havde,
som Mach selv gør opmærksom paa, lært, at Talbegrebets Gyldighed
beror paa Analogi. --- Geometrien undersøger forskellige
Maalestokkes indbyrdes Forhold og spare os derved direkte
Maaling. Synsrummet og Berøringsrummet ere forskellige; men
en Forskydning i det ene svarer til en Forskydning i det andet.
--- Kvantitative Ordninger ere simplere og mere omfattende
end kvalitative Ordninger, og de lette Oversigten og Behand%
% --- p. 90 ---
% running head: Analogiens Betydning.
\opage{90}lingen af store Grupper af Erfaringer. Saadanne Ordninger
muliggøres i Naturvidenskaben ved Hjælp af Begreber som
Kraft, Masse, Atom. Disse Begreber ere kun Tankemidler, hvis
Betydning ene bestaar i at genkalde økonomisk ordnede Erfaringer,
skønt de fleste Naturforskere tillægge dem Realitet
udenfor Tænkningen. --- Aarsagsforholdet betyder kun den
nøjeste Forbindelse. nogen Beskrivelse kan udtrykke. Det har
kun logisk, ikke fysisk Betydning. Det Samme gælder Kontinuitetsprincipet:
vi statuere Kontinuitet, hvor det viser sig, at
Forskelligheden voxer med Afstanden mellem Leddene i en
Række, medens den aftager, jo mere Leddene nærme sig til
hverandre, og tilsidst forsvinder for os. Fordelen ved at anvende
Kontinuitetsprincipet overalt, hvor det er muligt, er, at
vi endnu ved de mindste skelnelige Dele af vedkommende System
kunne anstille lignende Betragtninger som ved de større
Dele. Kun Erfaring kan afgøre, hvor langt dette kan fortsættes.


Tænkningens Økonomi under Tillempelsen efter Erfaringen
medfører, som det let ses af det Foregaaende, en stadig
Anvendelse af \textit{Analogier.} I sin Afhandling \textit{Die Aehnlichkeit
und die Analogie als Leitmotiv der Forschung} (i første Bind
af Ostwalds »Annalen der Naturphilosophie«) behandler Mach
dette Begreb. Analogi betegner et saadant Forhold mellem to
Begrebssystemer, at baade Forskellen mellem to tilsvarende Begreber
i de to Systemer og Overensstemmelsen i deres indbyrdes
Forhold tydeligt træder for Dagen. Den gør en enhedsmæssig
Opfattelse af uensartede Kendsgerninger mulig og er
derfor af stor biologisk og erkendelsesteoretisk Betydning, en
Betydning, om hvilken hele Naturvidenskabens Historie vidner.
Men Analogi er ikke Identitet, og Forskningens Midler maa
vel sondres fra Forskningens Resultater.

\textit{Den mekaniske Naturopfattelse} hviler paa en omfattende
Analogi mellem Massernes Bevægelse i Rummet og Tingenes
kvalitative Forandringer (i Temperatur, Elektricitetsforhold o. s. v.).
At denne Analogi i saa stor Udstrækning kan gennemføres,
giver ikke Ret til at antage, at alle fysiske Processer »egentligt«
ere mekaniske. De mekaniske Love kunne tjene som for%
% --- p. 91 ---
% running head: Fornemmelsernes Funktionsforhold.
\opage{91}male Mønstre og Fingerpeg. Legemernes Bevægelser i Rummet
ere de simpleste og anskueligste Processer, --- de, som det er
lettest at forfølge i vor Fantasi. Desuden har enhver fysisk
Proces, hvad enten det nu er Varme, Elektricitet, Klang o. s. v.,
en mekanisk Side, der fremtræder som Udvidelse, Zitren, Tiltrækning
o. s. v. Derfor kan man klart fremstille og tydeliggøre
fysiske Processer ved Hjælp af mekaniske Analogier, skønt
man ikke har Ret til at sige (med Wundt f. Ex.), at alle fysiske
Aarsager ere Bevægelsesaarsager. Rent mekaniske Processer
gives der slet ikke, da magnetiske og elektriske Virkninger
og Varmeprocesser stedse gøre sig gældende foruden den
blotte Bevægelse. Enhver Proces hører egentligt til alle Fysikens
Omraader. ---

Denne Opfattelse af vor Erkendelse er ingen Opløsning af
Videnskaben. Den berøver os ingen virkeligt værdifulde Synspunkter,
kun de overflødige. Og den fjerner det Dogmatiske,
der har indsneget sig hos mange Forskere.

Den Verdensanskuelse, som bliver mulig paa Grundlag af
en saadan Erkendelsesteori, kan ikke være Materialisme. Fornemmelsernes
kvalitative Beskaffenhed kan ikke afledes af rent
kvantitative Forhold og Forbindelser, det Psykiske ikke af det
Fysiske. Fysik og Psykologi angaa i Virkeligheden det Samme,
nemlig Fornemmelseskomplexer; --- kun behandle de det fra
forskellige Synspunkter.

Men Verden er ikke blot en Sum af Fornemmelser. Der
gælde bestemte Funktioner mellem disse, og Erkendelsen af
disse Funktionsforhold er en Erkendelse af Virkeligheden. Forskellen
mellem Idealisme og Realisme er teoretisk set ligegyldig.
Tilværelsens Elementer kalde vi Fornemmelser, naar vi
betragte dem som Elementer i en psykisk Verden. Videnskabens
Opgave er kun den at finde den lovmæssige Sammenhæng i
Alt, hvad der sker; den tilstræber derfor nødvendigvis en \textit{Monisme.} ---


Maaske kommer Mach her for let fra store Vanskeligheder.
De for Fysik og Psykologi fælles »Elementer« staa noget ubestemt
og mystisk, som en Taagemasse, der ikke er formet og
artikuleret. Det er vel saa, at Forskellen mellem Fysik og
% --- p. 92 ---
% running head: Heinrich Hertz.
\opage{92}Psykologi skyldes en Arbejdsdeling, der har vist sig hensigtsmæssig.
Men en saadan Arbejdsdeling er ikke let at ophæve
igen; den faar Karakteren af en »Fordom«, der gør Modstand,
saalænge ikke aldeles bestemte Erfaringer nødvendiggøre en ny
intellektuel Tillempelse. Og selve den Kendsgerning, at Arbejdsdelingen
har vist sig nødvendig og hensigtsmæssig, tyder paa
kosmologiske Betingelser. Tilværelsen kan neppe være saa
simpel, som Mach forudsætter, naar forskelligartede Synspunkter
ere nødvendige for at forstaa den. Monismen maa vindes
gennem Kamp med større Modsætninger end dem, Mach anerkender.


Dertil kommer, at Forholdet mellem Elementerne og deres
funktionelle Relationer henstaar uklart. Det kan neppe være
ligegyldigt for Elementerne, at de netop staa i disse bestemte
Forhold, og omvendt maa vel Relationerne være betingede ved
de Elementer, for hvilke de gælde. Men hvori bestaar saa
Virkeligheden? I de blotte Funktionsforhold alene kan den
ikke bestaa. Her rejser sig da en hel Række Spørgsmaal om
det Enkeltes Forhold til Sammenhængen, Spørgsmaal, hvis Besvarelse
bliver af Betydning for den hele Verdensanskuelse.

At Mach ikke har forfulgt Problemerne saa langt, forringer
ikke Værdien af hans Arbejde, der har virket vækkende
og frugtbringende i sjælden Grad og synes at tiltrække sig stigende
Opmærksomhed.

3. I sine erkendelsesteoretiske Ideer slutter \textsc{Heinrich}
\textsc{Hertz} (f. 1857 i Hamburg, d. 1894 som Professor i Fysik i
Bonn) sig til Mach. Den geniale, saa tidligt bortgangne Forsker
er især berømt ved sin experimentale Paavisning af Identiteten
af Lys og Elektricitet, en Opdagelse, der bekræftede Ørsteds
Anelse, Faradays Ideer og Maxwells Beregninger. I en
Tale ved det tyske Naturforskermøde 1889 har han i klar og
tilgængelig Form gjort Rede for sine Opdagelser \textit{(Ueber die
Beziehungen zwischen Licht und Elektrizität).} Hans Lærer
Helmholtz har i Indledningen til tredie Bind af Hertz's Skrifter
givet en Karakteristik af den udmærkede Forsker og af
hans videnskabelige Udviklingsgang. Her have vi kun at gøre
med den Maade, paa hvilken Hertz opfattede sin Videnskabs
% --- p. 93 ---
% running head: Symbolistisk Erkendelseslære.
\opage{93}første Grundsætninger, og den Karakteristik, han giver af Erkendelsesarbejdet.
Han har udtalt sig herom i en klar og tankevægtig
Afhandling, Indledningen til hans »Principien der Mechanik«
(Werke III).

Videnskabens Opgave er at udlede det Fremtidige af det
Forbigangne. Forat kunne gøre dette danne vi os Billeder eller
Symboler (Scheinbilder oder Symbole) af den Beskaffenhed, at
deres tankenødvendige Følger stedse igen ere Billeder af de
forestillede Genstandes \textit{naturnødvendige} Følger. Erfaring viser,
at dette er muligt, hvad der tyder paa en Overensstemmelse
mellem Naturen og vor Aand. Men det maa vel fastholdes, at
vore Billeder kun i denne ene Henseende, nemlig Overensstemmelsen
mellem Tankenødvendighed og Naturnødvendighed, behøve
at stemme med det Virkelige. Om de ellers stemme med
det, kunne vi aldrig faa at vide.

Dersom det nu er muligt at danne flere, lige logiske og
konsekvent anvendelige Billeder, foretrække vi det simpleste.
Ganske vist vil det --- da vor egen Aands Natur stedse lægger
sig for Dagen i Billeddannelsen --- ikke være muligt at udelukke
unødvendige og uanvendelige Elementer. Simpelheden
kan stedse blive større, og Anvendeligheden maa stedse prøves
paany. Om Naturen selv følger de simpleste Veje, kan ikke
godtgøres.

Til sin Fremstilling af Mekaniken mener Hertz ikke at behøve
andre Symboler end Tid, Rum og Masse, hvorimod saadanne
Begreber som Kraft (Energi) og Atom ere overflødige.
Hverken Kræfter eller Atomer kunne vi iagttage. Hvad vi se,
er stedse Masser i Bevægelse. Nu frembyder det, vi se, ikke
fuld Lovmæssighed og Sammenhæng. Vi maa derfor antage,
at den virkelige Verdens Mangfoldighed er større end Mangfoldigheden
af den Verden, der aabenbarer sig for vore Sanser.
Ville vi have et afrundet, lovmæssigt Verdensbillede, maa
vi bag de Ting, vi se, formode andre, usynlige Ting. Men det
Simpleste vil da være at forestille os dette Skjulte som Masser
i Bevægelse, kun forskellige fra de synlige Masser ved at være
utilgængelige for vore Sanser og vore sædvanlige Iagttagelsesmidler.
Om Kraft (Energi) tale vi kun, naar det gælder at
% --- p. 94 ---
% running head: Kan Kraftbegrebet elimineres ?
\opage{94}bestemme Forholdet mellem to Bevægelser; ved Kraft forstaas
dels Aarsag til en Bevægelse, dels Virkning af en Bevægelse:
det er altsaa et afledet Begreb. Ogsaa Begrebet Atom er en
ren Abstraktion; hvad der er givet, er stedse kun et materielt
System, aldrig det enkelte Massepunkt. ---

Det maa selvfølgeligt være Mekanikens Sag, om den kan
nøjes med de Forudsætninger, Hertz stanser ved. Her interesserer
det os at se, hvorledes en fremragende Naturforsker
har søgt at gøre sig Rede for sit intellektuelle Arbejde. At
denne Redegørelse er bestemt ved hans eget specielle Arbejdsomraade,
kan jo ikke være Andet. Hertz's store Opdagelse
var, at de elektriske Svingninger frembyde ganske samme Forhold
som Lyset og Varmestraalerne, saa at der i Rummet maa
foregaa langt flere Bevægelser, end vore Sanser direkte vise
os. De usynlige Masser og Bevægelser kom derved ganske naturligt
til at spille en Hovedrolle i Hertz's videnskabelige
Fantasi; i denne Henseende havde han en stor Forgænger i
Faraday, der allerede stræbte at udelukke Alt, hvad der ikke
kunde anskues, og derfor tvivlede om det tomme Rums og
Fjernkraftens Realitet. --- Om det er muligt ganske at udelukke
Kraftbegrebet (Energibegrebet), er tvivlsomt; og egentligt vil
Hertz ogsaa kun nedsætte det til et underordnet, afledet Begreb.
Maaske vilde han være kommen til at se anderledes paa
Sagen, dersom han var gaaet nærmere ind paa Aarsagsproblemet.
Aarsagssætningen opstiller han blot som en nødvendig
Forudsætning, --- i Form af den Grundantagelse, at der bestaar
faste, lovmæssige Forhold mellem materielle Systemers forskellige
Stillinger. Men skarpt og punktuelt stilles Aarsagsproblemet
først, naar vi stanse ved det enkelte Øjeblik, i hvilket den
forudgaaende Tilstand er ved at forsvinde og den følgende
endnu ikke indtraadt. Her betegner Begrebet Kraft en Forventning
om, hvad der vil ske, en Tillid til, at det Tilkommende
vil staa i et lovmæssigt Forhold til det Forbigangne.
Hertz siger ogsaa selv, at vi tale om Kraft, \textit{medens} den Proces,
ved hvilken vi udlede de fremtidige Erfaringer af de forudgaaende,
staar paa. (Einleitung p. 14). Men i en vis Forstand
leve vi stedse i saadanne Øjeblikke --- og »Hjælpebegrebet« er
% --- p. 95 ---
% running head: Wilhelm Ostwald.
\opage{95}der saaledes stadigt Brug for, netop jo mere intensivt det intellektuelle
Arbejde er\textsuperscript{30}).

Fra Mach adskiller Hertz sig ved mere at betone det
Symbolske end det Økonomiske i vore Grundbegreber. Det
er Trangen til anskuelige Billeder, der især gør sig gældende
hos ham, og naar han helst vilde undvære Begreber som Kraft
og Atom, var det maaske, fordi de ikke gjorde klar Anskuen
mulig. Enten har Hertz været udpræget visuel, eller ogsaa har
det for ham været bestemmende, at visuelle Symboler have det
Fortrin for de motoriske (muskulære) Symboler, at de tilstede
direkte Maaling. Kraft er et Symbol, der viser tilbage til Erfaringen
om Muskelanspændelse; men denne kan kun maales
gennem selve Bevægelsen, den frembringer\textsuperscript{31}).

Af størst Interesse forekommer det mig at være, at Hertz's
Definition af de videnskabelige Symbolers Sandhed nøje stemmer
med den eneste Definition, man kan give af vore Sansekvaliteters
Sandhed, og som hans Lærer \textsc{Helmholtz} klart havde
fremsat i sin »Physiologische Optik« (§ 26) i følgende Ord:
»Forsaavidt vor Fornemmelses Kvalitet underretter os om Ejendommeligheden
af den ydre Indvirkning, ved hvilken den er
fremkaldt, kan den gælde som \textit{Tegn} paa denne, men ikke som
\textit{Afbillede} af den ... Et Tegn behøver ikke at have nogen
Art af Lighed med det, hvis Tegn det er. Forholdet mellem
dem er kun det, at samme Objekt under samme Betingelser
fremkalder samme Tegn.... Vi kalde vore Forestillinger
om Yderverdenen \textit{sande,} naar de give os tilstrækkelig Anvisning
om Følgerne af vore Handlinger i Yderverdenen, og lade
os drage rigtige Slutninger om de Forandringer, der kunne
ventes i denne.«\textsuperscript{32})

4. I en ejendommelig Modsætning til Hertz staar \textsc{Wilhelm}
\textsc{Ostwald}, den berømte Kemiker, idet han netop vil gøre
Energien, det Begreb, Hertz vil udskyde, til Grundbegrebet.
Maaske hænger denne Modsætning sammen med, at Ostwald
er motorisk, medens Hertz synes at være visuel. Men det kan
ogsaa hænge sammen med deres Studiefag, idet Kemikerne have
større Anledning end Fysikerne til at tænke paa Energiforholdene,
medens Fysikerne have lettere ved at gennemføre visu%
% --- p. 96 ---
% running head: Alt er Energi !
\opage{96}elle Skemaer end Kemikerne, for hvem de faa en mere udpræget
hypotetisk Karakter.

Ostwald (f. 1853, fra 1887 Professor i Kemi i Leipzig)
fører en Krig mod al Materialisme, idet han vil reducere al
Materie til Energi. Denne Tanke kommer første Gang frem i
hans Tale paa det tyske Naturforskermøde 1895 \textit{(Die Ueberwindung
des wissenschaftlichen Materialismus),} og den er udførligt
udviklet i hans, til Mach dedicerede \textit{Naturphilosophie}
(1902). Ostwald har fremdeles søgt at virke for Drøftelse af
de Spørgsmaal, der ligge paa Grænsen mellem Naturvidenskab
og Filosofi, ved at grundlægge et nyt Tidsskrift \textit{Annalen der
Naturphilosophie.}

Alt er Energi, og der existerer ikke Andet end Energi!
Det er Ostwalds Hovedsætning. Alle Materiens Egenskaber
kunne opløses til Energi: Masse er Kapacitet for Bevægelsesenergi;
Rumsopfyldelse er Volumenergi; Vægt er en vis Art
Stillingsenergi; de kemiske Egenskaber ere forskellige Arter af
den Energi, som fremtræder ved Stoffernes Forvandlinger. Energi
er Arbejde, eller Alt, hvad der opstaar af Arbejde eller forvandles
til Arbejde. Det er det mest omfattende Begreb, Videnskaben
hidtil har dannet. Det omfatter baade Bestaaen (Substantialitet)
og Virken (Kausalitet). Hvis man ved Substans mener det, der
bestaar, er Energien den sande Substans. Den udstrømmer fra
Solen over Jorden, anvendes umiddelbart eller ophobes som
kemisk Energi for siden at forbruges i Livsvirksomheden.

Begreber som Materie og Legeme have vi kun Anledning
til at danne, naar forskellige Energiens Forbindelser bevirke
en sammensat Ligevægtstilstand. Skal der \textit{ske} Noget, maa der
være Styrkeforskelligheder til Stede mellem de forskellige Energier.
I Virkeligheden gives der kun Tilnærmelser til fuldstændig
Ligevægt. --- Sætningen om Materiens eller de materielle
Elementers Bestaaen betyder kun, at Elementerne stedse kunne
frembringes paany efter enhver Sammensætning, i hvilken de
ere indtraadte.

Da Ostwald har villet give en Naturfilosofi, ikke blot en
Erkendelsesteori, søger han at bringe Bevidsthedslivet ind under
samme Synspunkt, som det, fra hvilket han betragter hele den
% --- p. 97 ---
% running head: Fysisk og psykisk Energi.
\opage{97}organiske og uorganiske Natur. Ogsaa Bevidsthed er efter ham
Energi. At vi forestille alt Ydre som Energi, finder netop sin
Forklaring ved, at alle Bevidsthedsprocesser ere af energetisk
Art. Dette søger Ostwald at godtgøre ved en Undersøgelse af
Opmærksomhed, Erindring, Sammenligning og Villen. Og han
henviser til Bevidsthedens syntetiske Karakter, som Kant saa
stærkt fremhævede, og som ytrer sig i, at alle psykiske Oplevelser
staa i indre Sammenhæng som opstaaede i en og samme
Hjerne eller Aand. Bevidsthedsenergien ytrer sig i denne Sammenfatten.
Da vore Forestillinger om Yderverdenen ere bestemte
ved vor Bevidstheds Natur, er det intet Under, at de
bære et lignende Præg som Bevidstheden selv. Bevidsthedsfænomenernes
Betydning bestaar i, at de gøre Opsamling af
Erfaringer mulig, saa at tidligere Oplevelser kunne benyttes til
Sammenligning med nye.

Hermed mener Ostwald at have løst Problemet om Forholdet
mellem det Fysiske og det Psykiske. Men naar han
definerer aandelig Energi som »ubevidst \textit{eller} bevidst Nerveenergi«,
er det klart, at Problemet stedse stiller sig paany, idet
der maa spørges, hvorledes Overgangen sker fra ubevidst til
bevidst Nerveenergi. Det nytter ikke at svare, at det sker ved
Opmærksomheden, ti dersom denne forudsættes forbunden med
Bevidsthed, bevæge vi os i en Kreds, og dersom den ikke forudsættes
bevidst, er det ikke godt at forstaa, hvorledes Bevidsthed
kan opstaa ved Samvirken af to ubevidste Processer.
Begrebet Nerveenergi henstaar hos Ostwald som et rent mystisk
Begreb.

Dette hænger sammen med, at Ostwalds Energilære først
og fremmest er bygget paa Erfaringer fra Omverdenen. Energien
definerer han som det, der følger af Arbejde og omsættes
til Arbejde. Men det Arbejde, her er Tale om, er Overvindelse
af Modstand mod en Bevægelse. Kun tilsyneladende
kan Ostwald derfor udelukke de geometriske Egenskaber fra
sit Energibegreb; de ere ligesaa oprindelige som de rent dynamiske.
Masse og Energi ere korrelative Begreber, der gensidigt
maa bestemmes ved hinanden. Massen kendes ved det
konstante Forhold mellem Energi og Hastighedsforøgelse, og
% --- p. 98 ---
% running head: Richard Avenarius .
\opage{98}Energien kendes, naar man forudsætter Massen og Hastighedsforøgelsen
kendte. Hastighedsforøgelsen kendes ved Maaling af
Rums- og Tidsforholdene. Vi komme altsaa ved det fysiske
Energibegreb ikke ud over det Rumlige. Ligesaa lidt som det
lykkedes Hertz og Mach at udelukke Kraftbegrebet, lykkedes
det Ostwald at udelukke Materiebegrebet. --- Det er hidtil hverken
lykkedes ham eller nogen Anden at opstille et Energibegreb,
der kunde lægges til Grund baade for Begrebet om
den psykiske Energi, der ytrer sig i Sammenfatten og Udsondren
af Bevidsthedselementer, og for Begrebet om fysisk
Energi, der ytrer sig i Overvindelsen af en Modstand i Rummet.

\chapterhead{II.}{PROBLEMERNES NATURHISTORIE.}

Ad filosofisk Vej kom \textsc{Richard Avenarius} til en Opfattelse
af Erkendelsen, der er beslægtet med den, som Maxwell
og Hertz kom til ved at tænke over deres Videnskabs
Principer, og som Mach kom til ved Studiet af Naturvidenskabens
Historie.

Avenarius (f. 1843) studerede i Leipzig, hvor han af Fysiologen
Ludwig indførtes i den strenge mekaniske Naturopfattelse
i dens Anvendelse paa de organiske Fænomener, og af Filosofen
Drobisch i den herbartske Filosofi med dens Lære om
Forestillingerne som Udtryk for Sjælens Selvopholdelsestrang.
Studiet af Spinoza fik stor Betydning for ham ved at stille ham
overfor et storartet Forsøg paa streng systematisk Reduktion af
alle Tanker til en eneste. Hans første Skrift behandlede Spinozas
Filosofi (\textit{Ueber die beiden ersten Phasen des Spinozischen
Pantheismus.} 1868). Han søgte her paa skarpsindig Maade
at vise, hvorledes Spinozas System er blevet til ved Sammensmeltning
af tre forskellige Tankerækker, en religiøs, der udtrykkes
i Gudsbegrebet, en naturvidenskabelig, der udtrykkes i
Naturbegrebet, og en abstrakt metafysisk, der udtrykkes i Substansbegrebet.
De tre Begreber udtrykke tilsidst for Spinoza Et
og det Samme (deus = natura = substantia). --- Ved Reduktionen
af det Organiske til Mekanik, ved Opfattelsen af Forestillingerne
% --- p. 99 ---
% running head: Den rene Erfaring.
\opage{99}som Udtryk for Selvopholdelse og ved Studiet af det mest
gennemførte Forsøg paa en Identitetslære forberedtes hos Avenarius
den Opfattelse af Erkendelsen, han første Gang udtalte
i sin Afhandling \textit{Philosophie als Denken der Welt gemäss dem
Princip des kleinsten Kraftmasses.} Prolegomena zu einer Kritik
der reinen Erfahrung. (1876). Til Grund ligger her det
Faktum, at Bevidstheden ikke raader over en uendelig Forestillingskraft;
derfor maa den anvende Økonomi i sin Tænkning.
Den stræber da at føre det Ubekendte tilbage til det
Bekendte; det sker i al Genkendelse og i al Begriben. Og den
stræber overhovedet efter at reducere Forskelle og Mangfoldighed
til det mindst Mulige. Filosofi er netop den videnskabelige
Stræben efter at tænke Indbegrebet af det i Erfaringen
Givne med den mindst mulige Kraftanvendelse. Den stræber
efter at naa til en ren Erfaring ved at udelukke alle de Forestillinger
fra Erkendelsen, som ikke ere indeholdte i selve det
Givne; kun naar dette er naat, anvendes der ikke mere Kraft
paa at tænke det Givne, end dette selv fordrer.

1877 blev Avenarius Professor i Filosofi i Zürich, og her
udarbejdede han sit Hovedværk \textit{Kritik der reinen Erfahrung}
(1888---1890). Det er en Problemlære, der ad rent biologisk
og psykologisk Vej søger at paavise Betingelserne for Problemers
Opstaaen og Forsvinden. Avenarius vil overfor Problemerne
stille sig paa den psykofysiske Experimentators eller
Sindssygelægens Standpunkt og altsaa studere Problemerne rent
naturhistorisk. Han mener at kunne konstatere en stadig Tilnærmelse
til »ren Erfaring« ved stigende Udrensning af alle
Tankeelementer, som ikke indeholdes i selve det Givne. Det
er en Tilnærmelse til et rent empirisk Verdensbegreb. Hvorledes
et saadant Verdensbegreb vilde være beskaffent, har han
udviklet i sit Skrift \textit{Der menschliche Weltbegriff} (1891), som
dog især supplerer Hovedværket ved at søge Grunden til Erfaringens
»Urenhed« i Animismen og overhovedet i Trangen
til at opfatte Tingene i Analogi med os selv.

Avenarius' Værk er Frugt af alvorlig Tænkning og bygget
paa en Fylde af Iagttagelser og Exempler. Men han har
skadet dets Virkning ved Indførelsen af en meget kunstig og
% --- p. 100 ---
% running head: Problemernes Naturhistorie.
\opage{100}unødvendig Terminologi, der virker afskrækkende paa mange
Læsere. Der er noget Tragisk i, at hans betydelige Tankearbejde
derved hindres i at vinde den Anerkendelse, det fortjener.
Hans Helbred nedbrødes under hans ivrige Studier.
Forgæves søgte han Lindring paa Kursteder, saaledes paa Skodsborg,
hvor jeg --- efterat jeg tidligere havde besøgt ham i Zürich ---
paa Spadsereture med ham i Dyrehaven første Gang
indviedes i »den rene Erfaring«. Et Par Aar efter døde han
(1896). En sjelden Tankeenergi i Forening med kunstnerisk
Interesse og et aabent, vennesælt Sindelag er Hovedtrækkene
i det Billede, de, der kom i personlig Berøring med ham, bevare
i deres Erindring.

I Karakteristiken af Avenarius' Filosofi vil jeg ikke følge
hans egen Fremstilling, ej heller anvende hans Terminologi. Ved
gentaget Studium af hans Hovedværk er det blevet mig klart,
at hans Grundtanker kunne fremstilles paa en langt simplere
Maade, end han har gjort det, --- og at de kun derved komme
til deres fulde Ret.

Det er Problemernes Naturhistorie, Avenarius vil give, idet
han ad fysiologisk, psykologisk og historisk Vej vil vise, under
hvilke Betingelser de opstaa, skærpes, løses eller forsvinde.
Stoffet til denne Undersøgelse kan hentes allevegne fra, ti Barnet,
den Vilde, det praktiske Menneske, den religiøst Troende
have ligesaa vel deres Problemer som Videnskabsmanden og
Filosofen. Alle staa her ens, da det gælder om Noget, som
følger af de naturlige, almenmenneskelige Forhold. Forskellen
er kun den, at den videnskabelige Problemstilling og Problemløsning
nærmer sig mest til ren Erfaring, d. v. s. til Genfremstilling
af det Givne uden subjektive Tilføjelser.

Et Problem er altid --- hvor det saa opstaar --- Tegn paa
et Spændingsforhold (en »Vitaldifferens«) mellem Individet og
dets Omgivelser, et Spændingsforhold, der opstaar ved, at de
Paavirkninger, der komme fra Omgivelserne, enten kræve større
eller mindre Energi, end Individet raader over. --- Naar Paavirkning
(P) og Energi (E) ganske svare til hinanden (hvad vi
kunne udtrykke ved P=E), er den højeste Grad af Livsopholdelse
(»vitales Erhaltungsmaximum«) til Stede. Der finder Genkendelse
% --- p. 101 ---
% running head: Problemernes forskellige Udspring.
\opage{101}af Omgivelserne Sted; Individet føler sig hjemligt til Mode og
har Tillid til sine Iagttagelser og sine Forestillinger. --- Naar
der nu formedelst Forandringer i Omgivelserne kræves nyt og
større Arbejde af Individet end før, uden at Energien er forøget
eller dens Anvendelsesmaade ændret, opstaar der et Problem.
(Dette Forhold kan udtrykkes ved P > E). Individet
finder Afvigelser, Undtagelser og Modsigelser i det Givne; dette
staar som »noget Andet«, end det var før, maaske som noget
»Vidunderligt«; Trygheden er ophævet, og Individet føler sig
fremmed i Verden; Genkendelse er foreløbigt umulig. Ethvert
ægte Problem er en Art Hjemve og udløser en Higen efter at
ophæve Fremmedheden. Med stigende Kultur, Dannelse og
Kundskab stiger Muligheden for saadan Fremmedhed. --- Hvis
omvendt Energien stiger mere end Paavirkningerne, eller disse
tage af (hvad der kan udtrykkes ved E > P), opstaar et Problem
af den modsatte Grund: der mangler Anvendelse for
Energien; denne udløses da frit, i usædvanlige Retninger, som
ikke helt bestemmes ved det Givne; den higer mod det Fjerne;
Fare og Lidelse opsøges; det bliver Emancipations- eller Stormog
Trængselstider, overhovedet den praktiske Idealismes Tider.
--- Denne sidste Form har Avenarius meget lidet dvælet ved.
Som Repræsentanter kunde Rousseau og andre »subjektive«
Tænkere være anførte. Hume har i sin psykologiske Erkendelseslære
lagt stor Vægt paa dette Punkt, og --- som vi ville
faa at se --- Guyau's og Nietzsche's hele Filosofi er bygget paa
dette Forhold.

Det Forhold, Avenarius især holder sig til, er det, hvor
Paavirkningerne kræve en anden eller større Arbejdsydelse,
end der idetmindste foreløbigt svarer til Individets Energi
(P > E). I den nærmere Undersøgelse af, hvad dette Forhold
medfører, anser Avenarius det for mest videnskabeligt, d. v. s.
mest nærmende sig til ren Erfaring, at anlægge en fysiologisk
Betragtningsmaade, altsaa betragte de Processer, Problemstillingen
udløser, som Processer i Centralnervesystemet. Hvad et Individ
udsiger om sine Tilstande under Problemstilling, Problembehandling
og Problemløsning, betragtes da som Symptomer
paa, hvad der foregaar i Centralnervesystemet. Vi faa saaledes
% --- p. 102 ---
% running head: Den objektive og den subjektive Vitalrække .
\opage{102}to Rækker (»Vitalrækker«), en subjektiv: Individets Tilstande,
udtrykte i det Udsagte, og en objektiv: Forandringerne i Centralnervesystemet.
Hin er en matematisk Funktion af denne,
og den subjektive Række kan derfor betegnes som den afhængige
Vitalrække, den objektive Række som den uafhængige
Vitalrække. Det er blot af Hensigtsmæssighedsgrunde, at den
objektive Række staar som den »uafhængige«. Forholdet mellem
de to Vitalrækker kunde i og for sig vendes om, og Avenarius'
definitive Standpunkt er betegnet ved en mundtlig
Ytring af ham: »Ich kenne weder Physisches noch Psychisches,
sondern nur ein Drittes«. --- Ved nøjere Undersøgelse af hans
Fremstilling viser det sig dog at være en Illusion, naar han
mener at lægge den objektive Vitalrække til Grund. Første
Bind af hans Værk skal fremstille den, og andet Bind den
subjektive Række. Fremstillingen i første Bind har nu den
Ejendommelighed, at der ingen specielle Exempler findes;
disse komme først i andet Bind og ere alle af »subjektiv« Art.
Den største Rigdom af Kendsgerninger findes altsaa i den subjektive
Række, og det er kun en meget mager og rent skematisk
Fremstilling, der kan gives af den objektive Række; selve
dette magre Skema bliver endda kun forstaaeligt, naar de psykiske
Korrelata opsøges. Dette Forhold er karakteristisk for
Forholdet mellem Psykologi og Fysiologi og indeholder en Advarsel
mod at antage, at man naar større »Videnskabelighed«
ved at proklamere den »biologiske« Metode som den ene rette
ved Behandlingen af Spørgsmaal som dem, der her er Tale
om\textsuperscript{33}). ---

En Vitalrække er en Proces, ved hvilken der opnaas en
ny Ligevægt, efterat en tidligere Ligevægtstilstand er ophævet
ved en Problemstilling. Den har tre Stadier, der (naar vi holde
os til den subjektive Række) kunne betegnes som Trang, Arbejde
og Befrielse. Det første Stadium kan igen ytre sig paa
forskellig Maade: som Savn og Længsel, som Begær, som Tvivl,
som Angst, Smerte, Anger eller Skuffelse. Det maa erindres,
at der her er Tale om det praktiske Livs Problemer og Spændingsforhold,
saavel som om den rene Teoris. Det andet Stadium
kan fremtræde som Handlen eller Stræben, som Kamp
% --- p. 103 ---
% running head: Individuelle, sociale og universelle Løsninger.
\opage{103}eller jevn Virksomhed, som en Voven, hvor Personligheden
sættes ind, eller som en taalmodig Udreden af Betingelserne
for at løse en Opgave. Det tredie Stadium er Sejrsfølelsens og
Tilfredsstillelsens Tid.

Det tredie Stadium (Løsningen, Deproblematisationen) kan
naas paa forskellig Maade. Dersom Spændingen er mindre betydelig,
kan den ophæves ved Restitution, idet der enten igen
sker en Forandring i Omgivelserne, ved hvilken den gamle
Tilstand paany indtræder, eller Individet foretager en momentan
Tillempelse. Hvis Spændingen er større, maa der ske en Substitution
ved Vanens eller Øvelsens Hjælp eller ved Udvikling
af nye Tilstande eller Evner. Den fuldstændige Løsning kræver,
at der udvikles Funktioner, som svare til det Konstante
og Universelle i de Forhold, Omgivelserne kunne frembyde,
saaledes at Individet for Fremtiden bliver uafhængigt af de
væsentligste Forandringer. Ved Restitutionen fører Hjemveen
tilbage til det gamle Hjem; ved Substitutionen fører den til, at
der bygges et nyt Hjem. Det kan være, at denne Løsning er
rent individuel (som i en Livsanskuelse, den enkelte Personlighed
er ene om); men den kan ogsaa være social (som i en
Folkereligion), eller universel (som i Videnskaben). Den formelle
Logiks Indhold er de Former, i hvilke det har vist sig,
at Vitalrækker kunne komme til universel Løsning. Forskellighederne
ere her reducerede til det mindst Mulige, og det saakaldte
Identitetsprincip angiver egentligt blot en Fordring om
en saadan Reduktion, saa at der naas et »heterotisk Minimum«,
ɔ: den mindst mulige Forskelligartethed i Erfaringens Indhold.
Det gælder om en Udsondring af de Elementer, der kun have
individuel og forbigaaende Betydning, og som skrive sig fra
Arvelighed eller Tradition, med Fastholden af de Elementer,
som stadigt finde Anvendelse. Kaldes hine Elementer \textit{α}, disse
a, saa gælder det i y = f (\textit{α}, a) at faa \textit{α} til at blive o. Derved
gaar den individuelle eller sociale Løsning over til at blive
universel, og vi nærme os til den rene Erfaring. Dette sker,
jo flere Iagttagelser der foreligge, og jo mere de gennemtænkes.
Logisk holdbare ville tilsidst kun de Elementer være, der
have universel Betydning (altsaa a, ikke \textit{α}).

% --- p. 104 ---
% running head: Biologisk og logisk Holdbarhed .
\opage{104}Men der er Forskel paa logisk og biologisk Holdbarhed.
Det kan være, at Individer, Folkeslag og Tidsaldre ikke kunne
undvære visse Forestillinger, skønt disse ikke have universel
Gyldighed (altsaa afgjort høre til Klassen \textit{α}, ikke til Klassen a).
Der vil da bestaa »Beskyttelsesformer«, som unddrages Erfaringens
Kontrol. Herhen høre mange Elementer i individuelle Livsanskuelser,
i Folkereligioner, i filosofiske Systemer og i naturvidenskabelige
Hypoteser og Principer. Saadanne Beskyttelsesformer
ville naturligvis ikke være uforanderlige; det biologisk
Uundværlige vil variere med Forholdene. Maaske støttes deres
Berettigelse ligefrem paa deres Uundværlighed; deres Gyldighed
»postuleres« da.

Hvad der er Problem for den Ene, er det ikke altid for
den Anden, enten fordi denne mener at have en Løsning, eller
fordi Betingelserne for Problemets Opstaaen ikke ere til Stede
hos ham. Og hvad der er Problem for den Ene, er maaske
netop Løsning for den Anden (f. Ex. Skabelsen af Intet). Det
rent subjektive Symptom (»Vished«, »Selvfølgelighed«) er ingen
Garanti for, at en virkelig Løsning (en absolut Deproblematisering)
er naat, og at en ny Problemstilling (Problematisering)
ikke vil indtræde. Der kan stedse udløses nye Vitalrækker,
indtil de stadigt gentagne Paavirkninger føre til fuldstændig
Udskillen af de unødvendige og uholdbare Biforestillinger (\textit{α}),
saa at der ikke udløses anden og mere Energi, end der netop
behøves (a). Der foregaar stadigt en saadan Udskillelsesproces,
ved hvilken vi nærme os til et heterotisk Minimum. Den rene
Erfarings Standpunkt vil være karakteriseret ved, at Erkendelsen
væsentligt bestaar i en Beskrivelse, saavidt muligt ved Hjælp
af kvantitative, ikke kvalitative Bestemmelser, med stadig Paavisning
af Ækvivalenter og med Udledning af det Følgende fra
det Foregaaende. I praktisk Henseende foregaar en analog Udvikling,
idet det fasteste Samfundsliv viser sig at være det,
indenfor hvilket de sociale Forskelligheder ere reducerede til
det mindst Mulige. Derved bliver det muligt, at den enkelte
Del af Samfundet kan hævde sig --- ikke paa de Andres Bekostning,
men netop ved at arbejde for deres Opholdelse. ---

Avenarius indrømmer selv, at den rene Erfaring kun til%
% --- p. 105 ---
% running head: Økonomisk og symbolistisk Erkendelsesteori.
\opage{105}nærmelsesvis kan naas. Han burde sikkert have lagt langt
større Vægt paa dette Approximative; derved vilde hans Opfattelse
af Problemerne være bleven mere frugtbar. Ti det
kommer hos Avenarius let til at se ud, som om det er en Fejl
hos os, at vi ikke staa paa den rene Erfarings Standpunkt; og
det kan det dog ikke være, naar det er et Standpunkt, der aldrig
fuldt kan naas. Den rene Erfaring staar hos Avenarius
som den rene Ide hos Platon. Vor virkelige Erkendelse er
stedse i Arbejde paa at naa den rene Ide eller den rene Erfaring.
Men vi maa gaa et Skridt videre. Kan selv den rene
Erfaring undvære Hjælpebegreber (\textit{α}), naar den dog skal vedblive
at være Erfaring? Vi kunne dog ikke gøre os til Et
med »Omgivelserne«; disse vedblive, selv om vor Erfaring
bliver nok saa »ren« for subjektive Elementer, dog at »omgive«
os; Forskellen mellem dem og os falder ikke bort. Og allerede
det, at Noget er \textit{givet,} indeholder, at der er subjektive Elementer
til Stede. Avenarius har oprindeligt betegnet sin Opgave
som Forsøg paa at give en Filosofi som Tænken af Verden med
mindst muligt Kraftforbrug; det er et saadant Forsøg, der efter
hans Mening fører til »ren Erfaring«. Men Tænkningens Økonomi
fører nødvendigvis til Symbolisme, da man ikke kan godtgøre,
at det økonomiske Hensyn (Simpelhedens Princip) er et
Princip af rent universel, ikke delvis subjektiv Gyldighed. I
hvert Tilfælde fører Principet om Tænkningens Økonomi ikke
uden videre til slet ingen Hjælpebegreber at bruge, ligesaa lidt
som en økonomisk Husholdning er en saadan, der slet ingen
Gryder og slet intet Brændsel bruger. Her er et Punkt, hvor
Avenarius ikke til sin Fordel afviger fra de filosoferende Naturforskere,
som ovenfor ere omtalte; disse saa klart Sammenhængen
mellem Økonomi og Symbolisme, skønt de lagde forskellig
Vægt paa disse to Sider ved Erkendelsen. --- Hermed
hænger det ogsaa sammen, at Avenarius indtager et noget dogmatisk
Standpunkt i Forholdet mellem Fysiologi og Psykologi.
Han ser ikke, at erkendelsesteoretisk kende vi stedse den »uafhængige
Vitalrække« gennem den »afhængige Vitalrække«: vi
vilde ingen Forestilling have om Centralnervesystemet og dets
Tilstande, naar vi ikke overhovedet havde Evne til at danne
% --- p. 106 ---
% running head: Avenarius og Vurderingsprobelmet.
\opage{106}Forestillinger; og selv vore nuværende Forestillinger om Centralnervesystemet
skyldes Problemstillinger og Problemløsninger,
der endnu ikke tilfredsstille den »rene Erfarings« Fordringer.


Ogsaa i en anden Henseende lægger Avenarius ikke tilstrækkelig
Vægt paa, at vi ere stedte i en Udvikling, og at
Forestillingen om en ideal Afslutning kun kan have Betydning
som Maalestok for Udviklingens Gang og Retning. Han holder
sig overvejende til Tilfældet P > E, altsaa til saadanne Tilfælde,
hvor Arbejdskravet stilles udefra. Men en Betingelse for Livets
--- baade Erkendelses- og Villieslivets --- Fremgang er det, at
Arbejdskravet ogsaa kan blive stillet indefra, idet der er et
Overskud af Kraft, der søger Anvendelse (E > P). I hele Naturen
begynder Livet med et saadant Overskud; det er det,
der gør Væxt mulig og fører til nye Spirer. Netop naar Livet
rører sig kraftigt, higer det ud over det Givne, opsøger eller
fremkalder nye Erfaringer, sætter nye Begyndelser\textsuperscript{34}). Det sker,
naar der voves, og det sker særligt, naar der vurderes. Dersom
Avenarins, til hvis Planer det hørte at skrive en »Friluftsetik«,
var kommen til at behandle Vurderingsproblemet, vilde han
sikkert have lagt Vægt paa, at der i al Vurdering gaas ud over
det blot Givne. Her kunne de subjektive Elementer allermindst
udskilles, selv om de stadigt maa berigtiges gennem objektive
Kriterier. Hos Guyau og Nietzsche ville vi komme til at beskæftige
os med dette Spørgsmaal. ---

Disse kritiske Bemærkninger skulde kun bidrage til at udskille
nogle af de subjektive Biforestillinger (\textit{α}), der have hindret
den energiske Forsker i at se klart paa Sagerne. Hans
udmærkede Ide om en Problemernes Naturhistorie\textsuperscript{35}) vilde uden
dem --- og uden den afskrækkende Form, i hvilken den i hans
Bøger er fremstillet --- ganske sikkert have gjort større Indtryk
paa Tidsalderen. Det er en Ide, der vil blive tagen op
igen, og det Arbejde, der er sat ind paa den, vil ikke være
spildt.

% --- p. 107 ---
\opage{107}\grouphead{TREDIE GRUPPE.}{VURDERINGSFILOSOFIEN.}

De tre Forhold mellem Erkendelsens indre og ydre Betingelser,
som Avenarius skelnede imellem, kunne bruges til at
karakterisere Forskellen mellem de tre Grupper af filosofiske
Forsøg, der fremtræde i den nyeste Tid. Den systematiske
Gruppe er væsentligt karakteriseret ved, at der forudsættes Ligevægt
mellem Trang og Tilfredsstillelse eller mellem Evne og
Arbejde (altsaa E = P). Den erkendelsesteoretisk-biologiske
Gruppe lægger Hovedvægten paa, at der fra Tilværelsen kan
stilles andre og større Fordringer til vor Evne, end denne fuldt
ud kan tilfredsstille (altsaa E < P). Den tredie Gruppe betoner
endeligt, at der kan være et Overskud af Trang og Evne
(altsaa E > P), hvad der fører til en Vurdering af Tilværelsen
og til en Higen ud over det Givne. Som Repræsentanter for
denne sidste Synsmaade vil jeg her betragte Guyau og Nietzsche.
Det vilde være berettiget ogsaa at nævne andre moderne Vurderingsfilosofer
her, navnligt William James og Rudolph Eucken;
men dem har jeg nyligt i et andet Skrift \textit{(Om nogle religionsfilosofiske
Arbejder fra den nyeste Tid.} Universitetsprogram
1903) havt Lejlighed til at omtale.

Guyau og Nietzsche staa begge paa Udviklingslærens Grund
og haabe paa højere Livsformers Opstaaen. De bygge dette
Haab paa Livets Kraft og Sundhed, --- paa den Overbevisning,
at der er en overstrømmende Fylde af Energi, som ikke finder
% --- p. 108 ---
% running head: Den indre Energis Overvægt over de ydre Paavirkninger.
\opage{108}Udtryk og Afløb i vor nuværende Erfaring og vore nuværende
Livsforhold. De stille et stort Ja mod Pessimismens Nej,
--- den Ene i inderligt bevæget Stemning og mild Resignation
overfor Værdiernes Skiften og Forgængelighed, den Anden i
Trods og Foragt overfor Fortiden og i ubændigt, tilsidst krampagtigt
Haab om Fremtiden.

Et Hovedproblem hos begge er Forholdet mellem Instinkt
og Reflexion, --- mellem Livets udelte Energi paa dets tidligere
Trin og dets delte Virken under Kulturens og Reflexionens
Fremskriden. Siden dette Problem først toges op mod Slutningen
af det 18. Aarhundrede, af Rousseau, Lessing og Kant,
er det ikke blevet stillet saa energisk som af de to Forfattere,
vi her beskæftige os med. De polemisere begge mod den
ensidige Intellektualisme og støtte sig til Følelses- og Villieslivet
som det, der aldrig ganske kan udformes i klare, rationelle
Forestillinger. Medens det hos den anden Gruppe er den
udefra givne Fylde, der kan true med at overvælde Tanken,
er det her den indefra, fra Sindets indre Verden opstigende
Fylde, der driver Tanken til dens Grænse.

Hermed hænger det sammen, at deres Fremstillinger staa
paa Grænsen af Filosofi og Poesi. Stemning og Lidenskab virke
med paa ethvert Punkt, ikke altid til Fordel for Undersøgelsens
Klarhed og Konsekvens, men til Fordel for den literære,
eller endog agitatoriske Virkning, deres Skrifter øve. Deres
Tanker faa for en stor Del mere Betydning som Symptomer
paa, hvad der bevæger sig i Sindet hos dem, der ere stærkest
grebne af Livet i vore Dage, end som aktive Bidrag til Behandling
af Problemer. Dette gælder i ganske særlig Grad for
Nietzsches Vedkommende.

Endnu et Lighedspunkt er der: de ere begge syge, og
deres Tanker og Værker ere for en stor Del blevne til under
stadig Kamp med Sygdom. Dette afgør Intet med Hensyn til
Tankernes Værdi. Der kunde jo være Tanker, der kaste Lys
over Livet, og som netop kun kunde blive til under saadanne
Forhold. Grænseterritoriet mellem Sundhed og Sygdom kunde
jo være ganske særligt frugtbart. Vi maa da først og fremmest
drøfte Betydningen og Gyldigheden af deres Tanker, og først
% --- p. 109 ---
% running head: Jean Marie Guyau .
\opage{109}naar denne Prøvelse er foretagen, kunne vi benytte, hvad vi
vide om deres individuelle Forhold, til at forklare os deres
Aandsretning.

\chapterhead{I.}{JEAN MARIE GUYAU.}

Guyau er et Exempel paa tidlig Modning. Allerede tyve
Aar gammel (han blev født 1854) vandt han det franske Videnskabsselskabs
Præmie for et Arbejde over Nyttemoralens Historie
fra Epikur til vore Dage. Hans første Studier angik især
Platon og Kant, og han paavirkedes tillige stærkt af sin Stedfader
Fouillée. En glimrende Fremstillingsevne og en indtrængende
kritisk Energi vakte Opmærksomhed for hans Arbejde,
og lod det finde Anerkendelse selv hos dem, der bleve stærkt
kriticerede i det, som f. Ex. Spencer. Under de fortsatte Studier
nærmede han sig noget mere til de Synsmaader, han havde
angrebet, og han har i sine positive Fremstillinger ikke sjeldent
Vanskelighed ved at slippe forbi Indvendinger, han selv havde
opstillet overfor den engelske Skole. Det er æstetiske, etiske
og religionsfilosofiske Problemer, han især har beskæftiget sig
med. Nagende Tvivl om vore Værdiers Betydning i Tilværelsen
forfølge ham fra hans Ungdom og faa Udtryk ikke blot i
filosofiske Skrifter, men ogsaa i Digte \textit{(Vers d'un philosophe).}
Hans oprindelige, platoniserende Tro paa Verdens Rationalitet
blegnede under Reflexionens og Brystsygdommens forenede Indflydelse.
Vel synes det, at han senere igen fastholdt Muligheden
af Tilværelsens ideale Betydning, men det klare, filosofiske
Blik for Problemets Vanskelighed mistede han ikke. Man
har fundet et fransk Træk i den Maade, paa hvilken hos ham
Stemningens Inderlighed og Tankens Klarhed saa ofte gaa sammen.
Han véd, hvor mange af hans Syner og Forhaabninger
der ere Illusioner; men han fastholder dem, naar de blot ere
frugtbare Illusioner, der formaa at sætte Tanken og Villien i
energisk Arbejde. (Digtet Illusion féconde). Han ser Undergangen
klart for Øje, da hans Sygdom er skreden frem; men
han ønsker at lægge al sin Kærlighed til Mennesker ind i sin
% --- p. 110 ---
% running head: Kritik af den engelske Etik.
\opage{110}sidste Sang (La cicade). Neppe er der nogen anden Forfatter,
der som han har vist, hvilken Inderlighed og Højhed Livet kan
have, selv om alle Dogmer betragtes som Illusioner. Han
levede iøvrigt --- afset fra Sygdommen
--- under lykkelige Forhold.
Forskende og digtende, elskende og elsket nærmede han
sig Døden med fuld Bevidsthed om Sygdommens Udfald. Han
døde i Menton 1888. Paa hans Grav blev der sat følgende
Linier af et af hans Skrifter: »Vore højeste Bestræbelser synes
netop at være de mest unyttige; men de ere at ligne med Bølger,
der have kunnet naa til os og derfor ogsaa nok ville kunne
naa længere. Jeg er sikker paa, at det Bedste, jeg ejer, vil
overleve mig. Nej, ikke en af mine Drømme maaske vil være
spildt; Andre ville tage dem op, ville drømme dem igen efter
mig, indtil de engang gaa i Opfyldelse. Det er ved de stadigt
hendøende Bølgers Hjælp, at Havet formaar at forme sin Bred
og det uhyre Leje, i hvilket det bevæger sig.«

\sectionhead{A. KRITIK AF DEN ENGELSKE ETIK. \textit{(La morale\\ anglaise contemporaine.} 1879).}

Allerede i Guyaus kritiske Arbejde finde vi Antydninger
af hans definitive Opfattelse. Hvad han indvender mod den
engelske Skoles Etik, er især tre Ting.

Englænderne lære, at Moralen væsentligt skyldes klog Beregning
(Bentham), eller Forestillingsassociation (Hartley, de to
Mill'er), eller naturlig Udvælgelse (Darwin), eller Tillempelse
efter Livsforholdene (Spencer). De bygge Moralens Gyldighed
og Bestaaen paa bestemte indre og ydre Udviklingsforhold. Den
moralske Følelse er et Naturfænomen, som Mennesket kun betragter
som en objektiv Magt, fordi det ikke har mærket, hvorledes
det er blevet til. Den uvilkaarlige og ubevidste Væxt af
den moralske Følelse er den engelske Etiks Hovedemne. Men
her indvikler den sig nu, mener Guyau, i en Selvmodsigelse
og arbejder sig selv imod. Ti selve disse engelske Teorier maa
jo aabne Menneskets Øje for, hvorledes det hænger sammen
med den hemmelighedsfulde og moralske Følelses Natur. Den
% --- p. 111 ---
% running head: Trangen til Expansion og Virksomhed.
\opage{111}kommer til at staa som et Bedrag, man søger at befri sig for.
Reflexionen vil opløse, hvad der rent uvilkaarligt har formet sig.

Men --- er der nu her foregaaet en virkelig Udvikling?
Englænderne forudsætte som sidste Grundlag den egoistiske
Selvopholdelsestrang, og bliver den ikke ved at være den samme,
om der saa kommer nok saa mange Beregninger, Associationer
og Tillempelser til? Evnen til uinteresseret Hengivelse kan
ikke forklares gennem saadanne Faktorer. Derfor er det ikke
sagt, at uinteresseret Sympati er en oprindelig Tendens. Men
hvad Englænderne opfatte som oprindelig Egoisme, er ifølge
Guyau noget Andet og Mere, som ligger dybere end baade
Egoisme og Sympati. Det er en Trang til Virksomhed, til Udfoldelse
og Udvidelse, som er til Stede overalt, hvor Livet er
sundt og kraftigt. Denne Trang behøver ikke at gaa ud over
Andre, men kan tvertimod føre til Forbindelse med Andre og til
Arbejde for dem. Det er da ikke først gennem mange Mellemled,
men allerede i sin første Spire, at den menneskelige Trang
indeholder en Mulighed for Hengivelse til et mere omfattende
Liv end det rent individuelle. --- Guyau udvikler her en Tanke,
der allerede var kommen frem hos Vauvenargues og Rousseau,
uden at han dog synes at have den fra dem. ---

Naar der saaledes lægges Vægt paa den indre, af ydre
Foranledninger til Dels uafhængige Trang til Virken og Udvikling,
saa faar vor Stræben ikke --- som det i Nyttemoralen
bliver Konsekvensen --- Karakteren af at være et blot Middel.
Det kan da gælde, at vi leve for at ville og handle, ikke blot
omvendt. Il faut vivre pour vouloir et agir! Nutiden skal
hverken ofres for Fremtiden eller for Fortiden. I selve Villiesakten
ligger der et Formaal: il y a dans la volonté quelque
chose de définitif! ---

\sectionhead{B. ETIK. \textit{(Esquiss d'une morale sans obligation ni\\ sanction.} 1885).}

Til Grund for sin Etik lægger Guyau Begrebet Liv i Ordets
videste Betydning, hvorefter Modsætningen mellem det Bevidste
og det Ubevidste og mellem det Egoistiske og det Sym%
% --- p. 112 ---
% running head: Appel til det Ubevidste og Uvilkaarlige.
\opage{112}patiske falder bort. Livet ytrer sig som en Trang til Væxt,
Bestaaen, Forplantning, Udvidelse. Det er Biologien, ved hvis
Hjælp han vil raade Bod paa de Vanskeligheder, i hvilke Englænderne
indviklede sig ved deres psykologiske Begrundelse:
dog indrømmer han, at Darwin og Spencer egentligt allerede
have paavist et biologisk Grundlag. En Hovedsag er det for
ham, at Formaalene tilsidst ikke frembringes af Bevidstheden,
naar denne sætter sig Formaal: hvad der sker, er, at man
bliver sig bevidst eller gør til Genstand for sig, hvad der allerede
ubevidst har gjort sig gældende. Formaalene komme altsaa
ikke helt udefra.

Etiken bliver da Læren om de Midler, ved hvilke det af
Naturen selv satte Maal Livets Væxt og Udfoldelse --- kan
naas. Og det etiske Hovedbud bliver: Udfold dit Liv i alle
Retninger! Vær saa rigt et Menneske som muligt, baade hvad
Inderlighed og hvad Omfang af dine Bestræbelser angaar! Og
derfor --- lev i Samfund med Andre og uddan din Evne til
at leve i Samfund saa meget som muligt! --- Pligten er kun
en Overstrømmen af Liv, der kræver at blive brugt og at give
sig hen. Den skyldes ikke Tvang og ydre Nødvendighed, men
er Udtryk for en Kraft. I Sympati eller Altruisme træder
denne overstrømmende Kraft især frem. Egoismen derimod er
Tegn paa en Indsnevring, en Isolation, og den beror tilsidst
paa en Illusion. Den højeste og tilmed den i Naturen mest
grundede Dyd er Højsindet (générosité). ---

Guyau har søgt at undgaa at falde for den Kritik, han selv
har rettet mod Englænderne, ved at gaa tilbage til det Ubevidste
og Uvilkaarlige. Men han overser de forskellige Muligheder
--- og dermed de forskellige Problemer --- der ligge ved
Overgangen fra det Ubevidste til det Bevidste og fra det Uvilkaarlige
til det Vilkaarlige. Det er ikke sagt, at det Bevidste
udtømmer eller ganske svarer til det Ubevidste, og et lignende
Forhold er der mellem det Vilkaarlige og det Uvilkaarlige.
Midler forudsætte et Formaal, og Etiken kan ikke søge Midler,
naar den ikke véd, hvilket Formaal de skulle tjene. Hvad vi
gøre til Formaal, beror igen paa, hvad vi tillægge Værdi.
De middelbare Værdier (Midlerne) forudsætte de umiddelbare
% --- p. 113 ---
% running head: Expansionen og Eftertanken.
\opage{113}Værdier. Med disse maa man derfor begynde, og de kunne
ikke findes uden Psykologiens Hjælp. En empirisk Lære om
Værdier maa danne Grundlaget, og dersom der er flere fundamentale
Værdier, der ikke kunne reduceres til en enkelt, saa
ligger her en Vanskelighed for al Etik lige ved den første Begyndelse.
Det er kun, saalænge man holder sig i vag Almindelighed,
at en Appel til Biologien er tilstrækkelig. Saasnart der
stilles bestemte Spørgsmaal, dukker Værdibegrebet op. Det har
Guyau selv afgivet et Exempel paa. Overfor Bourgets Paastand,
at det er vilkaarligt og dogmatisk at skelne mellem naturlige
og unaturlige Følelser og Drifter, hævder Guyau (i sit Skrift
\textit{L'art au point de vue sociologique.} p. 375), at der meget vel
kan være en Forskel i Værdi, selv om de forskellige Tilstande
opstaa med lige stor Nødvendighed. Der gives, siger han, et
Kriterium for naturlig Værdi (valeur naturelle), som maa søges
i Livets Kraft og Omfang og i den Bevidsthed og den Lystfølelse,
som er den indre Aabenbaring deraf. --- Det er klart,
at selv om »Livet« var nok saa kraftigt og omfattende, --- naar
der ingen »Bevidsthed og Lystfølelse« var knyttet til det, vilde
der intet Værdibegreb kunne opstilles, og ingen etisk Diskussion
finde Sted\textsuperscript{36}).

Men den Trang, der fyldestgøres ved Livets Tiltagen i Kraft
og Omfang, er jo uvilkaarlig og instinktsmæssig; vil den da ikke
svækkes ved at drages frem som Genstand for Opmærksomhed
og Eftertanke? --- Det var den Hovedvanskelighed, Guyau havde
fundet hos Englænderne, og som han derfor maa se at befri
sin egen Teori for. Han udtaler, at Fornuften kun lammer et
Instinkt, naar den gaar i en anden Retning end dette, eller
naar det er en Fordel, at den træder i dets Sted. Men ingen
af Delene er Tilfældet her. Den uvilkaarlige Hengivelse og Udfoldelse
har sin store praktiske Betydning. Lysten til at vove
udspringer af Livets Overskud paa Kraft, der føder Haab og
Tro, og derved gør Noget muligt, som ellers ikke vilde være muligt.
I dristig Tankebevægelse, i begejstret Opofrelse ytrer denne
Livets Kraft og Sundhed sig. Den eneste Sanktion, som Fremtidens
Etik anerkender, vil være Tilfredsstillelsen ved at vove (le
plaisir de risque). --- Jeg kan ikke se, at Englænderne ikke
% --- p. 114 ---
% running head: Kunsten som Livsudfoldelse.
\opage{114}ligesaa vel kunde beraabe sig paa disse Grunde; faktisk gøre de
det ogsaa, og i højere Grad end Guyau, da deres Etik er mere empirisk
end hans. Hvis Klogskab, Association og Tillempelse ikke
kunde føre til noget Værdifuldt, vilde de ikke være de stadige
Veje, ad hvilke Udviklingen skrider frem. Og desuden: hvis
Reflexionen kan opløse Instinktet, kan den saa ikke ogsaa hæmme
Expansionen?

\sectionhead{C. ÆSTETIK. \textit{(Les problèmes de l'estétique\\ contemporaine.} 1884).}

Guyaus Æstetik frembyder et interessant Exempel paa,
hvorledes et Problem kan stille sig forskelligt efter de forskellige
Tiders Vilkaar. For tidligere Forskere, som Kant og Schiller,
og endnu langt senere for Spencer, stod det som Opgaven
at skaffe Plads for Kunsten midt i Livets Kamp og trods denne.
De toge deres Tilflugt til den samme Tanke, der spiller saa
stor en Rolle hos Guyau, idet de henviste til, at naar der kommer
et Overskud af Energi, benyttes det naturligt til lignende
Funktioner som dem, der anvendes i Livskampen, kun at det
nu er Billede og Leg, det gælder, ikke selve Livets Virkelighed.
Guyau indvender mod hele denne Retning i Æstetiken:
vi søge i Kunsten netop mere og fyldigere Liv, end det sædvanlige
Liv kan bringe; Kunsten er en Udvidelse af Livet, og
det er kun Skade, at vi maa ty til Billedet og Legen forat opnaa
den. Det Uvirkelige er ingenlunde en Betingelse for Kunstskønheden,
men netop en Begrænsning, den maa finde sig i.
Under kunstnerisk Skuen opstaar Ønsket om at blive til det,
vi betragte, og det er Digterens Sorg, at han ikke kan gøre sig
til Et med Livets hele Fylde (Digtet Le mal du poète). Derfor
staar efter Guyaus Opfattelse Naturskønhed over Kunstskønhed.


Hermed hænger det sammen, at Guyau paastaar, at alle
Sanserne virke med ved det Skønne. En Drik Mælk ved en
Bjerghytte kan virke som en hel Pastoralsymfoni. Alle Sider i
vort Væsen virke med --- gennem Sansning, gennem Erindring,
gennem Fantasi. Skøn kalder Guyau en Paavirkning, der paa
én Gang sætter Følelse, Tanke og Villie i Bevægelse og vækker
% --- p. 115 ---
% running head: Fremtidens Irreligion.
\opage{115}Lyst ved denne samlede Virken. Mere og mere vil al Nydelse
og Lystfølelse faa Skønhedens Præg, saa at den Forskel, der
paa et Mellemstadium af Udviklingen maa gøres mellem det Behagelige
og det Skønne, vil falde bort; denne Forskel betinges
kun ved, at det Dyriske endnu er saa stærkt i Mennesket.
Da ville ogsaa det Æstetiske og det Etiske falde sammen i Hævdelsen
af Livets naturlige Værdier. Tilsidst vil det vise sig, at
ingen Kunstner egentligt kan fremstille Andet end sit eget Liv
--- og Værkets Værdi vil bero paa Værdien af dette Liv!
--- Læserne af Guyaus Skrifter ville ved denne sidste Ytring tænke,
at den passer paa Guyaus egen Forfattervirksomhed, besjælet
som denne er af den store Eros, der var Sjælen i hans Liv.

\sectionhead{D. RELIGIONSFILOSOFI. \textit{(L'irreligion de l'avenir.} 1887).}

Guyaus mærkeligste og betydeligste Bog er den, der handler
om »Fremtidens Irreligion«. Denne Bog er paa én Gang
skreven med den mest indgaaende Kritik og med den inderligste
Følelse, med den fuldeste Forvisning om, at Religionens Tid er
forbi, og med den frejdigste Forhaabning om, at en høj og ideel
Livsanskuelse aldrig vil mangle Betingelser for at udvikle sig. Den
er i sin Titel rettet mod Forsøgene paa at konstruere en Fremtidsreligion,
men tillige mod »Antireligionen«, idet den vil hævde
den indre Sammenhæng mellem det religiøse Standpunkt og det,
der skal afløse det.

Det religiøse Problem hænger for Guyau nøje sammen med
det etiske og det æstetiske. Alle tre Omraader mødes for ham
i Livets Begreb. I det koncentrerede personlige og sociale Liv
fremtræde alle tre Elementer, og først derved komme de hvert
for sig til deres Ret. Religionen begynder vel som en Slags
mytologisk Fysik; men de intellektuelle Elementer ere dog ikke
de overvejende. Hvad der ligger til Grund, er Menneskets Erfaringer
i det Livsfællesskab, i hvilket det lever. Dette Livsfællesskab
udstrækker han uvilkaarligt til hele Tilværelsen, saa
Guder og Mennesker danne et stort Samfund. Religionen er
ikke blot Antropomorfisme, men særligt Sociomorfisme. Religionernes
Forskel beror især paa, hvilken Art af Samfund man tænker
sig Universet at være, en Familie eller et Rige? --- og i
% --- p. 116 ---
% running head: Det Evige i Religionen.
\opage{116}sidste Tilfælde: styret af en Fader eller en Tyran? --- Religionens
sociale Karakter faar sit mest udprægede Udtryk i Kultus,
gennem hvilken Mennesket træder i Samfund med Guderne.

Efterhaanden sublimeres og idealiseres baade den religiøse
Kultus og de religiøse Forestillinger. Indre Kultus tenderer til
at afløse ydre Kultus, og Mystiken afløser Mytologien. Højdepunktet
er naat, naar Gud staar som det personificerede moralske
Ideal. --- Som væsentlige Træk i al Religion fremhæver Guyau:
mytologisk Naturforklaring, der i de højere Religioner fortsættes
i Antagelsen af Mirakler, --- en Kreds af Dogmer som absolute
Sandheder, --- og Kultus som en Kreds af Handlinger med
overnaturlige Virkninger. Hvor disse tre Ting ikke findes, vil
Guyau ikke tale om Religion.

En Religions Opløsning sker ikke direkte og udefra, men
sker ved, at dens indre Livsbetingelser falde bort. Og dette sker
igen langsomt, alt eftersom Industri, Videnskab, Individualisme
og selvstændig Etik udvikle sig. Af disse Faktorer lægger Guyau
mest Vægt paa den uafhængige individuelle Overbevisning: det
er denne Faktor, der især har ført til en Religions Sejr over
en anden; den vil tilsidst føre til den gradevise Opløsning af
ethvert dogmatisk System, og derved til Religionens sluttelige
Ophør. Denne Opløsningsproces kan vare længe; men Ideerne
virke uafbrudt og umærkeligt; --- de skride frem, selv om de
--- ligesom trætte Soldater --- undertiden sove paa Marschen.

Men den sluttelige Mangel af Religion (l'absence finale de
religion), Irreligionen, er ikke det Samme som Antireligion, Religionens
fuldstændige Modsætning. Den skal tvertimod blive
en højere Grad af Religion og af Civilisation, idet det Bedste i
det religiøse Liv efter Dogmets Bortfalden vil voxe i Styrke og
i Omfang. Det Evige i Religionen er selve den Stræben, der
har frembragt den, --- Trangen til at gaa ud over det nøgne
Faktum og finde en større Sammenhæng. Den menneskelige
Aand er som Svalen: dens lange Vinger egne sig ikke til en
Flugt langs Jorden, men til høj og dristig Flugt i den frie Luft;
det gælder blot om, at den kan hæve sig op fra Jorden, og det
falder ofte svært, men den uforgængelige Higen mod Idealet
vil atter og atter skaffe den Luft under Vingerne. Denne Higen
% --- p. 117 ---
% running head: Følelsen ved Livet.
\opage{117}vil røre sig stærkere end før, naar den er udskilt fra Religionen.
Kun vil man føle mindre Trang til at finde et bestemt
Sprog, i hvilket de evige Gaader kunne udtrykkes; de ville tale
for sig selv. De sidste Hypoteser ville stedse beholde en udpræget
individuel Karakter. Ingen Religion kan udtrykke Alt,
hvad den enkelte Personlighed i sin hele Ejendommelighed indeholder
eller trænger til. Enhver maa frembringe sin egen Gud
eller sin egen Bibel. Og forskellige Overbevisninger kunne meget
vel bestaa Side om Side mellem Mennesker, ligesom forskellige
Planter kunne trives i samme Jordbund. --- Særligt ville de
etiske Elementer i Religionen blive bevarede og frigjorte. Allerede
nu tenderer Religionen hos de ædleste Mennesker til at
blive Et med Kærligheden. --- Stemningen overfor Livet og Tilværelsen
vil ikke være enkelt og stadig; den vil skifte og svinge.
Udviklingen har jo varet en Evighed --- og har saa kun ført
til en Verden som denne med alle dens Disharmonier og med
al Uvisheden i dens Løb! Jo mere Mennesket ved en Art ideal
Sociologi lever sig ind i hele Tilværelsen, saaledes som han kan
kende den, des mere vil Sørgmodigheden voxe; Tanken er Smerte,
ikke blot Lys. Gud selv maatte, hvis han existerede, føle den
største Smerte, da han som et uendeligt Væsen bitrest maatte
føle sin Afmagt til at hjælpe. Men Helhedsstemningen vil have
Ophøjethedens Præg, saa at Sorgen hos den Stræbende kun bliver
et Element i Livsfølelsen. Kraftoverskuddet vil afføde Haab,
Tro og Lyst til at vove. Selv overfor Døden, »den største Nyhed
i det intellektuelle Liv næst efter Fødselen«, staar den Stræbende
trods al Smerte med urokket intellektuel Fasthed: Notre
dernière douleur reste aussi notre dernière curiosité! ---

Denne sidste Ytring er karakteristisk for Guyaus utrættelige
Tankearbejde, der gik Haand i Haand med hans dybe Alvor,
vemodige Stemning og store Resignation. Man har sagt om ham,
at han sejler ud paa det uendelige Hav uden nogensinde at
kaste Anker. Men kan man nogensinde kaste Anker én Gang
for alle? --- Hvad det religiøse Problem angaar, har han gravet
saa dybt som Faa, og han har klart set Nødvendigheden af at
spørge, om den Livskoncentration, Religionerne have betinget i
deres klassiske Tider, kan forsvinde, uden at der kommer Ækvi%
% --- p. 118 ---
% running head: Friedrich Nietzsche.
\opage{118}valent derfor. Han har her søgt at hævde Kontinuiteten paa de
aandelige Værdiers Omraade, skønt han mener, at Begrebet Religion
kun gælder for disse Værdier, naar de fremtræde i en
vis bestemt Form.\textsuperscript{37}) Kontinuiteten mellem Religion og Irreligion
betoner han i hvert Tilfælde saa stærkt, at det Bedste i
Religionen (og tillige det, der har været den egentlige Tendens
ved dens Tilblivelse) skal bevares i Irreligionen.

\chapterhead{II.}{FRIEDRICH NIETZSCHE.}

\sectionhead{A. \textit{Karakteristik og Biografi.}}

Om det væsentligt er en Digter eller en Tænker, vi her
have for os, kan være Genstand for Strid, og har været Genstand
for Strid. Det kan i hvert Tilfælde paastaas, at det vilde
have været heldigt for den blivende Værd af Nietzsches Værker,
og at der vilde være bleven større Harmoni mellem Digteren
og Tænkeren i ham, dersom han havde anvendt Dialogens
eller Dramaets Form, eller maaske en lignende Form som Kierkegaards
Pseudonymer, til at fremstille sine Ideer. Ti det er forskellige,
modstridende Tanker, der røre sig hos ham, og røre sig med
samme Lidenskab; det er, som om det er forskellige Personligheder,
der fremtræde, ikke blot i hans forskellige Skrifter, men
ofte i et og samme Skrift. Men selve sin Hovedtanke --- i hvert
Tilfælde det, jeg vil forsøge at paavise som. hans Hovedtanke
--- iklædte sig for ham poetisk Form, idet han kun i store, billedlige
Træk kunde faa Luft for, hvad der inderst inde laa ham
paa Hjerte. Den betydeligste Mand blandt hans Ungdomsvenner,
Erwin Rohde, betragtede ham væsentligt som Digter, og
glædede sig derfor, da »Zarathustra« begyndte at udkomme.
»Jeg har,« siger han, »længe havt den Følelse, at Nietzsche lider
under en Fylde af Poesi, der ikke vil udfoldes til egentlig Digtning
og nu volder Feber og Nød i hans Indre.« Paa en meget
klar Maade har Nietzsche selv udtalt sig om Forholdet mellem
det digteriske og det filosofiske Element hos ham, og om begges
Betydning for ham, i en Fortale (af 1886) til en ny Udgave af
% --- p. 119 ---
% running head: Videnskab som Kunst, og Kunsten som Liv.
\opage{119}hans Ungdomsskrift »Die Geburt der Tragödie«, idet han betegner
det som sin stadige Opgave at se Videnskaben under Kunstens
Synspunkt og Kunsten under Livets Synspunkt. Denne
Udtalelse viser tillige, at det er med Rette, jeg betegner hans
Retning, sammen med Guyaus, som Vurderingsfilosofi. Det er,
ligesom hos Guyau, Livets Kraft og Fylde, der er hans inderste
Grundtanke; og Nietzsche har da ogsaa med stor Interesse læst
flere af Guyaus Skrifter. Hvad han tilstræber, er en ny, positiv
Livsvurdering paa Grundlag af en kulturhistorisk Betragtning.
I Forhold hertil er det, der i større Kredse er faldet mest
i Øjnene, Læren om Herre- og Slavemoral, den sociale Dualisme
og Foragten for Mængden (»Hjorden«), afledet og underordnet,
--- og, som jeg skal forsøge at vise, galt afledet. Det er
paa dette Punkt, i Forholdet mellem hans egentlige Grundtanke
og hans specielle Synsmaader og Antipatier, at de største Modsigelser
findes hos ham. Men de af hans Skrifter, i hvilke disse,
i Forhold til Grundtanken afledede (og galt afledede) Udtalelser
findes, ere blevne mest læste. Hans literære Virkning beror
paa denne Side af hans Forfattervirksomhed. Uden her at ville
forny den Strid, jeg for en Del Aar tilbage havde i Anledning
af Nietzsche med Georg Brandes, kan jeg dog ikke undlade at
bemærke, at det var uheldigt, at Nietzsche blev indført i vor
Literatur fra denne Side. Vægten blev lagt paa det mest Iøjnefaldende
hos ham, men ikke paa det Fundamentale. Nu, da man
--- især efter Udgivelsen af saa mange af Nietzsches efterladte
Skrifter og Udkast --- nogenlunde kan overse hans Produktion,
viser der sig for mig det Humoristiske, at omtrent det Eneste,
min daværende Modstander og jeg vare enige om, var Noget,
hvori vi begge havde Uret. Vi vare enige om, at Nietzsche var
Modstander af Velfærdsmoralen, men det viser sig nu, at Nietzsche
i Virkeligheden var Velfærdsmoralist i stor Stil. Dette maa fremhæves
her, hvor vi ikke have at gøre med, hvad der literært
maaske er mest epokegørende, men med det, der filosofisk er
det Afgørende: Grundtankerne, deres Begrundelse og deres
Kamp.

Den Betydning, det literære og poetiske Element har hos
Nietzsche, medfører større Vanskeligheder for Fremstilling og
% --- p. 120 ---
% running head: Ungdomsstudier .
\opage{120}Karakteristik end hos andre Filosofer. Forat give et nogenlunde
rigtigt Billede af ham, maa min Fremstilling være udførligere,
end den behøvede at være, hvis det blot gjaldt et rent Tankeopgør.
Dog vil det have stor psykologisk Interesse at studere
Forholdet mellem Tænkning, Kunst og Liv hos et Geni, hvis
Higen var at lade al Tænken og Digten tjene Livets Uddybelse
og Forhøjelse. Et psykologisk Udbytte her vil ogsaa være af
filosofisk Betydning. ---

Friedrich Nietzsche er født i preussisk Sachsen den 15.
Oktober 1844 af en gammel Præsteslægt. Selv vilde han gerne
være af fornemmere Oprindelse, og mente, at hans Oldefader
havde været en polsk Adelsmand. Nogle af hans Kritikere have
derimod netop ment at finde Præstesønnen i hans Tankegang.
Meget tidligt førtes han ind under andre Indflydelser end de
kirkelige, idet han i Skolen og ved Universitetet med Iver dyrkede
klassiske Studier. Førstehaandskendskabet til græsk Aandsliv
og Begejstringen over det blev afgørende for hans hele Livs
Digten og Tragten. Men ved Siden deraf blev Studiet af Schopenhauer,
som han dyrkede samtidigt med sine filologiske Studier,
af dybt indgribende Betydning. Af andre Filosofer have
Dühring og Albert Lange havt Indflydelse paa ham; men Schopenhauer
blev hans egentlige »Opdrager«. I Anledning af sit
Schopenhauerstudium skrev han (1869) til sin Ven Paul Deussen:
»En Filosofi, som vi kun antage af ren Erkendelsestrang,
bliver aldrig ganske vor Ejendom, fordi den aldrig \textit{var} vor Ejendom.
Enhver Enkelts rette Filosofi er \textit{ἀνάμνησις} [Erindring].«
Erkendelsestrangen ytrede sig hos Nietzsche stedse i uadskillelig
Sammensmeltning med Kunsttrangen og Livstrangen. Det
er hans Styrke, og det er hans Svaghed --- som Tænker, især
da hans Tanker paa Grund af hans impulsive og disharmoniske
Natur saa ofte dikteredes af Øjeblikkets Stemning eller af en
Reaktion mod Tendenser i Omverdenen.

Nietzsches anden store Opdrager blev Richard Wagner.
Allerede under sine Studieaar i Leipzig havde han gjort personligt
Bekendtskab med den store Mester, og da han 1869 blev
filologisk Professor i Basel, omgikkes han fortroligt med Wagner
og hans Hustru, der i disse Aar boede i Nærheden af Luzern.

% --- p. 121 ---
% running head: Nietzsche og hans Opdragere.
\opage{121}Endnu en tredie Opdrager fik han: Sygdommen, der i
mange Aar tærede paa hans Kraft og tilsidst blev uhelbredelig
Sindssygdom. Han skal have paadraget sig den, medens han som
Sygepasser deltog i den fransk-tyske Krig. Den var paadraget
ved Smitte, og \textsc{Møbius} \textit{(Ueber das Pathologische bei Nietzsche)}
hævder derfor, at den Hjernesygdom, hvoraf han kom til at lide,
ikke kan bringes i Forbindelse med hans oprindelige Personlighed,
skønt den fik bestemmende Indflydelse paa dennes senere
Udvikling, og skønt han maa anses for at have været arveligt
prædisponeret. Ifølge Møbius forværredes Nietzsches Tilstand ved
gal Lægebehandling og ved hans eget Misbrug af Kloral. Naar
Møbius tillige mener at kunne pege paa et bestemt Sted i et
af Nietzsches Skrifter, hvor Sygdommen skulde have begyndt at
ytre sin Indflydelse, tror jeg, det er aldeles ubegrundet; men
derom senere. --- Sygdommen gjorde ham det efter lang Kamp
umuligt at fortsætte sin Universitetsvirksomhed; tillige var Filologien
bleven ham for snever, efterhaanden som de store Problemer
tog Fart i ham og satte Fantasien og Tanken i Bevægelse.
1879 tog han sin Afsked og levede nu i de følgende
Aar mest i Engadin om Sommeren og ved Rivieraen om Vinteren.
I den vidunderlige Bjergnatur i Engadin, med dens lette
Luft og dens mærkelige Belysninger, højt hævet over det sædvanlige
Menneskeliv i Dalene, --- heroppe »paa Vidderne« har
han efter eget Udsagn undfanget nogle af sine betydningsfuldeste
Ideer.

Ejendommeligt for Nietzsche er det, at alle hans tre Opdragere
virkede mest afgørende paa ham gennem den Reaktion,
de udløste, eller den Kontrastvirkning, de betingede. I Sygdommen
hærdedes hans Villle, og Tapperheden i Lidelser blev hans
Hoveddyd. Han lærte Livets Smerte paa første Haand, men tillige
Livets overstrømmende Kraft til at hævde sig trods Alt.
Sygdommens Opdragelse blev han aldrig fri for; de to andre
Læremestre kunde han skille sig fra, i hvert Tilfælde tilsyneladende.
Med den levende Opdrager, Richard Wagner, blev
Bruddet smerteligt. Hos Schopenhauer var det Pessimismen,
der paa et vist Punkt af hans Udvikling frastødte ham; hos
Wagner frastødtes han baade af Pessimismen og af den Ven%
% --- p. 122 ---
% running head: Nietzsche og Hjorden.
\opage{122}ding til middelalderlig Kristendom, han mente at finde hos den
store Tonekunstner,\textsuperscript{38}) hvis menneskelige Svagheder han efterhaanden
fik et skarpt Blik for. Siegfried havde han beundret
og elsket; allerede ved 
Wotan --- Guden som Pessimist!
--- var han bleven betænkelig; Parsifal hadede han. Han havde
fra først af beundret Wagners Kunst som »dionysisk«; senere
stod den for ham som et Udtryk for Dekadence.

Ved Studiet af Nietzsches Biografi\textsuperscript{39}) leder man naturligt
efter Oplevelser, der kunne have bragt ham i Berøring med
Folket, den store Masse, »Hjorden«, som han senere kaldte den.
Det var jo naturligt at tænke sig, at han særligt ofte havde
havt Lejlighed til at se Mennesker optræde en masse. Kun to
saadanne Lejligheder forefindes, og begge Gange vaktes hans
Foragt. Den ene Gang var det øldrikkende tyske Studenter;
den anden Gang var det Teaterpublikummet i Bayreuth, der
ikke svarede til hans Forventninger om den ideale Tilhørerkreds
ved den store Kunsts Genfødelse. --- Derimod har Nietzsche
aldrig havt Lejlighed til at leve i Folket og lære dets Stræben,
dets Evne til Udvikling og Organisation at kende; han stod fuldstændigt
fremmed overfor de største sociale Begivenheder i det
19. Aarhundrede. ---

I Sygdommens Gang mener Møbius at kunne skelne flere
Perioder. En Stigning af Sygeligheden finder Sted fra 1881
til 1884; derefter en Synken i de følgende Aar, og saa en ny
Stigen, indtil der i Begyndelsen af 1889 kom en Katastrofe.
Han sendte Breve om til flere Bekendte, snart undertegnede
»Dionysos«, snart »den Korsfæstede«. En Dag faldt han i Afmagt
om paa Gaden i Turin. En Ven bragte ham til Basel,
hvorfra han senere bragtes til sin Hjemstavn. De sidste Aar henlevede
han i en fuldstændig Sløvhedstilstand i Weimar, omhyggeligt
plejet af sine Nærmeste. Den 25. August 1900 døde han.
--- Da der ingen Obduktion fandt Sted, har Sygdommens Art
og Sæde ikke med fuld Sikkerhed kunnet konstateres. Møbius
antager, at den har været progressiv Paralyse.

Nietzsche var en Modsætningernes Mand. Den Modsætning,
som især fremtræder hos ham, prægende hele hans Færd og
skæbnesvanger for hans Liv og hans Betydning, er Modsætnin%
% --- p. 123 ---
% running head: Nietzsches Karakter.
\opage{123}gen mellem høj Begejstring og dyb Modbydelighed. Under
sunde Forhold ere disse to Modsætninger to forskellige Sider
af samme Grundstemning; hos Nietzsche vare de fjendtlige Magter,
der kæmpede om Herredømmet i hans Sind og i hans Villie.
I denne Kamp gik han til Grunde. Beundring for det Store
og Høje, Længsel efter det, Ønsket om at virke for det var vel
det Inderste i hans Sjæl; men Antipatien og Foragten for det
Lave og Svage overvældede, hypnotiserede ham, saa at det, han
vilde kæmpe for, kom til at træde i Skygge for det, han vilde
kæmpe mod. De »reaktive« Stemninger --- som han dog selv
foragtede hos Andre --- fik atter og atter Overvægten hos ham.
Han har selv følt, hvilken Fare der her truede hans aandelige
Liv. I Aaret 1874 skriver han til Frøken v. Meysenbug: »Wie
wird mir zu Muthe sein, wenn ich erst alles Negative und Empörte,
was in mir steckt, aus mir herausgestellt habe, und doch
darf ich hoffen, in fünf Jahren ungefähr diesem herrlichen Ziele
nahe zu sein.« Dette Haab gik ikke i Opfyldelse, selv efterat
han var begyndt at arbejde paa den positive Fremstilling af
sine Ideer. Endnu i det sidste Aar før Katastrofen afbrød han
sit sammenhængende Arbejde for at faa Luft for sine Antipatier,
skønt han havde aabnet sine Ventiler tilstrækkeligt mange
Gange forud.

Den fremherskende Plads, Antipatierne kom til at spille i
hans Produktion, staar ikke blot i Modsætning til hans Begejstring,
men til hans kærlige og inderlige Sindelag som det
ytrede sig i hans nærmeste Forhold. Ogsaa dette følte han selv.
Han skriver til Frøken v. Meysenbug (1878): »I dette Efteraar
satte jeg mig for at begynde enhver Dag. med at spørge mig
selv: Gives der Ingen, du idag kunde gøre noget Godt for?
. . . . Med mine Skrifter volder jeg saa mange Mennesker Fortræd,
at jeg naturligt maa ønske at gøre det godt igen paa en
eller anden Maade.« Og nogle Aar efter (1883) skrev han til
Erwin Rohde, at han vel vidste, at det Billede, hans Bøger nu
[i den mest polemiske Del af hans Forfattervirksomhed] gave
af ham, ikke var det Billede, hans Ven bar af ham i sit Hjerte.
Han havde, fortsætter han, virkeligt en »anden Natur« [end den,
% --- p. 124 ---
% running head: Nietzsches Tapperhed.
\opage{124}der paa denne Tid ytrede sig i hans Bøger]; ved den »første«
alene vilde han for længe siden være gaaet til Grunde!

Ikke blot mod Begejstringen og Kærligheden stred den
stærke Fremhersken af Antipatierne, men ogsaa mod den Optimisme,
Nietzsche saa lidenskabeligt arbejdede for i sine senere
Aar. Han forlangte i stigende Grad og i overvættes Udtryk en
Bekræftelse af Livet, en Anerkendelse af dets Værdi. Men
hvorledes kunde han forlange dette, naar Livet var saa opfyldt
af Væmmelighed og Usselhed, som han atter og atter udtalte,
og ikke kunde blive færdig med at udtale?

De stridende Stemninger kunde hos ham ikke forsones,
hverken i den store Humors eller det dybe Vemods Form. Om
Shakespeare er det blevet sagt, at hans Humor er et Udtryk
for hans Troskab mod Livet. Paa denne Maade formaaede
Nietzsche ikke at vise sin Troskab; maaske hang det sammen
med, at han hverken var Digter eller Tænker paa afgjort Maade,
men at ogsaa her to Naturer brødes hos ham. Hans Samtidige
og i mange Henseender hans Tankefrænde Guyau fandt i Vemodet
og i Tankens evige Stræben den harmoniske Udligning.
Ogsaa denne Form var Nietzsche negtet. En Klangfarve, der i
sig samlede de stridende Toner, naaede han ikke. Svingninger
mellem dionysisk Rus og fortærende Disharmoni førte ham stedse
nærmere til Opløsningen. Her vil det være berettiget at finde
en Virkning af Sygdommen. Han har kæmpet redeligt mod
den. Ti han var en tapper Mand, et Villiesmenneske. Han følte
dybt »die Vehemenz der inneren Schwingungen«; men naar hans
Kritikere kaldte ham »excentrisk«, sagde han med Stolthed:
»Disse Herrer, der intet Begreb have om mit Centrum, om den
store Lidenskab, i hvis Tjeneste jeg lever, ville vanskeligt kunne
have Blik for, hvor jeg virkeligt var udenfor mit Centrum.«
(Brev til Deussen 1888).

Sin Tapperhed viste han overfor Arbejde, Lidelse, Miskendelse
og Ensomhed. »Ingen Smerte«, siger han i et Brev, »har
formaaet og skal formaa at forføre mig til et falsk Vidnesbyrd
om Livet, saaledes som jeg erkender det.« »Min Existens,« skriver
han omtrent samtidigt, »er en frygtelig Byrde, som jeg for
længe siden havde kastet bort, naar jeg ikke netop i denne Til%
% --- p. 125 ---
% running head: Den sidste Fjende .
\opage{125}stand af Lidelse og næsten absolut Forsagelse gjorde de lærerigeste
Prøver og Experimenter paa det aandelige og sædelige
Omraade; denne efter Erkendelse higende Glæde fører mig op
paa Højder, hvor jeg sejrer over al Pinsel og al Haabløshed.
I det Hele er jeg dog lykkeligere end nogensinde før.« »Agt
vel derpaa,« skrev han senere, »at de Aar, i hvilke jeg ophørte
at være Pessimist, netop var de Aar, i hvilke min Livskraft var
ringest.«

Kun én Ting frygtede han: Medlidenheden. Den havde
han i sin pessimistiske Tid og under Schopenhauers »Opdragelse«
forset sig paa. Denne Frygt, der staar i en underlig
Modsætning til hans Tapperhed, finder tildels sin Forklaring ved
hans egen Haardførhed, men hang ogsaa sammen med hans
eget bløde Sind, for hvilket Medlidenheden kunde være en Fristelse.
Ved Medlidenheden tænker han nærmest paa den passive,
sentimentale Form af Medfølelse, der svækker baade den,
der har den, og Genstanden for den. Det endte med, at Medlidenheden
stod for ham som den sidste store Synd, man endnu
kunde falde i, naar andre Synder vare blevne umulige. Den
ejendommelige hypnotiske Virkning, Nietzsches Antipatier havde
paa ham, ytrer sig ogsaa her. Tilsidst mærkede han Sygeluft
rundt om i Europa, hvor han gik og stod.

Formen for Nietzsches Forfattervirksomhed bestemtes ved
hans Temperament og ved hans Sygelighed. Han havde en Trang
til Koncentration, og der er i hans Produktion visse positive
Hovedtanker, hvis fuldstændige Gennemarbejdelse han higede
efter. Men han formaaede ikke at bøje de enkelte Øjeblikkes
overvættes Stemning ind under, hvad Grundtanken og Grundstemningen
krævede. Den Begrænsning, der gør Mesteren, stod
ikke i hans Magt. Deraf Splidagtigheden i hans Udtalelser; det
ene Øjebliks Rus stemmede ikke altid med det andets. Aforismen
blev hans Yndlingsform. Hver enkelt Aforisme lader et
Synspunkt komme til fyldig, ofte slaaende Udtryk, for saa at
give Plads for andre Aforismer, der uden videre sige det Modsatte.
Det egentlige Tankearbejde bliver ikke gjort. Med Rette
havde man betegnet Aforismeformen som farlig for Nietzsche.
Han kom ind paa den, dels ved Efterligning af ældre franske
% --- p. 126 ---
% running head: Nietzsches Skrifter.
\opage{126}Forfattere, især den af ham højt beundrede Larochefoucauld,
dels paa Grund af Sygeligheden, der gjorde ham sammenhængende
Arbejde umuligt. Han producerede paa Spadsereture og
nedskrev sine Tanker i løsreven Form; andre maatte saa redigere
dem. »Ich kritzele auf meinen Wegen hie und da etwas
auf ein Blatt; ich schreibe nichts am Schreibtisch; Freunde entziffern
meine Kritzeleien.« (Brev til Eiser 1880). En stor Del
af Nietzsches Popularitet skyldes vistnok denne Form, der gør
det saa let at faa i hvert Tilfælde et »Hilsebekendtskab« med
hans Bøger. Selv siger han om en af sine i denne Form forfattede
Bøger: »Ein Buch, wie dieses, ist nicht zum Durchlesen
und Vorlesen, sondern zum Aufschlagen, namentlich im
Spazirengehn und auf Reisen .« (Morgenröthe. Aph. 454). Det
er, som han selv har sagt, ikke let at finde hans egentlige Centrum;
men hans egen Fremstillingsform har gjort det vanskeligere,
end det i sig selv er.

\sectionhead{B. \textit{Nietzsches Skrifter.}}

Det er en Følge af Nietzsches disharmoniske og svingende
Temperament, at hans Skrifter ikke kunne ordnes rent kronologisk,
men maa deles i Grupper, der maa karakteriseres ved
det forskellige Forhold, der fremtræder mellem de Grundtanker,
som i det Hele ere de raadende hos ham.

1. \textit{Første Gruppe.} Hovedproblemet stilles, og den definitive
Løsning antydes. Karakteristisk for denne første Gruppe er
det kulturhistoriske Grundlag for Problemstillingen, som Nietzsche
henter fra sine filologiske Studier. Han støtter sig til Sociologi,
medens Guyau væsentligt støtter sig til Biologi. \textit{Die Geburt der
Tragödie} (1872) er ikke blot Nietzsches første betydelige Skrift,
men ---
hvad virkelig Problembehandling angaar --- maaske det
betydeligste af alle hans Skrifter. Hans Hovedtanke om en ny
Livsvurdering kommer allerede frem her. I et af sine sidste
Skrifter (Götzendämmerung) siger han selv: »Die »Geburt der
Tragödie« war meine erste Umwertung aller Werte.« Den tragisk-patetiske
Livsopfattelse --- symboliseret i Dionysos og Apollon
--- stilles i Modsætning til den intellektuelle Optimisme, der
% --- p. 127 ---
% running head: Første og anden Gruppe af Skrifter.
\opage{127}repræsenteres af Sokrates. Forholdet mellem Videnskab, Kunst
og Liv fremtræder her mellem disse tre Skikkelser --- og Dionysos
staar højest, skønt det apollinske, formende og begrænsende
Element anerkendes mere end i de senere Bøger.

Det saaledes vundne Synspunkt anvendes i \textit{Unzeitgemässe
Betrachtungen} (1873---1876) paa samtidig tysk Aandskultur. Af
de fire Betragtninger indeholder den ene en skarp Kritik af
Strauss, den anden en Polemik mod Overvurderingen af den
historiske Betragtningsmaade, den tredie og den fjerde forherlige
Schopenhauer og Wagner som de store Opdragere. Det
er de higende, tragisk stemte Personligheder, der stilles i Modsætning
til den ved sine kritiske Resultater selvtilfredse Forsker
og til den rent objektive Syslen med Fortiden. --- Senere
mente Nietzsche, at han havde lagt sin egen haabende Stemning
ind hos Schopenhauer og Wagner, der i Virkeligheden vare
Dekadenter, og i hans senere Produktion bliver han ikke færdig
med at reagere mod sin egen tidligere Beundring. Han brænder,
hvad han har tilbedt, og tænder stadigt Baalet paany.

I denne første Gruppe komme --- foruden Grundforholdet
mellem Videnskab, Kunst og Liv --- endnu andre af Nietzsches
Hovedtanker allerede frem: den radikale Aristokratisme, der i
de store Mennesker ser Historiens Maal, og den dermed sammenhængende
sociale Dualisme mellem Herrer og Slaver.

2. \textit{Anden Gruppe.} Her træde Modsætningerne indenfor
Nietzsches Tankegang stærkere frem, og tillige hans Modsætning
til andre Standpunkter. Bruddet med Schopenhauer og med
Wagner blev her især af Betydning. Han taler senere om en
antiromantisk Kur, han underkastede sig, fordi han led under
den farligste Form for Romantik. Det gjaldt at forsvare Livet
mod de Slutninger, man plejer at drage fra Smerte, Skuffelse
og Isolation. Han mener sikkert her det Samme, som naar han
i et Brev af 1883 taler om en lang og tung aandelig Askese,
han har underkastet sig i de sidste sex Aar. Den Kur eller
Askese, han omtaler, bestod i realistiske Studier. Han stiftede
Bekendtskab med Naturvidenskab, med ældre fransk Filosofi og
med moderne engelsk Filosofi. Der kom derved et nyt Element
ind i hans Tankegang, som var vanskeligt at forene med
% --- p. 128 ---
% running head: Det etiske Problem og Gentagelsen .
\opage{128}de tidligere; og der blev Anledning til ny Polemik og ny Reaktion.
Det var paa dette Tidspunkt, han kom ind paa Aforismeformen,
som han ikke igen formaaede at befri sig for. Dens
Herredømme maatte naturligvis blive des stærkere, jo flere forskellige
Tendenser, Sympatier og Antipatier der kom frem hos
ham og hver for sig fordrede Udtryk.

Det var dog ikke blot en forbigaaende Kur, han her underkastede
sig. Det etiske Problem arbejdede sig for ham bestemtere
frem, som en særlig Form for det almindelige Kulturproblem,
de første Skrifter havde stillet. En grundig Tvivl om de
i Livet og i Tænkningen hidtil gældende moralske Forudsætninger
(»Vurderinger«) opstod hos ham. »Ich begann, unser
Vertrauen zur Moral zu untergraben.« (Fortale af 1886 til
»Morgenröthe«).

Herhen høre Skrifterne \textit{Menschliches, Allzumenschliches}
(1878---80), \textit{Morgenröthe} (1881). \textit{Fröhliche Wissenschaft} (1882).

I Slutningen af »Fröhliche Wissenschaft« findes for første
Gang en Tanke, der i Nietzsches senere Aar spillede en stor
Rolle hos ham, Tanken om, at Tilværelsen var rytmisk i sin Udvikling,
saa at det, der én Gang var sket, vilde komme igen
paa samme Maade i alle Enkeltheder i en følgende Periode. I
Forbindelse dermed træder Zarathustraskikkelsen frem som Livsbekræftelsens
store Profet, der vil stige ned til Menneskene for
at forkynde dem, at Livet er saa skønt, at dets Gentagelse kan
og maa villes. (Fröhl. Wiss. Aphr. 341---342). Kort efter udarbejdede
Nietzsche første Del af sin Zarathustrabog. Han skriver
til en Ven (Juni 1883): »Hvad jeg endnu har tilbage af
Livet (kun Lidet, tror jeg), skal anvendes til helt og fuldt at udtrykke
det, for hvis Skyld jeg overhovedet har holdt Livet ud.
Tiden til at tie er nu forbi; min Zarathustra, som i Løbet af
nogle Uger vil blive sendt til dig, vil forraade dig, hvilken høj
Flugt min Villie tager. Lad dig ikke bedrage af denne Bogs
Legendeform; bag alle de jevne og sælsomme Ord ligger mit
dybe Alvor og min hele Filosofi.« --- Jeg forstaar ikke, hvorfor
Möbius tror, at paa hint Sted i »Fröhliche Wissenschaft« er
Nietzsches Sindssygdom begyndt at angribe hans Forestillingsliv,
efterat den forud kun havde angrebet hans Følelsesliv.
% --- p. 129 ---
% running head: >>Hinsides Godt og Ondt« .
\opage{129}Zarathustraskikkelsen havde beskæftiget Nietzsche fra hans Barndom
af, og den forberedes desuden ved Dionysosskikkelsen i
»Geburt der Tragödie«. Tanken om alle Tings Genkomst havde
Nietzsche fra sine græske Studier, idet den forekommer hos
Pytagoræere og Stoikere; den kommer nu ret naturligt frem
hos ham som en Prøvesten for, hvor stor Trangen til Livsbekræftelse
er. Ideen kan forekomme grel; men den er psykologisk
forstaaelig hos Nietsche netop paa dette Tidspunkt. Han
siger selv i en autobiografisk Optegnelse: »Grundkompositionen
til »Zarathustra«, den evige Genkomsts Tanke, denne højeste
Formel for Bekræftelse, som overhovedet kan naas, hører hen
til August 1881; den er henkastet paa et Blad med Paaskrift:
6000 Fod over Mennesker og over Tiden! I de Dage gik jeg
gennem Skovene ved Silvaplanasøen og stansede ved en mægtig,
pyramidalsk optaarnet Klippeblok i Nærheden af Surlei.
Der fik jeg denne Tanke.«

Selve Zarathustrabogen regner jeg ikke med til denne Gruppe.
Arbejdet paa den blev atter og atter afbrudt ved Trangen til i
nye Aforismer at faa Luft for polemiske og antipatiske Stemninger.
Han naaede ikke til at faa sin »store Alvor og hele sin
Filosofi« udtalt. \textit{Jenseits von Gut und Böse} (1886) skulde være
en Kommentar til Zarathustrabogen (som endnu ikke var færdig
--- og aldrig blev det). Men det er en Kommentar, der forudsætter,
at man forstaar selve Bogen, den ufærdige Bog, der
skal kommenteres. Nietzsche skriver i et Brev (Oktbr. 1886):
»Hast du dich in meinem »Jenseits« umgethan? Es ist eine
Art von Commentar zu meinem »Zarathustra«. Aber wie gut
müsste man mich verstehen, \textit{in wie fern} es zu ihm ein Commentar
ist!« --- Man kan ikke lade være at spørge, hvormange
af dem, der som gode Nietzscheanere trøstigt begyndte at leve
»hinsides Godt og Ondt«, der virkeligt have fyldestgjort denne
Fordring, deres Mester stiller. Men det er heller ingen let Opgave,
der stilles dem. Og det er en ny Form for Tragedien i
Nietzsches Liv og Tænkning, at han kom til at skrive Kommentarer
til Værker, der ikke udførtes, Kommentarer, der alligevel
selv behøvede Værkerne til --- Kommentar! --- \textit{Die Genealogie
der Moral} (1887) skulde igen være et Supplement til »Jenseits«.
% --- p. 130 ---
% running head: Tredie Gruppe af Skrifter.
\opage{130}Den er et Forsøg paa historisk Begrundelse af den sociale Dualisme,
der i Nietzsches aforistiske Skrifter staar saa stærkt i Forgrunden
hos ham.

I disse to Skrifter kommer Teorien om Slaveopstande i Moralen
og om Nødvendigheden af absolut Nyskaben af Værdier
tydeligt frem, skønt Nietzsche selv fører dem tilbage til »Menschliches---Allzumenschliches«
og til sit Ophold i Syditalien i Aaret
1876. Det er de Ideer hos Nietzsche, som ere faldne mest i
Øjnene i større Kredse. Stærkest fremtræder her Nietzsches
store Trang til at føle sig i Afstand fra Mængden, --- hvad han
kalder »das Pathos der Distanz«, --- og hans Kastemoral udtales
i den krasseste Form. Men det maa vel fastholdes, at de
to Skrifter kun ere Udsnit af et allerede fra Aaret 1881 af paatænkt
Værk, og at de derfor ikke kunne betragtes som Nietzsches
sidste Ord. Den Stemning, der kommer til Orde i dem, staar
netop tilsidst i en afgørende Strid med hans dybeste Tragten,
en Strid, han undertiden selv har følt. --- Fru Elisabeth Førster-Nietzsches
Oplysninger i Fortalen til femtende Bind af Nietzsches
Værker viser os den literære Sammenhæng i hans Produktion
fra 1881 til 1888, og særligt de forskellige Værkers Forhold
til Hovedværket og dettes positive Afslutning. Deraf ses ogsaa,
at en Række voldsomt polemiske Værker fra hans sidste sunde
Aar (1888) ogsaa ere Udsnit af Hovedværket. Saaledes: \textit{Der
Fall Wagner, Götzendämmerung, Nietzsche contra Wagner} og
\textit{Antichrist.}

3. \textit{Tredie Gruppe.} Sine Antipatier og sine polemiske Stemninger
havde Nietzsche ondt ved at faa fuldt Udtryk for. Han
afbryder lige tilsidst sit Hovedværk for med feberagtig Hast at
udsende sine Protester. Og dog maatte det ligge nær for den,
der vil foretage en Omvurdering (endog af \textit{alle} Værdier!), at
gøre det klart for sig og Andre, paa hvilket Grundlag og efter
hvilken Maalestok denne Omvurdering skal ske. Et saadant
Grundlag har Nietzsche ogsaa søgt, skønt han ikke fik det fuldt
fremstillet. Men han saa ikke klart, at den radikale Aristokratisme
og den sociale Dualisme, han i sine Skrifter havde forkyndt,
og som han skyldte sine største literære Triumfer, stode
i uforligelig Modstrid med hans Hovedtanke, hans Grundværdi:
% --- p. 131 ---
% running head: Kulturhistorisk Udgangspunkt .
\opage{131}Livets Kraft, Sundhed og Lykke. --- I to ufuldendte Værker har
han arbejdet paa at faa denne Hovedtanke frem som sidste Værdimaaler.
Det ene er \textit{Also sprach Zarathustra. Ein Buch für
Alle und Keinen,} hvoraf fire Dele udkom 1883---1886. Med
disse fire Dele maa sammenholdes det overordentligt betydningsfulde
Udkast til en afsluttende Del, som er trykt i tolvte Bind
af Værkerne (p. 321). Det andet Værk, der i mere filosofisk
Form skulde give, hvad »Zarathustra« gav i poetisk Form, er
\textit{Der Wille zur Macht. Versuch einer Umwerthung aller Werthe;}
det existerer kun i Fragmenter, der nu ere udgivne i femtende
Del af Værkerne. Det skulde slutte ligesom »Zarathustra«, idet
den sidste Bog skulde have til Overskrift: »Dionysos. Philosophie
der ewigen Wiederkunft«. --- Men Nietzsche havde givet
sine Antipatier i den Grad Luft, at han nu manglede Kraft til
at forme sit eget positive Ideal i Billede og i Tanke. Og særligt
var han saa uvant med al sammenhængende Tænkning og
saa mistænksom overfor den, og overhovedet overfor enhver bestemt
Form, at tilsidst Apollon saalidt som Sokrates kunde udtrykke,
hvad Dionysos mente. Og i sin Extase sønderrev Dionysos
da sig selv.

\sectionhead{C. \textit{Kulturhistorisk Udgangspunkt.}}

Nietzsche vilde en ny Renaissance, og ligesom den italienske
Renaissance begynder han med at drage Grækerne frem. Hans
filologiske og hans filosofiske Interesse forenedes her. Da han
som ganske ung Mand var bleven filologisk Professor i Basel,
frygtede han at blive Filister og »Heerdenmensch«, hvis han
gik op i Specialiteter. Han vilde gennemtrænge sin Videnskab
med nyt Blod ved at benytte den til at belyse de store Livsproblemer.
Og ved dette Forsøg skulde Schopenhauer og Richard
Wagner være hans Læremestre.

I \textit{Geburt der Tragödie} fremhæves som det vigtigste Moment
i den græske Kulturs Historie Foreningen af dionysisk
Begejstring med apollinsk Harmoni. Det bakkantiske Sværmeri,
der var Udtryk for overstrømmende Livsfylde, blev ved Delfi's
Indflydelse bragt ind i klare Former ved at faa billedligt og
% --- p. 132 ---
% running head: Dionysos —1 Apollon - Sokrates.
\opage{132}kunstnerisk Udtryk. Den tragiske Kunst opstod af de dionysiske
Satyrkor; Korsangenes Indhold tolkede fra først af Korets
Visioner, idet det saa sin Guds Forfølgelse, Sønderrivelse og
Genoprejsning for sig i store Billeder. Saaledes blev Tragedien
født af Musikens Aand, som Tolkning af Livets store Skæbner.
Den herskende Opfattelse af Græciteten, som især skyldes Winckelmann,
Goethe og Schiller, oversaa den titaniske Baggrund
for den antike Harmoni: Dionysos som Forgænger for Apollon!
Billedet, Ordet og Tanken ere sekundære, afledede i Forhold til
den store Livstrang.

Den græske Tragedie gik til Grunde ved Sokratismen, som
indskærpede, at det kom an paa Viden, paa Forstaaelse. Under
dens Indflydelse blev Euripides den første nøgterne Tragiker;
Tragedien maatte efterhaanden vige Pladsen for den nøgterne,
borgerlige Komedie, og Platon ironiserer over den digteriske Begejstring,
der ikke selv véd, hvad den gør. En ny Type træder
her ind i Menneskeheden: det teoretiske Menneske, Forskeren,
for hvem maaske tilsidst det at søge staar over det at finde.
Fra den Tid af er Kampen mellem den tragiske og den teoretiske
Verdensbetragtning staaende, indtil Befrielsen kommer --- gennem
Goethes Faust, Kants Fornuftkritik, Schopenhauers Livsanskuelse
og Wagners Musik. Men højest blandt disse befriende
Magter staar Musikens Kunst, som Wagner har øvet den: i den
lægge vi Øret til Tilværelsens Hjertekammer; i den rører sig
det Dionysiske, den bacchantiske aldrig hvilende Higen, som
inderste Kraft og Grund i hele Tilværelsen. Det er af Betydning.
at Trangen til Billeder vækkes; men der higes ud over
ethvert Billede. Kunsten er ikke blot Efterligning af Naturen,
men Supplering og Overvindelse af den. ---

Nietzsche benytter Kulturhistorien i store Træk; det gælder
ikke blot hans Udtalelser om græsk Kultur, men ogsaa hans
Omtale af Kristendommen, Reformationen og Revolutionen i
senere Skrifter. I Dionysos, Apollon og Sokrates har han personificeret
tre Tendenser, der stedse kæmpe og brydes i Menneskelivet.
Rent historisk har han ikke Ret i at se Sokratismen
som Hovedaarsag til den græske Kulturs Forfald. Hovedaarsagen
laa i, at det græske Liv droges ind i de store Verdensbegiven%
% --- p. 133 ---
% running head: Radikal Aristokratisme.
\opage{133}heders Hvirvel, og at det ved sin (om man vil: apollinske) Udformning
i smaa Samfund manglede Betingelsen for at hævde
sig under de nye Forhold. Og psykologisk er det ogsaa ensidigt
at lade Energien formes af Fantasien og begge oplæres af Tanken.
Der foregaar uafbrudt i den Enkelte og i Slægten ogsaa
den modsatte Strømning, fra Tanke gennem Billede til Følelse
og Villie. Og Sokrates var selv en positiv Kraft, Forbilledet for
de store, helstøbte Karakterer i den klassiske Oldtids Slutningstid,
den Tid, da det ikke længere gjaldt om at synge i Kor,
men da den Enkeltes Stemme (ligesom i den fuldt dramatiserede
Tragedie) skulde vidne om, hvad han havde faaet ud af Livet.\textsuperscript{40})

\sectionhead{D. \textit{Historiens Maal og den sociale Dualisme.}}

Overfor den overdrevne Betydning, historisk Dannelse efter
Nietzsches Mening har faaet i den nyere Tid, udraaber han:
Husk dog paa at leve Livet! (Memento vivere!) Lad Fortiden
ikke tynge for stærkt paa Eder, saa det gaar ud over Instinkt,
Personlighed, Kunst og Tænkning! Ellers vil der --- som allerede
Hesiodos frygtede --- komme en Tid, da Menneskene blive
fødte graahaarede!

Historiens Betydning ligger ikke der, hvor man i Reglen
finder den, i den endeløse Udviklingsproces, eller i de store
Massers Skæbne. I de enkelte, i Sandhed store Mennesker koncentrerer
Historiens hele Værdi sig; i dem har den naat sit
Maal, og den lange historiske Proces har kun den Betydning
at give den Anledning og den Kraft, som behøves til saadanne
Menneskers Opstaaen. Menneskehedens Endemaal (Ziel) kan ikke
ligge ved Historiens Ende, men kan kun ligge i dens højeste
Exemplarer. Den megen Tale om historisk Udvikling er af det
Onde; den spilder vor Kraft, og den ophidser Masserne. De
store Masser af Mennesker ere kun Midler eller Hindringer eller
Kopier --- im Uebrigen hole sie der Teufel und die Statistik!
--- Storhed beror ikke paa Massevirkning; det Ædleste og Bedste
virker slet ikke paa Masserne.

Her udtaler Nietzsche sin radikale Aristokratisme (aristokratisk
Radikalisme er en urigtig Benævnelse). (Unzeitgemässe
% --- p. 134 ---
% running head: To Motiveringer.
\opage{134}Betrachtungen. Zweites Stück).\textsuperscript{41}) Den er radikal Aristokratisme
i den Betydning, at den er gennemført Aristokratisme.
Substantivet Radikalisme vilde Nietzsche neppe vedkende sig, da
han tverimod er meget konservativ i Alt, hvad der angaar »Hjorden«.
Her maa man vælge; man kan ikke være Tilhænger af
Nietzsche og samtidigt være »radikal« i social og politisk Henseende
--- egentligt heller ikke i kirkelig Henseende. Af hans
Tilhængere har Pseudonymen Peter Gast, der har redigeret, udgivet
og kommenteret flere af hans Bøger, klart indset dette.

Den første Motivering af Aristokratismen hos Nietzsche
hænger saaledes sammen med hans Opfattelse af Historien, hvis
Betydning han vil reducere til saa Lidet som muligt, hvilket jo
unægteligt lykkes, naar det store Menneske staar der, hævet
over Fortid og Fremtid, ikke Virkning af hin, ikke Aarsag til
eller Middel for denne. En Inkonsekvens er der dog, da den
store Masse jo foruden at være Hindring eller Kopi skal kunne
være Middel. Og hvilken Interesse Statistiken »im Uebrigen«
skal kunne have af at beskæftige sig med Massen, naar denne
ikke engang ses som Middel eller Hindring eller Kopi, --- det
er ikke godt at se; heller ikke, hvad Fornøjelse Djævelen kan
have af den, naar den ikke ses som Hindring. --- Schopenhauers
Indflydelse, fra hvilken Nietzsche aldrig ret frigjorde sig,
ligger bag ved hans Historiefilosofi. Det er forunderligt, at
Nietzsche ikke har set, hvor pessimistisk denne Opfattelse er,
skønt han arbejdede hele den sidste Del af sin Virketid paa at
befri sig fra Pessimismen. En skarp Modsigelse kommer derfor,
som vi i det Følgende ville se, frem hos ham paa dette
Punkt.

Senere hen giver han en noget anden Motivering af sin
Aristokratisme end denne historiefilosofiske og pessimistiske. I
\textit{Menschliches---Allzumenschliches} gaar han ud fra Livets Kraft
og Kulturens Oprindelighed som Maalestok: Skal Livet overhovedet
gaa fremad og opad, kan dets Fremgang og Opadstigen
kun være repræsenteret af en højere Kaste, som kan øve frit
Arbejde, og som forudsætter en lavere Kaste, der som en
Slags Kykloper --- kan øve det materielle Arbejde, der stedse
mere eller mindre er et Tvangsarbejde. Kun hin højere Kaste
% --- p. 135 ---
% running head: Den højere Kaste Maal, ikke tillige Middel.
\opage{135}har Mulighed for Aandsfrihed. »Den højere Kaste«, siger Nietzsche
i et af sine senere Skrifter \textit{(Antichrist,} § 57), »kalder jeg:
de Færreste, og den har, da den er den fuldkomne, ogsaa de
Færrestes Forrettigheder, hvortil hører det at fremstille Lykken,
Skønheden og Godheden paa Jorden«. De to Kaster bør holdes
vel ude fra hinanden. Deres Antagonisme er nødvendig. Den
lavere Kaste skal saavidt muligt udstyres med Maskindyder eller
Hjordedyder. Religionen, den gængse Moral og de borgerlige
Dyder passe for Hjorden, ikke for de Udvalgte. »Ich möchte«,
siger Nietzsche \textit{(Der Wille zur Macht.} Aphor. 472), »die liebenswürdigen
Tugenden nicht unterschätzen; aber die Grösse der
Seele verträgt sich nicht mit ihnen.«

Den højere Kaste er kun Maal, ikke tillige Middel. Det er
ifølge Nietzsche Tegn paa Korruption, naar et Aristokrati ikke
længere betragter sig som Samfundets egentlige Mening og Retfærdiggørelse
og ikke stoler paa sin Ret til at gøre utallige
andre Mennesker til ufuldstændige Mennesker, til Slaver og Redskaber.
Samfundet er kun til som Underbygning og Stillads for
de udsøgte Væsener, af hvilke Aristokratiet bestaar --- ligesom
en Slyngplante paa Java støtter sig til et Egetræ forat naa op
i Lyset med sin Krone og der bærer sin Lykke til Skue!
\textit{(Jenseits von Gut und Böse.} Aphor. 258). Ved sin blotte Existens
udtrykker dette Aristokrati, hvad der giver Livet Værdi.
Derimod ligger dets Betydning ikke (som det f. Ex. gælder for det
Aristokrati, Platon og Comte skildrer i deres Statsidealer) i, at
det skal lede og fremhjælpe de andre Stænder. Tvertimod.
Endnu lige tilsidst \textit{(Wille zur Macht.} Aphor. 12) erklærer Nietzsche
det for et »Hovedsynspunkt«, at den højere Races Opgave
ikke bestaar i Ledelsen af den lavere, men at den lavere er
den Basis, paa hvilken en højere Race kan leve for sin egen
Opgave. Dog har selv denne højere Race et Fremtidsmaal: at
arbejde for Overmenneskets Komme. »O, mine Brødre,« siger
Zarathustra \textit{(Also sprach Zarathustra.} III), »jeg indvier Eder
til en ny Adel: I skulle være Frembringere og Opfostringsmænd,
Sædemænd for Fremtiden! ... Eders Børns Land skulle
I elske, og denne Kærligked skal være Eders nye Adel!« ---
Altsaa er der dog endnu et Maal i Fremtiden, og Historiens
% --- p. 136 ---
% running head: Kykloperne og deres Herrer.
\opage{136}Maal er ikke naat i selve Aristokratiet; en ny Racedannelse
indenfor den højere Race kræves. Problemet viser sig vanskeligere,
end det fra først af syntes.\textsuperscript{42})

Den lavere Race, Demokratiet, tjener som materielt Grundlag,
som Udfører af det materielle Arbejde. Demokratiske Indretninger
have --- trods deres Kedsommelighed --- den Nytte at holde
Tyranniet borte. Men hvorledes det er muligt at holde den
lavere Race indenfor bestemte Grænser, faa vi ikke at vide.
Nietzsche siger selv i sit første Skrift, at der intet mere Frygteligt
gives end en barbarisk Slavestand, der har lært at betragte
sin Existens som en Uret og bereder sig til at tage Hævn.
Men kan man da hindre Kykloperne i at blive vidende og i
at sammenligne? Det vil sikkert altid blive lettere at overbevise
Herrerne om, at de trænge til Slaver, end Slaverne om, at de
trænge til Herrer. Nietzsche er her altfor barnligt romantisk
--- eller maaske for barbarisk. Opretholdelsen af Militærstaten
erklæres af ham for det sidste Middel til at opretholde den store
Tradition om den høje og stærke Mennesketype \textit{(Wille zur Macht.}
Aphor. 327) --- en Ytring, der kaster et Lys over Nietzsches
indre Sammenhæng med moderne tyske Forhold. Hans Ven
Peter Gast priser da ogsaa Militæret som (Sligt maa siges paa
Tysk!): die herrlichste, straffste, männlichste Einrichtung unserer
plebejisch- und mercantilistisch effeminirten Zeit! --- Folkeoplysning
er da af det Onde, da det gør den lavere Race utilbøjelig
til at finde sig i sin Skæbne --- som Middel eller Kopi!

Dog kræves der af den højere Race et strængt Arbejde og
stor Selvbeherskelse, forat den skal kunne bevare sin Stilling.
Adelen maa regulere sit Liv paa bestemt Vis, hvis Respekten
for den skal bevares. De højere Klasser maa ikke selv leve som
Pøbel. De maa leve maadeholdent og uden til Skue stillet Yppighed.
Forgyldt Raahed og naragtig Opblæsthed nedbryder Respekten
for Kulturen. Ligeledes ere pludselig Berigelse uden
Arbejde og slet Fordeling af Ejendommen mislige. (\textit{Menschliches---
Allzumenschliches}). Her kræves dog altsaa en Hensyntagen,
som gør, at den højere Kaste eller Race ikke uden videre
kan leve Livet ud af sig selv. Den maa »regere« sig selv, og
maaske ogsaa de andre Kaster (nemlig naar Militæret hidkaldes).

% --- p. 137 ---
% running head: Nietzsche, Guyau og Carlyle.
\opage{137}Det er ikke godt at se, hvorfor Storheden skulde blive mindre
ved at øve en Funktion. Som om Kraftens Værdighed led
ved, at Kraften bruges! Hvor der virkeligt er et Overskud af
Kraft, vil der være Evne til at være Formaal og Middel paa
én Gang, og begge Dele i stor Stil. Og i Virkningens Rækkevidde
have vi nu engang den eneste Maalestok for Kraft.

Paa dette Punkt kommer Grundforskellen mellem Nietzsche
og Guyau frem. De stemme overens i at lægge stor Vægt paa
Livets Kraft og overensstemmende Fylde. Det er Genstand for
deres fælles Tro. I sine Exemplarer af Guyaus Bøger har Nietzsche
forsynet hans Udtalelser i denne Retning med Randbemærkninger,
der udtrykke Anerkendelse og Tilslutning. Men naar
Guyau i den uvilkaarlige Expansion finder et Grundlag for Hengivelse
og Sympati baade overfor Mennesker og overfor Ideer,
protesterer Nietzsche i sine Noter, og naar denne Tanke vender
tilbage hos Guyau, sætter Nietzsche i Marginen: »Fix Ide!«
Selv fortolker han nemlig den uvilkaarlige, af Kraftoverskud betingede
Virketrang som en »Villie til Magt«, og nærmere som
en Trang til at lade sin Magt gaa ud over Andre. Med Larochefoucauld
vil han føre alle Følelser og Drifter tilbage til Egoisme;
og Egoismen opfatter han særligt som Trang til Udladning
af Magt (Macht auslassen).\textsuperscript{43}) Utvivlsomt har Guyau her
set rigtigt, naar han betragter hin uvilkaarlige Trang som liggende
dybere end Modsætningen mellem Egoisme og Altruisme.
Nietzsche er her mere dogmatisk end den franske Filosof. ---

En stor Modsigelse er der i den sociale Dualisme, naar der
af den lavere Kaste kræves Beundring og Ærbødighed for en
Kultur, den selv skal staa ganske udenfor. Hvor kan Hjorden
have Respekt for de store Fænomener, naar der aldeles ingen
aandelig Forbindelse er mellem dem? --- \textsc{Carlyle} har her set
rigtigere end Nietzsche, naar han \textit{(Past and Present} I, 5) i Anledning
af Slagordet om »Talentets Aristokrati« bemærker: »Det
sande Øje for Talentet forudsætter den sande Respekt for det, og
--- o Himmel! --- forudsætter saa mange Ting!« Disse mange
Ting lader Nietzsche roligt ligge, optagen som han er, af »Afstandens
Pathos«.

% --- p. 138 ---
% running head: 
\opage{138}\sectionhead{E. \textit{Slaveopstanden i Moralen.}}

Den sociale Dualisme er ifølge Nietzsche ikke blot det Rette
og Ønskelige, men ogsaa det Naturlige. Den finder sin Forklaring
ved Moralens historiske Oprindelse \textit{(Genealogie der Moral).}
De moralske Begreber og Navne skrive sig fra de Herskende,
Kraftige og Lykkelige og udtrykke disses Lykkefølelse
og Selvfølelse i Kontrasten til de Svage og Afmægtige. De »Gode«
vare oprindeligt de Fornemme, de Mægtige, de Høje. »God« betyder
»af første Rang«; derimod har det fra først af Intet at
gøre med uegoistiske Handlinger. Den højere Kultur begynder
ofte med, at en kraftig Race overvælder en svagere, hvorved et
Afstandsforhold, ved hvis Hjælp Fortrinene kunne føles, bliver
til. Det politiske Forrangsbegreb bliver til et sjæleligt Forrangsbegreb.


Et ivrigt Theognisstudium, Nietzsche drev i sin første Ungdom,
har her sat sine Spor. Hos Theognis, en indædt Aristokrat,
der var fordreven fra sin Fødeby Megara, betyder »de
Gode« Aristokratiets Mænd og »de Onde« Folkets Mænd. Iøvrigt
er det Folkevandringen, Renaissancen og Napoleon, der tjene
ham som historiske Exempler. Muligvis have ogsaa Barndomsindtryk
virket med. Hans Fødeby Naumburg var et bureaukratisk
Knudepunkt i Provinsen, og Geheimeraadscirklernes Selvfølelse
har været den første Art »Fornemhed«, den vordende
Apostel for Aandsaristokrati har lært at kende. Karakteristisk er
det i hvert Tilfælde, at han kun betragter og undersøger Afstandsforholdet
fra den ene Side og aldeles ikke tænker paa,
hvorledes den samme Afstand virker paa den anden Part,
--- hvilke Følelser af Frygt og Ærefrygt, Tillid og Beundring de
store og straalende Skikkelser kunne vække. Det er dog snarest
fra denne Side af Forholdet, de moralske Forestillinger have
deres historiske Oprindelse.\textsuperscript{44}) Og det vilde ogsaa for Nietzsche
selv have været mest konsekvent at antage dette, da de »fornemme«
Naturer efter hans Opfattelse skulle hvile i sig selv,
leve ud af sig selv, uden reaktive Følelser. Saa er det ikke
konsekvent at lade dem danne Forholdsbegreber, eller overhovedet
blive sig Afstanden bevidst. --- Iøvrigt maa den herskende
% --- p. 139 ---
% running head: En lang Række af » Slaveopstande<< .
\opage{139}Race allerede før Erobringen have dannet et Samfund, og indenfor
dette vil der allerede have dannet sig moralske Forestillinger
paa Grundlag af en Vurdering af de Egenskaber (f. Ex.
Tapperhed), der for vedkommende Samfund vare af største
Værdi. ---

Den naturlige Vurdering, der skriver sig fra de Kraftige
og Lykkelige, hos hvem Villien til Magt kan ytre sig frit, er
nu ifølge Nietzsche bleven omstyrtet ved en moralsk Slaveopstand.
under hvilken vor Kultur endnu lider. Den skyldes
især Jøderne, --- dette Folk, der kaldte sig selv det udvalgte,
men som Tacitus med større Ret har betegnet som født til
Slaveri; --- det var deres Profeter, der gjorde det at være rig
og mægtig til Et med at være ugudelig og ond. »Verden« blev
et Skældsord, og Fattigdom og Lidelse lovpristes. Fortsættelsen
og den fulde Udvikling heraf kom i Urkristendommen. I Johannes'
Aabenbaring spores især Afmagtens Had. Der erklæredes
Krig mellem Judæa og Rom, mellem Svagheden og Styrken.
Frihed, Stolthed og aandelig Selviskhed ofredes. Især har Paulus
gennemført denne Tankegang. Apostlene forstod ikke Jesu
Død, der netop var en stor Frihedshandling; de gjorde den til
et Sonoffer, til en Forherligelse af Lidelsen. --- I sine senere
Skrifter opdynger Nietzsche en saadan Mængde Haansord i sin
Omtale af Kristendommen, at hans Søster kun kan forklare sig
det ved den overdrevne Brug af Kloral i det sidste Aar før
Katastrofen.

Der er dog ifølge Nietzsche ikke blot foregaaet én moralsk
Slaveopstand, men flere. Før Kristendommen ere allerede Buddhismen
og Sokratismen saadanne Opstande, kun af mere fornem
Art. Senere er Reformationen en Slaveopstand overfor den katolske
Kirkes og Renaissancens fornemme Verdslighed; --- ja,
ogsaa Fritænkeriet, Revolutionen, Demokratiet og --- Naturvidenskaben
(der er en Demokratisering af Naturen ved sit Princip
om almen Lovmæssighed!) ere saadanne Opstande. I stigende
Fart udvider Nietzsche Begrebet »moralsk Slaveopstand«; efterat
han én Gang har faaet Øje for det, bliver det som en fix Idé,
der anvendes paa højst forskellige, indbyrdes modstridende Bestræbelser.
--- Om det noget Steds er berettiget at finde Vidnes%
% --- p. 140 ---
% running head: De Lykkelige og de Ulykkelige.
\opage{140}byrd om Nietzsches Sygelighed i hans Ideer, saa er det paa
dette Punkt, hvor han oversaver den Gren, paa hvilken han selv
sidder. ---

Men hvad Nietzsche ikke har set, er, at »Livet« netop rører
sig rigt og kraftigt hos dem, han kalder Slaver. Urkristendommen
er netop, hvad Nietzsche kalder en dionysisk Bevægelse,
et aandeligt Gennembrud, tildels i Extasens Form, et Udtryk for
den store Længsel, et Excelsior! Og en lignende Rolle, som Delfi
spillede overfor Bacchantismen, har Hierarkiet og den kirkelige
Organisation spillet, inddæmmende den voldsomme Strøm,
--- ikke altid paa apollinsk Vis, trods den Fornemhed, hvoraf det
organiserede Hierarki præges i Nietzsches Øjne. --- Og analoge
Bemærkninger kunne gøres med Hensyn til de andre »Slaveopstande«. ---


En Begrundelse af sin Teori om Slaveopstandene har Nietzsche
fundet i den formentlige Iagttagelse, at de Lykkelige ere
bedre Mennesker end de Ulykkelige. De Lykkelige leve ud af
sig, udfolde deres Virksomhed uvilkaarligt, som indefra bestemt
Aktivitet. De behøve hverken Had eller Fordrejelse, da de ikke
ere afhængige. De Ulykkeliges Virksomhed er bestemt udefra,
er af reaktiv Art; Misundelse, Mistænksomhed, Had og Svig
ligge dem nær, afhængige som de ere. De Lykkelige sige Ja,
de Ulykkelige Nej. --- Negtes kan det naturligvis ikke, at naar
de aktive eller lykkelige Naturer »udlade deres Magt« saaledes,
at det volder Lidelse omkring dem, saa have de lettere ved at
glemme end de, hvem Magtudladningen er gaaet ud over. Men
inkonsekvent er det, naar Foragt skal spille saa stor en Rolle
i de Lykkeliges Psykologi; ti det er dog en afgjort reaktiv
Følelse. ---

Dog denne Begrundelse er ikke den definitive. Den afgørende
Begrundelse for den sociale Dualisme og for Teorien
om Slaveopstanden finder Nietzsche --- mærkværdigt nok --- i
Velfærdsprincipet.

% --- p. 141 ---
% running head: 
\opage{141}\sectionhead{F. \textit{Filosofisk Begrundelse af den sociale Dualisme.}}

Nietzsche kalder sig ofte Immoralist og erklærer, at han
vil ophæve al Moral, og han har som bekendt kaldet en af sine
Bøger »Hinsides Godt og Ondt«. Det maa dog af det Foregaaende
være klart, at det kun er Slavemoralen, han vil udrydde.
Og et Sted siger han ogsaa (i Analogi med et bekendt
Epigram af Schiller), at han vil udrydde Moralen --- »aus Moralität«.
Omvurdere vil han --- men dertil kræves en Maalestok, et Princip, der nødvendigvis maa udtrykke en Grundværdi.
Spørge vi om, hvilket dette Vurderingsprincip er, finde vi det
Overraskende, at det er Velfærdsprincipet (das menschliche Gedeihen,
se Fortalen til »Geneal. d. Moral«). Hans Grundspørgsmaal
er, om den tidligere moralske Vurdering har tjent til Menneskeslægtens
Velfærd eller ikke. Hans Hovedsynspunkt er, om
en Handling eller en Personlighed udtrykker stigende eller synkende
Livskraft. Og han stiller sig --- ganske vist noget sent
--- den Opgave at ordne Værdierne videnskabeligt i en Trinrække,
der bestemmes ved Kraftens Styrke og Omfang (Zahl- und
Maass-Scala der Kraft. \textit{Wille zur Macht.} Aphor. 353).
Desværre har han været saa optagen af at deklamere mod
Slavemoralen og mod Hjorden, at han ikke er naat til for Alvor
at behandle denne Opgave og derved i bestemtere Form
udvikle, hvad han kalder sin \textit{moralistiske Naturalisme} (W. z.
M. Aphor. 192). Det er en dionysisk Vurdering, der kræves,
--- en Livskraftens og Livsglædens Etik. Alt ligefra den Glædens
Højde, hvor Mennesket helt igennem føler sig som en guddommeliggjort
Form for og Selvretfærdiggørelse af Naturen og
ned til sunde Bønders og sunde Halvmenneskedyrs Glæde
--- hele denne lange, uhyre Lys- og Farveskala af Lykke kaldte
Grækerne, ikke uden den taknemmelige Ærefrygt, der følger
med Indvielsen i et Mysterium, Dionysos! (W. z. M. Aphor.
482). Det gælder altsaa at tilvejebringe den Lykke, der svarer
til det givne Livstrin. Velfærdsetiken, som Nietzsche saa ofte
haaner, kan --- visse »Fornemheder«, der skjule over uløste Problemer,
fraregnede --- gaa ind paa denne Opfattelse, især da
det udtrykkeligt erklæres, at det ikke gælder nogen Enkelts
% --- p. 142 ---
% running head: Nietzsche som Utilitarianer.
\opage{142}eller nogen Kastes Lykke, men at den kraftige Livsudfoldelse
og Livslykke, der kræves, kræves for Slægtens Skyld
--- nur des Gesammt-Lebens willen. Kun for Slægtens Skyld i det Hele
kan der tilkendes den Enkeltes Selvopholdelse Værdi; egentligt
grunder sig Begrebet om Individ, om den Enkelte, saaledes som
man hidtil forstod det, paa en Vildfarelse. \textit{(Götzendämmerung.}
Aphor. 33).

Den sociale Dualisme og hele Herremoralen finder da heri
sin sidste Begrundelse. Herrerne og Overmenneskene have da
ikke blot deres Maal i sig selv; den sidste Vurdering af dem
beror paa, om de ligge paa Slægtslivets opstigende Linie eller
ikke, --- beror altsaa paa deres Bidrag til Menneskelivets Udvikling.
--- Nietzsches Discipel og Fortolker Peter Gast indrømmer,
at Mesteren i en vis Forstand er Utilitarianer; kun skiller
han sig fra den sædvanlige Utilitarianisme, der kun tager Hensyn
til den allernærmeste Nytte. (Hvor denne sædvanlige Utilitarianisme
findes, oplyses desværre ikke. Utilitarianismen forlanger
i alle sine Former, at Handlingernes Virkninger skulle
forfølges saa langt som muligt).

Nu var det i sig meget vel muligt, at Slægtens Velfærd
kunde kræve en skarp Modsætning mellem Herrer og Slaver,
skønt det er vanskeligt at se, hvorledes en saadan Modsætning
vil kunne hævdes og gennemføres paa det vesteuropæiske Kulturtrin.
Men et alvorligt Forsøg paa Paavisning heraf er ikke
gjort. Stuart Mill har i sin Tid ud fra sin demokratiske Grundanskuelse
gjort et langt alvorligere Forsøg paa at vise Nødvendigheden
af individuel Frihed, Ejendommelighed og Storhed, end
Nietzsche ud fra sin aristokratiske Grundanskuelse har gjort paa
at vise Nødvendigheden af en Slavestand. Og dog haaner Nietzsche
Mill og de andre engelske Tænkere, hvis Skotvinger han
--- som streng, metodisk Forsker betragtet --- ikke er
værdig at løse.

En Ting staar i hvert Tilfælde fast: det var tænkeligt, at
en social Dualisme tjente Menneskeslægten bedst; men det er
utænkeligt, fordi det er selvmodsigende, at man paa én Gang
kan se det sidste Maal i enkelte Individer eller Kaster alene og
% --- p. 143 ---
% running head: Det højeste Udtryk for Villien til Magt.
\opage{143}dog i samme Øjeblik se hele Slægtens Velfærd som det sidste
Maal.

Men paa de sidste, afgørende Punkter er Nietzsche ikke
Filosof, men Digter. Hvad han ikke naaede til at gennemarbejde
med sin Tanke, det har staaet for ham i et stort Fremtidssyn,
og mod dette maa vi vende os for at høre hans sidste Ord.\textsuperscript{45})

\sectionhead{G. \textit{Det sidste Ja --- og Zarathustras Død.}}

Hvad der driver Nietzsche frem til en Tanke eller en Vision,
der staar i saa skarp Modsigelse til, hvad han forud eller
samtidigt hævder som Polemiker, er en stor Kontrastvirkning,
en vældig Trang til aldeles at faa Bugt med den Livsopfattelse,
han i sin første Ungdom var kommen ind paa under Schopenhauers
og Wagners Indflydelse. Som Polemiker reagerer han
mod sine Samtidige; i sine Syner reagerer han mod sig selv,
og han kan ikke faa denne Reaktion stærkt nok udtrykt. Tankens
Form slaar ikke til her; men til Gengæld staar en høj,
af mægtig Stemning baaren Poesi til hans Raadighed. Det er
Grundtanken fra sit Ungdomsskrift om Tragediens Fødsel, han
nu vender tilbage til. Livet skal forherliges. Livstrangen ytrer
sig som Villie til Magt, og denne Villie kan og skal være saa
stærk og modig, at den endog kan vælge Smerten --- og vælge
den ikke blot som Middel, --- vælge at gentage Livet, det selvsamme
Liv, den engang har lært at kende, med dets Møje og
Smerte.

Zarathustra, Profeten, har oppe i sin Bjerghule samlet de
højere Mennesker, dem, der have lidt under den demokratiske
Pøbels Miskendelse og Forfølgelse. Han forklarer dem, at de
stedse ville faa det værre og haardere; stedse flere af dem ville
gaa til Grunde. Det er nemlig ikke Menneskets Bestaaen, det
gælder om, men Overmenneskets Tilblivelse. Hvad vi kunne
elske hos Menneskene, er en Overgang og en Undergang. Der
er saa store, fjerne og høje Kræfter i Virksomhed, at det intet
Under er, om det skrøbelige Kar, i hvilket de udvikle sig, sønderbrydes.
Medlidenheden med disse højere Mennesker er Zarathustras
sidste Synd. \textit{(Also sprach Zarathustra.} IV). ---

% --- p. 144 ---
% running head: Overmenneskets Tilblivelse.
\opage{144}Hvad er da Overmennesket? Det er en Existensform, der
forholder sig til Mennesket, som dette til Aben. Dets Tilblivelse
er Betydningen af Jordens Existens (der Sinn der Erde). Det
er en ny Livstype, der skal virkeliggøres, og ved Ordet Overmenneske
udtrykkes denne Types Begreb eller Symbol. \textit{(Wille
zur Macht.} Aphor. 390). Denne Types Virkeliggørelse vil Menneskeheden,
naar den store Middagstime kommer midt paa Udviklingens
Vej, se som sit Haab og som sin Opgave. Maalet
kommer altsaa dog til at ligge i Fremtiden og i en fjern Fremtid,
hvad Nietzsche forhen protesterede imod. Han antager ikke
længere, at Historien har naat sit Maal i Fortidens Store; det
er ham ikke nok. Men hvad der især er af Betydning: de Antydninger,
Nietzsche i sine mere positive Udviklinger giver af
den nye Livstype, pege i en anden Retning end hans Skildringer
af Herremoralen. Det er ikke længere den blotte Kraft, men
den højsindede Menneskekærlighed, han kræver. Selv i Kristendommen
finder han nu for liden Kærlighed, idet den raabte Ve
over de Leende. --- Saasnart Nietzsche positivt skal udvikle sit
Ideal og holder de reaktive Stemninger borte, undergaar Begrebet
Overmenneske en Forskydning. Det ses først og fremmest
hos Zarathustra selv. Da han har faaet de højere Mennesker
til at gaa ind paa sin Tanke, aner han, at den store Middag
er nær. Nu kan han ikke blive i sin ophøjede Bjergensomhed.
Han trænger til ligesom Solen at lade sit Lys skinne for Verden.
»Du høje Stjerne,« siger han til Solen, »du store Lykkeøje,
hvad vilde al din Lykke være, naar du ikke havde Nogen
at lyse for!« (Also sprach Zarathustra. IV. Das Zeichen). »Endnu
engang vil jeg ned til Menneskene. Hos dem vil jeg gaa til
Grunde, døende vil jeg give dem min rigeste Gave. Det har
jeg lært af Solen, naar den gaar ned, den overvættes Rige:
Guld udøser den da i Havet af sin uudtømmelige Rigdom, saa
at den fattigste Fisker kommer til at ro med gylden Aare. Det
saa jeg selv engang, og jeg blev ikke mæt af Taarer derover!«
(Also spr. Zar. III. Von alten und neuen Tafeln). Nu hader
Zarathustra sit eget Had. »Jeg er bleven til En, der velsigner
og siger Ja! Længe kæmpede jeg og var en Stridsmand, forat
jeg kunde faa Hænderne fri til at velsigne!« (ibid. Vor Sonnen%
% --- p. 145 ---
% running head: Midnatssangen.
\opage{145}aufgang). --- Her ytrer Villien til Magt sig unægteligt paa en
anden Maade end hvor Foragten for Hjorden har Ordet. Og det
er selve Livstrangen, Trangen til at forjage Pessimismen ud
af alle Kroge, der fører Nietzsche til dette Højdepunkt. Det er
den konsekvente Udfoldelse af Nietzsches Grundtanke, der her
fuldbyrdes. ---

Nu, da Zarathustra har faaet saa klart Blik paa Maalet og
paa sin rette Opgave, istemmer han sin Midnatssang. Denne
Sang (Nietzsche kalder den »den ensomste Sang, som nogensinde
er skreven«) havde Nietzsche skrevet en Nat paa en Loggia
paa Piazza Barbarina i Rom; han henlægger den i sit Digt
til Midnatstimen foran Zarathustras Hule. Der er i denne Strofe
anslaaet en Helhedsstemning og antydet en sammenfattende
Tanke, som han ellers ikke formaaede at fastholde og udrede.
Han er her kommen saa nær som muligt til et stort, koncentreret
Udtryk for begejstret Resignation og tilkæmpet Haab. ---
Jeg gør intet Forsøg paa at oversætte den, men gengiver Tankerne,
saa godt jeg formaar.

»Midnatten forkynder: Verden er dyb, langt dybere, end
Dagen tænker! --- Dyb er Verdens Vé. Men Lyst er endnu
dybere end Hjertevé. Vé betyder Undergang. Men Lyst vil
Evighed, vil dyb, dyb Evighed!«

Gennem Smerten, Sorgen, Tvivlen og Foragten havde Nietzsche
arbejdet sig ind i en dyb Tro paa Livet. Han udtaler her
i digterisk Form den gamle biologiske Teori om Lystfølelsen
som Udtryk for Livets Kraft og Fremgang. Og han ser i enhver
Rørelse af Lyst en Villie til Livets Bestaaen og Fortsættelse.


Men denne sidste Tanke, Villien til Livets Bestaaen, antager
hos ham en mærkelig Form. Den bliver Et med en Villie
til at gentage Livet. Ti for Nietzsche var det kommet til at
staa som en Nødvendighed, at det, der sker, maa komme igen
paa samme Maade i Fremtiden. Han gaar ud fra, at Verden
bestaar af et bestemt Antal Elementer, saa at kun et bestemt,
endeligt Antal af Kombinationer bliver muligt. Er det udtømt,
maa Gentagelsen komme. Verdensløbet er et Kredsløb, der allerede
har gentaget sig uendeligt ofte og ligeledes vil gentage sig
% --- p. 146 ---
% running head: >>Wohlan, noch einmal ! <<<
\opage{146}i Fremtiden.\textsuperscript{46}) Da denne Tanke gik op for Nietzsche, greb den
ham med Forfærdelse. Han betragtede den som en naturvidenskabelig
Nødvendighed, og den satte hans Optimisme paa en
haard Prøve. Ikke blot Smerten skulde nu vælges og tages,
men Alt, hvad der var lidt, skulde lides atter og atter igen.
Des større blev Sejrsjubelen, da han overvandt Frygten for Gentagelsen,
og i denne Overvindelse fandt den højeste Grad af
Livsbekræftelse, som overhovedet var mulig. Paa en interessant
Maade mødes Nietzsche her med Kierkegaard, for hvem jo ogsaa
Gentagelsen var en Prøvesten paa Livsanskuelsens Kraft og
Alvor. Kierkegaards Sætning, at »den, der vil Gentagelsen, han
er en Mand«, er ogsaa en af Nietzsches Grundsætninger.
--- Denne Tanke er Grundtanken i Zarathustrabogen, hvor den dog
ofte utydeliggøres ved, at Aforismer, der staa i mindre klar
Sammenhæng med den, trænge sig frem. Dersom Nietzsche
havde faaet udarbejdet Bogens Slutning, vilde Bogens Ide være
kommen bedre til sin Ret.

Denne Tanke om Genkomsten forkynder Zarathustra først
for de højere Mennesker, han har samlet om sig ved sin Bjerghule.
Efter nogen Modstand faar han dem til at gaa ind paa
den, saa selv »det hæsligste Menneske« udraaber: »War \textit{Das}
--- das Leben? --- Wohlan, noch einmal!« Derefter er det saa,
Zarathustra stiger ned forat forkynde Livets Vilkaar for den
store Menneskemasse. Han samler Menneskene til Fest og giver
dem nye Love. Der fastsættes en Rangforordning, hvor de virkelige
Livsværdier lægges til Grund. Stændernes Kamp er nu
lykkeligt overstaaet. Det erkendes udtrykkeligt for Herskerkastens
Opgave at skaffe sig de Beherskedes dybe, ubetingede
Tillid. Hadet til den demokratiske Nivelleren har havt sin Tid,
som nu er forbi. Zarathustra skrider saa til sin egentlige Opgave.
Først forkynder han det store Haab om Overmenneskets
Komme, der bliver muligt ved den nye Vurderings Gennemførelse.
Og derpaa kommer det store, frygtelige Øjeblik, da
han forkynder Menneskene, at Alt vil komme igen. Nu kan
dog denne Tanke udholdes, ikke blot af ham selv, men ogsaa af
Menneskene. Ti da han spørger dem: »Ville I alt dette endnu
engang?« --- svare Alle ja, og Zarathustra dør af Glæde der%
% --- p. 147 ---
% running head: Kritik af Læren om Gentagelsen.
\opage{147}over. (Se Udkastene til 5. og 6. Del af »Also sprach Zarathustra«.
Werke XII p. 321 ff.). ---

I poetisk Form har Nietzsche her taget den sociale Dualisme
og Herremoralen tilbage. Allerede det, at Zarathustra
stiger ned forat forkynde Menneskenes Mængde de højeste Sandheder,
viser dette, foruden at det ligger i, at Stændernes Kamp
er afsluttet. Der er nu et fælles Maal for Slægten.

Man vil ikke kunne negte, at det er en haard Prøve,
Nietzsche sætter vor Livstro paa. Det er jo dog muligt at have
en stor Tro paa Livet, selv om man ikke ønsker igen at begaa
de samme Dumheder eller Slyngelstreger og lide de samme
Smerter og Sorger atter og atter. Men der er heller ikke
Grund til at antage absolut Gentagelse. Erfaringen viser os
den ikke, og Tilværelsen staar, jo mere vor Erkendelse fordyber
sig i den, som uudtømmelig; stedse afslører den større
Mangfoldighed og fjernere Muligheder. En stridende Livsanskuelse
--- og enhver Livsanskuelse, der vil se Virkeligheden
under Øjne, maa være en stridende --- maa have en Hovedplads
for Tanken om det Nye og Ubekendte, der stikker op
snart som Fare, der truer, snart som Opgave, der ægger, snart
som Haab, der vinker. Der er Intet, der taler for, at Tilværelsen
skulde være saa snever, at et begrænset Antal af Existenser
i al Evighed skulde afløse hverandre. Vor Tro paa Livet
behøver derfor ikke at gennemgaa den Skærsild, Nietzsche i
dionysisk Begejstring vil lade den gaa igennem.

Som han ellers har ladet Disharmoniens og Foragtens
Stemning potensere til det Yderste, saaledes har han i sine afsluttende
Tanker ladet Reaktionen mod al Pessimisme naa den
højeste Potens. Det er Følelsesexpansioner, som paa hvert Punkt
beherske hans Tanke. Ikke ved videnskabelig Problembehandling,
men ved den lidenskabelige Maade og den ofte geniale Pathos,
hvormed indbyrdes modstræbende Synspunkter komme til Orde
hos ham og derved stilles klart og skarpt frem, hører han
hjemme i Filosofiens Historie. Hans Betydning er symptomatisk.
Han har ladet os blive Vidne til et indre Drama i en
Sjæl, der følte Tidens og Livets Tendenser dybt og energisk.

% --- p. 148 ---
\backmatter
\chapter*{NOTER.}
\addcontentsline{toc}{chapter}{Noter}
\markboth{Noter}{Noter}
\opage{148}1) S. 8. Denne Frygt for Materialisme kan føre til betænkelige Følger,
idet Wundt for at hævde Psykologiens Selvstændighed ganske naturligt
søger at begrænse »Elementerne« til det mindst Mulige, da det jo
kun er ved Elementerne, ikke ved deres Forbindelse, at et fysiologisk
Korrelat skal være nødvendigt. Saaledes har han fundet det betænkeligt,
at jeg har kaldt den nye Egenskab, en Fornemmelse faar, naar den genkendes,
for en Kvalitet (»Bekendthedskvaliteten«), idet han mener, at jeg
derved føres til at sætte en fysiologisk Forklaring i Stedet for en psykologisk.
Og han finder heri et Tidens Tegn: »Dieser Forscher ist durchaus
nicht den principiellen Vertretern der materialistischen Psychologie zuzurechnen.
Gerade dies aber scheint mir bezeichnend für das Ueberwuchern
hypothetischer und psychologisch unfruchtbarer physiologischer Erklärungsversuche
in der heutigen Literatur, dan selbst so ... unbefangene Psychologen
wie Høffding sich nicht ganz von denselben freihalten können .«
(Philos. Studien. X. p. 61). Hertil har jeg at svare, at efter min Opfattelse
maa man ved alle psykiske Processer spørge om, hvorvidt det er
muligt at finde en fysiologisk Forklaring; --- det gælder for Forbindelser
og Sammenligninger, saavel som for de enkelte Elementer og Kvaliteter.
Og ifølge Identitetshypotesen, som det jo ogsaa er Wundts Mening
at fastholde, udelukker en fysiologisk Forklaring ikke den psykologiske;
ethvert Fænomen bør, saavidt muligt, belyses fra begge Sider. Jeg opfatter
psykologisk den umiddelbare Genkendelse i Analogi med Lighedsassociation
og Sammenligning, som et Grænsetilfælde i Forhold til dem,
nemlig som den mest elementære Form, i hvilken Lighedsforholdet gør
sig gældende i Bevidstheden. Den psykologiske Forklaring, Wundt giver,
udvisker Forskellen mellem de mere elementære og de mere sammensatte
psykiske Processer, idet han antager, at frie Forestillinger medvirke ogsaa
i de simpleste Tilfælde af Genkendelse. --- En lignende Stilling som overfor
den umiddelbare Genkendelse indtager Wundt ved Drøftelsen af Webers
% --- p. 149 ---
% running head: Noter til Side 9--12 .
\opage{149}Lov, idet han skelner mellem selve Fornemmelserne og deres »Sammenligning«.
\textit{(Physiologische Psychologie.} 5. Udg. I. p. 541). Det er et stort
Spørgsmaal, om denne Distinktion er holdbar, hvad de mest elementære
Fornemmelsesprocesser
angaar. --- Af Literatur om Wundt vil jeg især
henvise til \emph{Edmund König}: \textit{W. Wundt.} Seine Philosophie und Psychologie.
Stuttgart 1901. (Frommanns Klassiker der Philosophie). \emph{Allan}
\emph{Vannérus}: \textit{Vid studiet af Wundts psykologi.} (Stockholm 1896), Nytte har
jeg ogsaa havt af en utrykt Konferensafhandling om Wundt som Psykolog
af min unge Ven Mag. art. \emph{Aug. Bjarnason}.

2) S. 9. Smlgn f. Ex. \emph{Wundt}: \textit{System der Philosophie.} (1. Udg.) p.
533: »Da den aandelige Energis Væxt kun bestaar i de organiske Formers
\textit{kvalitative} Fuldkommengørelse, forbliver den kvantitative Bestemmelse
af de fysisk-kemiske Energier fuldstændig uberørt deraf .... Ved alt
Fremskridt i organisk Udvikling formeres Naturkræfternes Kvantitet ligesaa
lidt som Materiens Kvantitet. Men i Værdi ere Naturkræfterne og deres
Substrater voxede umaadeligt ved det organiske Livs Udvikling. Ja, Værdibestemmelser
ere jo overhovedet først bleven mulige, og tillige nødvendige,
ved Opstaaen af formaalstjenlige Villieshandlinger og af de til dem
knyttede Forestillinger og Følelser. Smlgn. \textit{Physiol. Psychol.} 5 III. p. 781.

3) S. 9. Det aktuelle Sjælebegreb er især blevet gjort gældende af
Fichte, men ligger allerede i Kants Begreb om Bevidstheden som Syntese.
Hos os er det hævdet af F. C. Sibbern og af mig. (Se især mine \textit{Psykologiske
Undersøgelser.} p. 67---81).

4) S. 12. Wundt mener, at jeg i mine \textit{Psykologiske Undersøgelser} (p.
60) har misforstaaet hans Udtalelser om Apperceptionen i tredie Udgave
af »Physiologische Psychologie«, naar jeg deraf, at Forskellen mellem Association
og Apperception skal bero paa en Abstraktion, har sluttet, at
han ikke antager dem for forskellige Processer; der skal deri kun ligge,
»dass sie thatsächlich \textit{immer zusammen} vorkommen«. (Philos. Studien. VII
p. 229 f.), Men i den senere udkomne fjerde Udgave af »Physiologische
Psychologie« hedder det (II. p. 447) om de simpleste Forestillingsforbindelser:
»Diese Vorgänge unterscheiden sich durch die völlige Einflusslosigkeit
des Willens auf die Art ihres Eintritts auf das bestimmteste von
den nachher zu erörternden apperceptiven Verbindungen der Vorstellungen«,
--- og fremdeles (p. 479): »Die Association bringt die Vorstellungen nur
in diejenigen Verbindungen, in die sie vermöge ihrer eigenen Beschaffenheit,
unbeeinflusst von jeder inneren Willensthätigkeit, sich ordnen.« Og
naar Wundt stadigt (ogsaa i femte Udgave af sit psykologiske Hovedværk)
antager et særeget Apperceptionscentrum i Hjernen, synes ogsaa deri at
ligge, at Apperception og Association ikke altid behøve at forekomme sammen.
--- Den Aktivitet, der ytrer sig i Apperceptionen, blive vi os ifølge
Wundt umiddelbart bevidste ved en »Thätigkeitsgefühl«, som er fælles
for al Apperception. (Physiol. Psychol.\textsuperscript{4} II. p. 266, 270, 279). Andre Steder
(f. Ex. Philos. Stud. X. p. 109) fremstiller han Apperceptionen, eller
% --- p. 150 ---
% running head: Noter til Side 13-25.
\opage{150}overhovedet Villiesvirksomheden som Genstand for umiddelbar Anskuen.
(Smlgn. ogsaa Physiol. Psychol.\textsuperscript{4} II. p. 560: »Der Begriff der Thätigkeit
stammt zunächst nur aus unseren eigenen Willenshandlungen und wurde
erst von ihnen auf äussere bewegte Gegenstände übertragen«). I sine tidligere
Fremstillinger satte Wundt Bevidstheden om Aktivitet i nøje Forbindelse
med den umiddelbare Fornemmelse af Innervation, som han senere
har forkastet. --- Om de Vanskeligheder, der ere forbundne med enhver
Art af Aktivitetsbevidsthed, taler Wundt mærkværdigvis ikke. Smlgn. om
dette Spørgsmaal min Psykologi VII B. 4 og Psykologiske Undersøgelser
p. 68---81.

5) S. 13. Smlgn. Wundt's Afhandling \textit{Ueber die Definition der Psychologie}
(Philos. Studien. XII) p. 517. \emph{Friedrich Paulsen}: \textit{Einleitung
in die Philosophie.} Berlin 1892. p. 116---132 (Intellektualistische und voluntaristische
Psychologie).

6) S. 14. I min Fremstilling af Psykologien gik jeg allerede i første
Udgave (1882) ud paa at paavise en psykisk Aktivitet i Sansning, Forestillingsassociation,
Lyst og Ulyst saavel som i hvad der specielt kaldes
Villen, og jeg fandt i Villien Sjælelivets egentligste og inderligste Natur
(IV, 7 e). Men dersom jeg nu kunde omarbejde den fra Grunden af, vilde
jeg ogsaa i Fremstillingsformen gøre den til en Villiespsykologi.

7) S. 19. Karakteristiske Udtaler af \emph{Wundt} om Metafysiken findes
i  \textit{Philos. Studien} XIII. p. 80; 428. \textit{Einleitung in die Philosophie.} p. 348---352.

8) S. 20. Wundts Udvikling af dette Punkt findes i hans \textit{System der
Philosophie} IV, 1, 3. (1. Udg. p. 361---367). Jeg finder ikke hans Tankegang
ganske tydelig, men har fremstillet den efter min Tydning af den.

9) S. 23. Smlgn. om dette hele Spørgsmaal min \textit{Psykologi.} p. 66.
\textit{Religionsfilosofi.} p. 66–86. \textit{Filosofiske Problemer} (Universitetsprogram 1902).
p. 62---65. --- I sin fortræffelige Mindetale over Fechner karakteriserer
Wundt med Rette Fechners Spekulation som bygget paa Analogi, i Modsætning
til den naturvidenskabelige Metode, der kræver Induktion. Forholdet
mellem Analogi og Induktion er efter Wundt det, at hin begrænser
sig til et lille Antal af Kendemærker, medens denne kræver saa mange
som muligt. \textit{(Gustav Theodor Fechner.} Rede zur Feier seines hundertjährigen
Geburtstages. Leipzig 1901. p. 70---73). Men det, der ved Analogien
først og fremmest maa tages Hensyn til, er Udgangspunktet (hvad
Wundt kalder »Angriffspunkt«) eller Lighedspunktet. Analogien er en Forholdslighed,
og det Afgørende vil være, i hvilken Grad denne Forholdslighed
kan betragtes som Tegn paa en dybere liggende Identitet. I Metafysik
og Religion tydes Tilværelsen som Helhed ved Forestillingen om en
enkelt af dens Dele eller Sider, og Analogien maa her stedse vedblive at
være ufuldkommen.

10) S. 24. Smlgn. min \textit{Religionsfilosofi.} p. 217–246; 334–339.

11) S. 25. Ved »Formaalenes Heterogoni« forstaar Wundt det Samme,
% --- p. 151 ---
% running head: Noter til Side 26-38 .
\opage{151}jeg kalder objektiv Værdiforskydning (se min \textit{Etik} XIII, 4). Han synes
at antage, at Motivforskydningen (den subjektive Værdiforskydning) stedse
forudsætter, at der i Virkningen er Mere end i Hensigten. Dette er dog
ikke altid nødvendigt, da dels Indlevelsen i Handlingens Resultat, dels
andre tiltrædende Forhold kunne muliggøre nye Motivers Opstaaen. --- I
rent psykologisk Form er Motivforskydningen allerede beskreven af Spinoza,
Hartley og James Mill. I Historiens Filosofi har Hegel lagt stor
Vægt paa »Formaalenes Metamorfose«.

12) S. 26. Smlgn. mine kritiske Bemærkninger om Wundts Etik i
\textit{Etiske Undersøgelser.} p. 28 (hvor jeg har vist, at Wundt i sin Polemik
mod Velfærdsmoralen glemmer sin egen Lære om »Formaalenes Metamorfose«)
og i \textit{Etik} VIII, 4; XIII, 4.

13) S. 31. Om den italienske Filosofis Historie før Ardigò henviser
jeg til \emph{Louis Ferri}: \textit{Histoire de la philosophie en Italie an 19e siècle.}
(Paris 1869). \emph{Alfred Espinas}: \textit{La philosophie experimentale en Italie.}
(Paris 1880). I Ueberwegs Haandbog findes en af \emph{Luigi Credaro} forfattet
Oversigt. Med Hensyn til Fordømmelsen af Kriticismen smlgn.
\emph{Rudolf Eucken}: \textit{Thomas von Aqvino und Kant,} ein Kamp zweier Welten.
Berlin 1901.

14) S. 36. Disse sidste Bemærkninger om Ækvivalensen forekomme i
Skriftet \textit{L'unità della coscienza} (p. 410---413), som er langt senere end Afhandlingen
\textit{La formazione naturale.} De forekomme mig ikke ganske nøjagtige,
ti naar det indrømmes, at der foruden de indre Dispositioner maa
komme ydre Betingelser til, forat en »distinzione« skal indtræde, saa have
de specielle Former dog egentligt ikke deres Ækvivalenter i Indistinto
alene, men i dettes Dispositioner + de ydre Betingelser. Dette ligger nu
egentligt ogsaa i Ardigòs Lære, naar hau erklærer ethvert Indistinto for
begrænset og Forskellen mellem indistinto og distinto for relativ.

15) S. 36. Quando it pensiero perdi di vista l'infinito, fissandosi nel
distinto finito, esso infinito lo assiste inosservato, e costituisce la stessa
forza della logica del suo discorso. \textit{Form. nat.} p. 136. (Opere filos. II).

16) S. 37. Ardigò er klar over. at Gyldigheden ikke beror pua Tankens
Lighed med Genstanden, som tænkes, men paa dens egen indre Natur.
\textit{(La psicologia come scienza positiva.} p. 228) (Opere filos. I). Men dermed
er Gyldigheden ikke forklaret. Et andet Sted (se \textit{L'unità della coscienza.}
p. 441) (Opere filos. VII) kalder han Sandhed for en Egenskab ved Tanken
(una qualità del peusiero). Men naar Tanken er et Fænomen ved
Siden af andre Fænomener, hvad har saa hin Egenskab ved Tanken at
gøre med disse? --- En vis Dogmatisme ytrer sig her hos Ardigò.

17) S. 38. I en senere Udgave af \textit{La psicologia} (1882) føjer Ardigò
følgende Note til Omtalen af »den psykofysiske Realitet«: Questa espressione,
sostanza psicofisica, non è altro che l'Indistinto naturale precedente
e sottostante ai due fenomeni distinti del mondo della psiche e di quella
% --- p. 152 ---
% running head: Noter til Side 42-53.
\opage{152}della materia. (p. 387). --- Begrebet Indistinto indførtes først ved Afhandlingen
La formazione naturale (1877).

18) S. 42. Ardigò polemiserer her specielt mod Mamiani. Se Afhandlingen
\textit{La religione di T. Mamiani} i Opere philos. Vol. II. ---

19) S. 43. \emph{William James} har ganske vist nyligt i varme Ord udtalt
sin Tilslutning til den gamle engelske Skole og erklæret, at han bygger
paa dens Grund. (Smlgn. herom min Afhandling \textit{Om nogle religionsfilosofiske
Arbejder fra den nyeste Tid.} Universitetsprogram 1903. p. 71). Men
denne Tilslutning angaar kun visse almindelige Grundsætninger, der nu
anerkendes af de Fleste, som filosofere, og som iøvrigt hos James selv
fremtræde med Modifikationer, der ikke kunne udledes af selve den ældre
engelske Skoles Principer.

20) S. 46. Foruden de anførte Skrifter har Bradley ogsaa udgivet en
hel Række Afhandlinger i Tidsskriftet \textit{Mind.} I den følgende Fremstilling
ere ogsaa nogle af disse Afhandlinger benyttede.

21) S. 53. Bradley staar i bestemt Modstrid mod Voluntarismen. Det
følger af hele hans Standpunkt, at han ikke i en Appel til Villien kan
finde nogen Løsning af sine Gaader. »Hvad vi kende som Villie, forudsætter
en Relation og en Proces, og forudsætter tillige en uforsonet Disharmoni
af Elementer. Det Samme gælder Energi eller Aktivitet, eller
andet Lignende.« \textit{(Appearance and Reality.} p. 483). Bradley gaar her ud
fra, at det Begreb, paa hvilket vi bygge vor Verdensanskuelse, selv maa
være afsluttet og ikke føre til nye Spørgsmaal. Han gør, ligesom Goethe,
»Urfænomenet« til Noget, der ikke kræver Forklaring. (Smlgn. mine
\textit{Filosofiske Problemer.} p. 56 f.). At Villien er et Begreb, der sætter nye
Spørgsmaal i Bevægelse, deri har Bradley utvivlsomt Ret. Schopenhauer,
Voluntarismens Stamfader i det 19. Aarhundrede, har overset det. Men
gælder det Samme ikke om alle Begreber, man vil prøve at opstille som
afsluttende? Og gælder det ikke først og fremmest om Begrebet Erfaring,
som Bradley vil lægge til Grund? --- Jeg ser her bort fra den rent psykologiske
Side af Sagen. Bradley har endnu ikke offentliggjort sin Psykologi.
Dog ses det af hans Afhandlinger i »Mind« \textit{(On active attention.}
1902. --- \textit{The definition of will.} 1902---1903), at hans Psykologi har en
overvejende intellektualistisk Karakter.

Overfor James Ward, der havde kritiseret Bradley fra Voluntarismens
Standpunkt, bemærker A. E. \emph{Taylor}, en Tænker, der staar Bradley nær,
at det dog maa være klart, at Aktivitet ikke kan være Udtryk for det
Højeste eller det Absolutc, da Aktivitet stedse forudsætter en Modstand.
(Mind. 1900. p. 258). (Smlgn. min Religionsfilosofi p. 64; 81). Overfor
denne Indvending gives der sikkert ingen anden Udvej end den, filosofiske
Teister ofte have grebet, at antage Guddommen for begrænset. Saaledes
\emph{Rashdall} i Afhandlingen \textit{Personality, human and divine} (i »Personal
Idealism«). Problemet om det Onde har især ført flere Tænkere til dette
Resultat. Da ser man ganske vist med afgjort Tydelighed, hvorledes Pro%
% --- p. 153 ---
% running head: Noter til Side 59-83.
\opage{153}blemerne vedblive at stille sig paany. »Det Absolute« er ifølge Rashdall
ikke Gud, men »et Samfund, der omfatter Gud og alle Aander«. --- Det
nytter aabenbart ikke at føre Problembehandlingen over paa det mytologiske
Omraade.

22) S. 59. Enkelte Træk til Taines Karakteristik har jeg hentet fra
\emph{Boutmy} \textit{(Taine, Renan, Schérer.} Paris 1900) og fra \emph{Albert Sorels}
Mindetale over Taine ved sin Optagelse i Akadémie française (1895).

23) S. 63. \emph{Charles Renouvier}: \textit{Esquisse d'une classification systematique
des doctrines philosophiques.} Paris 1886. II. p. 395. --- Smlgn.
iøvrigt om Renan G. \emph{Séailles}: Ernest Renan. Essai de biographie
psychologique. Paris 1894. Mme \emph{Darmstetter}: \textit{La vie d'Ernest Renan.}
Paris  1898. \emph{Ed. Platzhoff}: \textit{Ernest Renan.} Ein Lebensbild. Leipzig
1900.

24) S. 65. Smlgn. mine \textit{Psykologiske Undersøgelser} (1889), 76. ---
\textit{Psykologi} VII B, 4.

25) S. 71. Renouvier har faaet denne Tanke fra sin tidligt afdøde
Ven \emph{Jules Lequier}, hvem han ofte henviser til. I Afhandlingen \textit{Doute
ow croyance} (L'année philosophique 1896) p. 44---51 anfører han en længere
Udtalelse af ham. --- Den hele Opfattelse minder om \emph{Fichtes} Filosofi.
Striden mellem Dogmatisme og Idealisme maa jo ifølge Fichte afgøres
ved et Valg, og Valget vil afhænge af, hvad Slags Menneske man
er. Se om dette Valg især Fichtes Skrift \textit{Die Bestimmung des Mensches.}
(Berlin 1800). p. 186; 193.

26) S. 76. Se nærmere herom min \textit{Psykologi} V D, og \textit{Filosofiske Problemer.}
p. 60; 62. --- Efterat Texten er gaaet i Trykkeriet, faar jeg Efterretning
om, at Renouvier er afgaaet ved Døden.

27) S. 83. Se Indledningen til Maxwell's Afhandling om Faradays
Kraftlinier (tysk Oversættelse ved Boltzmann i Ostwald's Klassiker der
Naturwissenschaft Nr. 69 p. 4). Ideen om, at Tallærens Anvendelse paa
Naturfænomen beror paa Analogi, er siden udført af \emph{Helmholtz} i hans
Afhandling i Festskriftet til Zeller (1887) \textit{Zählen und Messen, erkenntnisstheoretisch
betrachtet,} og af \emph{Ernst Mach} i hans Skrift \textit{Die Principien der
Wärmelehre, historisch-critisch erläutert} (1896). --- Hvor optagen Maxwell
har været af Begrebet Analogi, ses af, at der fra den samme Tid, i hvilken
han skrev Afhandlingen om Faradays Kraftteorier, haves et Udkast
fra hans Haand, som drøfter det Spørgsmaal, om der findes reale Analogier
i Naturen. (Det er trykt hos \emph{Campbell and Garnett}: \textit{The
Life of James Clerk Maxwell.} London 1882. p. 235---244). Ganske klart
forekommer det mig ikke at være, hvad Slutningsresultatet angaar. Det
er egentligt kun en Skizze, der antyder forskellige Analogier, som kunne
faa Betydning for Erkendelsen. Et Vidnesbyrd om den kritiske Filosofis
Indflydelse paa Maxwells Tankegang er hans Bemærkning om Analogien
mellem Grund og Aarsag: Reasons, when spoken of with relation to
% --- p. 154 ---
% running head: Noter til Side 86-102 .
\opage{154}objects, get the name of causes, which are reasons, analogically referred
to objects instead of thoughts. (p. 238).

28) S. 86. Paa et Møde af University College Christian Association
i Maj 1903 udtalte Lord Kelvin, at man, hvad Livets Oprindelse angaar,
ikke kunde sige, at Videnskaben hverken hævdede eller benægtede en
skabende Kraft; nej, Videnskaben \textit{hævdede positivt} en skabende og ledende
Kraft, da de rent fysiske Kræfter ikke kunde forklare Livet. Denne Udtalelse
foranledigede en livlig Diskussion, i hvilken Kelvins Paastand blev
bestemt imødegaaet af Thiselton-Dyer, Karl Pearson, Ray-Lankester o. fl.
\textit{(Times weekly.} 8. 15. og 22. May). Det kan kun være den store Fysikers
Autoritet paa hans eget Omraade, der har vakt saa stor Opmærksomhed
for denne ukritiske Udtalelse.

29) S. 88. Kirchhoffs Opfattelse fremstiller \emph{Boltzmann} \textit{(Gustav Robert
Kirchhoff.} Festrede 1888. p. 25) saaledes: »Det gælder ikke om at
danne dristige Hypoteser om Materiens Væsen og ud af Molekulernes Bevægelser
at gætte sig til Legemernes Bevægelser, men at danne Ligninger,
der uden Hypoteser saa fuldstændigt som muligt og med kvantitativ
Nøjagtighed svare til Erfaringsverdenen, ubekymret om Tingenes Væsen
og Kræfter. Ja, i sin Bog om Mekanik vil Kirchhoff endog banlyse alle
metafysiske Begreber, som Kraft eller Bevægelsesaarsag; han søger blot
de Ligninger, der med saa stor Nøjagtighed som muligt svare til de iagttagne
Bevægelser.«

30) S. 95. Smlgn. \textit{Filosofiske Problemer.} p. 46.

31) S. 95. Paa den filosofiske Kongres i Paris 1900 blev det under
Drøftelsen af en Afhandling, der sluttede sig til Mach's og Kirchhoff's
Standpunkt, udtalt, at naar man i Videnskaben sætter Accelerationen i
Stedet for Kraften, er det, fordi man vil sætte visuelle Symboler i Stedet
for muskulære, hvilket igen har sin Grund i, at hine ere nøjagtigt
maalelige.

32) S. 95. Smlgn. om dette dynamisk-symbolske Sandhedsbegreb \textit{Filosofiske
Problemer.} p. 37 f.

33) S. 102. Smlgn. \textit{Filosofiske Problemer} p. 15; 18. --- Avenarius'
Discipel Fr. \emph{Carstanjen} har givet et Resumé af første Bind af »Kritik
der reinen Erfahrung« i sit Skrift \textit{Richard Avenarius' biomechanische Grundlegung
der neuen allgemeinen Erkenntnisstheorie} (München 1894). Paa flere
Punkter bliver Hovedværkets Fremstilling tydeligere ved Carstanjen's Fortolkning
--- men Tydeligheden naas kun ved, at der benyttes psykologiske
Exempler, som hentes fra Værkets anden Del, der egentligt slet ikke
skulde være benyttede, da den »uafhængige« Række dog burde belyse
den »afhængige«, ikke omvendt. Desuden føjer Carstanjev selv flere »Symptomer«
til, som jeg i Texten har benyttet. --- Den i Texten anførte mundtlige
Ytring af Avenarius findes i Afhandlingen \textit{Richard Avenarius. Ein
Nachruf} von Fr. Carstanjen (Vierteljahrschrift für wissenschaftliche Philosophie
XX). --- Det er uberettiget, naar \emph{Wundt} i sin iøvrigt fortrinlige
% --- p. 155 ---
% running head: Noter til Side 106-122.
\opage{155}Kritik af Avenarius \textit{(Ueber naiven und kritischen Realismus.} Philos. Studien.
XIII. p. 331) kalder hans Lære Materialisme.

34) S. 106. Smlgn. min \textit{Religionsfilosofi.} p. 240.

35) S. 106. Medens Avenarius har givet en Problemernes Naturhistorie,
satte jeg mig i min Afhandling \textit{Filosofiske Problemer} den Opgave at finde
det for de forskellige Hovedproblemers Indhold Karakteristiske. Min Undersøgelse
bekræfter forsaavidt Avenarius', som det viser sig at være Forholdet
mellem Diskontinuitet og Kontinuitet, der kommer igen indenfor
hvert saakaldt Problem. En væsentlig Forskel er det, at jeg betoner Erfaringens
(Tilværelsens) Uudtømmelighed, medens Avenarius opererer med
Forestillingen om en Slutningstilstand. Jeg betoner derfor stærkere den
blivende Irrationalitet i vort intellektuelle Forhold til Tilværelsen.

36) S. 113. Fra et lignende Synspunkt som her er Guyaus Teori
bleven kritiseret i den skarpsindige Afhandling af \emph{Ch. Christophle}: \textit{Le
principe de la vie comme mobile moral selon J. M. Guyau.} Étude critique.
(1901). (Særtryk af Revue de Metaphysique et de Morale).

37) S. 118. Smlgn. herom min \textit{Religionsfilosofi.} p. 103. --- \emph{Fouillée}
(\textit{Le mouvemeut idealiste.} p. XXII) kalder Guyaus »Irreligion« en Fremtidsreligion,
skønt Guyau selv ikke vilde have den kaldt saa. F. \emph{Buisson}
\textit{(L'èducation morale et l'èducation religieuse,} i Samlingen »Questione de Morale«.
Paris 1900. p. 329) mener, at den fremfor Alt har Krav paa Navnet
Religion. --- Navnet har naturligvis ikke Stort at sige. --- En god
Karakteristik af Guyau og en Oversigt over hans Ideer har hans Stedfader
Fouillée givet i Skriftet \textit{La morale, l'art et la religion d'après Guyau.}
Paris 1889. Den amerikanske Religionsfilosof \emph{Josiah Royce} har ogsaa
givet en meget sympatetisk Karakteristik af Guyau (i en Afhandling i
hans \textit{Studies in Good and Evil.} New York 1898. p. 379---384).

38) S. 122. I sin Bog \textit{Richard Wagner, poète et peuseur} (3. éd. p.
301; 429---432) har \emph{Henri Lichtenberger} søgt at vise, at Nietzsche
gør Wagner Uret, naar han mener, at denne har »omvendt« sig. Der er
--- trods forskellige Nuanceringer --- en Kontinuitet i Wagners Livsanskuelse.
At han senere stiller Kristendommen over Buddhismen, medens
han tidligere stillede dem ved Siden af hinanden, er væsentligt et Udtryk
for, at han --- ligesom paa sin Vis Nietzsche selv --- reagerede
mod Pessimismen.
Konfessionel blev han aldrig i sin religiøse Opfattelse.

39) S. 122. Nietzsches Søster, Fru \emph{Elisabeth Förster-Nietzsche}
har begyndt Udgivelsen af en Biografi, af hvilken hidtil to Bind (\textit{Das
Leben Friedrich Nietzsches.} I og II, 1. Leipzig 1895---1897) ere udkomne.
De gaa til 1880. Dertil kommer Nietzsches \textit{Briefe} (2. Bde. 1900---1902). I
\emph{Crusius'} Biografi af Erwin Rohde findes mange Bidrag til Oplysning
om og Forstaaelse af Nietzsche; ligeledes i Paul Deussen's smukke
Bog \textit{Erinnerungen an Nietzsche} (Leipzig 1901). Nietzsches moderlige Veninde
\emph{Malwida von Meysenbug} har i sin Bog \textit{Individualitäten} (1901)
p. 1---41 givet en interessant Skildring af Nietzsches Personlighed.

% --- p. 156 ---
% running head: Noter til Side 133-146.
\opage{156}40) S. 133. Mod Sokrates er Nietzsche overhovedet streng og uretfærdig.
Han kan ikke finde Haansord nok imod ham. I »Geburt der Tragödie«
er det foruden Intellektualismen den optimistiske Tro paa Livet,
der bebrejdes ham. Senere \textit{(Fröliche Wissenschaft.} Aphor. 340) maa han
høre ilde for den Pessimisme, der for Nietzsche ligger i hans sidste Ord
hos Platon, Anmodningen om at ofre Arklepios en Hane. Han har betragtet
Livet som en Sygdom! --- Nietzsche følger her den traditionelle, nyplatoniske
Fortolkning af den bekendte Ytring. Professor \emph{Heiberg} har
vist, at denne kan fortolkes paa en langt simplere, mindre mystisk Maade,
naar man ser paa hele den Sammenhæng, i hvilken den forekommer. (Se
Oversigten over det danske Videnskabernes Selskabs Forhandlinger 1902.
\emph{J. L. Heiberg}: \textit{Sokrates' sidste Ord).}

41) S. 134. I sit sidste Skrift \textit{Wille zur Macht} betegner Nietzsche
selv sin Anskuelse som Aristokratisme (Werke XV. p. 427). Paa det citerede
Sted synes Nietzsche at polemisere mod Guyau og Fouillée, der
lægge saa stor Vægt paa den sociologiske Betragtning og dens Betydning
for Livsanskuelsen.

42) S. 136. Der er her indtraadt en Ændring i Nietzsches Opfattelse,
ti i \textit{Unzeitgemeine Betrachtungen} (Drittes Stück § 5) er det Opgaven for
\textit{Alle} at arbejde hen til det store Menneskes Opstaaen. Her er Dualismen
endnu ikke saa skarp, og Problemet ikke saa indviklet.

43) S. 137. Nietzsches Exemplarer af Guyaus Bøger med hans Randbemærkninger
findes i Nietzsche-Arkivet i Weimar, hvorfra \emph{Fouillée}
har faaet de vigtigste af disse i Afskrift. Han har benyttet dem i et interessant
Kapitel i sin Bog \textit{Nietzsche et l'Immoralisme} (1902). Livre III:
Les jugements de Nietzsche sur Guyau d'après des documents inédits.

44) S. 138. Smlgn. min \textit{Etik.} X, 3 (mod Slutningen) om Forbilledernes
Betydning --- i Analogi med Variationernes Betydning for Biologien.

45) S. 143. Jeg gaar ikke ind paa de Antydninger i erkendelsesteoretisk
og metafysisk Retning, Nietzsche har givet, især i tredie Bog af
\textit{Wille zur Macht.} I sin Opfattelse af Erkendelsens Principer minder han
om den økonomisk-biologiske Erkendelsesteori, og endnu mere afgjort minder
han i sin Villiesmetafysik om Schopenhauers Filosofi (kun at han taler
om »Villien til Magt« i Stedet for om »Villien til Livet«).

46) S. 146. Tanken om alle Tings Genkomst er ikke ny. Nietzsche
har sandsynligvis først lært den at kende gennem sine klassiske Studier.
Det er en antik Tanke, og den hænger nøje sammen med Forestillingen
om Verden som et begrænset Hele. Naar tillige Elementer og Kræfter
antages uforgængelige, er en rytmisk Gentagelse nødvendig. Hos Stoikerne
kom en saadan Gentagelse efter enhver Verdensbrand, baade for
Guder og Mennesker. Alt gaar da for sig paa samme Maade; Sokrates
ægter igen Zantippe o. s. v. Det var for Stoikerne (se især \emph{Marcus}
\emph{Aurelius}. Comm. XI, 1) en Støtte for deres Sindsro, at der intet Nyt
kom; hverken ville vore Efterkommere faa noget væsentligt Nyt at se, ej
% --- p. 157 ---
% running head: Note til Side 146.
\opage{157}heller have vore Forfædre set en anden Verden end vi. --- Den er i nyere
Tid oftere udtalt som formentlig Konsekvens af den naturvidenskabelige
Opfattelse. Saaledes af Blanqui, Le Bon, Nàgeli, Guyau. (Smlgn. \emph{Fouillée}:
Nietzsche et l'Immoralisme. IV, 3: Le retour éternel). Ogsaa hos
Dostojevski kommer Tanken frem --- som en Anfægtelse fra selve Djævelen!
(Smlgn. \emph{Merejkowsky}: \textit{Tolstoi et Dostojewsky.} Paris 1903. p. 300).
--- Gentagelsen er naturligvis kun nødvendig, naar man forudsætter Verdens
Begrænsning og tillige har en dogmatisk Tillid til de videnskabelige
Grundsætningers absolute Gyldighed. Dersom det derimod er os umuligt
at tænke nogen absolut Grænse for Verden i Tid, i Rum, i Energi, og
dersom hine Grundsætninger kun ere de ledende Hypoteser for vor Forsken,
udelukkes ikke Muligheden af nye Elementers og Kombinationers
Opstaaen. Forholdsloven fører os naturligt i denne Retning og stiller os
derved overfor Muligheden af Tilværelsens Irrationalitet. Smlgn. herom
mine \textit{Filosofiske Problemer.} Kap. II og III. Herefter vil snarest Goethe
have Ret, naar han siger: »Neu ist immer die Erfahrung«. --- Hos Nietzsche
er Læren om Gentagelsens Nødvendighed ikke konsekvent. Ti han
lægger i sin Erkendelsesteori stor Vægt paa, at Identitetsprincipet og de
andre Grundsætninger ere vor Villies Postulater, ndtrykke vort Ønske om
at bringe Naturen i vor Magt. Det er os, der lægge Identitet ind i Naturen.
»Man soll die Nöthigung, Begriffe, Gattungen, Formen, Zwecke,
Gesetze zu bilden (»eine Welt der identischen Fälle«), nicht so verstehen,
als ob wir damit die wahre Welt zu fixiren im Stande wären; sondern
als Nöthigung, uns eine Welt zurecht zu machen, bei der unsre Existenz
ermöglicht wird: --- wir schaffen damit eine Welt, die berechenbar, vereinfacht,
verständlich u. s. w. für uns ist.« \textit{(Der Wille zur Macht.} Aphor.
279). Naar Nietzsche antager, at der ikke i Virkeligheden gives identiske
Tilfælde, maa han konsekvent opgive Troen paa absolut Gentagelse. Erfaringen
viser jo ogsaa kun en Tilnærmelse til absolut Identitet og Gentagelse.
Nietzsches Metafysik og hans Erkendelsesteori stemme ikke overens.
Han polemiserer jo desuden undertiden mod Naturvidenskaben, fordi
den indfører Lighed for Loven i Naturen!

\end{document}
